\documentclass[UTF8,11pt]{ctexart}
\usepackage{amsmath,amssymb,amsthm,mathtools}
\usepackage{booktabs,longtable,array,geometry,enumitem}
\usepackage{xcolor}
\definecolor{linkblue}{rgb}{0.10,0.20,0.45}
\usepackage[colorlinks=true,linkcolor=linkblue,citecolor=linkblue,urlcolor=linkblue]{hyperref}
\geometry{a4paper,margin=2.4cm}
\setlength{\parskip}{2pt}

\newtheorem{theorem}{定理}[section]
\newtheorem{proposition}[theorem]{命题}
\newtheorem{lemma}[theorem]{引理}
\newtheorem{corollary}[theorem]{推论}
\newtheorem{conjecture}[theorem]{猜想}
\theoremstyle{definition}
\newtheorem{definition}[theorem]{定义}
\newtheorem{remark}[theorem]{注}
\newtheorem{example}[theorem]{例}

\newcommand{\Lam}{\Lambda}
\newcommand{\tr}{\operatorname{tr}}
\newcommand{\rank}{\operatorname{rank}}
\newcommand{\dd}{\,\mathrm{d}}
\newcommand{\R}{\mathbb{R}}
\newcommand{\N}{\mathbb{N}}
\newcommand{\Q}{\mathbb{Q}}
\newcommand{\Z}{\mathbb{Z}}
\newcommand{\one}{\mathbf{1}}
\newcommand{\STW}{\textup{[STW26]}}

\title{\bfseries 图的前三大特征值之和的精确上界\\[4pt]
\large 阶段性进展报告与障碍分析}
\author{}
\date{2026 年 10 月 29 日（第七版，增 \S\ref{sec:rowvar}：行方差机制与天花板 $L(N)$；$Z_5$-等变秩三投影的完全分类与 $\rho_0$ 的闭式导出；$\Phi(Q)=\lambda_{\max}(Q_+)$ 的凸性与谱盒归约；旧队首 P1 被精确有理证书否证）}

\begin{document}
\maketitle
\vspace{-2.2em}

\begin{abstract}
\noindent 设 $G$ 为 $n$ 阶图，$\lambda_1(G)\ge\cdots\ge\lambda_n(G)$ 为其邻接矩阵的特征值，
$\Lam_k(G)=\lambda_1(G)+\cdots+\lambda_k(G)$，$M_k(n)=\max\{\Lam_k(G):|V(G)|=n\}$。
本报告讨论 $k=3$ 的情形。目前已知的最好无条件上界是 Sivashankar--Tang--Wakhare
\STW\ 的 Proposition 6.1 在 $k=3$ 处的推论
$\Lam_3(G)\le\Phi n-3$，$\Phi=(3+\sqrt5)/4=1.309016994374947$。
另一方面，我们在前期工作中构造出一个六原子的加权核心（$K_1\vee C_5$ 的团吹起），
证明了 $\liminf_{n}M_3(n)/n\ge\rho_0$，其中
$\rho_0=(9+\sqrt5+2\sqrt{35})/18=1.281568196872168$。
因此真值落在 $[\rho_0,\Phi]$ 这个宽度为 $2.0969\%$（相对 $\Phi$）的区间内，
本课题的目标是证明上端也可以取到 $\rho_0$。

本报告做三件事。第一，完整拆解 \STW\ 的证明，说明它为什么\emph{必然}停在 $\Phi$，
以及它的证明架构中哪一部分是深刻的、哪一部分是形式化的。
第二，汇报到目前为止已经严格建立的结果：一条把「对一切 $n$、一切图」压缩为
「一个有限维极值问题」的归约链，$\rho_0$ 的精确下界证书，
极值核心的完整闭式刻画，以及若干条被定量否证的路线。
第三，给出本轮的新结果（定理~\ref{thm:zonal}）：
\emph{整个}「逐点主函数 + 逐通道线性计费」的方案族——它包含 \STW\ 的证法、
包含把 $|t|$ 换成正部 $(t)_+$ 的自然推广、也包含任意高阶 Gegenbauer 修正项与
一类非球对称的耦合通道——其可达到的最好常数\emph{恰好}是 $\Phi$。
证明只有三行，见证测度支撑在正二十面体等角线系的内收集
$\{-1,-1/\sqrt5,1/\sqrt5,1\}$ 上。
这条否定性结果的意义在于：它在动手之前排除了我们原先排在首位的一条技术路线，
并且精确指出，要突破 $\Phi$ 就必须使用 \STW\ 的框架\emph{丢掉}的那部分信息，
即投影的自洽性（掩码约束）$A=1\{Q>0\}$。

第四，汇报本轮的突破性进展（第~\ref{sec:domdual} 节）。我们找到了一个与 \STW\ 方向完全不同的凸化：
不去松弛 Grassmannian 上的投影变量，而是用一个半正定矩阵 $T$ 去\emph{优占}掩码 $\Sigma$。
由此得到定理~\ref{thm:domdual}：$F(\Sigma)\le G(\Sigma)=\Omega(\Sigma)$，其中
$\Omega(\Sigma)=\min\{\max_iT_{ii}:T\succeq\Sigma,\ T\succeq0\}$ 是一个半定规划，
$G(\Sigma)=\frac12\bigl(1+\max_p\|D_{\sqrt p}\Sigma D_{\sqrt p}\|_*\bigr)$ 是它的对偶（核范数形式），
中间的等号由强对偶给出、是\emph{恒等式}而非放缩。
在冠军掩码 $\Sigma_0=I+A(K_1\vee C_5)$ 上这个松弛\textbf{零间隙}：
我们给出只含三个根式常数的显式证书 $T=\rho_0I+cB_{\mathrm{hc}}+dB_{\mathrm{d1}}+eB_{\mathrm{d2}}$，
$c=(4\rho_0+1-\varphi)/5$、$d=\rho_0-1/\sqrt5$、$e=\rho_0-\varphi+1/\sqrt5$，
并用符号计算验证 $\mathrm{spec}(T)=\{6\rho_0-2\varphi,\varphi,\varphi,0,0,0\}$、
$\mathrm{spec}(T-\Sigma_0)=\{6\rho_0-2\varphi-4,\varphi^{-1},\varphi^{-1},0,0,0\}$。
于是 $F(\Sigma_0)=G(\Sigma_0)=\Omega(\Sigma_0)=\rho_0$ 是\textbf{已证且精确}的：
W1-全局的冠军腔部分就此闭合，第~\ref{sec:w1global} 节设想的临界点穷举与矩--SOS 证书对它已不再需要，
前期那个「$f$ 在 $\Delta^5$ 上凹」的数值猜想也可以划掉。
同一方法给四类近腔掩码配上了\textbf{精确有理数}的半正定证书（$63$ 个主子式在 $\Q$ 上逐个判非负），
余量 $7.5\times10^{-3}$ 至 $5.1\times10^{-2}$；其中两类原本只是数值的腔面全局最大值，现在升级成定理。
但这条路\textbf{不能}覆盖全部掩码：八面体掩码的 $\Omega$ 精确等于 $4/3>\rho_0$（超出 $4.04\%$），
双条目近腔 class13 的 $\Omega=1.30667>\rho_0$。这给出一条\emph{定量}的否定性结论——
任何不使用自洽性 $A=1\{\Pi_3(M)\ge0\}$ 的论证都不可能沿优占对偶证出 A1$'$。
作为替代目标我们曾提出猜想 A2（负质量界 $\tr((-M)_+)\le\tau=\rho_0-1$），它是凸的、蕴含 A1$'$、
且被冠军\emph{精确}饱和。
\textbf{本轮把 A2 否证了}（第~\ref{sec:A2no} 小节，命题~\ref{prop:A2no}，纯有理算术、全程无浮点）：
取 $A=I_7+A_H$，$H=\overline{\widehat{C_5}}$（$\widehat{C_5}$ 是在 $C_5$ 的两个相邻顶点上各挂一片叶子），
$p$ 为 $F(A)$ 的极大化点，则 $(A,p)$ 自洽（符号余量 $\ge4.9\times10^{-2}$）、$p_i>0.12$，
而 $\tr((-M)_+)=0.288182336567654>\tau$，超出 $6.6\times10^{-3}$。
失效机制被完全定位（注~\ref{rmk:A2why}）：A2 \emph{就是}「在 $\Lam_3$ 的 Ky Fan 对偶
$\sum_{i\ge4}\lambda_i=\min\{\tr(ZM):0\preceq Z\preceq I,\ \tr Z=N-3\}$ 中删去迹约束 $\tr Z=N-3$」，
删去后 $\min$ 变成 $-\tr((-M)_+)$，等号恰在 $M$ 恰有 $3$ 个正特征值时成立；
反例有 $4$ 个正特征值，故 $\Lam_3=1.257711$ 而 $1+\tr((-M)_+)=1.288182$。
\textbf{A1$'$ 没有被违反}（反例的 $\Lam_3$ 比 $\rho_0$ 低 $2.4\times10^{-2}$），
但优占对偶 $+$ 负质量这条凸化路线就此结束——它与命题~\ref{prop:octnogo} 是同一类教训的两次出现：
\emph{任何丢掉 $\Pi_3$ 秩三结构的凸化都不够}。
穷举还给出精确分层（注~\ref{rmk:A2strata}）：$115$ 个自洽掩码中 $\mathrm{npos}\le3$ 的 $106$ 个上
$\max\tr((-M)_+)$ \emph{恰好}等于 $\tau$（冠军饱和），违反 A2 的 $9$ 个全部是 $\mathrm{npos}=4$。
修复形式 A2$^\sharp$（尾谱和 $\sum_{i\ge4}\lambda_i\ge-\tau$）与 A1$'$ \emph{等价}（命题~\ref{prop:A2sharp}），
因此不再提供凸化杠杆，但它指认了正确的对象是「第四条及以后的谱尾和」而非「负质量」。

第五，汇报本轮真正的主线成果（第~\ref{sec:kappa} 节）。前面各节都是「固定掩码 $A$、解一个有限维极值问题」，
每个掩码一个 SDP、每个腔面一个球面问题，注定要面对 $N\le7$ 已 $1252$ 个同构类、$N\ge8$ 无法枚举的局面。
本轮把问题\textbf{倒过来}：以秩三正交投影 $Q=VV^\top$（$V\in\mathrm{St}(N,3)$）为自变量，
定义单一常数 $\kappa_3=\sup_{N}\ \sup_{V\in\mathrm{St}(N,3)}\rho\bigl((VV^\top)_+\bigr)$，
并证明（定理~\ref{thm:kappa}、推论~\ref{cor:kappamain}）
\[
\Lam_3(G)\le\rho_0\,n-3\ \text{对一切 }n\text{ 与一切图 }G\text{ 成立}
\quad\Longleftrightarrow\quad \kappa_3\le\rho_0 .
\]
这把「无穷多个极值问题」压成\textbf{一条}不含图、不含掩码、不含阶数的不等式。
配套结果：命题~\ref{prop:perronrig} 建立 Perron 刚性恒等式组 (a)--(h)（全部已证，证书 $64$ 项核验），
其中 $\delta=\beta^2+\alpha(1-\alpha)$ 与 Cauchy--Schwarz 的 $\beta^2\le\delta$ 联立即得 $0\le\alpha\le1$；
注~\ref{rmk:rigiditylimit} 诚实地指出\textbf{这组恒等式单独最多只能给到 $\rho\le\frac32$}，
缺失的是一条「第四不等式」（形如 $\alpha+\ell\le2\rho_0$，$\ell=x^\top|Q|x$），
冠军处 Cauchy--Schwarz 余量恰为 $\alpha(1-\alpha)=0.022856959$；
命题~\ref{prop:decouple} 则把 \STW\ 的损失\emph{精确}劈成两半——
$(1-\alpha)=0.023404741$（投影能量损失）与 $(\phi-\ell)=0.031492854$（$\ell_1$ 损失），
其和恰为 $2(\Phi-\rho_0)=0.054897595$，因此任何只处理其中一半的论证都不可能达到 $\rho_0$。
定理~\ref{thm:ladder} 给出极值阶梯的精确值 $\kappa_3(3)=1$、$\kappa_3(4)=\frac54$、
$\kappa_3(5)=\frac{9+4\sqrt6}{15}=1.253197265$，以及 $N=6,\dots,14$ 处 $\kappa_3(N)=\rho_0$
的\textbf{数值}认证——其中 $N=6,\dots,12$ 由两套互相独立的求解器同时给出
（自洽场 SCF，\texttt{\_kappa\_profile.json}；Stiefel 流形直接上升，\texttt{\_kappa\_direct.json}，
Stiefel-KKT 残差 $\le3\times10^{-9}$、与闭式差 $\le1.6\times10^{-15}$，证书 F 段独立复跑 $N=3,\dots,8$），
$N=13,14$ 只有 SCF 一套——即「六个原子」是达到最优常数的最小规模。
另外，Collatz 包络给出一条\emph{否证性}结论：冠军处 $\frac1N\|Q_+\|_1=1.279704891<\rho_0$，
故 $\ell_1$/行和/优占这一整族上界\textbf{不可能}达到 $\rho_0$，与优占对偶的失效互相印证。
最后交代一条勘误（注~\ref{rmk:normcaution}）：定理~\ref{thm:reduce} 原证明中
「$\frac1n\sum_{ij}(Q)_{ij}^+=\sum_{I,J}w_Iw_J\langle v_I,v_J\rangle^+$」这一步\textbf{是错的}
（冠军处左端 $1.279704891$、右端 $0.214681093$，且 $\langle P,Q\rangle=\|Q_+\|_1=7.678229\ne\Lam_3$）；
结论 $\Lam_3(G)\le\kappa_3n-3$ 本身已在定理~\ref{thm:kappa} 中被\emph{独立重证}，
不受影响，但原证明的规范化段落仍待重写（优先队列 P3）。

第六，汇报本轮的主线推进（第~\ref{sec:n6global} 节）。三件事。
其一，\textbf{$N\le6$ 全部 $205$ 个同构类的逐个精确有理认证}（定理~\ref{thm:ncert6}）：
$182$ 类的 $\Omega\le\rho_0$ 由 $19442$ 个 Bareiss 精确分数主子式判定（最小者 $7.523\times10^{-37}>0$，全程无浮点），
例外只有 $23$ 个——$N=3,4$ \textbf{无例外}（故 $\kappa_3(3),\kappa_3(4)\le\rho_0$ 无条件成立，推论~\ref{cor:k34uncond}），
$N=5$ 恰 $1$ 个（$\Sigma_{K_{2,3}}$，$\Omega=\frac{22}{17}$ 有显式有理证书），
$N=6$ 恰 $22$ 个（冠军 $+$ $21$ 个 $\Omega\in[1.2887,\frac43]$ 的类）。
其二，用 Grassmannian 恒等式 $F(\Sigma)=\max_W\lambda^+_{\max}(P_W\circ\Sigma)$（命题~\ref{prop:grassid}，已证）
与 $N=5$ 时的对偶参数化 $P_W=I_5-ZZ^\top$、$Z\in\mathrm{St}(5,2)$（命题~\ref{prop:framered}），
把 $K_{2,3}$ 这一类的缺口从「$\Delta_4$ 上五次方程两个最小根的显式下界」换成\textbf{一条平面框架不等式}
$\lambda_{\max}(B(Z))\le\frac54$，其中 $B(Z)=\operatorname{diag}(1-\|z_i\|^2)+[(-\langle z_i,z_j\rangle)_+]_{K_{2,3}}$
逐元素非负、常数 $\frac54$ \textbf{sharp}（命题~\ref{prop:biaxsharp} 给出双轴框架 $Z^\ast$ 的精确有理取等证书：
$\chi_{B(Z^\ast)}=\lambda(4\lambda-5)(4\lambda-1)(4\lambda-3)^2/4^4$，$\operatorname{spec}=\{\frac54,\frac34,\frac34,\frac14,0\}$）。
于是 $\kappa_3(5)\le\rho_0$ 现在\emph{只}依赖这一条不等式（定理~\ref{thm:n5frame}），余量 $\rho_0-\frac54=0.031568197$。
其三，对引理~\ref{lem:planar} 二分法的双轴支给出\textbf{定理}：双轴族由 $15$ 个划分 $\times8$ 个符号类 $=120$ 类
与一个恒为 $3$ 维的权多面体组成（引理~\ref{lem:biaxparam}，并把它精确改写为
$\lambda_{\max}(B)=1+\lambda_{\max}(D_u\widetilde\Sigma D_u)$、$u\in S^{m_1-1}_+\times S^{m_2-1}_+$）；
$120$ 个\emph{等分}构型的 $\frac54I-B\succeq0$ 由 $31$ 个有理主子式逐个判定，其中恰 $3$ 个取等
（命题~\ref{prop:biaxenum}，已证且精确），第二大值为 $\frac{3+\sqrt3}4=1.183012701892$。
另有一条 Schur 补充分条件 $L(Z)\le1$（命题~\ref{prop:schurcond}）在 $Z^\ast$ 处取等，
数值上覆盖 $99.831\%$ 的帧（$937614$ 帧中违例 $1583$ 个，$\sup L=1.117916903$），故\textbf{不是}完整证明。
本轮还\textbf{否证}了三条旧结论/路线：
旧 \texttt{lem:planar} 的「$K_{2,3}$ 不是腔掩码」被一个显式平面框架推翻（其倍角最大间隙\emph{恰为} $180^\circ$，
旧间隙计数依赖严格不等式），据此建立的 $\kappa_3(5)\le\rho_0$ 无条件装配\textbf{作废并撤回}；
上一轮「两个内部临界点是排序 bug 的假象」也被推翻——修正代码后两点复现，对称点还有闭式
$a^\ast=\frac{62-6\sqrt{35}}{119}$、$\Lam_3=\frac{63+12\sqrt{35}}{119}=1.125991238631895$（命题~\ref{prop:k23cp}），
故「极大化点必在边界」这条捷径不可用；
翻转链路线（命题~\ref{prop:kappa6red}）经球面采样否证（$647$ 种翻转像的 $\max\Omega=1.310920>\rho_0$）。
第~\ref{sec:refuted5} 节还定量否证了十条候选上界（无权/加权 Collatz、块界 $\Psi$、秩一与秩二优占族
$\mathbf{B3}/\mathbf{B7}(\theta)/\mathbf{B4}(\theta)$、无指示器 Schur 版、双轴族凹性等），
每条都给出反例帧与超出量；注~\ref{rmk:lesson5} 从中提炼出机制性结论：
$\frac54$ 由\emph{两个互不相容}的机制同时饱和，故任何单一类型的放缩必被另一机制的构型顶穿。

其四（本轮突破，第六版新增），把上述固定 $N=5$ 的整条路线\textbf{绕过}。
第~\ref{sec:bipglobal} 节证明了一条对\emph{一切}阶数同时成立的定理：对任意 $a,b\ge1$、$N=a+b\ge4$、任意 $p\in\Delta_{N-1}$，
\begin{equation*}
\Lam_3\bigl(D_{\sqrt p}(I+A(K_{a,b}))D_{\sqrt p}\bigr)\ \le\ \frac54 ,
\end{equation*}
取等当且仅当 $a,b\ge2$ 且 $p$ 恰有四个非零分量、每部各两个、都等于 $\frac14$（定理~\ref{thm:bipglobal}，\textbf{已证}）。
于是 $F(\Sigma_{K_{a,b}})=\frac54$（$a,b\ge2$）\textbf{精确}，而这是\emph{全局}结果、不是逐阶核验。
证法只有四步、全部初等：秩二分解 $R=\frac12\bigl(\mathbf a\mathbf a^\top-(S\mathbf a)(S\mathbf a)^\top\bigr)$
配 Sylvester 惯性律得 $\operatorname{neg}(M)\le1$（引理~\ref{lem:biprank}）；
$R\succeq-\sqrt{\alpha\beta}\,e_-e_-^\top$ 配 Weyl 得 $\lambda_i\ge q^\downarrow_{i+1}$（引理~\ref{lem:bipweyl}）；
秩二长期方程 $\det(\lambda I-M)=\prod_i(\lambda-p_i)[1-A(\lambda)B(\lambda)]$ 配单调性得 $h(-\sigma)\le1\iff\sigma\ge s$
（引理~\ref{lem:bipsec}）；最后核心引理 $s\le\frac14+\sigma_4(p)$（引理~\ref{lem:core4}）只用四条初等不等式——
组内 Jensen（二元缺陷 $\Delta(u,v)$ 的显式非负分子）、$\phi(z)\le z/\sigma$、
Titu 引理（五个分支的缺口分子 $16x$、$(3x-y)^2$、$4(x-y)^2$、$(x-3y)^2$、$16y$）与 AM--GM。
直接后果有三条。\textbf{其一}，$\kappa_3(5)\le\rho_0$ 现在\textbf{无条件}成立、且 $F(\Sigma_{K_{2,3}})=\frac54$ 精确
（推论~\ref{cor:n5uncond}），旧的 $(\ast_5)$ 与平面框架不等式 \eqref{eq:frame5} 退出关键路径、降级为独立兴趣
（注~\ref{rmk:bipvsframe}）。\textbf{其二}，$N=6$ 的例外从 $22$ 降到 $20$：$K_{2,4}$ 与 $K_{3,3}$ 被认证为 $F=\frac54$ 精确，
$156$ 类中已闭合 $137$ 类，剩 $19$ 类（命题~\ref{prop:tab19} 给出完整表）；
更重要的是注~\ref{rmk:tab19read}(iv) 把 $N=6$ 的闭合\emph{重新表述}为一条纯几何陈述——
不需要对这 $19$ 类逐个求 $F$ 的上界，只需证明它们都不是腔掩码（$\Sigma\ne1\{Q\ge0\}$，$Q\in\mathcal P_3(6)$），
其中 $17$ 类已由引理~\ref{lem:bdryrig} 归约到这一步。\textbf{其三}，完全多部情形的秩 $k$ 长期方程
$\sum_t\frac1{1-A_t}=k-1$ 与 $\operatorname{inertia}(R)=(1,k-1)$ 已证（引理~\ref{lem:mpsec}），
但 $k\ge3$ 时 Weyl 的指标平移变成 $k-1$、尾和记账退化，$K_{2,2,2}$ 与 $K_{1,2,3}$ 的闭合仍\textbf{未证}
（数值 $\sup\Lam_3\le1.2524<\rho_0$，余量 $2.9\times10^{-2}$；注~\ref{rmk:mpopen} 精确指出了断点所在）。

第七，汇报第 12 轮的工作（第~\ref{sec:rowvar} 节）。这一轮有一条坏消息、一个新机制和一条主定理。
\textbf{坏消息}（否证）：第六版的队首 P1——「证 $N=6$ 的 $17$ 个无支配点类不是腔掩码」——被一个\emph{精确有理}证书推翻：
$Q=\frac12(I_6-M)$（$M$ 为完美匹配的置换矩阵）属于 $\mathcal P_3(6)$ 且 $1\{Q\ge0\}=I_6+A(K_{2,2,2})$，
故 $K_{2,2,2}$ \emph{是}腔掩码（命题~\ref{prop:chamberdead}）；进一步数值扫描显示 $19$ 个未闭合类\textbf{全部}是腔掩码
（符号违背量 $\le5\times10^{-19}$），即「纯几何不存在性」这条路线能关闭的类数是 $0$。
\textbf{新机制}（已证）：我们找到一组作用在\emph{每一行}上的恒等式与不等式——
行能量恒等式 $a_i+r_i=g_i(1-g_i)$、Perron 行方程 $s_i=(\rho-g_i)x_i$、三角形恒等式 \eqref{eq:tri}
（引理~\ref{lem:rowid}），以及由它们导出的四条 Cauchy--Schwarz 不等式 $[f],[g],[h_1],[h_2]$（命题~\ref{prop:fourcs}）。
冠军处这四条\textbf{同时取等}（残差 $\le10^{-16}$），只有五个锥行的 $[f]$ 留有 $1.8\times10^{-3}$ 的余量。
由此得到一个四参数族 $(g,\pi,r,w)$ 上的松弛 REL，$\kappa_3(N)\le\mathrm{REL}(N)$ 无条件成立，
并且在\emph{双均匀层}（$g_i\equiv3/N$、$\pi_i\equiv1/N$）上 REL 可以精确求解：
\textbf{定理~\ref{thm:LN}}：$\rho\le L(N)=\frac{2+\kappa_N}{1+\kappa_N}$，$\kappa_N=\sqrt{\frac{3(N-1)}{N-3}}$。
这一个闭式同时重导了两条已知结果——$\lim_{N\to\infty}L(N)=\frac{1+\sqrt3}2$ 是 Mohar 的常数，
$L(6)=\frac{3+\sqrt5}4=\Phi$ 是 \STW\ 的常数——并且 $L(4)=\frac54$ \textbf{取等}，给出 $\kappa_3(4)=\frac54$ 的第二条证明。
命题~\ref{prop:fracdeg} 解释了它为什么不能更好：取等\emph{强制}钝角图正则、度为 $n^*=N\kappa_N/(3(1+\kappa_N))$，
并强制 $Q$ 两值，故 $n^*$ 必须是整数；$n^*\in\Z$ 当且仅当 $\frac{3(N-1)}{N-3}$ 是 $\Q$ 中的平方，
这只对一支\emph{无穷但极稀疏}的 Pell 例外集 $\mathcal E=\{4,9,28,99,364,1353,5044,18819,\dots\}$ 发生
（$n^*(6)=\frac{5-\sqrt5}2$ 无理）。所以 $\Phi-\rho_0$ 这个缺口的本质是\textbf{度的整性}，不是 Cauchy--Schwarz 太粗。
定理~\ref{thm:unattain} 进一步处理例外集：可达性会强制钝角图是\emph{强正则图}（或等团不交并），
而强正则图的重数整性对 $\mathcal E$ 中 $9\le N\le7.1\times10^8$ 的每一个成员都失败，
故 \textbf{$L(N)$ 对一切 $5\le N\le7.1\times10^8$ 不可达，$N=4$ 是唯一取等的阶数}。
（附带一条自我修正：命题~\ref{prop:fracdeg} 的第一版断言「$n^*$ 对一切 $N\ge5$ 无理」，
$n^*(9)=2$ 是反例，见注~\ref{rmk:fracdegfix}。）$N=6$ 的整数化阶梯（命题~\ref{prop:ladder6}）把双均匀层压到 $\frac{13}{10}$，
再用两值刚性（命题~\ref{prop:twovalue}，符号谱计算）排除其可达性。
\textbf{主定理}（已证）：$N=6$ 的\emph{全部} $Z_5$-等变秩三投影恰好构成一条一维代数弧（引理~\ref{lem:z5class}），
弧上的 Perron 特征值有闭式 $\rho(s)=\frac{6+\sqrt5-2s+\sqrt{P(s)}}{10}$，
它在弧内部有唯一临界点 $s^*=\frac{9+4\sqrt5-\sqrt{35}}{18}$，而 $\rho(s^*)=\frac{9+\sqrt5+2\sqrt{35}}{18}=\rho_0$
\textbf{精确相等}（定理~\ref{thm:conefam}）。于是 $\rho_0$ 的地位从「一个被验证过的下界见证」升级为
「一条显式代数弧的唯一极大值」。意外的副产品是：定理~\ref{thm:zonal} 的见证——正二十面体等角线系——
就在\emph{同一条弧}上，位于 $s=\frac12$，值 $\frac{5+\sqrt5+\sqrt{30}}{10}=1.271329355$，
即本报告两个最重要的六阶对象是同一条弧的中点与顶点。
另外两条一般性结果：$\Phi(Q)=\lambda_{\max}(Q_+)$ 在全体对称矩阵上是\textbf{凸}的，
故 $\operatorname{conv}(\mathcal P_3(N))=\mathcal S_N$ 且秩三约束在求上确界时无损（定理~\ref{thm:convex}）；
余秩一层的符号分类给出 $\kappa_2(3)=\frac87$ 的一条三行新证明（推论~\ref{cor:k23}）。
现状的量化阶梯是
$1.366025$（Mohar）$>1.309017$（$\Phi=$ REL 自由上确界）$>1.300000$（$\frac{13}{10}=$ 加整性）
$>1.284781$（加 $Z_5$ 对称）$>1.281568$（$\rho_0$）；
剩余的 $3.2\times10^{-3}$ 已定位到 REL 丢掉的秩三信息（\eqref{eq:tri} 与四点 Gram 行列式），见第~\ref{sec:diag} 小节。

第八，汇报第 13--14 轮的工作（第~\ref{sec:harm} 节）。这一轮换了两次语言，目的都是让阶数 $N$ 从问题里消失。
\textbf{调和恒等式}（定理~\ref{thm:harmonic}，已证）：$\rho H=3$、$H=\sum_ic_i/(G_+c)_i$，
于是主猜想等价于\emph{下界} $H\ge3/\rho_0=2.340882059435$（猜想~\ref{conj:Hform}）；
$\Psi\le\sqrt3$ 同时获得精确缺陷式 $3-\Psi^2=\sum_{i<j}c_ic_j(p_i-p_j)^2/(p_ip_j)$（定理~\ref{thm:psi}），
Mohar 常数 $\frac12(1+\sqrt3)$ 得到两个新证明，「永不取等」的最短理由是：取等会强制 $|Q|$ 秩一，
而秩一与 $\rank Q=3$ 矛盾（定理~\ref{thm:projmohar}）。
\textbf{$N$-自由重构}（定理~\ref{thm:nfree}，已证）：$\kappa_3$ 等于球面紧框架测度类
$\{\mu\ge0:\mu(S^2)=3,\ \int uu^\top\dd\mu=I_3\}$ 上核为 $(\langle u,u'\rangle)_+$ 的二次型上确界；
该集合凸紧，其极端点由 Carathéodory 至多 $6$ 个原子（命题~\ref{prop:extreme}）。
\textbf{二阶变分公式}（定理~\ref{thm:secondvar}，已证，$45$ 组数值核验到 $4\times10^{-8}$）：
沿矩核直线 $\rho''=2\rho\sum_{k\ge2}\frac{\rho_k}{\rho-\rho_k}\langle a,w_k\rangle^2$，
故 $Q_+\succeq0$ 即凸（推论~\ref{cor:psdconvex}）、内部极值必须满足 $\nu_-(Q_+)\ge N-6$（推论~\ref{cor:nuneg}），
并得到\emph{条件性}六原子归约定理（定理~\ref{thm:atomreduce}）：沿线段恒有 $Q_+\succeq0$ 时，
任意 $N\ge7$ 构型可化约到 $\le6$ 原子而不减小 $\rho$，主猜想遂\textbf{完全化归}为 $N=6$ 的有限问题
（这正是注~\ref{rmk:ladderinfo}(2) 一直缺少方法的那一步）。
数值上：$96/96$ 条矩核线段的最大值在端点、$18$ 条杀原子链全部单调，
冠军外加 $1\sim6$ 个原子后归约精确回到 $\rho_0$（差 $\le2\times10^{-15}$）；
$\kappa_3(N)=\rho_0$ 的数值验证由 $N=14$ 推进到 $N=20$。
\textbf{族 $F_n$}（定理~\ref{thm:Fn}，已证）：枢纽 $+$ 正 $n$ 边形族可逐 $n$ 精确求解，
$n=3,4,5,6$ 的峰值为 $\frac65$、$\frac{9+4\sqrt6}{15}$、$\rho_0$、$\frac{13+4\sqrt{11}}{21}$，
对应驻点 $\frac45$、$\frac{24-\sqrt6}{30}$、$\frac{45+4\sqrt5-\sqrt{35}}{90}$、$\frac{30-\sqrt{11}}{63}$；
$n\ge7$（数值）与 $n\to\infty$（极限 $1.248749$）均低于 $\rho_0$，
故\textbf{族内峰值唯一在 $n=5$、即 $N=6$}（推论~\ref{cor:Fnpeak}）。
$\rho_0$ 由此第一次有了不依赖数值搜索的解析来源，冠军的质量、夹角与全谱都有闭式（推论~\ref{cor:champclosed}），
而 $F_5$ 与第 12 轮的 $Z_5$-等变弧是同一条曲线（$s=5g-2$，注~\ref{rmk:F5arc}）。
诚实交代：无条件版归约仍是\textbf{猜想}；$\lambda_{\max}(|Q|)\le2\rho_0-1$ 被冠军\textbf{否证}
（$1.602534>1.563136$），挫折界差 $1.17\times10^{-2}$，
REL2 的 $1.314006274436$ 是伪值（$\rho H=3.078911\ne3$，被调和恒等式当场识破），
REL4 的数值求解不可靠（$N=4$ 只得 $1.249796<\frac54$）已放弃。

\medskip
\noindent\textbf{关键词：} 图特征值；Ky Fan 范数；投影常数；Gegenbauer 多项式；图极限；极值谱图论

\medskip
\noindent\textbf{阅读说明：} 全文对每一条结论标注其证明状态。
「已证」指本报告或所引文献中含有完整证明；「精确」指该结论的数值部分已由
\texttt{sympy} 在相应代数数域上判定为恒等式（差化简后恰为 $0$）；
「数值」指只有浮点证据、\emph{不}构成证明；「否证」指该命题已被反例推翻。
所有标注「数值」的条目都不应被当作成果引用。
\end{abstract}

\tableofcontents
\newpage

\section{问题、记号与已知结果}

\subsection{问题}

对 $n$ 阶图 $G$，记邻接矩阵特征值为 $\lambda_1(G)\ge\lambda_2(G)\ge\cdots\ge\lambda_n(G)$，并记
\[
\Lam_k(G)=\sum_{i=1}^k\lambda_i(G),\qquad
M_k(n)=\max\{\Lam_k(G):\ |V(G)|=n\},\qquad
c_k^{\Sigma}=\lim_{n\to\infty}\frac{M_k(n)}{n}.
\]
极限 $c_k^{\Sigma}$ 存在（$M_k(n)$ 关于 $n$ 几乎可加，且 $M_k(n)\le \frac{1+\sqrt k}{2}n$ 给出线性上界）。
$k=1$ 时 $c_1^{\Sigma}=1$（完全图取到）；$k=2$ 时 $c_2^{\Sigma}=8/7$，
这是 Ebrahimi--Mohar--Nikiforov--Ahmady 的猜想，
由 Kumar--Liu--Monterde--Pragada--Tait 用图极限与凸分析方法证明。
本报告讨论 $k=3$。

\begin{conjecture}[主猜想，记作 A1$'$]\label{conj:A1}
$c_3^{\Sigma}=\rho_0$，其中
\[
\rho_0=\frac{9+\sqrt5+2\sqrt{35}}{18}=1.2815681968721679\ldots
\]
等价地，对每个 $n$ 与每个 $n$ 阶图 $G$，$\Lam_3(G)\le\rho_0 n-3$。
\end{conjecture}

下界 $\ge$ 的一半已经严格证明（见定理~\ref{thm:lower}）；上界 $\le$ 的一半是本课题的全部内容，目前未证。

\subsection{三条已发表的上界}

\begin{enumerate}[leftmargin=1.6em,itemsep=2pt]
\item \textbf{Mohar (2009)}\ \ $\Lam_k(G)\le\frac{1+\sqrt k}{2}n$。$k=3$ 给
$\frac{1+\sqrt3}{2}n=1.366025403784439\,n$，且不带 $-k$ 修正项。
Mohar 在同一篇文章中明确提问：这个界离真值有多远。
\item \textbf{Das--Mojallal--Sun (2019)}\ \ 含负惯性指数 $\theta$ 的界
$\Lam_k(G)\le\frac{n}{2(\theta+1)}\bigl(\theta+\sqrt{\theta(k\theta+k-1)}\,\bigr)$。
$k=3$ 时，该界优于 $\Phi n$ 当且仅当 $\theta\le\theta^*=9+4\sqrt5=17.944$；
当 $\theta=\Theta(n)$ 时精确退化回 Mohar 的 $\frac{1+\sqrt3}{2}n$。
而我们构造的极值图（$K_1\vee C_5$ 的团吹起）恰好有 $\theta=\Theta(n)$，
所以这一条对稠密渐近情形完全无用。
\item \textbf{Sivashankar--Tang--Wakhare (2026)，即 \STW}\ \ Proposition 6.1：
对一切 $n$ 阶图与一切 $1\le k\le n$，
\[
\Lam_k(G)\ \le\ \frac{1+\gamma(k)}{2}\,n-k
\ \le\ \frac12\Bigl(1+\frac{k+\sqrt{k+2}}{1+\sqrt{k+2}}\Bigr)n-k ,
\]
其中 $\gamma(k)$ 是 $k$ 维绝对投影常数。$k=3$ 时 $\gamma(3)=\varphi=\frac{1+\sqrt5}{2}$，故
\[
\Lam_3(G)\ \le\ \Phi n-3,\qquad \Phi=\frac{1+\varphi}{2}=\frac{3+\sqrt5}{4}=1.309016994374947 .
\]
这是目前的无条件纪录。该文同时解决了 Nikiforov (2015) 的 Conjecture 4.2。
\end{enumerate}

于是
\[
\rho_0=1.281568196872168\ \le\ c_3^{\Sigma}\ \le\ \Phi=1.309016994374947,
\qquad \Phi-\rho_0=0.027448797502779 ,
\]
可改进空间的精确上限是 $\frac{\Phi-\rho_0}{\Phi}=2.0969\%$。

\subsection{一个值得注意的定量平行现象}

\STW\ 在 $k=2$ 处给出的界是 $\frac76 n-2=1.166667\,n-2$，而真值是 $\frac87 n=1.142857\,n$；
两者之比恰为 $\frac{49}{48}=1.020833$，即 \STW\ 的一般方法在 $k=2$ 处高出真值 $2.083\%$。
在 $k=3$ 处，$\Phi/\rho_0=1.021418$，即高出 $2.142\%$。
两个相对偏差几乎相同，这不是巧合的证据，但它与「$\rho_0$ 就是 $k=3$ 的真值」这一猜想高度相容：
\STW\ 的架构在 $k=2$ 处损失了 $2.08\%$，在 $k=3$ 处损失同量级是自然的。
\STW\ 自己在其 \S6.2 中承认，对 $k=2$「专门的工作给出了远为尖锐的答案」，
而他们的一般估计「只给出 $\frac76n-2$」。

\subsection{记号约定（必须钉死，否则会与 \STW\ 冲突）}

\STW\ 用 $c_k$ 表示\emph{单个}特征值的常数（$\lambda_k(G)\le c_k n-1$）。
本报告用 $c_k^{\Sigma}$ 表示\emph{前 $k$ 个之和}的常数。两者完全不同：
\STW\ 的 Table 1 列出的是 $c_5,\dots,c_{17}$（单个特征值），
其中并\emph{没有} $k=3$ 求和的任何下界。
\STW\ 还定义了混合投影常数
\[
\eta_{\mathrm{mix}}(k)=\sup_{N\ge k}\frac1N\max_{Q\in P_k(N)}\bigl(\|Q\|_1+\one^\top Q\one\bigr),
\]
并指出 $\Lam_k(G)\le\frac{\eta_{\mathrm{mix}}(k)}2n-k$ 而 $\eta_{\mathrm{mix}}(k)\le 1+\gamma(k)$。
\textbf{他们从未计算 $\eta_{\mathrm{mix}}(3)$ 的真值。}
由定理~\ref{thm:reduce} 可知 $c_3^{\Sigma}=\eta_{\mathrm{mix}}(3)/2$，
所以猜想~\ref{conj:A1} 恰好等价于 $\eta_{\mathrm{mix}}(3)=2\rho_0=2.563136393744$，
而 \STW\ 只给出 $\eta_{\mathrm{mix}}(3)\le 1+\varphi=2.618034$。
这是一个明确写在已发表文献中、且未被任何人占据的缺口。

其余记号：$P_k(n)=\{Q\in\R^{n\times n}:Q^\top=Q,\ Q^2=Q,\ \tr Q=k\}$ 为秩 $k$ 正交投影全体；
$\|Q\|_1=\sum_{ij}|Q_{ij}|$ 是\emph{逐元素}的 $\ell_1$ 范数（不是核范数；对投影而言核范数等于 $\tr Q=k$）；
$(x)_+=\max(x,0)$；$P_\ell$ 为 Legendre 多项式（$P_\ell(1)=1$）；
$s=\|Q\|_1/n$，$t=\one^\top Q\one/n$。

\section{\STW\ 证明的解剖}\label{sec:stwanat}

这一节的目的有两个：说清楚 \STW\ 到底做了什么（它并没有对任何具体的图跑算法），
以及精确定位它的证明在哪一步失去信息。后一件事决定了我们能不能改进它。

\subsection{第一步：Ky Fan 对偶与掩码（完全形式化，无损失）}

设 $A$ 为 $G$ 的邻接矩阵。由 Ky Fan 极大原理，
$\Lam_k(G)=\max_{Q\in P_k(n)}\tr(AQ)$。固定 $Q$，
因为 $a_{ii}=0$ 而 $0\le a_{ij}\le1$（$i\ne j$），有
\[
\tr(AQ)=\sum_{i\ne j}a_{ij}q_{ij}\ \le\ \sum_{i\ne j}(q_{ij})_+ ,
\]
且等号在 $A=1\{Q_{ij}>0\}$（$i\ne j$）时取到。两个 $\max$ 可以交换，于是
\begin{equation}\label{eq:mask}
M_k(n)=\max_{Q\in P_k(n)}\sum_{i\ne j}(Q)_{ij}^{+}.
\end{equation}
这一步的实质是：\emph{最优图由投影的符号模式决定}。它把「对一切图」换成了「对一切秩 $k$ 投影」，
$n$ 依然存在，但图已经消失了。这是 \STW\ 的 \S2，也是整个课题的出发点。

\subsection{第二步：拆成两项，分别封顶（信息损失发生在这里）}

用 $(x)_+=\frac12(x+|x|)$，并注意 $\sum_i q_{ii}=\tr Q=k$（投影的对角元非负），得
\begin{equation}\label{eq:st}
\sum_{i\ne j}(Q)_{ij}^{+}=\frac12\Bigl(\|Q\|_1+\one^\top Q\one\Bigr)-k .
\end{equation}
\STW\ 随后对两项\emph{分别}使用各自的上确界：
$\one^\top Q\one=\|Q\one\|_2^2\le\|\one\|_2^2=n$（因为 $Q\preceq I$），
以及 $\|Q\|_1\le\gamma(k)n$（绝对投影常数的定义）。代入即得 Proposition 6.1。

\begin{remark}\label{rmk:loss}
记 $s=\|Q\|_1/n$、$t=\one^\top Q\one/n$，则 \eqref{eq:st} 给出 $\Lam_k(G)=(s+t)n/2-k$。
$k=3$ 时 $s\le\gamma(3)=\varphi=1.618034$，$t\le1$，于是 $s+t\le 1+\varphi=2.618034$，
即 $\Lam_3\le\Phi n-3$。\textbf{但两个上确界不能同时取到：}
$s=\varphi$ 的唯一取到者是正二十面体的 $6$ 条等角线构成的 ETF，
而它的 $t=(3+\sqrt5)/6=0.872678\ne1$。
我们的极值构造（见定理~\ref{thm:champ}）给出 $(s,t)=(1.586541134689,\ 0.976595259055)$，
两项都严格低于各自的上确界，但和更大。精确的赤字分解是
\[
(\varphi-s)+(1-t)=0.031492854061+0.023404740945=0.054897595006=(1+\varphi)-2\rho_0 .
\]
\end{remark}

注~\ref{rmk:loss} 是本报告最重要的一句结构性观察：
\textbf{任何逐项放缩的路线都不可能改进 \STW}。要往下走，必须证明一条\emph{权衡}不等式
\begin{equation}\label{eq:tradeoff}
s+t\ \le\ 2\rho_0=2.563136393744336 ,
\end{equation}
而不是分别改进 $s$ 与 $t$。第~\ref{sec:nogo} 节将给出这条观察的一个远为严格的版本。

\subsection{第三步：$\|Q\|_1\le\gamma(3)n$ 是怎么证的（这才是 \STW\ 的数学内容）}

这一部分与图完全无关，是一条关于秩 $r$ 正交投影的、维数无关的解析不等式。
设 $Q=VV^\top$，$V\in\R^{n\times r}$ 列正交，行为 $x_i=c_iu_i$，其中 $c_i=\|x_i\|_2\ge0$，
$u_i\in S^{r-1}$。列正交性 $V^\top V=I_r$ 写成行的形式就是\emph{各向同性条件}
\begin{equation}\label{eq:iso}
\sum_{i=1}^n c_i^2u_iu_i^\top=I_r,\qquad q_{ij}=c_ic_j\langle u_i,u_j\rangle .
\end{equation}
于是 $\|Q\|_1=\sum_{ij}c_ic_j|\langle u_i,u_j\rangle|$。令
\[
f_2^{r}(t)=t^2-\frac1r,\qquad f_4^{r}(t)=t^4-\frac{3}{r(r+2)},\qquad
\phi(u)=uu^\top-\frac1rI_r,
\]
\[
C=\sum_ic_i,\qquad X^2=\Bigl\|\sum_ic_i\phi(u_i)\Bigr\|_F^2 .
\]

\STW\ 的论证由四块拼成。

\begin{enumerate}[leftmargin=1.6em,itemsep=3pt]
\item \textbf{两个正定核}（Lemma 3.2、3.3）。
$\langle\phi(u),\phi(v)\rangle_F=f_2^{r}(\langle u,v\rangle)$，所以 $f_2^r$ 是 Gram 矩阵、故在 $S^{r-1}$ 上正定；
$f_4^r$ 的正定性由附录 A 用一个四阶张量的显式 Gram 表示证明。
按 Schoenberg 定理，$S^{r-1}$ 上的正定函数恰是 $\sum_{\ell\ge0}a_\ell G_\ell^{r/2-1}$（$a_\ell\ge0$），
而 $f_2^r=\frac{r-1}{r}\widetilde G_2^{r/2-1}$，
$f_4^r=\frac{(r-1)(r+1)}{(r+2)(r+4)}\widetilde G_4^{r/2-1}+\frac{6}{r+4}f_2^r$。
\item \textbf{一条关键的二次型不等式}（Lemma 4.2）：
\begin{equation}\label{eq:l42}
C^2+rX^2\ \le\ rn .
\end{equation}
证明是把 $C^2+rX^2$ 展开、用 $2c_ic_j\le c_i^2+c_j^2$、再用各向同性 \eqref{eq:iso} 收尾，共五行。
这是整个论证中唯一真正用到「$\sum c_i^2u_iu_i^\top=I_r$」的地方。
\item \textbf{标量主函数}（Lemma 4.3）。若存在 $a\ge0$、$\rho\ge0$、$b\le ra$ 使
\begin{equation}\label{eq:maj}
|t|\ \le\ a+b\,f_2^{r}(t)-\rho\,f_4^{r}(t)\qquad\forall t\in[-1,1],
\end{equation}
则 $\|Q\|_1\le arn$。证明：把 \eqref{eq:maj} 代入 $t=\langle u_i,u_j\rangle$、乘以 $c_ic_j\ge0$、对 $i,j$ 求和；
常数项给 $C^2$，$f_2^r$ 项给 $X^2$，$f_4^r$ 项由正定性非负而系数为 $-\rho\le0$，故可丢弃；
再用 $b\le ra$ 与 \eqref{eq:l42}。
\item \textbf{系数的确定}（Lemma 4.4 与 Remark 4.5）。令 $s=\sqrt{r+2}$，取
\[
a_r=\frac{s^2+s-2}{r(s+1)},\qquad b_r=\frac{s(s^2+2s+3)}{2(s+1)^2},\qquad \rho_r=\frac{s^3}{2(s+1)^2}.
\]
则 \eqref{eq:maj} 成立，且余项有\emph{显式因式分解}
\begin{equation}\label{eq:fac}
a_r+b_rf_2^r(t)-\rho_rf_4^r(t)-|t|=\frac{(1-t)(st-1)^2(st+s+2)}{2(s+1)^2}\ \ge\ 0 .
\end{equation}
系数不是猜的：Remark 4.5 说明，极值等角线系的公角是 $\alpha=1/\sqrt{r+2}=1/s$，
故要求在 $t=1$ 与 $t=1/s$ 两点接触、且在 $t=1/s$ 处导数也为零，
即 $H(1)=H(1/s)=H'(1/s)=0$（$H$ 为 \eqref{eq:fac} 的左端），三个条件唯一确定三个系数。
\end{enumerate}

$r=3$ 时 $s=\sqrt5$，$3a_3=\frac{3+\sqrt5}{1+\sqrt5}=\varphi$，于是 $\|Q\|_1\le\varphi n$，即 $\gamma(3)\le\varphi$。
下界由正二十面体的 $6$ 条等角线给出：取 $12$ 个顶点中每对对映点的一个代表 $x_1,\dots,x_6$，
令 $X=[x_1/\sqrt2\ \cdots\ x_6/\sqrt2]\in\R^{3\times6}$，则 $XX^\top=I_3$，$Q=X^\top X$ 是秩 $3$ 投影，
$Q_{ii}=\frac12$，$|Q_{ij}|=\frac1{2\sqrt5}$，故
\[
\|Q\|_1=6\cdot\tfrac12+30\cdot\tfrac1{2\sqrt5}=3+3\sqrt5,\qquad
\gamma(3)\ge\frac{3+3\sqrt5}{6}=\frac{1+\sqrt5}{2}=\varphi .
\]
\begin{theorem}[$\gamma(3)=\varphi$，精确]\label{thm:gam3}
$\gamma(3)=\frac{1+\sqrt5}{2}$。上述六个恒等式（$XX^\top-I_3=0$、$Q^2-Q=0$、$\tr Q=3$、
$\|Q\|_1=3+3\sqrt5$、$\frac{3+3\sqrt5}{6}-\varphi=0$、$\frac{3+\sqrt5}{1+\sqrt5}-\varphi=0$）
已由 \texttt{sympy} 独立复核，化简后的差全部恰为 $0$。
\end{theorem}

\subsection{\STW\ 自己留下的两个缺口}

第一，\STW\ 的 Theorem 4.6 利用 \eqref{eq:fac} 的松弛给出了 $\gamma(r)$ 的严格改进，
但其假设是「$r\ge4$ 为偶数且 $r+2$ 不是完全平方」。\textbf{$r=3$ 是奇数，不在其中}；
而且 \STW\ 在第 515--517 行明确写出：Theorem 4.1 的界在 $r\in\{1,2,3,7,23\}$ 处是\emph{尖锐}的，
因为这些维数存在极值实等角紧框架。结合定理~\ref{thm:gam3}，
投影常数这条通道在 $k=3$ 处已经被榨干，天花板恰为 $\frac{1+\gamma(3)}2=\Phi$，不可能更低。

第二，\STW\ 定义了 $\eta_{\mathrm{mix}}(k)$ 并只给出 $\eta_{\mathrm{mix}}(k)\le1+\gamma(k)$，
在 \S6.1 中把它列为开放问题（他们的 Table 1 中若干 $r$ 处两侧下界不一致，
「suggesting that this inequality may be strict」，但明确声明这不构成分离的证明）。
$k=3$ 时 $\eta_{\mathrm{mix}}(3)$ 的真值无人计算。这正是猜想~\ref{conj:A1} 的位置。

\subsection{方法论上的一个澄清}

\STW\ 从头到尾没有对任何一阶的图做枚举或数值搜索。它的三步中，
第一步是恒等变形，第二步是一行放缩，第三步是一条关于球面上点系的解析不等式，
其证明依赖「正定核 + 标量主函数 + 一条二次型不等式」，
与 Cohn--Elkies--Viazovska 型的球面堆积线性规划界是同一类技术。
它的下界部分（\S5）确实用了确定性局部搜索来找候选核心，
但作者明确声明「we make no claim of global optimality」，
且每一个列出的构造都被直接验证。这个分工值得借鉴：
\textbf{上界必须是解析的，下界才允许是构造加数值搜索的}。

\section{已经严格建立的结果}

\subsection{归约链：把「对一切 $n$、一切图」变成一个有限维问题}

对有限图 $H$（$|V(H)|=N$）与权 $w\in\Delta^{N-1}$，记 $D_w=\mathrm{diag}(\sqrt w)$，定义\emph{核心}
\[
M(H,w)=D_w(A_H+I)D_w,\qquad
\psi_k(H)=\max_{w\in\Delta}\Lam_k\bigl(M(H,w)\bigr),\qquad
\psi_k=\sup_H\psi_k(H).
\]
$G=H[w;n]$ 表示团吹起：把 $H$ 的每个顶点换成 $n_v\approx w_vn$ 个顶点的团，每条边换成完全二部图。

\begin{theorem}[归约，已证]\label{thm:reduce}
以下三条成立。
\begin{enumerate}[label=(\roman*),leftmargin=2em]
\item 对每个 $n$ 与每个 $n$ 阶图 $G$，$\Lam_3(G)\le\psi_3\cdot n-3$；
\item $c_3^{\Sigma}=\psi_3=\eta_{\mathrm{mix}}(3)/2$；
\item 若 $(H,w)$ 满足 $\lambda_3(M(H,w))>0$ 且 $w$ 的分量全为有理数，则团吹起 $G_n=H[w;n]$ 满足
$\Lam_3(G_n)=\Lam_3(M(H,w))\cdot n-3$（对使 $n\cdot w$ 为整数的 $n$ 精确成立）。
\end{enumerate}
\end{theorem}

证明思路（(i)）。设 $Q\in P_3(n)$ 为 $G$ 的前三特征子空间投影，$Q=YY^\top$，$Y^\top Y=I_3$，行为 $y_i$。
令 $x_i=\sqrt n\,y_i$，$\mu=\frac1n\sum_i\delta_{x_i}$，则 $\int uu^\top\dd\mu=Q$。
把取值相同的 $x_i$ 归并为原子 $v_I$、权 $w_I$，并令 $P_{IJ}=1\{\langle v_I,v_J\rangle>0\}$。
$P$ 是对称 $0$--$1$ 矩阵且对角为 $1$，故 $P=I+A_H$，$H$ 为某个有限图。设 $U$ 的行为 $\sqrt{w_I}v_I$，则 $U^\top U=I_3$，且
\[
\frac1n\sum_{ij}(Q)_{ij}^{+}=\sum_{IJ}w_Iw_J\langle v_I,v_J\rangle^{+}
=\tr\bigl(U^\top D_wPD_wU\bigr)\le\Lam_3(D_wPD_w)\le\psi_3 ,
\]
倒数第二步用 Ky Fan（$U^\top U=I_3$）。$\square$

\begin{remark}
定理~\ref{thm:reduce} 的意义是：主定理 $\Lam_3(G)\le\rho_0n-3$ \textbf{完全等价于} $\psi_3\le\rho_0$，
其中 $n$ 与图都已消失。而且由 (iii)，一旦上界证出，$\rho_0$ 自动是\emph{最优}常数，
不需要另外证明尖锐性。这与 \STW\ 的第一步是同一件事，但多走了两步：
\STW\ 停在「对一切秩 $3$ 投影」，我们进一步停在「对一切有限核心 $(H,w)$」。
\end{remark}

\begin{theorem}[三维几何形式，已证]\label{thm:P3}
A1$'$ 等价于：对一切 $m$、一切 $p\in\Delta^{m-1}$、一切 $z_1,\dots,z_m\in\sqrt3\,S^2$ 满足
$\sum_ip_iz_iz_i^\top=I_3$，有
\begin{equation}\label{eq:P3}
\tag{P$_3$}\sum_{i,j}p_ip_j\bigl(\langle z_i,z_j\rangle\bigr)_+\ \le\ \rho_0 .
\end{equation}
\end{theorem}

\begin{remark}
（记号勘误）前期稿件中把 \eqref{eq:P3} 的约束写成 $z_i\in S^2$，这是错的：
取迹得 $\sum_ip_i|z_i|^2=1$，而 $\tr I_3=3$。正确的写法是 $z_i\in\sqrt3\,S^2$，
与测度形式 $C_3=\{\nu:\int uu^\top\dd\nu=I_3\}$（自动 $\int|u|^2\dd\nu=3$）一致。
本报告的定理~\ref{thm:zonal} 使用 $S^2$ 上的单位向量与 \STW\ 的 $c_i$ 分离写法，不受此影响。
\end{remark}

\begin{theorem}[自洽刻画与支撑界，已证]\label{thm:selfcons}
\begin{multline*}
c_3^{\Sigma}=\sup\Bigl\{\Lam_3\bigl(D_p^{1/2}AD_p^{1/2}\bigr):
N\in\N,\ A\in\{0,1\}^{N\times N}\ \text{对称},\\
A_{ii}=1,\ p\in\Delta^{N-1},\
A=1\bigl\{\Pi_3(D_p^{1/2}AD_p^{1/2})\ge0\bigr\}\Bigr\},
\end{multline*}
其中 $\Pi_3(M)$ 是 $M$ 的前三特征投影；且上确界可在 $N\le24$ 上取到，
Fenchel--Bunt 型精化给 $N\le20$。
\end{theorem}

\textbf{自洽条件 $A=1\{\Pi_3(M)\ge0\}$ 是本课题与 \STW\ 的分水岭。}
\STW\ 的框架里，$Q$ 可以是\emph{任意}秩 $3$ 投影；而在我们的问题里，
$Q$ 必须同时是掩码矩阵 $M$ 的前三特征投影。这是一条组合性质的约束，
\STW\ 的解析机器完全用不到它。第~\ref{sec:methods} 节将说明，这正是唯一的出路。

\subsection{精确下界与极值核心}

\begin{theorem}[下界，已证且精确]\label{thm:lower}
$c_3^{\Sigma}\ge\rho_0=\frac{9+\sqrt5+2\sqrt{35}}{18}$，且等号由核心 $H=K_1\vee C_5$（六原子轮图）
与权 $w^*=(a^*,b^*,b^*,b^*,b^*,b^*)$ 取到，其中
\[
a^*=\frac{91-25\sqrt7}{126}=0.197271565264962,\qquad
b^*=\frac{7+5\sqrt7}{126}=0.160545686946568,
\]
\[
\frac{a^*}{b^*}=5\sqrt7-12=1.228756555322 .
\]
对应的团吹起图列给出 $\Lam_3(G_n)=\rho_0n-3$（精确，非渐近）。
\end{theorem}

\begin{theorem}[冠军的完整刻画，已证且精确]\label{thm:champ}
对冠军核心，以下全部成立。
\begin{enumerate}[label=(\alph*),leftmargin=2em,itemsep=1pt]
\item $M^*=M(H,w^*)$ 的特征值为
$\rho_0,\ 0.76203144\ldots$ 等，其中 $\Lam_3(M^*)=\rho_0$，第四大特征值 $\mu_4=-0.083123<0$；
\item 一阶条件可化为一条初等二次方程 $a^2+24ab-31b^2=0$，其正根给出 $a^*/b^*=5\sqrt7-12$；
\item 负图 $\Gamma^-=C_5$ 是奇圈，且权在 $C_5$ 的五个顶点上等值（等权刚性）；
\item 在 $O(2)$ 归约下，核 $V_{IJ}=A_IA_J+R_IR_J\cos(\theta_I-\theta_J)$ 只取五个不同值：
\[
\begin{array}{lcl}
V(\text{hub},\text{hub})=\dfrac{14}{-7\sqrt5-7+5\sqrt7+3\sqrt{35}}=+1.681778636549,\\[4pt]
V(\text{hub},\text{cyc})=\sqrt{7/5}=+1.183215956620,\\[2pt]
V(\text{cyc},\text{cyc})\ \text{自身}=-\frac72-\frac{7\sqrt5}{10}+\frac{3\sqrt{35}}{10}+\frac{5\sqrt7}{2}=+3.323954628342,\\[2pt]
V(\text{cyc},\text{cyc})\ C_5\text{边}=(4-\sqrt7)\sqrt{7/5}=+1.602368657980,\\[2pt]
V(\text{cyc},\text{cyc})\ C_5\text{距}2=-\sqrt{7/5}=-1.183215956620 ;
\end{array}
\]
\item 刚性签名：$V(\text{hub},\text{cyc})=-V(\text{cyc},\text{cyc};\text{距}2)=\sqrt{7/5}$（等值反号），
$V(C_5\text{边})/V(\text{hub},\text{cyc})=4-\sqrt7$。
由此，在 $s=\sum w_Iw_J|V_{IJ}|$ 中「hub--cycle 对」（权重 $10ab$）与「距 2 的 cycle 对」（权重 $10b^2$）
单位贡献完全相同，故冠军处存在能级简并，\textbf{对偶乘子不唯一}，最优对偶解构成一整面；
\item 精确饱和：$\|M^*_{-}\|_1=0.083122814311+2\times0.099222691280=0.281568196871=\rho_0-1$（$10^{-12}$ 内），
$(s,t)=(1.586541134689,0.976595259055)$，$s+t=2\rho_0$（偏差 $0.0$）。
\end{enumerate}
\end{theorem}

\begin{remark}[$k=3$ 为什么比 $k=2$ 难，已证]
$k=2$ 的极值核心是 $P_3=K_1\vee2K_1$，三原子，负图 $\Gamma^-=K_2$ 是\emph{二部}的，
$e^\perp$ 是一维的（角度只有 $0$ 与 $\pi$，离散），常数域是 $\Q$，
判别式 $\Delta=a^2+6ab+b^2=1$ 是\emph{完全平方}，一阶条件是 $a=3b/2$（有理）。
$k=3$ 的极值核心是 $K_1\vee C_5$，六原子，负图 $\Gamma^-=C_5$ 是\emph{奇圈}，
$e^\perp$ 是二维的（角度 $\theta\in S^1$ 连续），对称群是 $O(2)$ 加冠军的 $D_5$，
常数域是 $\Q(\sqrt5,\sqrt7)$（四次），$\Delta=a^2+14ab+9b^2=35/49$，
$\sqrt\Delta=\sqrt{35}/7$ 无理，一阶条件是 $a/b=5\sqrt7-12$（无理）。

$k=2$ 的两套已发表证明（Huang--Wei 的补图关联秩一法、KLMPT 的图极限凸分析法）
都依赖「$\Gamma^-$ 二部 $\Rightarrow$ 符号模式退化为切换（switching）
$\Rightarrow$ 目标函数成为完全平方 $\|\sum\varepsilon_Iw_Iv_I\|^2$」这一结构。
$C_5$ 是奇圈，$\pm1$ 标注无法使所有边为负，切换结构失效。
此外 $\Delta=35/49$ 不是完全平方，意味着\textbf{任何「配方成完全平方」的初等证法在 $k=3$ 原则上不可用}。
这既是难度来源，也说明 $k=3$ 的不等式不是一个平方和。
\end{remark}

\subsection{Case A / Case B 二分}

\begin{theorem}[已证且精确]\label{thm:caseab}
记 $s=\|Q\|_1/n$、$t=\one^\top Q\one/n$、$\alpha_0=2\rho_0-\varphi=\frac{9-7\sqrt5+4\sqrt{35}}{18}=0.945102404994441$。
\begin{itemize}[leftmargin=1.2em,itemsep=1pt]
\item \textbf{Case A}（$t\le\alpha_0$）：$s+t\le\varphi+\alpha_0=2\rho_0$，
即\textbf{\STW\ 的 $\gamma(3)=\varphi$ 单独就足够}，不需要任何新想法。
\item \textbf{Case B}（$t>\alpha_0$）：全部难度集中在此，且 $t$ 被压缩到窄带 $t\in(0.9451,0.995]$。
\end{itemize}
\end{theorem}

\begin{remark}
这条二分的含义是：\textbf{让 \STW\ 饱和的那个几何对象（正二十面体 ETF）根本不是障碍}。
它的 $t=0.872678<\alpha_0$，落在 Case A 内，已被自动关闭。障碍只可能在 Case B，
而 Case B 的特征是 $t$ 接近 $1$，即 $\one$ 几乎落在前三特征子空间内。
\end{remark}

\subsection{负行挤压恒等式组：掩码刚性的定量发动机}\label{sec:squeeze}

方向庚（第~7.7 节）需要一个能把「负图 $\Gamma^-$ 的组合形状」翻译成「定量不等式」的工具。
本子节给出这样一组恒等式与不等式。它们对形式 9（定理~\ref{thm:form9}）的\emph{任何}可行标架都成立，
证明各只有一行，但在冠军处逐条精确饱和或给出显式正 slack，
因此是后续掩码刚性论证的定量发动机。

记 $G=A^\top A$（秩 $3$ 正交投影，$p_k=G_{kk}$，$\sum_kp_k=3$），
$y>0$ 为 $(G)_+$ 的 Perron 向量，$\rho=\lambda_{\max}((G)_+)$，$Y_2=\|y\|^2$。
对每个原子 $k$ 定义其\emph{负行三量}：
\[
N_k=-\!\!\sum_{l:\,G_{kl}<0}\!\!G_{kl}\,y_l\ \ge0,\qquad
E_k=\!\!\sum_{l:\,G_{kl}<0}\!\!G_{kl}^2,\qquad
S_k=\!\!\sum_{l:\,G_{kl}<0}\!\!y_l^2 ,
\]
即 $N_k,E_k,S_k$ 分别是负行的质量、能量与权平方和，三者都只在 $\Gamma^-$ 的边集上求和。

\begin{theorem}[负行挤压与钉扎，已证]\label{thm:squeeze}
对一切可行 $(G,y)$ 与每个 $k$：
\begin{enumerate}[label=(\roman*),leftmargin=2em,itemsep=1pt]
\item \textbf{正行恒等式}：$\sum_{l\ne k,\,G_{kl}>0}G_{kl}y_l=(\rho-p_k)y_k$；
\item \textbf{负行挤压}：$(\rho-p_k)^2y_k^2\le(p_k-p_k^2-E_k)\,(Y_2-y_k^2-S_k)$；
\item \textbf{负行 Cauchy--Schwarz}：$N_k^2\le E_kS_k$；
\item \textbf{全局钉扎}：$\rho\,Y_2-\sum_ky_kN_k=y^\top Gy$。
\end{enumerate}
\end{theorem}

\begin{proof}
(i) 把 $(G_+y)_k=\rho y_k$ 按对角、正、负三块拆开，负块移到右端即得（注意负块在两边各出现一次，相消）。
(ii) 对正行用 Cauchy--Schwarz：左端是正行的平方；第一个因子 $\sum_{l:+}G_{kl}^2=(G^2)_{kk}-p_k^2-E_k=p_k-p_k^2-E_k$
用到 $G^2=G$；第二个因子 $\sum_{l:+}y_l^2=Y_2-y_k^2-S_k$。
(iii) 对负行直接用 Cauchy--Schwarz。
(iv) $y^\top Gy=\sum_ky_k(Gy)_k$，而 $(Gy)_k=(G_+y)_k+(G_-y)_k=\rho y_k-N_k$。
\end{proof}

\begin{remark}[冠军处的精确饱和表，已证且精确]\label{rmk:squeezechamp}
把定理~\ref{thm:champ} 的闭式代入（$y$ 归一化为 $y_{\rm hub}=1$、
$y_{\rm cyc}=(\rho_0-p_{\rm hub})/(5h)$），四条全部精确成立，且 slack 有闭式：
\[
Y_2=\frac{91+25\sqrt7}{31}=5.069154283116605,
\]
\begin{align*}
y^\top Gy&=\frac{251\sqrt5+133\sqrt{35}}{558}+\frac{25\sqrt7+91}{62}\\
&=4.950512040312239,
\end{align*}
\[
\sum_ky_kN_k=\frac{95\sqrt5+37\sqrt{35}}{279}=1.5459548739683362
=\rho_0Y_2-y^\top Gy,
\]
\begin{align*}
\text{cycle 行 (ii) 的 slack}&=\frac{46-16\sqrt7}{405}=0.009056738328312,\\
\text{hub 行 (ii) 的 slack}&=0,\qquad \text{(iii) 在 cycle 行取等}.
\end{align*}
\end{remark}

remark~\ref{rmk:squeezechamp} 中有两个信号值得单独记录。
第一，\textbf{hub 行的挤压恰好饱和}：Cauchy--Schwarz 取等意味着 hub 的正行满足
$G_{{\rm hub},l}\propto y_l$，这是一条\emph{等式}型结构信息，可以直接用于刚性论证
（任何与冠军竞争的对象，其「hub 型」原子必须保持这个比例关系）。
第二，\textbf{(iii) 在 cycle 行取等}：意味着每个 cycle 原子的两条负边
满足 $|G_{kl}|/y_l$ 为常数，即负边在归一化意义下等权——这与等权刚性是同一现象的另一种表述。
这两条取等条件合起来，正是方向庚所需要的「掩码被钉死」的起点：
子步 1b 的任务是证明，任何 $\Gamma^-\ne C_5$ 吹起的构型都至少破坏其中一条取等，
从而由 (ii)(iii) 的显式 slack 把 $\rho$ 压到 $\rho_0$ 之下。

\subsection{冠军的局部刚性：近段稳定性常数}\label{sec:locrig}

定理~\ref{thm:squeeze} 是逐点的代数工具；本子节补一个\emph{局部}工具，
它是方向戊「近段」的第一块砖，也直接服务于方向庚的子步 1b。

在 Grassmann 图表 $X\mapsto\Pi(X)=A^\top(AA^\top)^{-1}A$（$A=[I_3\mid X]\in\R^{3\times6}$）下，
形式 9 的目标函数成为 $\rho(X)=\lambda_{\max}((\Pi(X))_+)$，定义域是 $\R^9$ 的开集；
冠军的符号腔 $\Sigma^*=\{\text{sign}\,\Pi=\text{sign}\,\Pi^*\}$ 是其中一个开室
（腔半径由 $\min|\Pi^*|=0.189960218562$ 控制）。
在室内 $(\Pi(X))_+=\Sigma^*\odot\Pi(X)$ 是 $X$ 的解析函数，且 $\lambda_{\max}$ 简单
（冠军处 $(\Sigma^*\odot\Pi^*)$ 的谱隙为 $0.588930287016$），故 $\rho$ 在室内解析。

\begin{theorem}[冠军的局部刚性，已证（计算机辅助认证）]\label{thm:locrig}
在冠军图表点 $X^*$ 处 $\nabla\rho(X^*)=0$（认证精度 $3.4\times10^{-51}$），且
\[
\operatorname{Hess}\rho(X^*)\ \preceq\ -\kappa I_9,\qquad
\kappa=0.020962946762952704 .
\]
因此冠军是其符号腔内的\textbf{非退化严格局部最大}；定量地，腔内
$\rho(X)\le\rho_0-\frac\kappa2\|X-X^*\|^2+O(\|X-X^*\|^3)$，
从而 $\rho(X)\ge\rho_0-\varepsilon$ 迫使 $\|X-X^*\|^2\le2\varepsilon/\kappa+O(\varepsilon^{3/2})$。
\end{theorem}

\begin{proof}[证明要点（认证管线）]
Hessian 用二阶摄动的基无关公式解析计算：
$\partial_{ts}^2\rho=v^\top M_{ts}v+u_t^\top M_sv+u_s^\top M_tv$，
其中 $M_\bullet=\Sigma^*\odot\partial_\bullet\Pi$，而 $u_\bullet$ 由
$(\lambda_1I-M_0)u=P_\perp M_\bullet v$ 解出（谱隙保证可逆，且避开重特征值子空间的基选择问题）。
$\partial_\bullet\Pi$ 与 $\partial^2_{\bullet\bullet}\Pi$ 由 $\Pi=A^\top(AA^\top)^{-1}A$ 的
二阶截断幂级数（含 $3\times3$ 级数求逆）精确给出，全程 50 位十进制。
认证：把 $9\times9$ Hessian 舍入到 20 位得 $\tilde H$，
$\|H-\tilde H\|_F\le9\times4.85\times10^{-21}=4.36\times10^{-20}$，
再由 Weyl 单调性 $\lambda_{\max}(H)\le\lambda_{\max}(\tilde H)+4.36\times10^{-20}
=-0.0209629467629527041<-\kappa$。
全部步骤见随附脚本 \texttt{证书\_冠军局部刚性.py}；
该解析 Hessian 与独立的高精度有限差（$h=10^{-3},10^{-4}$）在 8 位有效数字上吻合。
\end{proof}

\begin{remark}[这条定理在三条战线上的用法]
\emph{(戊)} 它就是近段稳定性：常数 $\kappa$ 显式，$C=\sqrt{2/\kappa}=9.7630\ldots$
给出「值损失 $\varepsilon$ $\Rightarrow$ 距离 $\le C\sqrt\varepsilon$」的近段版本
（三次余项的显式界是下一张证书，见协议 P2-近段 b）。
\emph{(庚)} 腔内没有竞争者：任何与冠军同腔且值 $\ge\rho_0-\varepsilon$ 的构型必在 $C\sqrt\varepsilon$ 邻域内，
于是子步 1b 只需处理\emph{跨腔}情形——即 $\Gamma^-$ 偏离 $C_5$ 吹起的构型必须穿过某面
$\{\Pi_{kl}=0\}$，而在面上的上界是有限个低维极值问题（每面一个），可逐面认证。
这把 1b 从「分析任意偏离」化为「有限个面上的上界」，是本轮最重要的结构简化。
\emph{(己)} 非退化严格局部最大排除了冠军附近的连续族竞争者，
交替映射的局部收敛速率也由 $\kappa$ 与谱隙控制。
\end{remark}

\subsection{三阶余项与完整的近段稳定性不等式（W3）}\label{sec:w3}

定理~\ref{thm:locrig} 给出腔内的二次主导项
$\rho(X)\le\rho_0-\frac\kappa2\|X-X^*\|^2+O(\|X-X^*\|^3)$，其中三次余项只有阶的估计。
本子节把余项做成\emph{显式常数}，得到一条完整的、带明确适用半径的近段稳定性不等式——
这是方向戊「近段」的收官，也是协议 P2-近段 b 的交付物。

\paragraph{方法}
在图表 $X\mapsto\rho(X)=\lambda_{\max}(\Sigma^*\odot\Pi(X))$ 上，$\rho$ 在冠军符号腔内解析，
其三阶导张量 $T_{ijk}=\partial^3\rho/\partial x_i\partial x_j\partial x_k\big|_{X^*}$
用\textbf{三阶截断幂级数}精确求出：把 $\Pi(X^*+tD)$ 的每个条目写成 $t$ 的三次多项式
（系数即 $0,1,2,3$ 阶导），其中 $3\times3$ 矩阵 $AA^\top$ 的逆用级数求逆递推
$(\sum_k t^kc_k)^{-1}=\sum_n t^nr_n$、$r_n=-c_0^{-1}\sum_{k\ge1}c_kr_{n-k}$，全程 $50$ 位十进制、无浮点舍入。
对简单特征值 $\lambda(t)=\lambda_{\max}(M_0+tM_1+t^2M_2+t^3M_3)$（$M_k$ 为掩码后的级数系数），
逐阶摄动给出 $\lambda_3=\langle u_2,M_1v_0\rangle+\langle u_1,M_2v_0\rangle+\langle v_0,M_3v_0\rangle$、
$\rho'''(0)=6\lambda_3$，其中 $u_1,u_2$ 由加边正则化的投影方程解出（见注~\ref{rmk:bordered}）。
完整对称张量再由\textbf{极化}恢复：对 $210$ 个一般方向 $c$ 计算
$\rho'''(0;c)=\sum_\tau m_\tau T_\tau\,c^\tau$（$m_\tau\in\{6,3,1\}$ 为重数），
解一个 $165\times165$ 的 Vandermonde 型线性方程组即得全部 $165$ 个 $T_\tau$。

\begin{remark}[一个必须记录的数值陷阱]\label{rmk:bordered}
投影算子 $L=\lambda_1I-M_0$ 在冠军处\emph{奇异}（$\det L=2.1\times10^{-51}$，零空间为 $\mathrm{span}(v_0)$）。
直接 \texttt{lu\_solve} 会因除以近零主元而返回被 $v_0$ 分量污染的解。
在\textbf{二阶}（Hessian）情形，冠军是临界点、一阶导 $\lambda_1=\langle v_0,M_1v_0\rangle=0$，
污染项恰乘以 $\lambda_1$ 而消失，故定理~\ref{thm:locrig} 的 $\kappa$ 不受影响
（本轮独立复核：Hessian 二次型在奇异解与正则化解之间逐位相同，且与有限差吻合到 $10$ 位）。
但到\textbf{三阶}，污染项进入 $u_2$ 的方程且\emph{不再}消失。正确做法是\textbf{加边正则化}
$L_{\rm reg}=L+v_0v_0^\top$：对任何 $r\perp v_0$，$L_{\rm reg}^{-1}r$ 恰为 $Lu=r$ 的 $\perp v_0$ 解。
本轮最初未加边，三阶张量与有限差量级偏差近一倍；加边后二者吻合到 $2.3\times10^{-10}$
（即五点中心差分 $O(h^2)$ 截断的极限）。
\end{remark}

\begin{theorem}[完整的近段稳定性不等式，已证（三阶张量精确，常数经有限差与边界采样认证）]\label{thm:w3}
在冠军符号腔 $\Sigma^*$ 内记 $d=\|X-X^*\|_F$。三阶导张量满足
\[
\max_{ijk}|T_{ijk}|=0.313572464368744174\ \ (\text{在对角元 }T_{555}\text{ 处，图表坐标 }(1,2)),\qquad
\|T\|_F=\Bigl(\sum_\tau m_\tau T_\tau^2\Bigr)^{1/2}=2.521337009637 .
\]
取三次余项常数 $C_3:=\|T(X^*)\|_F=2.521337009637$；由 Cauchy--Schwarz，
$C_3$ 是算子范数 $\|T\|_{\rm op}=\max_{\|c\|=1}\bigl|\sum_{abc}T_{abc}c_ac_bc_c\bigr|=1.932652351321$ 的严格上界。
则对腔内一切 $0<d<R:=3\kappa/C_3=0.024942655$，有
\[
\rho(X)\ \le\ \rho_0-\frac\kappa2\,d^2+\frac{C_3}6\,d^3\ <\ \rho_0,
\qquad \kappa=0.020962946762952704 ,
\]
即冠军在半径 $R=0.024942655$ 的腔内邻域中是唯一最大值点，值损失至少 $\frac\kappa2d^2-\frac{C_3}6d^3$。
\end{theorem}

\begin{proof}[证明要点]
腔内 $\rho$ 解析，Taylor 带 Lagrange 余项：
$\rho(X^*+D)=\rho_0+\frac12D^2\rho[X^*](\hat D,\hat D)\,d^2+\frac16D^3\rho[X^*+sD](\hat D,\hat D,\hat D)\,d^3$
（$\hat D=D/d$，某 $s\in(0,1)$；一阶项因 $\nabla\rho(X^*)=0$ 消失）。
由定理~\ref{thm:locrig}，$D^2\rho[X^*](\hat D,\hat D)\le-\kappa$。
由 $|D^3\rho[X](\hat D,\hat D,\hat D)|\le\|T(X)\|_{\rm op}\le\|T(X)\|_F$，
且在球 $B(X^*,R)$ 上 $\|T(\cdot)\|_F$ 连续、经边界采样验证其方向三阶导 $\le C_3$
（最坏方向在边界处为 $2.079<C_3$，随机边界方向 $\le0.502$），故余项 $\le\frac{C_3}6d^3$。
合并即得；当 $d<3\kappa/C_3$ 时 $\frac\kappa2d^2>\frac{C_3}6d^3$，右端 $<\rho_0$。
张量的精确性由两条独立检验背书：$210$ 方向线性方程组的残差 $1.6\times10^{-48}$
（即 $\rho'''(\cdot;c)$ 确为 $c$ 的齐三次型），以及五点中心差分对随机方向的三阶导复核（吻合到 $2.3\times10^{-10}$）。
全部步骤见随附脚本 \texttt{证书\_三阶稳定性.py}。
\end{proof}

\begin{remark}[定位与诚实交代]
$R=0.0249$ 远小于符号腔半径 $0.189960218562$，故整条不等式在腔内成立、无需担心掩码翻转。
其中\textbf{精确}的是三阶张量本身（$50$ 位、线性方程组残差 $10^{-48}$）；
\textbf{数值验证}的是 $C_3$ 在球上的\emph{一致性}（边界采样，余量约 $18\%$）——
把一致性也做成区间算术证书只需再补一个四阶导数的模，属直接但繁琐的工作，本报告不主张已完成。
本定理只覆盖\emph{近段}；真正的远段（跨腔、原子数直至 $20$）由第~\ref{sec:walls}--\ref{sec:assemble} 节的
腔面程序与 W2 处理，不在其射程内。
\end{remark}

\subsection{腔面与邻腔的数值定位（W1 第一步）}\label{sec:walls}

定理~\ref{thm:locrig} 把腔内的竞争者排除后，子步 1b 只剩「跨腔」情形：
任何符号模式异于冠军的构型，其模式与冠军模式之间至少隔着一面 $\{\Pi_{kl}=0\}$。
按 $D_5$ 对称性，邻接面只有三个对称类：hub--cyc 条目、cyc--cyc 距 1 条目、cyc--cyc 距 2 条目。
对每一类，定义光滑函数 $F_\Sigma(X)=\lambda_{\max}(\Sigma\odot\Pi(X))$；
在相应的面与邻腔上 $F_\Sigma$ \emph{恰好等于}真目标（因为被掩掉的条目在那里非正），
故 $\max_X F_\Sigma$ 是该面与其邻腔上真目标的一个上界。
掩码的取法必须小心：若冠军侧该条目为\emph{正}（前两类），则 $\Sigma$ 取冠军掩码去掉该条目；
若冠军侧为\emph{负}（第三类），则 $\Sigma$ 必须把该条目\emph{加回}正掩码，
否则 $F_\Sigma$ 在冠军侧退化为冠军目标本身、余量恒为 $0$（本轮曾因此得到假结果，已修正）。

在整张 Gr(3,6) 图表上以 11 个起点做 Nelder--Mead 加 Powell 优化，得到：

\begin{center}
{\small
\begin{tabular}{@{}llrr@{}}
\toprule
类 & 面 + 邻腔 & $\max F_\Sigma$ & 余量 $\rho_0-\max$ \\
\midrule
1 & hub--cyc 面 + 翻转腔 & $1.254642464407$ & $0.026926$ \\
2 & cyc--cyc 距 1 面 + 翻转腔 & $1.236504556084$ & $0.045064$ \\
3 & cyc--cyc 距 2 面 + 转正腔 & $1.256547360473$ & $0.025021$ \\
\bottomrule
\end{tabular}}
\end{center}

（第 1 类另做了带约束 $\Pi_{01}=0$ 的纯面优化，得 $1.252100833359$、余量 $0.029467$，与无约束版一致地低于 $\rho_0$。）
三个余量都与主缺口 $\Phi-\rho_0=0.027449$ 同阶，说明\textbf{腔面程序在数值上是闭合的}：
符号模式与冠军相差一个条目的构型都够不到 $\rho_0$。
\textbf{本小节全部为数值证据，不是证明。}
认证计划（下一步）：对每类的 argmax 解 KKT 方程组并以区间算术认证其临界性与 Hessian 负定（局部严格最大），
全局性则由定理~\ref{thm:squeeze} 的度序列不等式补足；
两个多步差的模式（相差两个条目）用同样的掩码上界函数再跑一轮，构成 W1 的完整清单。

\subsection{腔面的局部认证与双条目模式扫描（W1 的加固）}\label{sec:wallcert}

第~\ref{sec:walls} 节给出三类单条目腔面上界函数 $F_\Sigma$ 的数值最大值，均低于 $\rho_0$。
本子节做两件事：把每类 argmax 的\textbf{局部}最优性升级为区间认证（$50$ 位 + Weyl），
并把清单扩到相差\textbf{两个}条目的模式类。

\paragraph{单条目类的局部认证（已证）}
对第~\ref{sec:walls} 节的三类，在其数值 argmax $X_\Sigma$ 处重算 $F_\Sigma$ 的梯度与 $9\times9$ Hessian
（方法与定理~\ref{thm:locrig} 相同：三阶截断幂级数 + 加边正则化投影求解，见注~\ref{rmk:bordered}），
把 Hessian 舍入到 $20$ 位后用 Weyl 单调性认证 $\lambda_{\max}(\operatorname{Hess}F_\Sigma)<0$：

\begin{center}
{\small
\begin{tabular}{@{}llrrr@{}}
\toprule
类 & 面 & $\max F_\Sigma$ & $\lVert\nabla F_\Sigma\rVert_\infty$ & $\lambda_{\max}(\mathrm{Hess})\le$ \\
\midrule
1 & hub--cyc          & $1.254642464407$ & $2.4\times10^{-9}$ & $-0.0108321609014856$ \\
2 & cyc--cyc 距 1     & $1.236504556084$ & $5.0\times10^{-9}$ & $-0.0164028320295551$ \\
3 & cyc--cyc 距 2$^+$ & $1.256547360473$ & $3.5\times10^{-9}$ & $-0.00150760991144429$ \\
\bottomrule
\end{tabular}}
\end{center}

三类 argmax 均为 $F_\Sigma$ 的\textbf{非退化严格局部最大}（谱隙 $\ge0.479$，Hessian 负定且最大特征值有显式负上界）。
脚本 \texttt{证书\_腔面局部认证.py} 输出上表全部数字。
\textbf{但这只是局部认证}：每类 $F_\Sigma$ 的\emph{全局}最大值是否仍在该 argmax 处取得、从而全局 $<\rho_0$，
尚未证明（需约束 SOS/对偶证书，或归入 W2 的度序列引理）。这是 W1 留下的第一个缺口，见第~\ref{sec:assemble} 节。

\paragraph{双条目模式扫描（数值）}
把与冠军模式相差两个条目的情形按 $D_5$ 对称性归为六个代表模式，对每个模式用相应的掩码上界函数 $F_\Sigma$
（负侧条目须加回正掩码，规则同第~\ref{sec:walls} 节），以 $11$ 个起点做 Nelder--Mead 加 Powell 优化：

\begin{center}
{\small
\begin{tabular}{@{}llrr@{}}
\toprule
模式 & 相差的两个条目 & $\max F_\Sigma$ & 余量 $\rho_0-\max$ \\
\midrule
P1 & 两个 hub--cyc         & $1.247213595500$ & $+0.034355$ \\
P2 & hub--cyc $+$ 距 1     & $1.221062274052$ & $+0.060506$ \\
P3 & 距 1 $+$ 距 1（相邻） & $1.221062755745$ & $+0.060505$ \\
P4 & 距 1 $+$ 距 1（相离） & $1.216091741502$ & $+0.065476$ \\
P5 & 距 2$^+ $+$\,$ 距 2$^+$ & $1.253197264742$ & $+0.028371$ \\
P6 & hub--cyc $-$ 距 2$^+$ & $1.251351066898$ & $+0.030217$ \\
\bottomrule
\end{tabular}}
\end{center}

六个余量全部为正，且与主缺口 $\Phi-\rho_0=0.027449$ 同阶或更大。
\textbf{本段为数值证据}（脚本 \texttt{证书\_双条目模式扫描.py}）。
连同第~\ref{sec:walls} 节的单条目数值与上面的局部认证，近腔（相差一、两个条目）程序
\textbf{在数值上闭合、在单条目类上局部认证}；把它做成全局证明是 W1 的剩余工作，
与 W2 一起构成主线的两个缺口。

\subsection{W1-全局的结构归约与认证缺口}\label{sec:w1global}

第~\ref{sec:wallcert} 节把三类腔面的 argmax 认证为\emph{局部}最大；把它们升级为\emph{全局}上界
（记作 W1-全局）是主线的两个缺口之一。本子节先给一条把问题\textbf{降维}的结构定理，
再诚实交代认证为何仍未完成。

\begin{theorem}[腔面全局上界的球面归约，已证]\label{thm:wallreduce}
设 $\Sigma\in\{0,1\}^{6\times6}$ 为对称掩码、对角全为 $1$。则上界函数
$F_\Sigma(X)=\lambda_{\max}(\Sigma\odot\Pi(X))$ 在整张图表上的最大值（因图表是 $\mathrm{Gr}(3,6)$
的稠密开胞腔、$F_\Sigma$ 连续、$\mathrm{Gr}(3,6)$ 紧，故等于整个 Grassmannian 上的最大值）满足
\[
\max_{P\in\mathrm{Gr}(3,6)}\lambda_{\max}(\Sigma\odot P)\;=\;\max_{\|v\|=1}\Lam_3\bigl(\Sigma\odot vv^\top\bigr)
\;=\;\max_{v\in S^5}\Lam_3\bigl(D_v\Sigma D_v\bigr),\qquad D_v=\mathrm{diag}(v),
\]
其中 $\Lam_3$ 为三个最大特征值之和。
\end{theorem}

\begin{proof}
$\lambda_{\max}(\Sigma\odot P)=\max_{\|v\|=1}v^\top(\Sigma\odot P)v=\max_{\|v\|=1}\langle\Sigma\odot vv^\top,P\rangle$，
故 $\max_P\lambda_{\max}(\Sigma\odot P)=\max_{\|v\|=1}\max_{P\in\mathrm{Gr}}\langle\Sigma\odot vv^\top,P\rangle
=\max_{\|v\|=1}\Lam_3(\Sigma\odot vv^\top)$，末步用 Ky Fan 极大化 $\max_{P\in\mathrm{Gr}(3,6)}\langle C,P\rangle=\Lam_3(C)$；
再由 $(\Sigma\odot vv^\top)_{ij}=\Sigma_{ij}v_iv_j=(D_v\Sigma D_v)_{ij}$ 得第二式。
\end{proof}

于是 W1-全局化为一个\textbf{五维球面}上的显式特征值和问题。两个直接推论：
\emph{(i)}（迹恒等式）因 $\Sigma$ 对角全 $1$，$\mathrm{tr}(D_v\Sigma D_v)=\|v\|^2=1$，故
$\Lam_3(D_v\Sigma D_v)=1+\Lam_3(-D_v\Sigma D_v)$，目标值 $>1$ 当且仅当 $D_v\Sigma D_v$ 有负谱；
\emph{(ii)}（鲁棒 LMI 形式）由 Ky Fan 对偶，
$\max_{v\in S^5}\Lam_3(D_v\Sigma D_v)\le t\iff \Sigma\odot W\preceq tI\ \ \forall W\in\mathcal S$，
其中 $\mathcal S=\{W:0\preceq W\preceq I,\ \mathrm{tr}W=3\}$ 是秩三投影的凸包（一个谱面体）。
即 W1-全局等价于一条\textbf{鲁棒线性矩阵不等式}在 $\mathcal S$ 上成立，目标 $t=\rho_0$。

\begin{proposition}[临界点的自洽刻画，已证]\label{prop:w1kkt}
在 $g(v)=\Lam_3(D_v\Sigma D_v)$ 的光滑点（前三特征值互异处）有 $\nabla g(v)=2(\Sigma\odot P(v))v$，
其中 $P(v)$ 为 $D_v\Sigma D_v$ 的前三特征投影。故 $v^*\in S^5$ 是临界点当且仅当
$(\Sigma\odot P^*)v^*=g(v^*)\,v^*$ 且 $P^*$ 是 $D_{v^*}\Sigma D_{v^*}$ 的前三特征投影——一条自洽特征方程。
\end{proposition}

\begin{remark}[为什么「松弛成 SDP」没用（否定性引理）]\label{rmk:tight}
最自然的凸化是把秩一约束 $Z=vv^\top$ 松弛为 $Z\succeq0,\ \mathrm{tr}Z=1$，转而求 $\max\Lam_3(\Sigma\odot Z)$。
但 $Z\mapsto\Lam_3(\Sigma\odot Z)$ 是\emph{凸}函数（线性映射复合凸的 $\Lam_3$），
而 $\{Z\succeq0,\ \mathrm{tr}Z=1\}$ 的极点恰是秩一投影，故凸函数在紧凸集上的最大仍在极点达到：
$\max_{Z\succeq0,\,\mathrm{tr}Z=1}\Lam_3(\Sigma\odot Z)=\max_{\|v\|=1}\Lam_3(\Sigma\odot vv^\top)$。
松弛\textbf{零对偶间隙}，不可能靠它得到 $<\rho_0$ 的上界。同理，双线性形式
$\max\langle\Sigma\odot vv^\top,W\rangle$ 在 $S^5\times\mathcal S$ 之积上的极大也在双方极点
（秩一 $vv^\top$、秩三 $P$）达到，本质非凸。这就把 W1-全局的认证逼到\textbf{必须直面非凸性}：
要么穷举命题~\ref{prop:w1kkt} 的全部临界点并区间认证，要么给鲁棒 LMI（推论 ii）构造矩--SOS（Lasserre）证书。
二者都是明确、有限的计算，但都不是初等的。
\end{remark}

\paragraph{数值确认与否证扫描}
定理~\ref{thm:wallreduce} 的两条独立实现（$9$ 维图表的 Nelder--Mead/Powell 与 $5$ 维球面同法）
给出\emph{逐位相同}的最大值：class1 $1.254642464407$、class2 $1.236504556084$、class3 $1.256547360473$
（差 $\le2.2\times10^{-16}$），确认归约无误。球面上 $2000$ 次投影梯度上升（解析梯度 $2(\Sigma\odot P)v$）无一超过 $\rho_0$；
$4000$ 个非秩一混合 $Z$（$\mathrm{tr}Z=1$）的 $\Lam_3(\Sigma\odot Z)$ 亦全部 $\le$ 秩一最大值
（$1.171/1.158/1.174$ 对 $1.255/1.237/1.257$），印证注~\ref{rmk:tight}。三类 argmax 满足
命题~\ref{prop:w1kkt} 的自洽方程（残差 $\sim10^{-7}$），其 $D_{v^*}\Sigma D_{v^*}$ 的前三特征值为
$(0.695,0.283,0.277)$、$(0.799,0.266,0.172)$、$(0.786,0.262,0.209)$；二、三特征值近简并，
解释了全局最大盆地的狭窄（$2000$ 次随机上升仅 $1\sim18$ 次命中）。

\textbf{诚实交代}：本子节\emph{证明}了归约（定理~\ref{thm:wallreduce}、命题~\ref{prop:w1kkt}、注~\ref{rmk:tight}），
但\emph{没有}给出 W1-全局的完整认证——即尚未\emph{证明} $\max_{v\in S^5}\Lam_3(D_v\Sigma D_v)<\rho_0$，
目前只有强数值证据（两法逐位一致 + 否证扫描干净 + 局部认证）。缺口已从「任意图的全局最优」精确收缩为
「一个五维球面上显式特征值和的全局上界」，认证路径（临界点穷举 + 区间认证，或鲁棒 LMI 的矩--SOS 证书）明确但尚未执行。
脚本见 \texttt{证书\_腔面全局归约.py}。

\textbf{（2026-10-28 更新）}本段交代的认证缺口，对\emph{冠军掩码}已由定理~\ref{thm:champglobal}
用一条完全不同的凸化（优占对偶，第~\ref{sec:domdual} 节）\textbf{完全填平}：
$F(\Sigma_0)=\rho_0$ 已证且精确，证书是三个根式常数。对三类单条目腔面与两类双条目腔面，
定理~\ref{thm:wallsexact} 用精确有理数的半正定证书给出 $<\rho_0$ 的严格上界。
因此「临界点穷举 $+$ 区间认证」与「鲁棒 LMI 的矩--SOS 证书」这两条路，
对已覆盖的掩码从\emph{必须执行}降为\emph{可选的独立复核}；仍未覆盖的是
$\Omega>\rho_0$ 的那批掩码（如 class13 与八面体，见命题~\ref{prop:octnogo}）以及 $N\ge8$ 的情形。

\subsection{远腔反例扫描：对猜想 A1$'$ 的否证性体检}\label{sec:scan}

腔面程序（第~\ref{sec:walls} 节）只覆盖与冠军模式相差一两个条目的近腔；
远腔（模式差多个条目、原子数直至 $20$）目前\emph{没有任何引理}。
为防止主线建立在假猜想上，按研究方法协议把算法用作\textbf{否证工具}做了一次体检：
在 Stiefel 流形 $V_3(\R^K)$ 上以投影梯度上升（极小化 $-\lambda_{\max}((A^\top A)_+)$ 的相反数，
梯度 $2AM$ 投影到切空间、乘性回缩）最大化目标，
$K=6,\dots,12$，每档 $30$ 个随机起点外加结构化起点（$K_1\vee C_{K-1}$ 标架），共 $217$ 次独立上升。

\begin{center}
{\small
\begin{tabular}{@{}crcrl@{}}
\toprule
$K$ & 最优 $\rho$ & $\rho_0-\rho$ & 活跃原子 & 负度序列 \\
\midrule
6  & $1.279479170566$ & $+2.09\times10^{-3}$ & 6  & $(2,0,2,2,2,2)$ \\
7  & $1.281567173506$ & $+1.02\times10^{-6}$ & 7  & $(3,2,3,2,0,2,2)$ \\
8  & $1.281531188085$ & $+3.70\times10^{-5}$ & 8  & $(2,3,4,0,2,2,3,4)$ \\
9  & $1.281565442446$ & $+2.75\times10^{-6}$ & 9  & $(2,3,4,4,2,3,3,0,3)$ \\
10 & $1.281567315816$ & $+8.81\times10^{-7}$ & 10 & $(4,4,4,0,2,4,4,4,4,4)$ \\
11 & $1.281567360228$ & $+8.37\times10^{-7}$ & 11 & $(3,3,0,3,3,3,5,3,6,3,6)$ \\
12 & $1.281134352249$ & $+4.34\times10^{-4}$ & 12 & $(4,5,6,5,4,3,6,4,5,0,3,5)$ \\
\bottomrule
\end{tabular}}
\end{center}

未发现任何 $\rho>\rho_0$ 的点。差为 $10^{-6}$ 量级的各档，其构型均为「冠军加零权冗余原子」
（活跃原子中含负度 $0$ 者），与第~6.2 节对 $\mathrm{SR}_K$ 病态摄动的诊断一致；
真正新的局部最大（如 $K=8$、负度含 $4$ 者）比 $\rho_0$ 低至少 $3.7\times10^{-5}$。
\textbf{本小节是数值体检，不是证明}，其作用只有两个：
把「主线建立在假猜想上」的风险压低一档；
并给 W2 的度序列引理提示目标形状——含负度 $\ge3$ 活跃原子的构型，观测上界应比 $\rho_0$ 低一个显式量。

\subsection{逻辑装配：覆盖引理与条件性主定理（W4/W5）}\label{sec:assemble}

前面各节把上界问题拆成了有限块。本子节把它们的\textbf{逻辑关系}写清楚：
一条覆盖引理（W4）说明归约是\emph{全覆盖}的（没有图漏网），
一条条件性主定理（W5）说明「若两个具名缺口补上，则 $\rho_0$ 上界对一切 $n$ 成立」。
这样读者能一眼看出主线离终点还差哪两块，而不是差一团模糊的「还需努力」。

\begin{lemma}[覆盖，已证（由定理~\ref{thm:reduce}、\ref{thm:selfcons}、\ref{thm:form9} 组合）]\label{lem:cover}
对每个 $n$ 顶点图 $G$，存在一个\textbf{可行构型}——标架 $A\in\R^{3\times K}$（$AA^\top=I_3$）、
权 $p\in\Delta^{K-1}$ 且 $p_k>0$、原子数 $K\le20$——使得
$\Lam_3(G)\le n\cdot\rho(A)$，其中 $\rho(A)=\lambda_{\max}\bigl(\Sigma(A)\odot(A^\top A)\bigr)$、
$\Sigma(A)=1\{\Pi_3\ge0\}$ 为自洽掩码。换言之 $\sup_G\Lam_3(G)/n$ 被
$\sup\{\rho(A):A\ \text{可行、自洽、}K\le20\}$ 控制，而后者是 $\R^{3\times K}$ 上一个有限维、
模 $O(3)$ 紧的极大化问题。
\end{lemma}

\begin{proof}[证明要点]
定理~\ref{thm:reduce} 把 $\Lam_3(G)$ 归约到某核心 $(H,w)$ 的 $\psi_3$；
定理~\ref{thm:selfcons} 说明极值核心满足自洽掩码且支撑可压到 $N\le20$（$N\le24$ 经 Fenchel--Bunt 精化）；
定理~\ref{thm:form9} 把核心写成标架--Perron 形式 $\rho(A)$。三者串联给出可行构型；
$K\le20$ 来自支撑界，$p_k>0$ 来自「删去零权原子不改变值」。
\textbf{注意}：这里用的是无条件的 $N\le20$，\emph{不是}未证的支撑界 (2a)（$N\le8$）；
故引理成立，但它把极大化放在 $K\le20$（而非 $K\le6$）的大得多的空间上——这正是难点所在。
\end{proof}

\begin{theorem}[条件性主定理，已证（其条件尚未 discharged）]\label{thm:assemble}
设以下两条成立：
\textbf{(W1-全局)} 第~\ref{sec:walls}--\ref{sec:wallcert} 节每一类腔面/邻腔上界函数的\emph{全局}最大值 $<\rho_0$；
\textbf{(W2)} 在远腔（模式与冠军相差 $\ge3$ 个条目、原子数 $\le20$、含负度 $\ge3$ 的活跃原子）上，
$\rho(A)$ 有一个比 $\rho_0$ 低\emph{显式量}的上界。
则猜想 A1$'$（$c_3^{\Sigma}=\rho_0$，即猜想~\ref{conj:A1}）成立，从而对一切 $n$ 与一切 $n$ 顶点图 $G$，
\[
\Lam_3(G)\ \le\ \rho_0\,n-3,\qquad \rho_0=\tfrac1{18}\bigl(9+\sqrt5+2\sqrt{35}\bigr)=1.281568196872\ldots
\]
这会把现役上界 $\Phi n-3$（\STW）改进 $\Phi-\rho_0=0.0274488$（相对改进约 $2.1\%$），并达到下界 $\rho_0$ 的精确阶。
\end{theorem}

\begin{proof}[证明要点]
由引理~\ref{lem:cover}，$\sup_G\Lam_3(G)/n\le\sup\{\rho(A):A\ \text{可行、自洽、}K\le20\}$。
把可行域按符号腔分层：冠军腔由定理~\ref{thm:locrig} 与~\ref{thm:w3} 控制（腔内 $\rho\le\rho_0$，等号仅冠军）；
与冠军相差一、两个条目的近腔由 (W1-全局) 控制（各 $<\rho_0$）；
相差 $\ge3$ 个条目的远腔由 (W2) 控制（$\le\rho_0-\delta$）。三层取上确界得 $\rho_0$，且仅冠军腔达到，
故 $c_3^{\Sigma}=\rho_0$。$-3$ 的常数项来自归约定理~\ref{thm:reduce} 的边界修正。
\end{proof}

\begin{remark}[主线离终点还差哪两块]\label{rmk:gaps}
把整条链摊开后，缺口原本是\textbf{两个}：\emph{(W1-全局)}——每类腔面上界函数的全局最大值；
\emph{(W2)}——远腔的度序列型上界（只有第~\ref{sec:scan} 节的否证性数值扫描，没有引理）。
第~\ref{sec:domdual} 节把 W1-全局\textbf{实质推进}了：冠军腔已闭合（定理~\ref{thm:champglobal}，精确），
三类单条目腔面与两类双条目腔面已有 $<\rho_0$ 的精确有理数证书（定理~\ref{thm:wallsexact}），
$N\le7$ 的 $1252$ 个同构类中有 $832$ 个（$66.5\%$）被 $\Omega(A)\le\rho_0$ 无条件排除（第~\ref{sec:exh7} 小节）。
剩下的缺口重新表述为一条：
\textbf{(W1$'$)} 对\emph{自洽}掩码类 $\mathcal A=\{A=1\{\Pi_3(M)\ge0\},\ N\le20\}$ 证明
$\Omega(A)\le\rho_0$。
\emph{（2026-10-28 深夜更新）}原本列在这里的「退一步证明猜想~\ref{conj:A2}（负质量界）」
\textbf{已被否证}：命题~\ref{prop:A2no} 给出一个 $7$ 原子的显式自洽反例，
$\tr((-M)_+)=0.288182>\tau=0.281568$（纯有理算术证明）。
修复形式 A2$^\sharp$（尾谱和）与 A1$'$ 等价，不再是「退一步」。
因此 (W1$'$) 现在只剩 $\Omega(A)\le\rho_0$ 这一条凸形式，
而它与 (W1$''$)（$\kappa_3\le\rho_0$）相比是\emph{更弱}的表述：后者才是主方向。
命题~\ref{prop:octnogo} 说明这一步\emph{必须}用自洽性：在自洽类之外 $\Omega$ 至少能到 $4/3$。
第~\ref{sec:kappa} 节把 (W1$'$) 再压缩一档，得到与主定理\emph{等价}的单一不等式
\textbf{(W1$''$)}：$\kappa_3=\sup_{N}\sup_{V\in\mathrm{St}(N,3)}\rho((VV^\top)_+)\le\rho_0$
（推论~\ref{cor:kappamain}）。它的优点是彻底摆脱了掩码、图与阶数 $n$，
一阶条件有闭式 \eqref{eq:stiefelkkt}，$N=6$ 时符号模式只有 $156$ 个同构类；
代价是它必须对\emph{一切} $N$ 成立，而 $N\ge7$ 的最大值点是退化的（注~\ref{rmk:ladderinfo}(2)）。
(W2) 仍在，但第~\ref{sec:exh7} 小节的穷举把它的可疑范围压缩到「$\Omega>\rho_0$ 且自洽」的掩码上，
而 $N\le7$ 的扫描中这样的高值掩码一个也没有出现。
\end{remark}

\subsection{本轮突破：优占对偶与冠军腔的精确全局上界}\label{sec:domdual}

第~\ref{sec:w1global} 节把 W1-全局化为「五维球面上一族显式特征值和的全局上界」，
并在注~\ref{rmk:tight} 中说明\emph{把秩一约束 $Z=vv^\top$ 松弛成 $Z\succeq0,\ \tr Z=1$} 是零对偶间隙的、
不可能奏效的。本轮换了一个方向完全不同的凸化：不去动球面变量 $v$，
而是去\textbf{优占掩码本身}——找一个 $T\succeq\Sigma$、$T\succeq0$ 的矩阵，用它的对角线封顶。
这条路走通了，而且通得比预期彻底：冠军腔的全局上界被\emph{精确}求出（恰好等于 $\rho_0$），
证书只有三个根式常数；三类单条目腔面与两类双条目腔面也一并得到了严格上界。
同时它也给出了一条\textbf{定量}的否定性结果，说明这条路本身还不足以覆盖全部掩码。

\subsubsection*{三个上界函数与优占对偶}

对对称掩码 $\Sigma\in\{0,1\}^{N\times N}$（对角全为 $1$）定义三个量：
\[
F(\Sigma)=\max_{p\in\Delta^{N-1}}\Lam_3\bigl(D_{\sqrt p}\,\Sigma\,D_{\sqrt p}\bigr),\qquad
G(\Sigma)=\max_{p\in\Delta^{N-1}}\tr\Bigl(\bigl(D_{\sqrt p}\,\Sigma\,D_{\sqrt p}\bigr)_+\Bigr),
\]
\[
\Omega(\Sigma)=\min\Bigl\{\max_iT_{ii}\ :\ T\succeq\Sigma,\ T\succeq0\Bigr\},
\]
其中 $(\cdot)_+$ 是谱正部，$\Delta^{N-1}$ 是概率单纯形，$D_{\sqrt p}=\mathrm{diag}(\sqrt p)$。
$F$ 就是定理~\ref{thm:wallreduce} 里的上界函数
（$=\max_{P\in\mathrm{Gr}(3,N)}\lambda_{\max}(\Sigma\odot P)$）；$G,\Omega$ 是本轮新引入的。
$\Omega$ 是一个\textbf{半定规划}（变量 $T$ 与 $t$，约束 $T-\Sigma\succeq0$、$T\succeq0$、$\mathrm{diag}(T)\le t\one$）。

\begin{theorem}[优占对偶，已证]\label{thm:domdual}
对每个上述 $\Sigma$：
\begin{enumerate}[label=(\roman*),leftmargin=2em,itemsep=1pt]
\item $F(\Sigma)\le G(\Sigma)=\Omega(\Sigma)$；
\item $\Omega(\Sigma)=\max\bigl\{\tr(Y\Sigma):\ Y\succeq0,\ \exists\mu\in\Delta^{N-1},\ Y\preceq D_\mu\bigr\}$，
      强对偶成立且两侧都达到；
\item 因 $\Sigma_{ii}=1$，$\tr(D_{\sqrt p}\Sigma D_{\sqrt p})=\sum_ip_i=1$，故
      \[
      G(\Sigma)=1+\max_{p\in\Delta}\tr\Bigl(\bigl(-D_{\sqrt p}\Sigma D_{\sqrt p}\bigr)_+\Bigr)
               =\frac12\Bigl(1+\max_{p\in\Delta}\bigl\|D_{\sqrt p}\Sigma D_{\sqrt p}\bigr\|_*\Bigr),
      \]
      其中 $\|\cdot\|_*$ 是核范数（对称阵即特征值绝对值之和）；
\item $\Lam_3(M)=\tr(M_+)$ 当且仅当 $M$ 至多有 $3$ 个正特征值。故 $G(\Sigma)-F(\Sigma)$
      恰好度量「$G$ 的极大点处多出来的第四个及以后的正特征值质量」，
      两者相等当且仅当存在一个公共极大点 $p^*$ 使 $D_{\sqrt{p^*}}\Sigma D_{\sqrt{p^*}}$ 至多有 $3$ 个正特征值。
\end{enumerate}
\end{theorem}

\begin{proof}
(i) 设 $T$ 对 $\Omega$ 可行。合同变换保 Loewner 序，$D_{\sqrt p}TD_{\sqrt p}\succeq D_{\sqrt p}\Sigma D_{\sqrt p}$；
由 Ky Fan 极大原理 $\Lam_3$ 关于 Loewner 序单调，故
$\Lam_3(D_{\sqrt p}\Sigma D_{\sqrt p})\le\Lam_3(D_{\sqrt p}TD_{\sqrt p})$。
又 $D_{\sqrt p}TD_{\sqrt p}\succeq0$ 使其第 $4$ 个及以后的特征值非负，于是
$\Lam_3(D_{\sqrt p}TD_{\sqrt p})\le\tr(D_{\sqrt p}TD_{\sqrt p})=\sum_ip_iT_{ii}\le\max_iT_{ii}$。
对 $p$ 取上确界、对 $T$ 取下确界即得 $F\le\Omega$。而 $\Lam_3(M)\le\tr(M_+)$ 逐特征值显然，故 $F\le G$。

(ii) 对约束 $T-\Sigma\succeq0$、$T\succeq0$、$t\one-\mathrm{diag}(T)\ge0$ 分别取乘子
$Y\succeq0$、$Z\succeq0$、$\mu\ge0$，Lagrange 函数为
\[
L=t\,(1-\one^\top\mu)+\tr(Y\Sigma)+\tr\bigl((\mathrm{diag}(\mu)-Y-Z)\,T\bigr),
\]
它关于 $(t,T)$ 下有界当且仅当 $\one^\top\mu=1$ 且 $Y+Z=\mathrm{diag}(\mu)$，即 $0\preceq Y\preceq D_\mu$；
此时对偶目标为 $\tr(Y\Sigma)$。原问题在 $T=\rho I$（$\rho$ 充分大，$t=\rho+1$）处\emph{严格}可行，
故 Slater 条件成立、强对偶成立，且两侧最优值都达到。

(iii) 在对偶内层固定 $\mu$：令 $Y=D_{\sqrt\mu}Y_0D_{\sqrt\mu}$（$\mu_i>0$ 时可逆，边界情形由连续性），
则 $0\preceq Y\preceq D_\mu\iff0\preceq Y_0\preceq I$，且 $\tr(Y\Sigma)=\tr\bigl(Y_0\,D_{\sqrt\mu}\Sigma D_{\sqrt\mu}\bigr)$。
而 $\max_{0\preceq Y_0\preceq I}\tr(Y_0C)=\tr(C_+)$ 对任意对称 $C$ 成立（取 $Y_0$ 为 $C$ 的正谱投影）。
于是 $\Omega(\Sigma)=\max_{\mu\in\Delta}\tr\bigl((D_{\sqrt\mu}\Sigma D_{\sqrt\mu})_+\bigr)=G(\Sigma)$——
注意这是\emph{恒等式}，不是不等式：$\Omega$ 的对偶可行域恰好参数化了 $G$ 的可行域。
末式用 $\tr(M_+)=\tfrac12(\|M\|_*+\tr M)$ 与 $\tr(D_{\sqrt p}\Sigma D_{\sqrt p})=1$。

(iv) 设 $M$ 的特征值为 $\lambda_1\ge\cdots\ge\lambda_N$。$\tr(M_+)-\Lam_3(M)=\sum_{i\ge4}\lambda_i^+\ge0$，
等号当且仅当 $\lambda_4\le0$。$\square$
\end{proof}

\begin{remark}[与注~\ref{rmk:tight} 并不矛盾：两种凸化，一个无用一个有用]\label{rmk:tworelax}
注~\ref{rmk:tight} 否定的是\emph{松弛球面变量}：$Z\mapsto\Lam_3(\Sigma\odot Z)$ 是凸函数，
在谱单形 $\{Z\succeq0,\tr Z=1\}$ 上的最大仍在极点（秩一）达到，零间隙。
定理~\ref{thm:domdual} 动的是\emph{另外两个}地方：把 $\Lam_3$ 抬成 $\tr(\cdot)_+$，
并把掩码 $\Sigma$ 抬成它的优占矩阵 $T$。这两步都有间隙（前者由 (iv) 精确刻画，
后者由 $T$ 的对角线控制），但间隙\textbf{可以为零}——冠军处就是零（定理~\ref{thm:champglobal}）。
把两条注记合起来读，结论是：$\mathrm{Gr}(3,6)$ 上的非凸性没法靠松弛 $v$ 化解，
但可以靠\emph{提升}到原对偶变量对 $(T,Y)$ 上化解。这也解释了为什么
第~\ref{sec:w1global} 节的「临界点穷举 / 矩--SOS」路线不是唯一出路。
\end{remark}

\begin{remark}[$\Omega$ 的一个平凡上界与它的失效范围]\label{rmk:aplus}
取 $T=(\Sigma)_+$（谱正部）即得可行解，故 $\Omega(\Sigma)\le\max_i\bigl((\Sigma)_+\bigr)_{ii}\le1-\lambda_{\min}(\Sigma)$。
这个粗界在第~\ref{sec:exh7} 小节的穷举中排除了 $177/1252$ 个掩码；剩下的必须真解 SDP。
\end{remark}

\subsubsection*{冠军腔：精确求解}

\begin{theorem}[冠军腔的精确全局上界，已证且精确]\label{thm:champglobal}
设 $\Sigma_0=I+A(K_1\vee C_5)$ 为冠军掩码（$6\times6$：轮毂行全 $1$，圈上距 $1$ 为 $1$、距 $2$ 为 $0$），
$\varphi=(1+\sqrt5)/2$，$\rho_0=(9+\sqrt5+2\sqrt{35})/18$。令
\[
c=\frac{4\rho_0+1-\varphi}{5},\qquad
d=\rho_0-\frac1{\sqrt5},\qquad
e=\rho_0-\varphi+\frac1{\sqrt5},
\]
\[
T=\rho_0I_6+c\,B_{\mathrm{hc}}+d\,B_{\mathrm{d1}}+e\,B_{\mathrm{d2}},
\]
其中 $B_{\mathrm{hc}}$ 是轮毂—圈的关联指示矩阵，$B_{\mathrm{d1}},B_{\mathrm{d2}}$ 分别是 $C_5$ 上距 $1$、距 $2$ 的指示矩阵。则
\begin{enumerate}[label=(\alph*),leftmargin=2em,itemsep=1pt]
\item $T\succeq0$，$T-\Sigma_0\succeq0$，且 $\mathrm{diag}(T)=\rho_0\one$；
      数值为 $c=0.9016477597477553$，$d=0.8343546013722100$，$e=0.1107478036222310$；
\item $\mathrm{spec}(T)=\{6\rho_0-2\varphi,\ \varphi,\ \varphi,\ 0,0,0\}$，
      $\mathrm{spec}(T-\Sigma_0)=\{6\rho_0-2\varphi-4,\ \varphi^{-1},\ \varphi^{-1},\ 0,0,0\}$；
      两者秩都是 $3$，且 $\tr T=6\rho_0$；
\item 从而 $F(\Sigma_0)=G(\Sigma_0)=\Omega(\Sigma_0)=\rho_0$，等号在冠军权 $p=w^*$ 处取得。
\end{enumerate}
\end{theorem}

\begin{proof}
$T$ 与 $\Sigma_0$ 都关于「轮毂 $+$ 五圈」上的二面体群 $D_5$ 不变。
该置换表示分解为平凡表示（$2$ 维：$\mathrm{span}\{e_0,u_1\}$，$u_1=5^{-1/2}\sum_{j\ge1}e_j$）
与两个 $2$ 维不可约表示，故半正定性只需检查三个块：
\[
\text{平凡块}=\begin{pmatrix}\rho_0&\sqrt5\,c\\[1pt]\sqrt5\,c&\rho_0+S\end{pmatrix},\quad
\lambda_A=\rho_0+2d\cos72^\circ+2e\cos144^\circ,\quad
\lambda_B=\rho_0+2d\cos144^\circ+2e\cos72^\circ,
\]
其中 $S=2d+2e$（$\lambda_A,\lambda_B$ 各带重数 $2$，因轮毂不与这两个不可约子空间耦合）。
对 $N=T-\Sigma_0$ 同理，把 $(\rho_0,c,d,e)$ 换成 $(\tau,c-1,d-1,e)$，$\tau=\rho_0-1$。
代入 $\cos72^\circ=(\sqrt5-1)/4$、$\cos144^\circ=-(\sqrt5+1)/4$ 与 $c,d,e$ 的表达式，
用 \texttt{sympy} 精确化简得\textbf{八个恒等式}：
\[
\begin{array}{llll}
\rho_0(\rho_0+S)-5c^2=0, & \lambda_B=0, & \lambda_A=\varphi, & 2\rho_0+S=6\rho_0-2\varphi,\\[2pt]
\tau(\tau+S-2)-5(c-1)^2=0, & \lambda_A(N)=0, & \lambda_B(N)=\varphi^{-1}, & 2\tau+S-2=6\rho_0-2\varphi-4,
\end{array}
\]
外加两个正性 $\rho_0+S=3.1717730068\ldots>0$、$\tau+S-2=0.1717730068\ldots>0$。
于是两个平凡块迹正、行列式为零（半正定且各含一个零特征值），
$\lambda_A=\varphi>0$、$\lambda_B=0$、$\lambda_A(N)=0$、$\lambda_B(N)=\varphi^{-1}>0$，(a)(b) 得证。

(c)：由 (a) 与定理~\ref{thm:domdual}(i)，$F(\Sigma_0)\le\Omega(\Sigma_0)\le\max_iT_{ii}=\rho_0$。
反向，定理~\ref{thm:lower} 与~\ref{thm:champ} 已证 $\Lam_3(M^*)=\rho_0$，
而 $M^*=D_{\sqrt{w^*}}\Sigma_0D_{\sqrt{w^*}}$ 是 $F(\Sigma_0)$ 的一个可行点，故 $F(\Sigma_0)\ge\rho_0$。
夹逼得 $F=\Omega=\rho_0$；再由 $F\le G\le\Omega$ 得 $G=\rho_0$。$\square$
\end{proof}

\begin{remark}[证书为什么「刚好」：三重简并与黄金比]\label{rmk:whyexact}
定理~\ref{thm:champglobal} 的推导只解了三个方程（$T$ 平凡块行列式 $=0$、$\lambda_B(T)=0$、
$N$ 平凡块行列式 $=0$）就已唯一确定 $(c,d,e)$（$c$ 取二次方程的小根），
但\emph{第四个}条件 $\lambda_A(N)=0$ 自动成立，且 $\lambda_A(T)$ 自动等于 $\varphi$、
$\lambda_B(N)$ 自动等于 $\varphi^{-1}$。三个未知数满足六个约束，说明这个证书不是凑出来的，
而是被冠军的 $D_5$ 对称性与 $\rho_0$ 的代数性质\textbf{刚性}决定的。
黄金比 $\varphi=\gamma(3)$（三维绝对投影常数，也就是 \STW\ 上界的核心常数）
以二重特征值的身份出现在优占矩阵 $T$ 的谱里，$\varphi^{-1}$ 出现在 $T-\Sigma_0$ 的谱里；
这不是装饰：它定量地解释了为什么 $\rho_0$ 会紧贴在 $\Phi=(1+\varphi)/2$ 下方
（$\Phi/\rho_0=1.021418$，见第~1.3 小节）。
另有恒等式 $10\rho_0(\rho_0-1)=(\rho_0-1+\varphi)^2$，它把 $c$ 从 $\rho_0-\sqrt{2\rho_0(\rho_0-1)/5}$
化成了线性的 $(4\rho_0+1-\varphi)/5$。
\end{remark}

\begin{corollary}[W1-全局的冠军腔部分已闭合]\label{cor:chamberclosed}
由定理~\ref{thm:wallreduce} 与~\ref{thm:champglobal}，
\[
\max_{P\in\mathrm{Gr}(3,6)}\lambda_{\max}(\Sigma_0\odot P)=\max_{v\in S^5}\Lam_3(D_v\Sigma_0D_v)=\rho_0 .
\]
第~\ref{sec:w1global} 节遗留的「五维球面显式特征值和的全局上界」对冠军掩码\textbf{已经证明}；
该节末段所设想的临界点穷举或矩--SOS（Lasserre）证书，对冠军腔已不再需要。
\end{corollary}

\begin{remark}[顺带了结的两笔旧账]\label{rmk:concav}
(i) 前期数值上我们猜测 $f(p)=\Lam_3(D_{\sqrt p}\Sigma_0D_{\sqrt p})$ 在 $\Delta^5$ 上凹
（$3000$ 对随机中点检验全部满足，最小凹亏损 $1.4\times10^{-4}$；但同法在 class1 掩码上有 $6/3000$ 个
轻微反例，最小亏损 $-3.2\times10^{-3}$，故猜想本身可疑）。若凹成立，
$f(p)\le f(w^*)+\langle\nabla f(w^*),p-w^*\rangle=\rho_0$ 会给同一结论。
定理~\ref{thm:champglobal} 用凸对偶\textbf{绕过}了这个凹性猜想：它现在只是一条未被使用、
也不再需要的旁证，可以从队列里划掉。
(ii) 定理~\ref{thm:locrig}（局部刚性，$\kappa=0.020962946762952704$）与定理~\ref{thm:w3}
（三阶余项，$R=3\kappa/C_3=0.024942655$）只把冠军腔内的上界控制在\emph{邻域} $d<R$ 上；
推论~\ref{cor:chamberclosed} 把它升级为\emph{整个} $\Delta^5$。两者互补而非重复：
前者给等号成立的唯一性与稳定阶，后者给全域上界。
\end{remark}

\subsubsection*{腔面与近腔：精确有理数认证}

\begin{theorem}[腔面/近腔掩码的严格上界，已证（精确有理数认证）]\label{thm:wallsexact}
在 $\Sigma_0$ 上改一至两个条目得到下列掩码。每个掩码都存在一个\textbf{精确有理数}矩阵 $T$
（元素为分母 $2^{40}$ 的二进有理数，对角再加 $2^{-20}$ 的安全裕度）使
$T\succeq0$、$T-\Sigma\succeq0$、$\max_iT_{ii}\le\rho_0-\delta$，$\delta$ 见下表「余量」列：

\begin{center}
{\footnotesize
\begin{tabular}{@{}llrrrrc@{}}
\toprule
掩码 & 改动 & $\Omega(\Sigma)$ & 余量 $\rho_0-\Omega$ & $F(\Sigma)$（数值） & $\Omega-F$ & 自洽\\
\midrule
$\Sigma_0$（冠军） & —                 & $1.2815681969$ & $0$（\textbf{达到}）   & $1.2815681969$ & $0$        & 是\\
class1 hc          & $(0,1)\to0$       & $1.2740660391$ & $7.502\times10^{-3}$  & $1.2546424644$ & $1.94\times10^{-2}$ & 是\\
class2 d1          & $(1,2)\to0$       & $1.2365045561$ & $4.506\times10^{-2}$  & $1.2365045561$ & $\approx0$ & 否\\
class3 d2$+$       & $(1,3)\to1$       & $1.2565473605$ & $2.502\times10^{-2}$  & $1.2565473605$ & $\approx0$ & 是\\
class12            & $(0,1),(1,2)\to0$ & $1.2301542359$ & $5.141\times10^{-2}$  & $1.2210629132$ & $9.09\times10^{-3}$ & 否\\
\midrule
\multicolumn{7}{@{}l@{}}{\emph{以下两类 $\Omega>\rho_0$：优占对偶在此失效，右两列仅为数值}}\\
class13            & $(0,1)\to0,(1,3)\to1$ & $1.3066713982$ & $-2.510\times10^{-2}$ & $1.2500000000$ & $5.67\times10^{-2}$ & 否\\
八面体 $\Sigma_{\rm oct}$ & $J_6-$对径匹配  & $1.3333333333$ & $-5.177\times10^{-2}$ & $1.2531972647$ & $8.01\times10^{-2}$ & 否\\
\bottomrule
\end{tabular}}
\end{center}

\noindent 认证方式：$6\times6$ 有理对称矩阵半正定 $\iff$ 全部 $2^6-1=63$ 个主子式 $\ge0$（Sylvester 判据），
在 $\Q$ 上精确计算，\textbf{不含浮点、不含区间放宽}；实际最小主子式量级为 $10^{-25}$（正），
说明 $2^{-20}$ 的裕度足以吸收 $2^{-40}$ 的舍入。
\end{theorem}

\begin{remark}[class2 与 class3：数值变成定理]\label{rmk:class23}
第~\ref{sec:w1global} 节报告的两法逐位一致值 class2 $=1.236504556084$、class3 $=1.256547360473$
与本轮的 $\Omega$ 值\textbf{完全相同}（十位有效数字）。这不是巧合：这两个掩码在 $F$ 的极大点处
$\lambda_4\le0$（表中「$\lambda_4>0$?」列为否），故由定理~\ref{thm:domdual}(iv) $F=G=\Omega$，
而 $\Omega\le\rho_0-\delta$ 是已证的。于是那两个原本只是\emph{数值}的腔面全局最大值，
现在升级成了\emph{已证}的上界，并且是精确值。class1 则不同：它的 $\lambda_4=+0.01674>0$，
$\Omega-F=1.94\times10^{-2}$，松弛有真实间隙，但间隙被余量吸收了。
\end{remark}

\subsubsection*{这条路的边界：class13 与八面体}

\begin{proposition}[松弛类的定量天花板，已证且精确]\label{prop:octnogo}
\begin{enumerate}[label=(\alph*),leftmargin=2em,itemsep=2pt]
\item 令 $\Sigma_{13}$ 为 $\Sigma_0$ 上把 $(0,1)\to0$、$(1,3)\to1$ 得到的掩码。则
      $\Omega(\Sigma_{13})=1.3066713982>\rho_0$（超出 $2.510\times10^{-2}$，但仍 $<\Phi$），
      而其真值 $F(\Sigma_{13})=5/4=1.25$（数值，$\lambda_4=0$，$p$ 的第六个分量为 $0$）。
      优占对偶\textbf{不能}覆盖这个双条目近腔。
\item 令 $\Sigma_{\mathrm{oct}}=J_6-(\text{对径完美匹配})$，即正八面体六向量
      $\pm\sqrt3e_1,\pm\sqrt3e_2,\pm\sqrt3e_3$ 的 Gram 矩阵 $G$ 的符号掩码 $1\{G\ge0\}$。则
      $\Omega(\Sigma_{\mathrm{oct}})=G(\Sigma_{\mathrm{oct}})=\dfrac43$（精确），超出 $\rho_0$ 达
      $5.177\times10^{-2}$（相对 $4.04\%$）；其真值 $F=1.2531972647<\rho_0$，松弛间隙 $1/12$。
      这六个向量取 $p$ 均匀时满足定理~\ref{thm:P3} 的归一化 $\sum_ip_iz_iz_i^\top=I_3$，
      且它们在 \eqref{eq:P3} 中的目标值只有 $\sum_{ij}p_ip_j\langle z_i,z_j\rangle_+=3\sum_ip_i^2=0.5\ll\rho_0$，
      故\textbf{不违反 A1$'$}。
\item $\Sigma_{\mathrm{oct}}$ 不是自洽掩码：在 $p$ 均匀处 $D_{\sqrt p}\Sigma_{\mathrm{oct}}D_{\sqrt p}$ 的谱为
      $\{5/6,1/6,1/6,1/6,-1/6,-1/6\}$，$1/6$ 三重简并使 $\Pi_3$ 不唯一；
      任取一支都有 $1\{\Pi_3\ge0\}\ne\Sigma_{\mathrm{oct}}$（第三对对径原子的掩码条目对不上）。
\end{enumerate}
\textbf{结论}：任何只对「$\R^3$ 中标架可实现的符号掩码」成立、而不使用自洽性 $A=1\{\Pi_3(M)\ge0\}$ 的论证，
都不可能经由 $\Omega(A)\le\rho_0$ 证出 A1$'$。这为第~\ref{sec:obstacles} 节
「自洽性还没有被真正用起来」那条定性判断补上了一个\textbf{定量}反例：
松弛类的一致上界至少是 $4/3$，比 $\Phi=1.30902$ 还高。
\end{proposition}

\subsubsection*{负质量形式：一个新的、凸的、且尖锐的目标}

\begin{proposition}[加权 $\eta_{\mathrm{mix}}$ 与负质量恒等式，已证]\label{prop:negmass}
设 $(A,p)$ 自洽（$A=1\{\Pi_3(M)\ge0\}$，$A_{ii}=1$，$p\in\Delta$ 且 $p_i>0$），
$M=D_{\sqrt p}AD_{\sqrt p}$，$\Pi=\Pi_3(M)$，$Y=D_{\sqrt p}\Pi D_{\sqrt p}$。则
\begin{enumerate}[label=(\roman*),leftmargin=2em,itemsep=1pt]
\item $\Lam_3(M)=\tr(YA)=\dfrac12\bigl(\|Y\|_1+\one^\top Y\one\bigr)$，其中 $\|\cdot\|_1$ 是逐元素 $\ell_1$ 范数；
\item $0\preceq Y\preceq D_p$，$\sum_iY_{ii}/p_i=\tr\Pi=3$，且 $\one^\top Y\one=\langle\sqrt p,\Pi\sqrt p\rangle\le1$；
\item $\tr M=1$，故 $\Lam_3(M)\le\tr(M_+)=1+\tr\bigl((-M)_+\bigr)$；
\item 冠军处（全部为精确值）：$\|Y^*\|_1=1.586541134689$，$\one^\top Y^*\one=0.976595259055$，
      二者之和的一半 $=\rho_0$（残差 $4.4\times10^{-16}$）；
      $\tr\bigl((-M^*)_+\bigr)=\tau=\rho_0-1=0.281568196872$（残差 $2.2\times10^{-16}$，\textbf{精确饱和}）；
      $\mathrm{spec}(M^*)=\{0.7620314404,\ 0.2597683782^{\times2},\ -0.0831228143,\ -0.0992226913^{\times2}\}$，
      恰有 $3$ 个正特征值，故 $\tr((M^*)_+)=\Lam_3(M^*)=\rho_0$；
      对偶证书 $Y^*$ 与定理~\ref{thm:champglobal} 的 $T$ 满足互补松弛
      $\|Y^*(T-\Sigma_0)\|=7.2\times10^{-17}$、$\|(D_{w^*}-Y^*)T\|=3.3\times10^{-16}$。
\end{enumerate}
\end{proposition}

\begin{proof}
(i) $\Lam_3(M)=\tr(\Pi M)=\tr(\Pi D_{\sqrt p}AD_{\sqrt p})=\tr(D_{\sqrt p}\Pi D_{\sqrt p}A)=\tr(YA)$。
自洽性给 $\mathrm{sgn}(\Pi)=2A-J$，而 $p_i>0$ 使 $\mathrm{sgn}(Y)=\mathrm{sgn}(\Pi)$，故
$\tr(YA)=\sum_{ij}Y_{ij}\tfrac{1+(2A_{ij}-1)}2=\tfrac12\bigl(\sum_{ij}|Y_{ij}|+\sum_{ij}Y_{ij}\bigr)$。
(ii) $\Pi\preceq I$ 合同得 $Y\preceq D_p$；$Y_{ii}/p_i=\Pi_{ii}$，求和为 $\tr\Pi=3$；
$\one^\top Y\one=\langle\sqrt p,\Pi\sqrt p\rangle\le\lambda_{\max}(\Pi)\|\sqrt p\|^2=1$。
(iii) $\tr M=\sum_ip_iA_{ii}=1$，$\tr(M_+)=\tr M+\tr((-M)_+)$，$\Lam_3\le\tr(M_+)$。$\square$
\end{proof}

\begin{remark}[这正是 \STW\ 的 $\eta_{\mathrm{mix}}$，但带权且带自洽性]\label{rmk:etamix}
命题~\ref{prop:negmass}(i) 说明
\[
c_3^{\Sigma}=\frac12\sup\Bigl\{\|Y\|_1+\one^\top Y\one\ :\ Y=D_{\sqrt p}\Pi_3(M)D_{\sqrt p},\ (A,p)\ \text{自洽}\Bigr\},
\]
而 \STW\ 的 $\eta_{\mathrm{mix}}(3)=\sup_N\frac1N\max_{Q\in P_3(N)}(\|Q\|_1+\one^\top Q\one)$
是同一个表达式\textbf{去掉权 $p$、去掉自洽性}后的版本。
于是 A1$'$ 等价于「带权自洽版 $\eta_{\mathrm{mix}}(3)=2\rho_0=2.563136393744$」，
而 \STW\ 只给出无约束版 $\le1+\varphi=2.618034$。改进量 $2.097\%$ 全部来自这两条被丢掉的约束。
另外由 (i)(iii)，$\|Y\|_1\le\sqrt3$ 是 Cauchy--Schwarz 一步（$\|Y\|_1\le\|\sqrt{p_ip_j}\|_F\|\Pi\|_F=\sqrt3$）
给出的，它恰好等价于 Mohar 的 $\frac{1+\sqrt3}2$；这也解释了为什么「只用 (ii) 的约束、不用自洽性」
最好只能到 Mohar，而用逐点主函数的通道型方案（定理~\ref{thm:zonal}）最好只能到 $\Phi$。
\end{remark}

\begin{conjecture}[负质量界，记作 A2；蕴含 A1$'$，且被冠军精确饱和——\textbf{已于本轮否证，见命题~\ref{prop:A2no}}]\label{conj:A2}
对每个自洽核心 $(A,p)$，$M=D_{\sqrt p}AD_{\sqrt p}$ 满足
\[
\tr\bigl((-M)_+\bigr)\ \le\ \tau=\rho_0-1=0.281568196872\ldots
\]
等价地（把 $M$ 修成半正定所需的\textbf{最小}迹扰动）：存在 $N\succeq0$、$\tr N\le\tau$ 使 $M+N\succeq0$；
最小解就是 $N=(-M)_+$。
\end{conjecture}

\noindent\textbf{（2026-10-28 深夜：本猜想已被否证。）}
命题~\ref{prop:A2no} 给出一个 $7$ 原子的显式自洽核心，其 $\tr((-M)_+)=0.288182336568>\tau$，
超出 $6.6\times10^{-3}$；证明完全在 $\Q$ 上完成（有理特征多项式 $+$ 自实现 Sturm 序列 $+$ \texttt{isqrt}）。
下面保留原始陈述与它当时的动机，作为一条\emph{被定量否证的路线}存档；
失效机制见注~\ref{rmk:A2why}，穷举分层见注~\ref{rmk:A2strata}，修复形式见猜想~\ref{conj:A2sharp}。

\medskip
A2 比 A1$'$ 强（命题~\ref{prop:negmass}(iii)），且由 (iv) 是\textbf{尖锐}的。
它的好处是把目标从非凸的 $\Lam_3$ 换成\emph{凸}的 $M\mapsto\tr((-M)_+)=\max_{0\preceq Z\preceq I}(-\tr(ZM))$；
坏处是它丢掉了「只取前三」的信息——命题~\ref{prop:octnogo} 说明这个损失在自洽类\emph{之外}是致命的。
因此 A2 的正确读法不是「换一个好证的猜想」，而是：
\textbf{把全部难度集中到自洽性这一个约束上}，其余部分已经凸化完毕。

\subsubsection{A2 被否证：一个精确的七原子反例，以及它的失效机制}\label{sec:A2no}

\textbf{（2026-10-28 深夜，本轮最重要的负面结果。）}在用修正后的 \texttt{\_core.py}
重跑第~\ref{sec:exh7} 小节的 $1252$ 个掩码时，我们顺带记录了每个自洽掩码的
$\tr((-M)_+)$。结果是：$115$ 个自洽掩码中\textbf{有一个}的负质量严格超过 $\tau$。
这不是边界情形，超出量 $6.6\times10^{-3}$ 比数值误差大十二个数量级。
下面把它写成一个完全无浮点的定理。

\begin{proposition}[A2 为假：显式反例，已证（精确）]\label{prop:A2no}
令 $\widehat{C_5}$ 为「在 $C_5$ 的两个\emph{相邻}顶点上各挂一片叶子」得到的 $7$ 顶点连通单圈图
（$7$ 条边，度序列 $(3,3,2,2,2,1,1)$），令 $H=\overline{\widehat{C_5}}$
（$14$ 条边，度序列 $(3,3,4,4,4,5,5)$，$\omega(H)=4$、$\alpha(H)=2$、$|\mathrm{Aut}(H)|=2$，
graph6 码 \texttt{>>graph6<<FLp|w}，边集
$\{03,04,06,12,14,15,23,25,26,35,36,45,46,56\}$），
取 $A=I_7+A_H$ 与
\[
\begin{aligned}
p=\bigl(&0.128039679340616,\ 0.128039681290661,\ 0.142678696325684,\ 0.142678695287162,\\
&0.215687432177922,\ 0.121437907008806,\ 0.121437908569149\bigr),
\end{aligned}
\]
$M=D_{\sqrt p}AD_{\sqrt p}$。则
\begin{enumerate}[label=(\alph*),leftmargin=2em,itemsep=1pt]
\item $(A,p)$ \textbf{自洽}：$A=1\{\Pi_3(M)\ge0\}$，且符号余量极大——
      $A_{ij}=0$ 处 $|(\Pi_3(M))_{ij}|\ge4.94\times10^{-2}$，
      $A_{ij}=1$ 处 $(\Pi_3(M))_{ij}\ge1.06\times10^{-1}$；
      谱隙 $\lambda_3-\lambda_4=0.228268$，故 $\Pi_3(M)$ 良条件；
\item $\mathrm{spec}(M)$（降序）为
      \[
      \begin{aligned}
      &0.723222722079137,\quad 0.275747829160462,\quad 0.258740012261893,\quad 0.030471773066162,\\
      &0,\quad -0.130700331946254,\quad -0.157482004621400,
      \end{aligned}
      \]
      即 $M$ 有\textbf{四个}正特征值、一个零特征值、两个负特征值；
\item $\tr\bigl((-M)_+\bigr)=0.288182336567654\ldots>\tau=0.281568196872168\ldots$，
      超出 $6.614\times10^{-3}$。
\end{enumerate}
于是猜想~\ref{conj:A2} 为\textbf{假}。同时 $p$ 的每个分量都 $>0.12$，
且 $p$ 是 $F(A)$ 的极大化点（通用 KKT 残差 $1.5\times10^{-10}$、
$|\mathrm{diag}(\Pi_3M\Pi_3)-\Lam_3p|_\infty=5.6\times10^{-11}$，$400$ 起点复核），
$\mathrm{Aut}(H)$ 的对合 $(0\ 1)(2\ 3)(5\ 6)$ 使 $p$ 恰按 $\{0,1\},\{2,3\},\{5,6\},\{4\}$ 四组配对，
故这个反例同时落在命题~\ref{prop:negmass} 与猜想~\ref{conj:A2} 的\emph{全部}假设之内。
\end{proposition}

\begin{proof}[证明（精确，全程无浮点）]
(c) 是唯一的实质内容，(a)(b) 是它的适用性前提。

\textbf{第一步：把 $M$ 换成有理矩阵。}把 $p$ 舍入到 $2^{-60}$ 网格再归一化，得精确有理向量
$\hat p\in\Q^7$，$\sum_i\hat p_i=1$，$|\hat p-p|_\infty\le2.8\times10^{-17}$。
记 $\hat M=D_{\sqrt{\hat p}}AD_{\sqrt{\hat p}}$。因为
$D_{\sqrt{\hat p}}^{-1}\hat M D_{\sqrt{\hat p}}=A\,D_{\hat p}$，
二者\textbf{相似}，故特征多项式相同；而 $A\,D_{\hat p}$ 是\emph{有理}矩阵，
其首一特征多项式 $\chi(t)$ 可在 $\Q$ 上精确算出
（清分母后为整数系数，首项系数 $99$ 位十进制）。这一步是全部论证的关键：
它把一个含 $\sqrt{p_i}$ 的对称矩阵的谱问题，化成了一个\emph{有理系数}多项式的实根定位问题。

\textbf{第二步：Sturm 隔离。}对 $\chi$ 自实现 Sturm 序列（纯 \texttt{Fraction} 算术，
不用 \texttt{sympy} 内部 API、不用浮点）：序列长 $8$，说明 $\chi$ 无重根；
在 Cauchy 界 $R=2$ 的两端处符号变差数之差为 $7$，与「$7$ 阶实对称矩阵必有 $7$ 个实根」相符。
二分隔离得 $7$ 个互不相交的有理区间，再逐个精化到宽度 $<10^{-32}$，
其中负根 $2$ 个、零根 $1$ 个、正根 $4$ 个（这就同时给出了 (b)）。
区间中点与 \texttt{numpy} 的谱逐位相符（最大差 $3.3\times10^{-16}$）。于是
\[
\tr\bigl((-{\hat M})_+\bigr)\in
\bigl[\,0.28818233656765407380007104620745,\ \ \ldots+1.3\times10^{-32}\,\bigr]
\]
是一个\textbf{严格有理区间}。

\textbf{第三步：$\tau$ 的严格有理区间。}$\tau=\frac{-9+\sqrt5+2\sqrt{35}}{18}$。
用整数开方 \texttt{isqrt} 取 $\sqrt5,\sqrt{35}$ 到 $10^{-45}$，得
\[
\tau\in\bigl[\,0.2815681968721678862799251419346546754241,\ \ \ldots+1.7\times10^{-46}\,\bigr],
\]
包含性由 \texttt{sympy} 在代数数域 $\Q(\sqrt5,\sqrt{35})$ 上\emph{精确}判定
（\texttt{bool(tau\_lo < tau < tau\_hi) == True}），不是十进制串比较。

\textbf{第四步：比较。}上面两个有理区间\textbf{不相交}，间距 $6.614\times10^{-3}$。
两个有理数之间的大小比较是初等的，故 (c) 对 $\hat p$ 成立是一条\textbf{定理}，
不含任何浮点成分。

\textbf{第五步：(a) 的自洽性认证。}$80$ 位 \texttt{mpmath} 下重算 $\hat M$ 的秩三谱投影：
特征方程残差 $3.2\times10^{-81}$、正交残差 $5.6\times10^{-82}$、幂等残差 $5.3\times10^{-82}$，
而需要的符号余量是 $4.94\times10^{-2}$——比误差大\emph{八十个}数量级，
故 $\mathrm{sign}(\Pi_3(\hat M))=2A-J$ 逐条目成立是可靠的。

\textbf{第六步：鲁棒性与量词。}对 $p$ 施加 $2000$ 次 $10^{-9}$ 量级的相对随机扰动
（每次重新归一化），自洽性被破坏 $0$ 次，最小符号余量仍为 $4.94\times10^{-2}$、
最小超出量仍为 $6.614\times10^{-3}$。也就是说反例在一个\emph{显式开邻域}内一致成立。
猜想~\ref{conj:A2} 的量词沿用命题~\ref{prop:negmass} 的假设——「$(A,p)$ 自洽、$p\in\Delta$、$p_i>0$」，
\emph{并不}要求 $p$ 是驻点，因此显式有理向量 $\hat p$ 本身即构成反例。
即便坚持「核心 $=$ 驻点」的更强读法，由 $|\hat p-p^*|_\infty\le2.8\times10^{-17}$
与上述邻域鲁棒性，真驻点 $p^*$ 同样违反该界（余量 $6.6\times10^{-3}\gg10^{-15}$）。
\end{proof}

\begin{remark}[失效机制：负质量数的是\emph{全部}负特征值，$\Lam_3$ 只数前三个]\label{rmk:A2why}
反例中 $\Lam_3(M)=1.257710564<\rho_0$（余量 $2.386\times10^{-2}$），
所以\textbf{A1$'$ 完全没有被违反}——被否证的只是那条凸的\emph{充分}条件。
机制非常清楚：命题~\ref{prop:negmass}(iii) 给的是
$\Lam_3(M)\le\tr(M_+)=\tr M+\tr((-M)_+)=1+\tr((-M)_+)$，
这一步在 $M$ 恰有 $3$ 个正特征值时是\emph{等式}，在有 $4$ 个（或更多）正特征值时是\emph{严格}不等式。
反例正好是后者：$1+\tr((-M)_+)=1.288182>\rho_0$，而真正的 $\Lam_3$ 只有 $1.257711$；
被多算进去的正是第四个正特征值 $\lambda_4=0.030472$ 与 $\lambda_5=0$。
换句话说，\textbf{A2 的凸化恰好把「只取前三」这条信息全部丢掉了}，
而这正是它当初被寄予期望的地方。这与命题~\ref{prop:octnogo}（优占对偶在自洽类之外只能到 $4/3$）
是同一类教训的第二次出现：\emph{任何丢掉 $\Pi_3$ 的秩三结构的凸化都不够}。
\end{remark}

\begin{remark}[穷举给出的定量分层]\label{rmk:A2strata}
把 $115$ 个自洽掩码按 $\#\{\lambda_i(M)>0\}$ 分档（$N\le7$，\texttt{\_scan7F.json}）：
\begin{center}
{\small
\begin{tabular}{@{}lrrrl@{}}
\toprule
档 & 个数 & $\max\tr((-M)_+)$ & 与 $\tau$ 之差 & $\max F(A)$ \\
\midrule
$\mathrm{npos}=1$ & $3$ & $0$ & $-0.281568$ & $1.000000$ \\
$\mathrm{npos}=3$ & $103$ & $0.281568196872169$ & $+6.1\times10^{-16}$（冠军饱和） & $\rho_0$ \\
$\mathrm{npos}=4$ & $9$ & $0.288182336532$ & $\mathbf{+6.614\times10^{-3}}$ & $1.257711$ \\
\bottomrule
\end{tabular}}
\end{center}
两点值得记下。第一，$\mathrm{npos}\le3$ 档上恒等式 $\Lam_3(M)=1+\tr((-M)_+)$ 成立
（$106$ 个掩码上最大偏差 $1.6\times10^{-14}$），且 $\max\tr((-M)_+)$ \emph{恰好}是 $\tau$，
由冠军 $K_1\vee C_5$ 取到——即 A2 在这个占 $92\%$ 的子类上是\textbf{尖锐}的。
第二，违反 A2 的掩码\emph{全部}落在 $\mathrm{npos}=4$ 档，而这一档的 $\max F(A)=1.257711$
比 $\rho_0$ 低 $2.4\times10^{-2}$，是所有档里\emph{最}宽松的。
也就是说：A2 恰好在「它本来就等价于目标」的子类上成立，在「它比目标松」的子类上失效——
这条凸化没有提供任何真正的杠杆。
\end{remark}

\begin{conjecture}[尾谱和形式 A2$^\sharp$：A1$'$ 的\emph{等价}重写]\label{conj:A2sharp}
对每个自洽核心 $(A,p)$（$N$ 个原子），$M=D_{\sqrt p}AD_{\sqrt p}$ 满足
\[
\sum_{i=4}^{N}\lambda_i(M)\ \ge\ 1-\rho_0\ =\ -\tau =-0.281568196872\ldots
\]
（$N\le3$ 时左边按空和为 $0\ge-\tau$，平凡成立。）
\end{conjecture}

\begin{proposition}[A2$^\sharp$ 与 A1$'$ 等价，已证]\label{prop:A2sharp}
在 $\tr M=\sum_ip_iA_{ii}=1$ 的前提下，
$\Lam_3(M)=\tr M-\sum_{i\ge4}\lambda_i(M)=1-\sum_{i\ge4}\lambda_i(M)$，
故 $\Lam_3(M)\le\rho_0\iff\sum_{i\ge4}\lambda_i(M)\ge1-\rho_0=-\tau$。
此外 $\sum_{i\ge4}\lambda_i(M)=-\tr((-M)_+)+\bigl(\sum_{i\ge4}\lambda_i(M)\bigr)_+$，
所以 A2 $\Rightarrow$ A2$^\sharp$，而反向不成立（命题~\ref{prop:A2no} 的反例中
$-\tr((-M)_+)=-0.288182<-\tau$，但 $\sum_{i\ge4}\lambda_i=0.030472+0-0.130700-0.157482=-0.257710>-\tau$）。
\end{proposition}

\textbf{A2 到底是哪一步松弛？}这一点在反例之后才看得清楚，值得单独写明。
由 Ky Fan 最大原理的互补形式，
\[
\sum_{i\ge4}\lambda_i(M)=\min\bigl\{\tr(ZM):\ 0\preceq Z\preceq I,\ \tr Z=N-3\bigr\},
\]
取最小的一点是 $Z=I-\Pi_3(M)$；而\textbf{丢掉} $\tr Z=N-3$ 这条约束后，
\[
\min\{\tr(ZM):0\preceq Z\preceq I\}=-\tr\bigl((-M)_+\bigr),
\]
取最小的一点是 $Z=(-M)_+$ 的谱投影。两者相减正是命题~\ref{prop:negmass}(iii)：
$\Lam_3(M)=1-\sum_{i\ge4}\lambda_i\le1+\tr((-M)_+)$，
而等号 $\Lam_3(M)=1+\tr((-M)_+)$ 成立当且仅当
$\lambda_1+\lambda_2+\lambda_3=\sum_{\lambda_i>0}\lambda_i$（即 $\tr(M_+)=\Lam_3(M)$）：
当 $M$ 恰有 $3$ 个正特征值时自动成立；当有 $4$ 个或更多正特征值时\textbf{严格}；
当正特征值少于 $3$ 个时，还要求被 $\Lam_3$ 多算进去的 $\lambda_2,\lambda_3$ 恰为 $0$
（注~\ref{rmk:A2strata} 的 $\mathrm{npos}=1$ 档正是这种情形：$\tr((-M)_+)=0$、$\lambda_2=\lambda_3=0$，等号成立）。
所以 A2 \emph{就是}「在 $\Lam_3$ 的 Ky Fan 对偶中删去迹约束 $\tr Z=N-3$」，
而注~\ref{rmk:A2strata} 的分档表恰好把这条松弛的代价量化了：
它在占 $92\%$ 的 $\mathrm{npos}\le3$ 档上零代价（且尖锐饱和），在 $\mathrm{npos}=4$ 档上代价 $6.6\times10^{-3}$。

A2$^\sharp$ 的价值不在「更好证」——它\emph{等价}于 A1$'$（命题~\ref{prop:A2sharp}）。
它的价值在于\textbf{指认正确的对象}：需要控制的是「第四条及以后的谱尾和」
$\sum_{i\ge4}\lambda_i$（等价地，一族线性不等式 $\tr(ZM)\ge-\tau$，$0\preceq Z\preceq I$、$\tr Z=N-3$），
而不是「负质量」$\tr((-M)_+)$——后者是删去迹约束后的松弛，恰在 $\mathrm{npos}\ge4$ 时失效。
注意这里 $M\mapsto\sum_{i\ge4}\lambda_i$ 是\textbf{凹}的（因为 $\Lam_3$ 是凸的、$\tr M\equiv1$），
所以 A2$^\sharp$ 要求给出一个\emph{凹}泛函的下界，不再具有 A2 那种「凸泛函上界 $+$ 谱面表示」的形式；
这是保留秩三结构必须付出的代价。
由此得到本轮对主线的方法论结论：
\emph{优占对偶 $+$ 负质量这条凸化路线到此结束}——它两次（命题~\ref{prop:octnogo}、命题~\ref{prop:A2no}）
都在「丢掉 $\Pi_3$ 的秩三结构」这一步上失效。
主方向回到第~\ref{sec:kappa} 节的单一常数不等式 $\kappa_3\le\rho_0$（(W1$''$)），
因为 $\kappa_3=\sup_V\rho((VV^\top)_+)$ 的定义里\textbf{保留}了秩三投影本身，
正是 A2 与优占对偶先后丢掉的那部分信息。

\subsubsection{穷举筛选：$N\le7$ 的全部掩码}\label{sec:exh7}

$N\le7$ 的掩码恰是 $A=I+A_H$、$H$ 为不超过 $7$ 个顶点的简单图；按同构类共 $1252$ 个
（\texttt{networkx.graph\_atlas\_g}）。对每个类解一次 $\Omega$ 的 SDP：
\begin{itemize}[leftmargin=1.6em,itemsep=1pt]
\item $177$ 个被粗界 $\max_i((A)_+)_{ii}\le\rho_0$（注~\ref{rmk:aplus}）直接排除，无需 SDP；
\item $1075$ 个解了 SDP，其中 $655$ 个满足 $\Omega(A)\le\rho_0$，$420$ 个不满足（最大 $1.3733$）。
\end{itemize}
合计 $832/1252=66.5\%$ 的掩码\textbf{已被证明}满足 $F(A)\le\Omega(A)\le\rho_0$，
且这个证明与自洽性无关（对一切 $p\in\Delta^6$ 一致成立）。
剩下 $420$ 个用解析梯度（$\partial\Lam_3/\partial\log p_k=(\Pi_3M\Pi_3)_{kk}$）加 $25$ 个随机起点、
$L$-$BFGS$-$B$ 加 Nelder--Mead/Powell 收尾计算 $F(A)$。
\textbf{（2026-10-28 修正解析梯度后用 \texttt{\_core.py} 重跑，勘误见第~\ref{sec:kappa} 小节；
全部 $1252$ 类已跑完，耗时 $5603$ 秒，原始输出 \texttt{\_scan7F.json}。）}
\[
\max_{1252\ \text{类}}F(A)=1.281568196872=\rho_0,\qquad
\#\{F>\rho_0+10^{-9}\}=\mathbf 0,\qquad \#\{\text{自洽}\}=115 .
\]
残差：全部掩码上通用 KKT $\max\|(Q\odot A)x-\Lam_3x\|\le1.22\times10^{-8}$、
$\max|\mathrm{diag}(\Pi_3M\Pi_3)-\Lam_3p|\le6.09\times10^{-9}$；
限制到 $115$ 个自洽掩码上，Perron 残差降到 $4.76\times10^{-9}$、KKT 残差 $4.76\times10^{-9}$、
$g$-残差 $2.25\times10^{-9}$（表中「Perron 残差最大 $0.151$」是\emph{对全部掩码}取的，
其中绝大多数并不自洽，故该数无意义）。
自洽掩码按阶数分布为 $N=1,2,3,4,5,6,7\mapsto1,1,1,2,8,26,76$。

$F$ 的前若干个不同取值（降序）为
\[
\begin{aligned}
&\underbrace{1.281568196872}_{\rho_0},\quad 1.260028797704,\quad 1.257710563501,\quad
 \underbrace{1.256547360473}_{\text{腔面 class3}},\\
&\underbrace{1.254642464407}_{\text{腔面 class1}},\quad 1.254385521093,\quad 1.253202195454,\quad
 \underbrace{1.253197264742}_{\frac{9+4\sqrt6}{15}=\kappa_3(5)},\\
&1.252461952269,\quad 1.251351066898,\quad
 \underbrace{1.250785674353}_{K_1\vee C_6},\quad
 \underbrace{1.250000000000}_{5/4=\kappa_3(4)},\quad\ldots
\end{aligned}
\]
其中 $1.256547360473$ 与 $1.254642464407$ 正是第~\ref{sec:walls} 小节两类腔面的全局最大值，
$\frac{9+4\sqrt6}{15}$ 与 $\frac54$、$K_1\vee C_6$ 的值与第~\ref{sec:ladder} 小节的极值阶梯逐值吻合；
$1.260028797704$（$N=7$、自洽、$\mathrm{npos}=3$、负质量 $0.260029$）暂无闭式。

\textbf{一条新的结构观察。}$F=\rho_0$ 的掩码共 $17$ 个（$N=6$ 的冠军 $1$ 个 $+$ $N=7$ 的 $16$ 个，
其中自洽 $5$ 个），它们\textbf{全部}满足三条：含诱导 $C_5$、$\mathrm{npos}=3$、
$\tr((-M)_+)=\tau=0.281568196872$（最大偏差 $1.6\times10^{-15}$）。
$N=7$ 的那 $16$ 个中，$14$ 个的活跃权数为 $6$（最小权 $\le1.8\times10^{-17}$，即一个原子权已归零），
另 $2$ 个活跃权数为 $7$（最小权 $9.9\times10^{-2}$）；$5$ 个自洽掩码是冠军本身（$N=6$）加 $4$ 个 $N=7$ 的扩张。
这与定理~\ref{thm:selfcons} 的「负图必为 $C_5$ 吹起」方向一致，
也与注~\ref{rmk:ladderinfo}(2) 的退化性一致：$N\ge7$ 处取到 $\rho_0$ 的方式是\emph{加零权原子}，
而不是找到新的极值结构。这是优先队列 P2（$N\ge7$ 有效原子化约）的一条正面数值证据。

\textbf{同时，这次重跑也否证了猜想~\ref{conj:A2}}：$115$ 个自洽掩码中有一个
（graph6 码 \texttt{>>graph6<<FLp|w}，$N=7$，$F=1.257710563501$）的负质量为
$0.288182336532>\tau$。完整论证见第~\ref{sec:A2no} 小节。

\subsection{本轮突破：三阶问题的全局重构 $\kappa_3$、Perron 刚性与极值阶梯}\label{sec:kappa}

本小节是这一轮的主线成果。前面各节都在「掩码 $A$ 固定 $\to$ 解一个有限维极值」的框架里工作，
每个掩码一个 SDP、每个腔面一个球面问题；这样的路线注定要面对 $N\le7$ 就已经 $1252$ 个同构类、
$N\ge8$ 无法枚举的局面。本轮我们把整个问题\textbf{倒过来}：不再以掩码为自变量，
而是以\emph{秩三正交投影}为自变量，得到一个不含掩码、不含图、也不含阶数 $n$ 的单一常数 $\kappa_3$，
并证明主定理与「$\kappa_3\le\rho_0$」\textbf{完全等价}。这一步把课题从「无穷多个极值问题」
压缩成「一个不等式」，是本轮最重要的结构性进展。

\subsubsection{一个必须先交代的错误：softmax 梯度写错，以及它的影响范围}

在重跑第~\ref{sec:wallcert}、\ref{sec:scan}、\ref{sec:exh7} 小节的数值扫描时，
我们发现共用优化例程中 $\Lam_3(D_{\sqrt p}AD_{\sqrt p})$ 对 softmax 变量 $z$（$p=\mathrm{softmax}(z)$）
的解析梯度写错了。设 $M=D_{\sqrt p}AD_{\sqrt p}$、$\Pi_3$ 为 $M$ 的前三特征投影、
$g=\mathrm{diag}(\Pi_3M\Pi_3)$，则 $\partial\Lam_3/\partial z_k=g_k-\Lam_3p_k$；写成向量形式是
\begin{equation}\label{eq:gradfix}
\nabla_z\Lam_3 \;=\; g-\Bigl(\textstyle\sum_j g_j\Bigr)p \;=\; g-\Lam_3\,p ,
\end{equation}
因为 $\sum_jg_j=\tr(\Pi_3M\Pi_3)=\Lam_3(M)$。旧代码用的是 $g-\langle g,p\rangle\mathbf 1$。
两者只在 $p$ 为均匀分布时相等，一般情况下是\textbf{两个不同的向量}。
用中心差商逐分量核验：旧式最大偏差 $0.2\sim0.6$，新式 $4\times10^{-10}$（步长 $10^{-7}$）。
修正后的例程集中在 \texttt{\_core.py}（\texttt{top3}/\texttt{fg}/\texttt{Fmax}/\texttt{kkt\_diag}），
本报告随附的证书脚本一律改为调用它。

\textbf{影响范围（必须诚实交代）。}
\begin{itemize}[leftmargin=1.6em,itemsep=1pt]
\item \emph{所有已证结论不受影响}。定理~\ref{thm:reduce}、\ref{thm:P3}、\ref{thm:selfcons}、
\ref{thm:lower}、squeeze、locrig、w3、wallreduce、w1kkt、tight、
第~\ref{sec:domdual} 节的优占对偶与三个根式常数证书，全部由符号推导或 SDP/区间证书给出，
与数值梯度无关。
\item \emph{所有以该梯度驱动的数值扫描必须重跑}，即第~\ref{sec:walls}、\ref{sec:wallcert}、
\ref{sec:scan}、\ref{sec:exh7} 小节中「$\max F$」「余量」「命中次数」这类数字。
我们用修正后的 \texttt{\_core.py} 重跑了全部四项；\textbf{结论方向没有变化}：
$N\le7$ 的 $1252$ 个同构类中 $\max F$ 仍为 $\rho_0$、超过 $\rho_0$ 的仍为 $0$ 个。
\item \emph{一个教训}写进研究方法协议：任何以解析梯度驱动的优化，上线前必须过一次差商核验。
本项目此前恰好因为冠军的 $p$ 接近均匀（hub 权 $0.197$ 对 rim 权 $0.161$）而长期没暴露这个错误。
\end{itemize}

\subsubsection{全局重构：常数 $\kappa_3$}

记 $\mathcal P_3(N)=\{Q\in\mathbb R^{N\times N}:Q^\top=Q,\ Q^2=Q,\ \tr Q=3\}$
为秩三正交投影之集（Grassmannian $\mathrm{Gr}(3,N)$ 的算子实现），
并约定 $Q_+$、$Q_-$ \textbf{逐元素}取正部、负部：$(Q_+)_{ij}=\max\{Q_{ij},0\}$、
$(Q_-)_{ij}=\max\{-Q_{ij},0\}$，于是 $Q=Q_+-Q_-$、$|Q|=Q_++Q_-$。定义
\[
\kappa_3(N):=\sup_{Q\in\mathcal P_3(N)}\rho(Q_+),\qquad
\kappa_3:=\sup_{N\ge3}\kappa_3(N)=\lim_{N\to\infty}\kappa_3(N),
\]
其中 $\rho(\cdot)$ 为谱半径（$Q_+\ge0$ 对称，故 $\rho(Q_+)=\lambda_{\max}(Q_+)$）。
极限存在是因为 $\kappa_3(N)$ 单调不减：给 $Q\in\mathcal P_3(N)$ 补一行一列零得到
$\widetilde Q\in\mathcal P_3(N+1)$，而 $\widetilde Q_+$ 是 $Q_+$ 加一行一列零，谱不变。

\begin{theorem}[全局重构]\label{thm:kappa}
以下四条成立（(i) 的规范化链条见注~\ref{rmk:normcaution} 的勘误与修正写法）。
\begin{enumerate}[label=(\roman*),leftmargin=2em]
\item 对每个 $n$ 与每个 $n$ 阶图 $G$：$\ \Lam_3(G)+3\ \le\ \|Q_{G,+}\|_1\ $，其中 $Q_G=\Pi_3(A_G)$；
进而 $\Lam_3(G)\le\kappa_3\,n-3$（推导见注~\ref{rmk:normcaution}）。
\item $\kappa_3=\sup\{\Lam_3(D_x\Sigma D_x)\}$，上确界取遍
$N\in\N$、$\Sigma\in\{0,1\}^{N\times N}$ 对称且对角全 $1$、$x\in\Delta^{N-1}$、
并要求 $\Sigma=\mathbf 1\{\Pi_3(D_x\Sigma D_x)\ge0\}$（自洽）。
即上确界可以\textbf{只在自洽构型上取}。
\item $\kappa_3(N)$ 关于 $N$ 单调不减，且 $\rho_0\le\kappa_3\le\frac{1+\sqrt3}{2}$；
下界由冠军（$N=6$）取到等号。
\item \STW\ 的常数是本框架取 $x=\mathbf 1/\sqrt N$ 的\textbf{特例}：
$\frac1N\mathbf 1^\top Q_+\mathbf 1=\frac1N\|Q_+\|_1\le\frac1N\|Q\|_1\le\gamma(3)=\frac{1+\sqrt5}{2}$，
加上对角修正即得 $\Phi=\frac{3+\sqrt5}{4}$。
\end{enumerate}
\end{theorem}

\begin{proof}[证明要点]
(i) 第一式：$\Lam_3(G)=\langle A_G,\Pi_3(A_G)\rangle$（Ky Fan 取等），
而 $A_G$ 是 $0$--$1$ 且对角为零，故
$\Lam_3(G)=\sum_{ij}(A_G)_{ij}(Q_G)_{ij}\le\sum_{ij}(A_G)_{ij}(Q_G)^+_{ij}
\le\sum_{ij}(Q_G)^+_{ij}=\|Q_{G,+}\|_1$。
第二式的规范化见注~\ref{rmk:normcaution}。

(ii) 「$\ge$」：设 $Q\in\mathcal P_3(N)$，$x\ge0$ 为 $Q_+$ 的 Perron 向量、$\|x\|=1$，
$\rho=\rho(Q_+)$，取 $\Sigma=\mathbf 1\{Q\ge0\}$、$p=x^2\in\Delta$，令 $M'=\Sigma\odot xx^\top$。
因为 $\Sigma\odot Q=Q_+$（$\Sigma$ 恰是 $Q$ 的正部支撑），有
\[
M'x=(\Sigma\odot xx^\top)x=(\Sigma\odot Q)\,x=Q_+x=\rho x,
\]
故 $\rho$ 是 $M'$ 的特征值且 $x$ 是对应特征向量；进一步 $\Pi_3(M')=Q$
（在 $\mathrm{range}(Q)$ 与 $\ker Q$ 上逐块核对 $M'$ 的作用即可），于是该构型自洽、
目标值 $\Lam_3(M')=\rho$。故 $\kappa_3\le$ (ii) 的上确界。
「$\le$」：在自洽构型上 $\Sigma\odot Q=Q_+$，$Q=\Pi_3(\Sigma\odot xx^\top)$，故
$\Lam_3(D_x\Sigma D_x)=\langle\Sigma\odot xx^\top,Q\rangle
=\langle xx^\top,\Sigma\odot Q\rangle=x^\top Q_+x\le\rho(Q_+)\le\kappa_3$。

(iii) 单调性见上。下界：冠军 $Q^\ast=\Pi_3(\Sigma^\ast\odot x^\ast x^{\ast\top})$
满足 $\rho(Q^\ast_+)=x^{\ast\top}Q^\ast_+x^\ast=\rho_0$（定理~\ref{thm:lower} 与本节证书 E/F 段）。
上界 $\frac{1+\sqrt3}{2}$ 为 Mohar 型估计，见第~\ref{sec:methods} 节。

(iv) 逐式写出即得。\qedhere
\end{proof}

\begin{remark}[规范化勘误与正确写法，本轮发现]\label{rmk:normcaution}
定理~\ref{thm:reduce} 的原证明稿有一处\textbf{归一化不一致}，本轮复核时发现，必须记录。
稿中先令 $x_i=\sqrt n\,y_i$、$\mu=\frac1n\sum_i\delta_{x_i}$，于是
$\int uu^\top\dd\mu=\frac1n\sum_ix_i x_i^\top=YY^\top=Q$（这一步是对的）；
但接着写「设 $U$ 的行为 $\sqrt{w_I}v_I$，则 $U^\top U=I_3$」，
其中 $v_I$ 是原子（即 $x_i$ 的取值）、$w_I$ 是频率。由 $\frac1n\sum_ix_ix_i^\top=Q$ 得
$\sum_Iw_Iv_Iv_I^\top=Q$，这是 $N\times N$ 矩阵；而 $\sum_Iw_Iv_Iv_I^\top$ 按 $v_I\in\R^3$ 理解是 $3\times3$ 矩阵。
两个含义在稿中被混用，随后 $\frac1n\sum_{ij}(Q)_{ij}^{+}=\sum_{IJ}w_Iw_J\langle v_I,v_J\rangle^{+}$
这一步\textbf{不成立}（左边不带权、右边带权 $w_Iw_J$）。
冠军处数值核验：$\frac1n\sum_{ij}(Q)_{ij}^+=1.279704891254$，
$\sum_{IJ}w_Iw_J\langle v_I,v_J\rangle^+=0.214681093438$，$\|Q_+\|_1=7.678229347524$，
而 $\Lam_3(D_x\Sigma D_x)=x^\top Q_+x=1.281568196872$。

\textbf{正确写法}（本小节已按此陈述）：取 $v_I$ 为 $Y$ 的行（不乘 $\sqrt n$），则
$\sum_I\frac1nw_I\,v_Iv_I^\top=\frac1nI_3$ 与 $\langle v_I,v_J\rangle=Q_{IJ}$ 同时成立；
带权的量应当写成
\[
\Lam_3(D_x\Sigma D_x)=\langle\Sigma\odot xx^\top,\,Q\rangle,\qquad Q=\Pi_3(\Sigma\odot xx^\top),
\]
在自洽点等于 $x^\top Q_+x$。由此得到的正确归约链是
\[
\Lam_3(G)+3=\langle I+A_G,Q_G\rangle\le\|Q_{G,+}\|_1
\quad\text{且}\quad
\sup_G\frac{\Lam_3(G)+3}{n}=\kappa_3 ,
\]
其中第二个等号由 (ii) 与吹起构造（定理~\ref{thm:reduce}(iii)）给出。
\textbf{待办}：按本注重写定理~\ref{thm:reduce} 的证明；结论（$\Lam_3(G)\le\kappa_3n-3$）不变，
但链条必须逐式核对，不能再沿用稿中带权/不带权混用的写法。
\end{remark}

\begin{corollary}[主猜想的等价形式]\label{cor:kappamain}
主定理「对一切 $n$ 与一切 $n$ 阶图 $G$，$\Lam_3(G)\le\rho_0n-3$，且 $\rho_0$ 最优」
等价于单一不等式
\[
\boxed{\ \kappa_3=\rho_0\ }\qquad\text{即}\qquad
\sup\{\rho(Q_+):Q^\top=Q,\ Q^2=Q,\ \tr Q=3\}=\frac{9+\sqrt5+2\sqrt{35}}{18}.
\]
「$\ge$」已由冠军取等；全部难度在「$\le$」，即
\textbf{任意秩三正交投影的逐元素正部，其 Perron 根不超过 $\rho_0$}。
\end{corollary}

\begin{remark}[为什么这条重构是「数学的」而不是「算法的」]
$\kappa_3$ 的陈述里没有图、没有掩码、没有阶数，只有一个秩三投影的正部的 Perron 根。
它与投影常数、等角线、框架设计、两图（two-graphs）这些经典对象直接接壤：
$Q\in\mathcal P_3(N)$ 就是 $\R^3$ 中一个 $N$ 原子紧框架的 Gram 矩阵
（$\sum_Ia_Iv_Iv_I^\top=I_3$、$a_I=Q_{II}$、$\|v_I\|=1$），
而 $Q_+$ 是「只看夹角余弦为正的那些对」所得到的非负对称矩阵。
因此 $\kappa_3\le\rho_0$ 是一条\emph{框架/投影几何}的不等式，
原则上可以用框架能量、Schur--Horn 优占、两图分类这些结构性工具去攻，
而不需要对每一阶图跑算法。这正是研究方法协议要求的「数学优先」。
\end{remark}

\subsubsection{Perron 刚性：精确恒等式组}\label{sec:perronrig}

定理~\ref{thm:kappa}(ii) 说明上确界可以在自洽构型上取。设 $Q\in\mathcal P_3(N)$、
$x\ge0$、$\|x\|=1$、$Q_+x=\rho x$（$\rho=\rho(Q_+)>0$），令
\[
\alpha:=x^\top Qx=\|Qx\|^2,\qquad \beta:=x^\top Q_-x,\qquad d:=Q_-x,\qquad
\delta:=\|d\|^2,\qquad \ell:=x^\top|Q|x .
\]

\begin{proposition}[Perron 刚性系统，已证]\label{prop:perronrig}
在上述记号下，下列各条\textbf{无条件成立}（不依赖自洽性；自洽性只在把 $\rho$ 认作目标值 $\Lam_3$ 时才需要）：
\begin{enumerate}[label=(\alph*),leftmargin=2.2em]
\item $\rho=\alpha+\beta$；
\item $\alpha=\|Qx\|^2=x^\top Qx$（因 $Q$ 是投影）；
\item $\delta=\beta^2+\alpha(1-\alpha)$；
\item $\langle d,x\rangle=\beta$，$\ \langle d,Qx\rangle=(\rho-1)\alpha$，
$\ \langle d,(I-Q)x\rangle=\rho(1-\alpha)$；
\item $Qd=(\rho-1)\,Qx$（由 $Q^2=Q$ 与 $Q_+x=\rho x$ 两行推出）；
\item $\|(I-Q)x\|^2=1-\alpha$，且 $Qx\perp(I-Q)x$；
\item $\beta^2\le\delta$（Cauchy--Schwarz，因 $x\ge0$、$d\ge0$），从而 $\alpha(1-\alpha)\ge0$，
即 $\boxed{0\le\alpha\le1}$；
\item $\ell=\rho+\beta=\alpha+2\beta$（因 $|Q|=Q+2Q_-$ 且 $x\ge0$）。
\end{enumerate}
冠军处数值：$\alpha=0.976595259$、$\beta=0.304972938$、$\delta=0.115865452$、$\ell=1.586541135$、
$\rho=\rho_0=1.281568197$。
\end{proposition}

\begin{proof}
(a)(b)：$\rho=x^\top Q_+x=x^\top(Q+Q_-)x=x^\top Qx+\beta=\alpha+\beta$；
$\alpha=x^\top Qx=x^\top Q^2x=\|Qx\|^2$。

(c)：由 $d=Q_-x=(Q_+-Q)x=\rho x-Qx$ 直接展开范数，
\[
\delta=\|d\|^2=\|\rho x-Qx\|^2=\rho^2-2\rho\,\langle x,Qx\rangle+\|Qx\|^2
=\rho^2-2\rho\alpha+\alpha ,
\]
末一步用了 (b)（$\langle x,Qx\rangle=\|Qx\|^2=\alpha$）。另一方面，把 $\beta=\rho-\alpha$（即 (a)）代入
\[
\beta^2+\alpha(1-\alpha)=(\rho-\alpha)^2+\alpha-\alpha^2
=\rho^2-2\rho\alpha+\alpha ,
\]
两式逐项相同，故 (c) 成立。证书 A/B 段以 $10^{-8}$ 核验，实测残差 $5.9\times10^{-12}$；
这一残差完全来自 Perron 方程的数值残差 $\|Q_+x-\rho x\|_\infty=1.8\times10^{-11}$
（$Q$ 本身满足 $\|Q^2-Q\|_{\max}=4.4\times10^{-16}$，即是精确投影）。

(d)：$\langle d,x\rangle=x^\top Q_-x=\beta$。
$Qd=Q(Q_-x)=Q(Q_+x-Qx)=Q(\rho x)-Qx=(\rho-1)Qx$，
与 $x$ 作内积得 $\langle d,Qx\rangle=\langle Qd,x\rangle=(\rho-1)\|Qx\|^2=(\rho-1)\alpha$
（$Q$ 对称、$Qx$ 在 $\mathrm{range}Q$ 内）。
第三式 $=\langle d,x\rangle-\langle d,Qx\rangle=\beta-(\rho-1)\alpha$；
把 $\beta=\rho-\alpha$ 代入即得 $\rho-\alpha-(\rho-1)\alpha=\rho(1-\alpha)$。
于是 (d) 的三条内积恒等式其实只有\textbf{两条}独立信息：$\langle d,x\rangle=\beta$
与 $Qd=(\rho-1)Qx$；(c) 是它们的范数版本。
\emph{勘误}：本条早期草稿曾把 $\delta-(\rho-1)\alpha$ 也写成 $\langle d,(I-Q)x\rangle$ 的等价形式，
这是错的——由 (c) 有 $\delta-(\rho-1)\alpha=\rho^2-3\rho\alpha+2\alpha$，
在冠军处取值 $-0.159113$，与 $\rho(1-\alpha)=+0.029995$ 毫无关系；
证书脚本已同步删除该子句，只核验 $\beta-(\rho-1)\alpha$ 与 $\rho(1-\alpha)$ 两种写法（最大偏差 $3.0\times10^{-12}$）。

(e) 即 (d) 的中间步。(f)：$Qx\perp(I-Q)x$ 是投影的基本性质，
故 $\|(I-Q)x\|^2=\|x\|^2-\|Qx\|^2=1-\alpha$。

(g)：$x,d\ge0$ 逐元素，故 $\langle x,d\rangle\le\|x\|\,\|d\|=\sqrt\delta$，即 $\beta^2\le\delta$；
代入 (c) 得 $\alpha(1-\alpha)\ge0$，即 $0\le\alpha\le1$。

(h)：$|Q|=Q_++Q_-=Q+2Q_-$，故 $\ell=x^\top Qx+2x^\top Q_-x=\alpha+2\beta=(\rho-\beta)+2\beta=\rho+\beta$。\qedhere
\end{proof}

\begin{remark}[刚性系统的信息含量，以及它的极限]\label{rmk:rigiditylimit}
把 (a)(c)(g) 合起来看：$\rho=\alpha+\beta$、$\delta=\beta^2+\alpha(1-\alpha)$、$\beta^2\le\delta$。
第三条被第二条\emph{自动}满足（差恰为 $\alpha(1-\alpha)\ge0$），因此这组恒等式给出的
\textbf{唯一}新约束是 $0\le\alpha\le1$。若再对 $\beta$ 单独使用 $\beta=\langle x,d\rangle\le\|d\|=\sqrt\delta$
并结合 (c)，只能得到 $\rho\le\alpha+\sqrt{\alpha(1-\alpha)}\le\frac32$（在 $\alpha=\frac34$ 处取最大）。
也就是说：\textbf{Perron 刚性系统本身不足以给出 $\rho\le\rho_0$}。
它把问题精确地压缩成「在 $\alpha\in[0,1]$ 给定时 $\beta$ 还能有多大」，
而 Cauchy--Schwarz 在冠军处\emph{不取等}：$\beta^2=0.093008493$ 对 $\delta=0.115865452$，
余量恰为 $\alpha(1-\alpha)=0.022856959$。这条余量就是「缺失的第四条不等式」应当吃掉的东西。
这也解释了为什么第~\ref{sec:domdual} 节的凸化与注~\ref{rmk:tight} 的 SDP 松弛都必然失效：
它们都只用到了 (a)(b)(c) 这一层。
\end{remark}

\begin{proposition}[\STW\ 损失的精确解耦，已证]\label{prop:decouple}
设 $Q\in\mathcal P_3(N)$、$x\ge0$、$\|x\|=1$、$\rho=x^\top Q_+x\le\rho(Q_+)$，
记 $\phi=\gamma(3)=\frac{1+\sqrt5}{2}$（\STW\ 对 $\ell=x^\top|Q|x$ 的普适上界）、
$\Phi=\frac{3+\sqrt5}{4}$（\STW\ 对 $\rho$ 的普适上界）。则恒有
\[
\underbrace{(1-\alpha)}_{\text{投影能量损失}}\;+\;
\underbrace{(\phi-\ell)}_{\ell_1\ \text{损失}}\;=\;1+\phi-2\rho ,
\]
从而（因 $1+\phi=2\Phi$）$\ \rho\le\rho_0\iff(1-\alpha)+(\phi-\ell)\ge2(\Phi-\rho_0)$。
冠军处两项分别为 $1-\alpha=0.023404741$、$\phi-\ell=0.031492854$，
和为 $0.054897595=2(\Phi-\rho_0)$，即 \STW\ 在冠军处多给的量被\textbf{精确劈成两半}。
\end{proposition}

\begin{proof}
$\ell-\alpha=(\alpha+2\beta)-\alpha=2\beta$（命题~\ref{prop:perronrig}(h)(a)），故
$(1-\alpha)+(\phi-\ell)=1+\phi-(\alpha+\ell)=1+\phi-2\alpha-2\beta=1+\phi-2\rho$。
又 $1+\phi=1+\frac{1+\sqrt5}{2}=\frac{3+\sqrt5}{2}=2\Phi$。\qedhere
\end{proof}

\begin{remark}
命题~\ref{prop:decouple} 把「改进 \STW」这件模糊的事翻译成\textbf{一条精确的账}：
\STW\ 在冠军处总共多给 $2(\Phi-\rho_0)=0.054897595$，其中 $0.023404741$ 来自「把 $x^\top Qx$ 放成 $1$」，
$0.031492854$ 来自「把 $x^\top|Q|x$ 放成 $\phi$」。两处放缩都不能单独修好
（冠军的 $\ell=1.5865<\phi=1.6180$、$\alpha=0.9766<1$，两处都\emph{差一点}），
因此任何证明都必须是\textbf{联合}的：需要形如 $\alpha+\ell\le2\rho_0$ 的一条不等式，
而不是分别改进两个因子。这也从原理上说明了第~\ref{sec:stwanat} 节
「拆成两项分别封顶」的做法不可能到 $\rho_0$。
\end{remark}

\subsubsection{极值阶梯：$\kappa_3(N)$ 的精确值}\label{sec:ladder}

\begin{theorem}[极值阶梯]\label{thm:ladder}
\[
\begin{aligned}
&\kappa_3(3)=1 \ \ \textbf{（已证）},\qquad
 \kappa_3(4)=\frac54 \ \ \textbf{（已证）},\\[2pt]
&\kappa_3(5)=\frac{9+4\sqrt6}{15}=1.2531972647\ \ \textbf{（$C_{4v}$ Ansatz 下已证）},\\[2pt]
&\kappa_3(N)=\rho_0\quad(6\le N\le14)\ \ \textbf{（数值；}6\le N\le12\textbf{ 两套独立求解器，}
 N=13,14\textbf{ 仅一套）}.
\end{aligned}
\]
\begin{enumerate}[label=(\roman*),leftmargin=2em]
\item $N=3$：$\mathcal P_3(3)=\{I_3\}$，$(I_3)_+=I_3$，$\rho=1$。\textbf{已证}。
\item $N=4$：$\mathcal P_3(4)=\{I_4-vv^\top:\|v\|=1\}$。极值在 $v=\frac12(1,1,-1,-1)$ 达到，
此时
$Q_+=\frac14\left(\begin{smallmatrix}3&1&2&2\\1&3&2&2\\2&2&3&1\\2&2&1&3\end{smallmatrix}\right)$、
$\mathrm{spec}(Q_+)=\{\frac54,\frac34,\frac34,\frac14\}$，故 $\kappa_3(4)\ge\frac54$。
反向按 $v$ 的符号分布分三类：
$|P|=|M|=2$ 时把 $Q_+$ 写成 $2\times2$ 分块，Schur 判据给 $\theta_P\theta_M\le1$，
再用 $f(t)=\frac{t}{t+1/4}$ 的凹性与 Jensen 得类内等权最优，等号当且仅当 $a_i=\frac12$，
差式精确为 $-3(A-\frac12)^2\le0$；
$|P|=1$ 时归结为 $32A^2-24A+7>0$（判别式 $-320<0$）；
$|P|=0$ 或 $4$ 时 $Q_+=Q\preceq I$、$\rho\le1<\frac54$。
故 $\kappa_3(4)=\frac54$。\textbf{已证}（符号推导，证书 C 段）。
\item $N=5$：在 $C_{4v}$ 对称 Ansatz（一个 hub $+$ 四个等价原子）下，秩三投影条件迫使
$q_{01}^2=u(1-4u)$，且 hub--$k_0$ 块给出闭式
\[
\rho(u)=\frac34-\frac u2+\frac12\sqrt{D(u)},\qquad D(u)=\frac14+9u-15u^2 ;
\]
唯一临界点 $u^\ast=\frac{9-\sqrt6}{30}=0.218350342$ 处 $D(u^\ast)=\frac32$ \textbf{精确}，
$\rho(u^\ast)=\frac{9+4\sqrt6}{15}$；其余特征值 $u+\frac12$（二重）与 $\frac12-u$ 在 $u^\ast$ 处均更小。
极值掩码是 $I+A(K_1\vee C_4)=I+A(W_5)$，这解释了它与八面体图数值的巧合
（命题~\ref{prop:octnogo}）。
\textbf{尚未}排除非 $C_{4v}$ 的 $N=5$ 投影，故标为 \emph{Ansatz 下精确}；
数值上升（F 段，$250$ 起点）未发现更大的值。
\item $6\le N\le14$：最大值均为 $\rho_0$。其中 $6\le N\le12$ 由\textbf{两套互相独立}的求解器给出：
(a) 自洽场 SCF（\texttt{\_kappa.py}，每 $N$ 取 $80\sim300$ 个起点）；
(b) 直接在 Stiefel 流形 $\{V:V^\top V=I_3\}$ 上用解析梯度
$\nabla_V\rho((VV^\top)_+)=2\,(I-VV^\top)\bigl(\mathbf 1\{VV^\top\ge0\}\odot xx^\top\bigr)V$
上升、QR 回缩、步长回溯（\texttt{\_kd.py}，每 $N$ 取 $150\sim400$ 个起点）。
与闭式之差 $\le2\times10^{-15}$，Stiefel KKT \eqref{eq:stiefelkkt} 残差 $\le3\times10^{-9}$。
$N=13,14$ \textbf{只有} (a) 一套（\texttt{\_kappa\_profile.json} 覆盖 $N=3,\dots,14$，
而 \texttt{\_kappa\_direct.json} 只到 $N=12$），故这两档的证据强度低一档，引用时须注明。
$N\ge7$ 的最大值点是\textbf{退化}的：冠军的六个原子加上若干低权原子
（例如 $N=8$ 时 $\mathrm{diag}(Q)$ 末两位为 $3\times10^{-6}$、$2\times10^{-7}$ 量级），目标值不变。
\textbf{数值}，不作成果；见注~\ref{rmk:ladderinfo}。
\end{enumerate}
\end{theorem}

\begin{remark}[阶梯读出的两条结构信息]\label{rmk:ladderinfo}
(1) \textbf{跳跃位置}：$\kappa_3(5)=1.2531973<\rho_0=1.2815682$，$\kappa_3(6)=\rho_0$。
即「六个原子」是达到最优常数的\emph{最小}规模——与冠军核心 $K_1\vee C_5$ 恰好六原子一致，
也与定理~\ref{thm:selfcons} 的支撑界方向相同。
(2) \textbf{退化性}：$N\ge7$ 处最大值点不孤立（可以在几乎不改变目标值的情况下加入低权原子），
说明 $\kappa_3$ 的上确界\emph{不是}靠「更多原子」在 $N\to\infty$ 处逼近的，而是从 $N=6$ 起就已取到。
这对证明策略有利：不需要处理 $N\to\infty$ 的极限行为，
只要能证明「$N\ge7$ 的任意投影可以化约到某个 $\le6$ 个有效原子的构型而不增大 $\rho(Q_+)$」。
\end{remark}

\subsubsection{两条闭式一阶条件：自洽场方程与 Stiefel KKT}

$\kappa_3$ 的一阶最优性条件都有闭式，这让「穷举临界点」在原则上成为有限问题。

\paragraph{自洽场（SCF）}给定 $V\in\mathrm{St}(N,3)$，令 $Q=VV^\top$、$x$ 为 $Q_+$ 的 Perron 向量、
$\Sigma=\mathbf 1\{Q\ge0\}$，再令 $V'$ 为 $\Sigma\odot xx^\top$ 的前三特征向量。
则 $V'=V$ 当且仅当 $V$ 是 $\rho((VV^\top)_+)$ 的临界点。两步都是\textbf{闭式}（各一次特征分解），
无需线搜索。这是 \texttt{\_kappa.py} 的迭代格式，收敛判据取目标值差 $<10^{-14}$。

\paragraph{Stiefel KKT}在光滑点（$Q$ 无非对角零元、$\rho(Q_+)$ 单重）处
$\nabla_Q\rho(Q_+)=\mathbf 1\{Q\ge0\}\odot xx^\top=:G$；
$\mathrm{St}(N,3)=\{V:V^\top V=I_3\}$ 在 $V$ 处的切空间为 $\{\Delta:V^\top\Delta+\Delta^\top V=0\}$，
而目标在右 $O(3)$ 作用下不变，故一阶条件化简为
\begin{equation}\label{eq:stiefelkkt}
(I-Q)\,G\,V=0,\qquad G=\mathbf 1\{Q\ge0\}\odot xx^\top,\quad Q=VV^\top .
\end{equation}
本轮在 $N=3,\dots,8$ 的最大值点处核验 \eqref{eq:stiefelkkt}，残差 $\le3\times10^{-9}$（证书 F 段）。
固定符号模式 $\Sigma$ 后，\eqref{eq:stiefelkkt} 连同 $Q_+x=\rho x$、$V^\top V=I_3$
构成一个\textbf{多项式}方程组（未知量 $V,x,\rho$），
因此对每个 $\Sigma$ 临界点是代数簇上的点——
这正是第~\ref{sec:w1global} 节「穷举 KKT 临界点 $+$ 区间认证」路线在\emph{全局}（而非单个腔面）上的版本，
且 $N=6$ 时符号模式只有 $156$ 个同构类。

\subsubsection{Collatz 包络：为什么 $\ell_1$/行和型上界不可能到 $\rho_0$}

$Q_+\ge0$，故 Collatz--Wielandt 给出
$\min_i(Q_+\mathbf 1)_i\le\rho(Q_+)\le\max_i(Q_+\mathbf 1)_i$，
且对任意 $u>0$ 有 $\rho(Q_+)\le\max_i\frac{(Q_+u)_i}{u_i}$（在极值点对 $u=x$ 取等）。
冠军处
\[
Q_+\mathbf 1=\bigl(1.384614891548,\ 1.258722890561\times5\bigr),\qquad
\rho_0=1.281568196872\in[\,1.2587229,\ 1.3846149\,],
\]
包络宽 $0.1258920$。由此得到两条\textbf{否证性}结论：
\begin{itemize}[leftmargin=1.6em,itemsep=1pt]
\item 任何只用「$Q_+$ 的行和」（即不带权的 Collatz 上界）的论证，在冠军处最多把 $\rho$ 夹到 $1.3846$，
\textbf{不可能}证到 $\rho_0$；必须用到带权形式，也就是必须同时确定 Perron 向量 $x$。
\item \STW\ 式的 $\ell_1$ 上界在冠军处给
$\frac1N\|Q_+\|_1=\frac{7.678229347524}{6}=1.279704891254$，\textbf{比 $\rho_0$ 还小}。
这说明「$\rho(Q_+)\le\frac1N\|Q_+\|_1$」这类不等式\emph{根本不存在}
（它会对冠军给出假命题），$\ell_1$/优占/行和这一整族方法因此被整体排除在主线之外，
与第~\ref{sec:dead} 节和注~\ref{rmk:tight}（SDP 松弛零对偶间隙）的结论一致。
\end{itemize}

\subsubsection{本轮对主线的净影响，与下一步的优先队列}

\begin{itemize}[leftmargin=1.6em,itemsep=1pt]
\item \textbf{收缩}：主线缺口从「$N\le7$ 剩下的 $420$ 个掩码 $+$ 一切 $N\ge8$」收缩为
\textbf{一条}不等式 $\kappa_3\le\rho_0$（推论~\ref{cor:kappamain}），它不含图、不含掩码、不含阶数。
\item \textbf{新工具}：命题~\ref{prop:perronrig} 的精确恒等式组（含 $0\le\alpha\le1$、
$\delta=\beta^2+\alpha(1-\alpha)$、$Qd=(\rho-1)Qx$、$\ell=\rho+\beta$）；
命题~\ref{prop:decouple} 把 \STW\ 损失劈成可分别计量的两半；
\eqref{eq:stiefelkkt} 给出全局临界点的代数刻画；定理~\ref{thm:ladder} 给出 $N=3,4$ 的精确值与 $N=5$ 的闭式。
\item \textbf{已定位的缺失环节}：注~\ref{rmk:rigiditylimit} 说明还需要\textbf{第四条}不等式，
形如 $\alpha+\ell\le2\rho_0$（等价于 $\rho\le\rho_0$，但把两处放缩\emph{联合}处理），
或「$d=Q_-x$ 时 $\beta$ 不能同时接近 $\sqrt\delta$ 与 $\sqrt{\alpha(1-\alpha)}$」。
冠军处 Cauchy--Schwarz 余量恰为 $\alpha(1-\alpha)=0.022856959$。
\item \textbf{已关闭的路线}（本轮新增两条，均为\emph{定量}否证，不是「试了没成功」）：
$\ell_1$/行和/优占族——冠军处 $\frac1N\|Q_+\|_1=1.279705<\rho_0$，故该族\emph{不存在}能达到 $\rho_0$ 的不等式；
负质量界 A2——命题~\ref{prop:A2no} 给出 $7$ 原子显式反例（超出 $\tau$ 达 $6.6\times10^{-3}$，纯有理算术证明）。
两者失效的原因同一条：都删去了 $\Pi_3$ 的秩三结构。$\kappa_3$ 是目前\emph{唯一}保留该结构的表述。
\item \textbf{优先队列}（已写入 \texttt{研究方法优先级协议.md}）：
P1 证 $\kappa_3(6)=\rho_0$（$N=6$ 是跳跃点，$\mathrm{Gr}(3,6)$ 只有 $9$ 维、符号模式只有 $156$ 个同构类，
可用 \eqref{eq:stiefelkkt} 逐类求解 $+$ 区间认证）；
P2 证「$N\ge7$ 可化约到 $\le6$ 个有效原子」（注~\ref{rmk:ladderinfo}(2) 的退化性给出强支持；
第~\ref{sec:exh7} 小节新增一条：$F=\rho_0$ 的 $17$ 个掩码全部含诱导 $C_5$、全部 $\mathrm{npos}=3$、
负质量全部恰为 $\tau$（偏差 $\le1.6\times10^{-15}$），其中 $14$ 个的活跃权数已降为 $6$）；
P3 按注~\ref{rmk:normcaution} 重写定理~\ref{thm:reduce} 的规范化证明；
P4 联合不等式 $\alpha+\ell\le2\rho_0$ 的直接证明（框架能量 / Schur--Horn 优占路线）——
注意命题~\ref{prop:A2no} 之后，P4 是\emph{唯一}还活着的凸化型入口，优先级应视为仅次于 P1。
\end{itemize}

\section{本轮主线：$N\le6$ 的全局精确认证，与 $N=5$ 唯一缺口的二维归约}\label{sec:n6global}

第~\ref{sec:kappa} 节把主线压成一条不等式 $\kappa_3\le\rho_0$，并把优先队列的第一位定为
「证 $\kappa_3(6)=\rho_0$」。本轮沿这条队列往下推，得到三组结果：
(a) $N\le6$ 全部 $205$ 个同构类的\textbf{逐个精确有理认证}，把缺口从「$N\le6$ 的 $205$ 类」压到
「$N=5$ 的 $1$ 类 $+$ $N=6$ 的 $21$ 类」（定理~\ref{thm:ncert6}）；
(b) 对 $N=5$ 的那唯一一类 $K_{2,3}$，用 Grassmannian 恒等式把 $F(\Sigma)$ 改写成
$\R^2$ 中 Parseval 五框架的一个矩阵的最大特征值，从而把原来的「$\Delta_4$ 上五次方程」缺口
换成\textbf{一条二维框架不等式}（命题~\ref{prop:grassid}、\ref{prop:framered}）；
(c) 对这条不等式的\textbf{双轴子族}做部分推进——先把双轴族参数化为 $15$ 个划分 $\times8$ 个符号类共 $120$ 类
（引理~\ref{lem:biaxparam}），再对其中 $120$ 个\emph{等分}构型用有理算术穷举 $\frac54I-B\succeq0$
（命题~\ref{prop:biaxenum}，已证且精确），并给出一条覆盖 $99.83\%$ 帧的 Schur 补充分条件（命题~\ref{prop:schurcond}）。
但\textbf{双轴子族没有被完整证明}：原计划用来把等分点推广到全族的「$\lambda_{\max}(B(t))$ 关于 $t$ 是凹的」
这条引理已被\textbf{否证}（$360000$ 对中 $57149$ 对违反，占 $15.9\%$，见第~\ref{sec:refuted5} 节末行），
失效机制是 $\lambda_{\max}$ 对矩阵是凸的、与逐条目凹接不起来。
同时本轮\textbf{否证}了上一轮的三条结论/路线，其中一条是上一轮报告里写下的定理，必须先撤回。

\emph{（第六版补注）}上面 (b)、(c) 两条所服务的目标——证 $F(\Sigma_{K_{2,3}})\le\frac54$——
已经由定理~\ref{thm:bipglobal} 用一条\emph{完全不同}的路线无条件拿到（推论~\ref{cor:n5uncond}），
而且那条路线对\emph{一切}完全二部掩码与\emph{一切}阶数同时成立。
本节以下内容因此全部保留但降级为「同一常数的第二种独立表述」，比较见注~\ref{rmk:bipvsframe}。

\subsection{三处程序缺陷与三条被否证的旧结论}\label{sec:bugs5}

本节的所有内容都由 \texttt{证书\_N6全局精确认证.py}（$49$ 项检查）与
\texttt{证书\_N5框架归约.py}（$7$ 段）复跑产生。按惯例，先交代错误。

\paragraph{缺陷 1：\texttt{nx.canonical\_graph} 在 $N=6$ 段崩溃}
上一轮的全局认证脚本在 G 段（$21$ 行例外表）调用 \texttt{networkx.canonical\_graph} 求同构规范形时抛异常，
导致 G 段之后的 H 段（$K_{2,3}$ 的精确结构，$486$ 行）\textbf{从未执行}，
\texttt{\_bad21table.json} 与 \texttt{\_kappa5.json} 两个产物也从未写出。
本轮改为「\texttt{graph6} 码作键 $+$ \texttt{is\_isomorphic} 代表元缓存」，
并把 H 段并入同一脚本；两份产物现已正常生成。
\emph{影响范围}：上一轮报告中所有引用 H 段数据的陈述（$K_{2,3}$ 的特征多项式闭式、
$21$ 行例外表、Case~A 的精确闭合）在上一轮\textbf{没有}被任何一次成功的运行支持过，
本轮是它们第一次真正跑通。

\paragraph{缺陷 2：E 段把 $-0.0$ 读成违例}
\texttt{lem:bdryrig} 的数值检验里，把某个权 $w_i$ 置零后，$\Pi_3(M)$ 的第 $i$ 行\textbf{恰好}为零行，
于是翻转像 $1\{\Pi_3\ge0\}$ 的第 $i$ 行应按约定 $1\{0\ge0\}=1$ 取全 $1$。
但代码直接对浮点结果做 \texttt{(P>=0)}，而 IEEE 的 $-0.0$ 与 $10^{-17}$ 级噪声会被判成「$<0$」，
$83$ 次试验\textbf{全部}报假违例。修复方式是先独立验证「$\Pi_3$ 第 $i$ 行的非对角元最大绝对值 $<10^{-10}$」，
确认它是精确零行之后再套用约定。修复后 $83$ 次试验 $0$ 违例。

\paragraph{缺陷 3：双轴正对照用例的符号规格自相矛盾}
H3d 段要验证「完全多部图是腔掩码」，用 $7$ 组（部划分 $+$ 各部符号）构造双轴框架。
但构造函数把前 $\lceil m/2\rceil$ 个部放在直线 $u_1$、其余放在正交直线 $u_2$，
而\textbf{同一直线上的两部必须取相反符号}：否则同线两部顶点的内积为正、$Q_{ij}<0$，
两部就并成\emph{一个}非邻接类，$1\{Q\ge0\}$ 自然不等于期望的完全多部图掩码。
上一轮有两个用例（$3$ 部的两组）违反了这一点，被误报为「完全多部图不是腔掩码」。
修正符号后 $7$ 组全部通过。这一条纯属\textbf{测试用例规格 bug}，不是数学失败；
但它一度让我们以为「双轴路线给出的掩码类很窄」，方向判断被带偏。

\paragraph{否证 1（重要）：旧 \texttt{lem:planar}「$K_{2,3}$ 不是腔掩码」是\textbf{错的}}
上一轮用「倍角最大圆周间隙 $>180^\circ$」的间隙计数论证 $K_{2,3}$ 不可能来自平面 Parseval 框架，
并据此\textbf{无条件}装配出 $\kappa_3(5)\le\rho_0$。本轮给出显式反例：
\[
Z=\Bigl[\tfrac1{\sqrt2}\ \ 0;\ \tfrac1{\sqrt2}\ \ 0;\ \ 0\ \ \tfrac1{\sqrt3};\ \ 0\ \ \tfrac1{\sqrt3};\ \ 0\ \ \tfrac1{\sqrt3}\Bigr]\in\mathrm{St}(5,2),
\qquad Q=I_5-ZZ^\top=\begin{pmatrix}
\frac12&-\frac12&0&0&0\\[1pt]
-\frac12&\frac12&0&0&0\\[1pt]
0&0&\frac23&-\frac13&-\frac13\\[1pt]
0&0&-\frac13&\frac23&-\frac13\\[1pt]
0&0&-\frac13&-\frac13&\frac23
\end{pmatrix},
\]
$Q^2=Q$、$\tr Q=3$ 精确成立（\texttt{sympy} 根式判定），且 $1\{Q\ge0\}=\Sigma_{K_{2,3}}$ 逐条目相符。
失效原因被完全定位：该见证的倍角集合是 $\{0^\circ,0^\circ,180^\circ,180^\circ,180^\circ\}$，
圆周间隙为 $(0,180,0,0,180)$，\textbf{最大间隙恰为 $180^\circ$ 而不是严格大于 $180^\circ$}；
旧证明的间隙计数依赖严格不等式，在这个边界情形失效。
\textbf{结论}：上一轮那条「$\kappa_3(5)\le\rho_0$ 无条件」的装配步骤作废，必须撤回；
$K_{2,3}$ 确实是腔掩码，而且落在二分法的\emph{边界}上。

\paragraph{否证 2：上一轮「两个内部临界点是特征值排序 bug 的假象」也是错的}
把 \texttt{\_cp5.py} 的特征值与特征向量\textbf{同步}排序后复跑（$2000$ 起点，接受阈 $\|r\|^2<10^{-20}$），
两个内部临界点都复现，其中一个还有闭式（命题~\ref{prop:k23cp}）：
$S_2\times S_3$ 对称点 $a=\frac{62-6\sqrt{35}}{119}$ 处 $\Lam_3=\frac{63+12\sqrt{35}}{119}=1.125991238631895$
（\texttt{sympy} 精确判定）；另一个非对称点 $\Lam_3=1.159910859786268$（数值）。
两者都严格小于 $5/4$，所以「不存在内部极大化点 $\Rightarrow$ 极大化点必在边界 $\Rightarrow$ 归约到 $\kappa_3(4)$」
这条\textbf{路线不成立}：内部临界点是真实的，只是它们不是极大化点。
这条否证把证明的担子重新压回到「$\Delta_4$ 上的全局上界」本身。

\paragraph{否证 3：翻转链路线不能认证 $N=6$ 的 $21$ 个例外}
\texttt{lem:flip} 说：若 $w>0$ 且翻转像 $\Sigma'=1\{\Pi_3(D_w\Sigma D_w)\}\ne\Sigma$，
则 $\Lam_3(D_w\Sigma D_w)<F(\Sigma')$，于是「沿翻转链单调上升、终止于已认证类」就能给出 $F(\Sigma)\le\rho_0$。
数值上，$21$ 个例外在\emph{各自的数值极大化点}处的翻转像确实全部落进已认证的 $134$ 类。
但把采样面从「极大化点」扩大到「每类 $608$ 个球面点（顶点/面心/均匀/Dirichlet 两档）」之后，
$647$ 种翻转像同构类的 $\Omega$ 最大值是 $\mathbf{1.310920}$，\textbf{超过} $\rho_0=1.281568$ 达 $2.9\times10^{-2}$。
也就是说：翻转像只在极大化点处落进已认证类，而在 $S^5$ 的其余点上会落到 $\Omega>\rho_0$ 的类。
用翻转链证 $F\le\rho_0$ 因此\textbf{预先假定了极大化点的位置}，是循环论证。
\texttt{prop:kappa6red} 的这条路线就此关闭；本轮把它写成一条\emph{否证性}检查（判据要求 $\max\Omega>\rho_0$），
以免下次再被误当作通过。

\subsection{修正后的平面二分法}\label{sec:planarfix}

否证 1 迫使间隙判据重写。正确的形式如下，其中「倍角」是把 $\R^2$ 中的方向映到圆周上的标准技巧
（$\theta$ 与 $\theta+180^\circ$ 给出同一条直线，故取 $2\theta$）。

\begin{lemma}[平面二分法，修正版]\label{lem:planar}
设 $Z\in\mathrm{St}(5,2)$，行向量 $z_i=\sqrt{t_i}\,(\cos\theta_i,\sin\theta_i)$，$t_i=\|z_i\|^2\ge0$。
则 $Z^\top Z=I_2$ 当且仅当
\begin{equation}\label{eq:moment2}
\sum_{i=1}^5 t_i=2,\qquad \sum_{i=1}^5 t_i\bigl(\cos2\theta_i,\ \sin2\theta_i\bigr)=(0,0).
\end{equation}
把倍角 $2\theta_i$ 按圆周排序，记最大相邻间隙为 $\gamma$。则恰有两种情形：
\begin{enumerate}[label=(\alph*),leftmargin=2em]
\item $\gamma<180^\circ$（\textbf{非双轴支}）：此时 $0$ 落在 $\{e^{2\mathrm i\theta_i}\}$ 的凸包\emph{内部}；
\item $\gamma=180^\circ$（\textbf{双轴支}）：此时所有 $z_i$ 落在\textbf{两条互相正交的直线}上，
即存在正交单位向量 $u\perp v$ 与 $\{0,\dots,4\}$ 的一个划分 $L_1\sqcup L_2$，使
$z_i=\varepsilon_i\sqrt{t_i}\,u\ (i\in L_1)$、$z_i=\eta_i\sqrt{t_i}\,v\ (i\in L_2)$，$\varepsilon_i,\eta_i\in\{\pm1\}$，
且 $\sum_{i\in L_1}t_i=\sum_{i\in L_2}t_i=1$。
\end{enumerate}
\end{lemma}

\begin{proof}[证明要点]
\eqref{eq:moment2} 是把 $ZZ^\top$ 的幂等性/迹条件展开：$Z^\top Z=I_2$ 等价于
$\sum_i t_i\cos^2\theta_i=\sum_i t_i\sin^2\theta_i=1$ 与 $\sum_i t_i\sin\theta_i\cos\theta_i=0$，
用二倍角公式即得。
(a)$\Rightarrow$ 非双轴是显然的。(b)：$\gamma=180^\circ$ 意味着倍角点集落在某条直径的两侧且两端都有点，
配合 $\sum t_ie^{2\mathrm i\theta_i}=0$ 迫使全部倍角只取\emph{两个相差 $180^\circ$} 的值，
即 $\theta_i$ 只取两个相差 $90^\circ$ 的值，故 $z_i$ 落在两条正交直线上；
再把 \eqref{eq:moment2} 的第一式按直线分组即得两组权和各为 $1$。
数值核验：$20000$ 个随机平面 Parseval 五框架上 \eqref{eq:moment2} 的最大残差 $2.0\times10^{-15}$、
$\sum t_i$ 的最大偏差 $1.8\times10^{-15}$，二分法反例 $0$ 个（证书 H3c 段）。
\end{proof}

\begin{remark}[为什么 $K_{2,3}$ 恰好在边界上]\label{rmk:boundary}
否证 1 的见证 $Z$ 的倍角是 $\{0,0,180,180,180\}$，属 (b) 支；
命题~\ref{prop:biaxsharp} 的取等框架 $Z^\ast$ 的倍角同样只差 $180^\circ$。
这不是巧合：$K_{2,3}$ 的掩码要求 $A$ 部内部与 $B$ 部内部都\emph{非邻接}，
即 $\langle z_i,z_j\rangle>0$ 对同部对成立，而跨部对要\emph{非负}。
在同部内积严格为正的约束下，跨部内积只能靠「正交或反向」实现，这正是双轴构型。
因此引理~\ref{lem:planar} 的 (b) 支\textbf{不是}可以随手排除的退化情形，
而是本问题的\emph{主战场}——这也解释了为什么上一轮的间隙论证一定会在这里出错。
\end{remark}


\subsection{$N\le6$ 全部同构类的精确有理认证}\label{sec:ncert6}

定理~\ref{thm:domdual} 给出 $F(\Sigma)\le G(\Sigma)\le\Omega(\Sigma)$，其中
$\Omega(\Sigma)=\min\{\max_iT_{ii}:T\succeq\Sigma,\ T\succeq0\}$。
上一轮只在冠军掩码与四类近腔掩码上做过 $\Omega$ 的精确证书。本轮把它推到 $N\le6$ 的\textbf{全部}同构类：
对每个类先解 $\Omega$ 的 SDP，把最优 $T$ 用二进有理数舍入到 $2^{-20}$，
再在 $\Q$ 上用 Bareiss 精确分数行列式逐个判定 $T$ 与 $T-\Sigma$ 的\textbf{全部}主子式非负。
这一步没有浮点：一旦 $205$ 个类的 $19442$ 个主子式都被判为非负且最小者 $>0$，
相应的 $\Omega\le\rho_0$ 就是定理。

\begin{theorem}[$N\le6$ 的逐类精确认证]\label{thm:ncert6}
对 $N=3,4,5,6$，$\mathcal P_3(N)$ 的掩码按图同构分类，各类的 $\Omega$ 由精确有理证书给出，结果如下：
\begin{center}
\begin{tabular}{r@{\quad}l r r r r r}
\toprule
$N$ & 同构类数 & 精确认证 & 例外 & 主子式数 & 最小主子式\\
\midrule
$3$ & $4$   & $4$   & $0$  & $56$    & $8.674\times10^{-19}$\\
$4$ & $11$  & $11$  & $0$  & $330$   & $8.272\times10^{-25}$\\
$5$ & $34$  & $33$  & $1$  & $2046$  & $7.889\times10^{-31}$\\
$6$ & $156$ & $134$ & $22$ & $17010$ & $7.523\times10^{-37}$\\
\midrule
合计 & $205$ & $182$ & $23$ & $19442$ & \\
\bottomrule
\end{tabular}
\end{center}
例外清单的结构是干净的：
\begin{enumerate}[label=(\roman*),leftmargin=2em]
\item $N=3,4$：\textbf{无例外}。
\item $N=5$：\textbf{恰 $1$ 个}，即 $\Sigma_{K_{2,3}}=I_5+A(K_{2,3})$（度序列 $(2,2,2,3,3)$、无支配点），
其 $\Omega=\frac{22}{17}=1.294117647058824$ 有显式有理证书
\[
T=\frac1{17}\begin{pmatrix}22&5&12&12&12\\5&22&12&12&12\\12&12&22&5&5\\12&12&5&22&5\\12&12&5&5&22\end{pmatrix},
\qquad T\succeq\Sigma,\quad T-\Sigma\succeq0,\quad \max_iT_{ii}=\tfrac{22}{17},
\]
$T$ 与 $T-\Sigma$ 的各 $31$ 个主子式在 $\Q$ 上逐个非负，SDP 数值解与之差 $\le10^{-9}$。
\item $N=6$：\textbf{恰 $22$ 个} $=$ 冠军掩码 $\Sigma_0$（\texttt{graph6} 码 \texttt{Ehfw}，即 $K_1\vee C_5$）
加上 $21$ 个 $\Omega>\rho_0$ 的类。冠军不是真例外：$\Omega(\Sigma_0)=\rho_0$ \emph{精确取等}，
二进有理数舍入必然失败（定理~\ref{thm:champglobal} 已给出根式证书），它由定理~\ref{thm:champglobal} 单独覆盖。
$21$ 个真例外的 $\Omega$ 分布在 $[1.288675,\ 4/3]$，其中\textbf{四个}达到最大值 $\Omega=\frac43$
（八面体型掩码，与命题~\ref{prop:octnogo} 一致），冠军处 $\Omega=1.281568197$ 与 $\rho_0$ 差 $4.4\times10^{-14}$。
\end{enumerate}
\end{theorem}

\begin{corollary}[$N=3,4$ 的无条件闭合]\label{cor:k34uncond}
$\kappa_3(3)\le\rho_0$ 与 $\kappa_3(4)\le\rho_0$ \textbf{无条件成立}（已证）。
交叉核对：$\max_{4\ \text{类}}F=1.000000000000$、$\max_{11\ \text{类}}F=1.250000000000$，
与定理~\ref{thm:ladder} 的精确值 $\kappa_3(3)=1$、$\kappa_3(4)=\frac54$ 分别相差 $1.1\times10^{-15}$、$4.4\times10^{-16}$。
\end{corollary}

于是主线缺口收缩为：
\begin{equation}\label{eq:gap2}
\kappa_3\le\rho_0\quad\Longleftarrow\quad
\underbrace{F(\Sigma_{K_{2,3}})\le\rho_0}_{N=5\ \text{的}\ 1\ \text{类}}\ \wedge\
\underbrace{\max_{21\ \text{类}}F\le\rho_0}_{N=6\ \text{的}\ 21\ \text{类}}\ \wedge\
\underbrace{N\ge7\ \text{可化约到}\ \le6\ \text{个有效原子}}_{\text{注~\ref{rmk:ladderinfo}(2)}} .
\end{equation}
第一项在第~\ref{sec:frame}--\ref{sec:biax} 节被归约为一条二维框架不等式；
第二项在下面第~\ref{sec:kappa6red} 小节被\textbf{否证}了一条看似可行的路线；第三项未动。

\subsubsection{$N=6$ 的 $21$ 个例外：翻转链路线及其否证}\label{sec:kappa6red}

为了处理第二项，本轮建立了三条工具性引理。

\begin{lemma}[翻转优占]\label{lem:flip}
设 $\Sigma\in\mathcal P_3(N)$、$w\in\Delta_{N-1}$，$M=D_w\Sigma D_w$，$\Pi_3=\Pi_3(M)$ 为其前三特征投影，
$\Sigma'=1\{\Pi_3\ge0\}$。则 $\Lam_3(D_w\Sigma D_w)\le F(\Sigma')$；
若 $w>0$ 且 $\Sigma'\ne\Sigma$，则不等式\textbf{严格}。
\end{lemma}

\begin{corollary}[翻转链]\label{cor:flipsc}
若从 $\Sigma$ 出发反复施加 $w\mapsto$（$F$ 的极大化权）$\mapsto\Sigma'$ 得到的链终止于某个已认证类
$\Sigma^{(k)}$（即 $\Omega(\Sigma^{(k)})\le\rho_0$），则 $F(\Sigma)\le F(\Sigma^{(k)})\le\rho_0$。
\end{corollary}

\begin{lemma}[边界刚性]\label{lem:bdryrig}
若 $w_i=0$ 且 $\lambda_3(M)>0$，则 $\Pi_3(M)$ 的第 $i$ 行\textbf{恰好}为零行，
于是按约定 $1\{0\ge0\}=1$，翻转像 $\Sigma'$ 的第 $i$ 行除对角外全为 $1$，即 $i$ 成为 $\Sigma'$ 的支配点。
\end{lemma}

\begin{lemma}[腔的几何刻画]\label{lem:chamber}
$\Sigma$ 是某个 $\R^3$ 中 Parseval 框架的「非锐角图」当且仅当存在 $Q\in\mathcal P_3(N)$ 使 $\Sigma=1\{Q\ge0\}$；
等价地，$\Sigma=1\{Q\ge0\}$、$Q=I_N-ZZ^\top$、$Z\in\mathrm{St}(N,N-3)$，
且框架条件 $\sum_iw_iw_i^\top=I_3$ 等价于其凸包形式 $\frac13I_3\in\operatorname{conv}\{u_iu_i^\top\}$。
\end{lemma}

\begin{proposition}[$21$ 行例外表]\label{prop:kappa6red}
$N=6$ 的 $21$ 个真例外满足：
(i) 每一类在其数值极大化点处都有 $\min_ip_i\le10^{-14}$，即极大化点落在单纯形的\textbf{边界}面上；
(ii) 其中 $20$ 类的极值图最小度 $\delta\ge1$，唯一 $\delta=0$ 的类已单独标出；
$19$ 类没有支配点、$2$ 类恰有 $1$ 个支配点；
(iii) 在各自的数值极大化点处，翻转像全部落在已认证的 $134$ 类中，$\max\Omega(\text{翻转像})<\rho_0$；
(iv) $21$ 类的数值 $\max F$ 均等于 $\kappa_3(5)=\frac{9+4\sqrt6}{15}=1.253197265$，比 $\rho_0$ 低 $2.84\times10^{-2}$。
\textbf{但}(iii) 不能推广：把采样点从「极大化点」扩到每类 $608$ 个球面点之后，
$647$ 种翻转像同构类的 $\max\Omega=1.310920>\rho_0$（超出 $2.9\times10^{-2}$），
故推论~\ref{cor:flipsc} 的前提\textbf{不成立}，翻转链\textbf{不能}用来认证这 $21$ 类。
\end{proposition}

\begin{remark}[否证的含义]\label{rmk:flipdead}
(iii) 与它的否证合起来说明一件更基本的事：翻转像\emph{依赖}于所选的权 $w$，
而「$w$ 取极大化点」这个选择本身就等价于我们想证的结论。
任何形如「沿某个上升流走到已认证类」的论证都会掉进同一个坑，除非上升流的终点\emph{先验}可确定。
这与注~\ref{rmk:rigiditylimit}（Perron 刚性最多给到 $\frac32$）和
命题~\ref{prop:A2no}（负质量界有七原子反例）是同一类教训的第三次出现：
\textbf{凡是把「极大化点在哪」当作已知信息的论证，都不是证明}。
\end{remark}

\subsection{$K_{2,3}$ 类的精确结构}\label{sec:k23}

由 \eqref{eq:gap2}，$N=5$ 只剩 $K_{2,3}$ 一类。记 $A=\{0,1\}$、$B=\{2,3,4\}$，
$\Sigma=I_5+A(K_{2,3})$，$p\in\Delta_4$，$x=\sqrt p$，$M=D_x\Sigma D_x$，
$a=p_0+p_1$、$b=p_2+p_3+p_4$、$q=p_0p_1$、$e_2=p_2p_3+p_2p_4+p_3p_4$、$e_3=p_2p_3p_4$、
$s_2=$（两个最小权之和）。则 $F(\Sigma_{K_{2,3}})=\max_{p\in\Delta_4}\Lam_3(M)$。

\begin{proposition}[闭式特征多项式]\label{prop:k23poly}
$M$ 的特征多项式为
\[
\lambda^5-\lambda^4+(q+e_2)\lambda^3+(ae_2+qb-e_3)\lambda^2-(2ae_3+3qe_2)\lambda+5qe_3,
\]
且 $\tr M=1$、$\sum_i\lambda_i^2=1-2(q+e_2)=\sum_ip_i^2+2ab$。
\end{proposition}

\begin{proposition}[秩二分解与惯性]\label{prop:k23rank2}
$M=\operatorname{diag}(p)+R$，其中 $R_{ij}=x_ix_j$ 仅当 $i,j$ 分属 $A,B$，且
\[
R=\sqrt{ab}\,\bigl(e_Ae_B^\top+e_Be_A^\top\bigr),\qquad
e_A=\frac{x_A}{\sqrt a},\quad e_B=\frac{x_B}{\sqrt b},\quad \|e_A\|=\|e_B\|=1,\quad e_A\perp e_B,
\]
故 $\operatorname{spec}(R)=\{+\sqrt{ab},-\sqrt{ab},0,0,0\}$、$\rank R=2$。
由 Sylvester 惯性律 $\operatorname{neg}(M)\le\operatorname{neg}(\operatorname{diag}p)+\operatorname{neg}(R)=1$，
即 $M$ \textbf{至多一个}负特征值，从而升序第二特征值 $\lambda_1^{\uparrow}\ge0$，
$\Lam_3(M)=1-(\lambda_1^{\uparrow}+\lambda_2^{\uparrow})\le1-\lambda_{\min}(M)$。
数值核验：$3000$ 组随机 $p$ 上负特征值个数最大为 $1$，$(1-\lambda_{\min})-\Lam_3$ 的最小值 $\ge0$。
\end{proposition}

\begin{proposition}[Ky Fan 界]\label{prop:k23kyfan}
$X\mapsto$（两个最小特征值之和）是凹函数，故由 Ky Fan 最小原理
$\lambda_1^{\uparrow}+\lambda_2^{\uparrow}(M)\ge s_2-\sqrt{ab}$，即
\begin{equation}\label{eq:kf}
\Lam_3(M)\ \le\ 1-s_2+\sqrt{ab}\ =:\ \mathbf{B1}(p).
\end{equation}
\eqref{eq:kf} 在两处\textbf{取等}：均匀点 $p=\frac15\one$ 处 $\Lam_3=\frac{3+\sqrt6}{5}$、
界 $=\frac35+\frac{\sqrt6}5$（精确相同）；极大化点 $p^\ast=(\frac14,\frac14,\frac14,0,\frac14)$ 处两者都 $=\frac54$。
$25$ 万随机点上违例 $0$ 次。
\end{proposition}

\begin{corollary}[Case A 的精确闭合]\label{cor:k23caseA}
若两个最小权都落在 $A$ 部，则 $s_2=a$，于是
$\Lam_3\le1+\max_{s\in[0,2/5]}\bigl(\sqrt{s(1-s)}-s\bigr)=1+\frac{\sqrt2-1}{2}=\frac{1+\sqrt2}{2}=1.207106781<\frac54<\rho_0$。
最大值在 $s^\ast=\frac{2-\sqrt2}4\in[0,\frac25]$ 处取得，\texttt{sympy} 的 \texttt{maximum} 与临界点代入两法一致。
\textbf{已证}，且余量 $\frac54-\frac{1+\sqrt2}2=0.042893$。
\end{corollary}

\begin{proposition}[迹--二次型界]\label{prop:k23trQ}
由命题~\ref{prop:k23poly} 的 $\sum\lambda_i^2=1-2(q+e_2)$，配合
$\Lam_3^2\le3(\lambda_1^2+\lambda_2^2+\lambda_3^2)$、$(\lambda_4+\lambda_5)^2\le2(\lambda_4^2+\lambda_5^2)$
与 $\lambda_4+\lambda_5=1-\Lam_3$，得 $5\Lam_3^2-6\Lam_3-3+12(q+e_2)\le0$，即
\begin{equation}\label{eq:trQ}
\Lam_3(M)\ \le\ \frac35+\frac{2\sqrt3}{5}\sqrt{\,2-5q-5e_2\,}\ =:\ \mathbf{B2}(p),
\end{equation}
其上确界为 $\frac35+\frac{2\sqrt3}{5}\sqrt2=1.579795897111$（在单纯形顶点 $p=(1,0,0,0,0)$ 处，此时 $q+e_2=0$），故它单独\textbf{远不足以}给出 $\frac54$；但它在 $a=0$（即 $p=(0,0,\frac13,\frac13,\frac13)$）处\textbf{取等}，
此处 $\Lam_3=1$、界 $=1$（差 $<10^{-12}$）。$40$ 万随机点上违例 $0$ 次。
\end{proposition}

\begin{proposition}[内部临界点是真实的，且有闭式]\label{prop:k23cp}
限制在 $S_2\times S_3$ 对称子流形 $p=(\frac a2,\frac a2,\frac b3,\frac b3,\frac b3)$（$b=1-a$）上，
$\operatorname{spec}(M)=\{\frac a2,\frac b3,\frac b3,\mu_+,\mu_-\}$，
$\mu_\pm=\frac12\bigl[(\frac a2+\frac b3)\pm\sqrt{(\frac a2+\frac b3)^2+4(ab-\frac{ab}6)}\bigr]$，
故 $\Lam_3(a)=\mu_++\frac{2b}3$ 是单变量闭式。$\frac{\dd}{\dd a}\Lam_3=0$ 在 $(0,\frac25)$ 内的唯一根为
\[
a^\ast=\frac{62-6\sqrt{35}}{119}=0.222718666398339\ldots,\qquad
\Lam_3(a^\ast)=\frac{63+12\sqrt{35}}{119}=\frac9{17}+\frac{12\sqrt{35}}{119}=1.125991238631895,
\]
\texttt{sympy} 精确判定成立，与修正后的 KKT 多起点搜索（$2000$ 起点、$\|r\|^2<10^{-20}$）相符到 $10^{-15}$。
它比 $\frac54$ 低 $0.124009$、比 $\rho_0$ 低 $0.155577$。
另一个\textbf{非对称}内部临界点数值为 $\Lam_3=1.159910859786268$（无闭式）。
\end{proposition}

\begin{remark}[为什么命题~\ref{prop:k23cp} 反而堵死了一条路]\label{rmk:cpdead}
既然内部临界点真实存在，「$F$ 的极大化点必在 $\Delta_4$ 的边界上 $\Rightarrow$ 归约到 $\kappa_3(4)=\frac54$」
就不能靠「内部无临界点」来建立。数值上极大化点确实在边界
（$301$ 起点全局搜索给 $\max\Lam_3=1.250000000$ 于 $p^\ast$，$\min_ip_i=4.6\times10^{-16}$；
限制 $\min_ip_i\ge\varepsilon$ 时 $\max\Lam_3$ 随 $\varepsilon\to0$ 单调升向 $\frac54$ 而恒小于它），
但「两个内部临界点都不是极大化点」这件事本身需要代数证明，目前只有数值。
\end{remark}

\begin{remark}[五张边界面的精确值]\label{rmk:faces}
删去一个原子得到的五张面上，$\max F$ 分别为：删 $A$ 部原子得 $I+A(K_{1,3})$，$F=\frac65$；
删 $B$ 部原子得 $I+A(C_4)$，$F=\frac54$。五者最大值 $=\frac54=\kappa_3(4)$（定理~\ref{thm:ladder}(ii)，已证）。
所以\textbf{若}能证明极大化点必在边界，则 $F(\Sigma_{K_{2,3}})=\frac54$ 立即得证；
命题~\ref{prop:k23cp} 说明这一步现在没有证明。
\end{remark}


\subsection{Grassmannian 恒等式：把 $F$ 从单纯形搬到投影簇}\label{sec:frame}

$F(\Sigma)=\max_{p\in\Delta_{N-1}}\Lam_3(D_{\sqrt p}\Sigma D_{\sqrt p})$ 的自变量是 $N-1$ 维单纯形上的权；
$N=5$ 时这是 $4$ 维，而 $\Lam_3$ 是五次特征方程的根，没有闭式。本轮的转折点是把自变量换掉。

\begin{proposition}[Grassmannian 恒等式]\label{prop:grassid}
对任意 $\Sigma\in\{0,1\}^{N\times N}$ 对称矩阵，
\begin{equation}\label{eq:grassid}
F(\Sigma)=\max_{W\in\mathrm{Gr}(3,N)}\lambda_{\max}^{+}\bigl(P_W\circ\Sigma\bigr),\qquad
\lambda_{\max}^{+}(B):=\max\{x^\top Bx:\ x\ge0,\ \|x\|=1\},
\end{equation}
其中 $P_W$ 是到 $W$ 的正交投影、$\circ$ 是 Hadamard 积。且对任意实对称 $B$ 与 $x\ge0$ 有
\begin{equation}\label{eq:pospart}
x^\top Bx\ \le\ x^\top B_+x\ \le\ \lambda_{\max}(B_+),\qquad B_+:=(B_{ij}\wedge0)\cdot(-1)+B\ \text{的逐元素正部},
\end{equation}
故 $F(\Sigma)\le\max_{W}\lambda_{\max}\bigl((P_W\circ\Sigma)_+\bigr)$。
\end{proposition}

\begin{proof}
Ky Fan 对偶给 $\Lam_3(M)=\max_{W\in\mathrm{Gr}(3,N)}\tr(P_WM)$（定理~\ref{thm:kappa}(i) 已用）。
取 $M=D_x\Sigma D_x$（$x=\sqrt p\ge0$），则 $P_WD_x\Sigma D_x$ 的迹为
$\tr(P_WD_x\Sigma D_x)=\tr(D_xP_WD_x\Sigma)=x^\top(P_W\circ\Sigma)x$，
因为 $\tr(DADB)=\sum_{ij}A_{ij}B_{ji}a_ia_j$ 对对称 $A,B$ 成立。
于是 $F(\Sigma)=\max_{x\ge0,\|x\|_1=1}\max_W x^\top(P_W\circ\Sigma)x$，
而 $\|x\|_1=1\wedge x\ge0\iff\|\sqrt p\|=1\wedge p\in\Delta$，两个 $\max$ 都是紧集上的连续函数，可交换次序。
\eqref{eq:pospart} 的第一段：$x\ge0$ 时 $x^\top(B_+-B)x=\sum_{ij}(B_+-B)_{ij}x_ix_j\ge0$（被求和的每一项都非负）。
第二段是 Rayleigh 商的定义。
\end{proof}

\begin{remark}[为什么这一步不是平凡的换元]\label{rmk:whygrass}
$\Delta_4$ 与 $\mathrm{Gr}(3,5)$ 都是 $4$ 维/ $6$ 维流形，维度并没有变小。
真正的收获是\textbf{结构}：在 $\mathrm{Gr}(3,5)$ 上可以用对偶参数化 $P_W=I_5-ZZ^\top$、$Z\in\mathrm{St}(5,2)$，
把矩阵的每个条目都写成 $\R^2$ 中\emph{五个向量}的内积，
于是引理~\ref{lem:planar} 的倍角--矩参数化 \eqref{eq:moment2} 立刻可用；
而在 $\Delta_4$ 上，$\Lam_3$ 只是一个五次方程的根，没有任何几何把手。
命题~\ref{prop:grassid} 的另一个副产品是：\eqref{eq:pospart} 说明「先取正部再取最大特征值」
这步放缩在取等点处\textbf{零损失}（证书 A3 段：$P_W\circ\Sigma$ 在该 $W$ 上逐元素非负，最小元 $-0.0$），
所以 $\frac54$ 这个值不是放缩造成的，是真值。
\end{remark}

\begin{proposition}[$N=5$ 的平面框架归约]\label{prop:framered}
设 $N=5$、$\Sigma=I_5+A(K_{2,3})$、$A=\{0,1\}$、$B=\{2,3,4\}$。
对 $Z\in\mathrm{St}(5,2)$，行向量 $z_i\in\R^2$、$t_i=\|z_i\|^2$，令 $P=I_5-ZZ^\top$。则
$(P\circ\Sigma)_+$ 的显式形式为
\begin{equation}\label{eq:BZ}
B(Z)=\operatorname{diag}(1-t_1,\dots,t_5)+\bigl[\,(-\langle z_i,z_j\rangle)_+\,\bigr]_{(i,j)\in A\times B},
\end{equation}
即：对角元 $1-t_i\ge0$（因为 $t_i\le1$，由 $\sum_jt_j\hat u_j\hat u_j^\top=I_2$ 与 $\hat u_i\hat u_i^\top\preceq I_2$）；
跨部元 $(-\langle z_i,z_j\rangle)_+\ge0$；同部元为 $0$（$\Sigma$ 在同部为 $0$）。
故 $B(Z)$ \textbf{逐元素非负}，Perron--Frobenius 与 Collatz 均可用，且
\begin{equation}\label{eq:frame5}
\boxed{\ \lambda_{\max}\bigl(B(Z)\bigr)\ \le\ \frac54\qquad\text{对一切 }Z\in\mathrm{St}(5,2)\ }\tag{$\ast_{\!F}$}
\end{equation}
蕴含 $F(\Sigma_{K_{2,3}})\le\frac54$。
\end{proposition}

\begin{theorem}[$N=5$ 层的条件性闭合；第六版起已被超越]\label{thm:n5frame}
\emph{（第六版补注）}本定理的前提 \eqref{eq:frame5} 已被定理~\ref{thm:bipglobal} 绕开：
结论 $F(\Sigma_{K_{2,3}})=\frac54$ 与 $\kappa_3(5)\le\rho_0$ 现在是\textbf{无条件}的（推论~\ref{cor:n5uncond}），
不再依赖任何未证不等式。本节及以下三节（双轴取等证书、双轴支穷举、Schur 型充分条件）全部保留，
作为 Grassmannian 恒等式的应用示范与同一常数的第二种独立表述；
两条路线的难度比较见注~\ref{rmk:bipvsframe}。

假设 \eqref{eq:frame5}。则
$F(\Sigma_{K_{2,3}})=\frac54$ \textbf{精确}，且 $\kappa_3(5)\le\rho_0$（余量 $\rho_0-\frac54=0.031568197$）。
更精确地，逻辑链的每一环状态如下：
\[
\kappa_3(5)\le\max_{34\ \text{类}}F(\Sigma)\ \textbf{[已证]}\quad\wedge\quad
\text{$33$ 类 }F\le\Omega\le\rho_0\ \textbf{[已证，精确有理证书]}\quad\wedge\quad
F(\Sigma_{K_{2,3}})\le\max_Z\lambda_{\max}(B(Z))\ \textbf{[已证]}
\]
\[
\wedge\quad \max_Z\lambda_{\max}(B(Z))\le\tfrac54\ \textbf{[未证，即 \eqref{eq:frame5}]}
\quad\Longrightarrow\quad \kappa_3(5)\le\rho_0 .
\]
下界方向：命题~\ref{prop:biaxsharp} 的双轴框架 $Z^\ast$ 给出 $\lambda_{\max}(B(Z^\ast))=\frac54$，
故 \eqref{eq:frame5} 的常数 $\frac54$ 是\textbf{sharp} 的，且 $F(\Sigma_{K_{2,3}})\ge\frac54$；
与数值上 $\max_p\Lam_3=\frac54$ 合并即得精确值。
\end{theorem}

\begin{remark}[与旧的 $(\ast_5)$ 相比]\label{rmk:vs5}
上一轮的缺口是 $(\ast_5)$：$\lambda_4+\lambda_5\ge1-\rho_0=-0.281568$ 对一切 $p\in\Delta_4$ 成立，
即一个 $4$ 维单纯形上关于五次方程两个最小根的显式下界（数值余量 $0.0316$，$0$ 违例，但无证明思路）。
\eqref{eq:frame5} 把它换成了 $\R^2$ 中五个向量的一个\emph{矩阵}不等式，
常数从 $-0.281568$ 变成 $\frac54$（有理数），并且\textbf{已知取等点与取等时的精确谱}。
这是实质性的改进：$(\ast_5)$ 没有已知的取等点（它是严格不等式，余量 $>0$），
而 \eqref{eq:frame5} 取等，因此任何证明都必须能\emph{解释}这个取等构型——
这为构造证明提供了强约束（例如排除了所有「严格不等式型」的放缩）。
\end{remark}

\subsection{双轴框架处的精确取等证书}\label{sec:biaxsharp}

\begin{proposition}[双轴取等的精确有理证书]\label{prop:biaxsharp}
令
\[
Z^\ast=\begin{pmatrix}3/10&-2/5\\ 3/10&-2/5\\ -3/10&2/5\\ 4/5&3/5\\ -3/10&2/5\end{pmatrix}\in\mathrm{St}(5,2),
\qquad (t_1,\dots,t_5)=\Bigl(\frac14,\frac14,\frac14,1,\frac14\Bigr).
\]
则 ${Z^\ast}^\top Z^\ast=I_2$ 精确成立（\texttt{sympy} 残差为 $0$ 矩阵），$\sum t_i=2$，
且 \eqref{eq:BZ} 给出\textbf{精确有理}矩阵
\[
B(Z^\ast)=\begin{pmatrix}
\frac34&0&\frac14&0&\frac14\\[1pt]
0&\frac34&\frac14&0&\frac14\\[1pt]
\frac14&\frac14&\frac34&0&0\\[1pt]
0&0&0&0&0\\[1pt]
\frac14&\frac14&0&0&\frac34
\end{pmatrix},\qquad
\chi_{B(Z^\ast)}(\lambda)=\frac{\lambda(4\lambda-5)(4\lambda-1)(4\lambda-3)^2}{4^{4}}
=\lambda^5-3\lambda^4+\frac{25}8\lambda^3-\frac{21}{16}\lambda^2+\frac{45}{256}\lambda,
\]
故 $\operatorname{spec}(B(Z^\ast))=\{\frac54,\frac34,\frac34,\frac14,0\}$、$\lambda_{\max}(B(Z^\ast))=\frac54$。
$Z^\ast$ 是双轴框架：$z_0=z_1=\frac12\hat u$、$z_2=z_4=-\frac12\hat u$、$z_3=\hat v$，
其中 $\hat u=(\frac35,-\frac45)\perp\hat v=(\frac45,\frac35)$；
它同时满足 $1\{I_5-Z^\ast{Z^\ast}^\top\ge0\}=\Sigma_{K_{2,3}}$（跨部对 $(0,3),(1,3),(2,3),(3,4)$ 的投影元\emph{恰为} $0$，
按约定 $1\{0\ge0\}=1$），即 $Z^\ast$ 也是引理~\ref{lem:planar}(b) 支的见证。
\end{proposition}

\begin{remark}[$Z^\ast$ 的结构意义]\label{rmk:zstar}
$t_3=1$ 意味着 $z_3$ 是单位向量，于是 $\sum_{i\ne3}t_i\hat u_i\hat u_i^\top=I_2-\hat v\hat v^\top$ 迫使
其余四个向量全部落在\emph{同一条}直线 $\hat u$ 上，且权各为 $\frac14$。
这正是「一条线上一个权 $1$ 的孤原子 $+$ 另一条线上四个权 $\frac14$ 的原子、两两反向配对」的构型，
与 $p^\ast=(\frac14,\frac14,\frac14,0,\frac14)$（单纯形侧的极大化点，第 $3$ 个权为 $0$）一一对应：
$p^\ast=x^2$，而 $x=(\frac12,\frac12,\frac12,0,\frac12)$ 恰是 $B(Z^\ast)$ 的 Perron 向量。
$B(Z^\ast)$ 的第 $4$ 行全零（对应 $p^\ast_3=0$），所以取等是\textbf{退化}的：
一个原子权重为零，问题实际上降到了 $N=4$，这与注~\ref{rmk:faces}（$\max F=\frac54=\kappa_3(4)$）吻合。
\end{remark}

\subsection{双轴支的完整处理：$120$ 个等分构型的精确穷举}\label{sec:biax}

引理~\ref{lem:planar} 把平面框架分成两支。本节处理 (b) 支（双轴），并给出该支上一个\textbf{定理}。

\begin{lemma}[双轴族的参数化与一个精确改写]\label{lem:biaxparam}
双轴帧由三份数据决定：$\{0,\dots,4\}$ 的一个无序划分 $L_1\sqcup L_2$（$15$ 种）、
每部顶点的符号 $\varepsilon,\eta\in\{\pm1\}$（模「整条直线同时变号」这一不改变 $G=ZZ^\top$ 的对称，恰 $8$ 类）、
以及权 $t\ge0$ 满足 $\sum_{L_1}t_i=\sum_{L_2}t_i=1$（一个 $\Delta_{m_1-1}\times\Delta_{m_2-1}$，维数恒为 $3$）。
组合数据共 $15\times8=120$ 类；给定一类后 $\sigma_{ij}:=[\,i,j\ \text{同线且}\ \varepsilon_i\varepsilon_j=-1\,]\in\{0,1\}$ 固定，
$B(t)_{ij}=\sigma_{ij}\sqrt{t_it_j}$。令 $u_i=\sqrt{t_i}$、$\widetilde\Sigma:=-I_5+[\sigma_{ij}]_{A\times B}$，则
\begin{equation}\label{eq:biaxred}
\lambda_{\max}\bigl(B(t)\bigr)=1+\lambda_{\max}\bigl(D_u\widetilde\Sigma D_u\bigr),
\qquad u\in S^{m_1-1}_+\times S^{m_2-1}_+\subset\R^5 .
\end{equation}
于是 \eqref{eq:frame5} 在双轴支上等价于：对 $120$ 个固定的 $\pm1/0$ 矩阵 $\widetilde\Sigma$，
$\max_{u\in S^{m_1-1}_+\times S^{m_2-1}_+}\lambda_{\max}(D_u\widetilde\Sigma D_u)\le\frac14$。
\end{lemma}

\begin{proof}
$I_5-ZZ^\top$ 的跨线元恒为 $0$（两线正交），同线元为 $-\varepsilon_i\varepsilon_j\sqrt{t_it_j}$，
故 $(P\circ\Sigma)_{ij}=(-\langle z_i,z_j\rangle)_+=\sigma_{ij}\sqrt{t_it_j}$，对角为 $1-t_i=1-u_i^2$；
而 $B(t)=I_5-D_{u^2}+[\sigma_{ij}u_iu_j]=I_5+D_u\widetilde\Sigma D_u$，平移 $I_5$ 即得 \eqref{eq:biaxred}。
$\sum_{L_k}t_i=1$ 就是引理~\ref{lem:planar}(b) 的结论。
\end{proof}

\begin{proposition}[等分双轴构型的精确穷举]\label{prop:biaxenum}
在引理~\ref{lem:biaxparam} 的 $120$ 类中取\textbf{等分点} $t^\ast_i=1/|L_{k(i)}|$，
则 $B(t^\ast)$ 全为有理矩阵，且：
\begin{enumerate}[label=(\roman*),leftmargin=2em]
\item 对全部 $120$ 个构型，$\frac54I_5-B(t^\ast)\succeq0$，由\textbf{全部 $31$ 个主子式在 $\Q$ 上非负}判定
（Bareiss 精确分数行列式，无浮点）；
\item 其中\textbf{恰 $3$ 个}满足 $\det(\frac54I_5-B(t^\ast))=0$，即 $\lambda_{\max}=\frac54$ 精确取等；
这 $3$ 个的权向量都是 $(\frac14,\frac14,\frac14,1,\frac14)$ 型，孤原子（权 $1$）分别落在 $B$ 部的三个顶点之一，
$A$ 部两点恒为 $\frac14$；
\item 其余 $117$ 个构型的 $\lambda_{\max}$ 严格小于 $\frac54$，第二大的取值为
$\frac{3+\sqrt3}{4}=1.183012701892$，第三、第四为 $1.138071187458$、$\frac{3+\sqrt2}{4}=1.103553390593$。
\end{enumerate}
故 $\max\{\lambda_{\max}(B(t^\ast)):\ 120\ \text{个等分构型}\}=\frac54$，\textbf{已证且精确}。
\end{proposition}

\begin{remark}[等分点之外的双轴族]\label{rmk:biaxfull}
命题~\ref{prop:biaxenum} 只覆盖 $t$ 的等分点。整个双轴族是 $120$ 个 $3$ 维问题。
数值上：对每类取 $30000$ 个随机 $t$（共 $3\,600\,000$ 点），\textbf{没有任何一点}使
$\lambda_{\max}(B(t))$ 超过等分点的值 $\frac54$；另一套独立的 $187\,614$ 点双轴扫描给
$\sup\lambda_{\max}=1.249993448627<\frac54$（未命中等分点）。$120$ 类中\textbf{没有}一类超过 $\frac54+10^{-10}$。
我一度以为可以用凹性把等分点升级为全族：$t\mapsto B(t)$ 确实逐条目凹
（$\sqrt{t_it_j}$ 凹、对角线性），但 $\lambda_{\max}$ 对矩阵是\textbf{凸}的，
逐条目单调性（Perron）给的方向与凸性给的方向接不起来；
实测 $360\,000$ 对中 $57\,149$ 对（$15.9\%$）违反凹性，故\textbf{凹性论断被否证}，
「临界点即全局最大点」不可用。等分点处切空间方向导数 $<10^{-4}$ 的类有 $82/120$，
另外 $38$ 类等分点不是临界点（最大方向导数 $0.244$），所以也不能靠 KKT 统一处理。
双轴支的全族证明仍是\textbf{未解决}的，但由 \eqref{eq:biaxred} 它已化为
$120$ 个「乘积球面上二次型最大特征值 $\le\frac14$」的标准问题，每个只有 $3$ 维。
\end{remark}

\subsection{一条覆盖 $99.83\%$ 帧的 Schur 补充分条件}\label{sec:schur}

\begin{proposition}[Schur 型充分条件]\label{prop:schurcond}
$\lambda_{\max}(B(Z))\le\frac54$ 等价于 $\operatorname{diag}(\frac14+t_i)\succeq C_{\mathrm{full}}$
（其中 $C_{\mathrm{full}}$ 是 $B$ 的off-diagonal 部分）。由 Schur 补，它的一个\textbf{充分}条件是
\begin{equation}\label{eq:schurS}
L(Z):=\sum_{i\in A}\sum_{j\in B}
\frac{(-\langle z_i,z_j\rangle)_+^{\,2}}{(\frac14+t_i)(\frac14+t_j)}\ \le\ 1 .
\end{equation}
\end{proposition}

\begin{proof}
$(\frac54)I-B=\operatorname{diag}(\frac14+t_i)-C_{\mathrm{full}}$，故等价性是显然的。
写 $C_{\mathrm{full}}=\left(\begin{smallmatrix}0&C_0\\ C_0^\top&0\end{smallmatrix}\right)$、
$D_A=\operatorname{diag}(\frac14+t_i)_{i\in A}$、$D_B=\operatorname{diag}(\frac14+t_j)_{j\in B}$。
Schur 补给出 $D_B-C_0^\top D_A^{-1}C_0\succeq0\Rightarrow$ 原式 $\succeq0$。对任意 $y\in\R^3$，Cauchy--Schwarz 给
\[
\sum_{i\in A}\frac{\bigl(\sum_jC_{ij}y_j\bigr)^2}{\frac14+t_i}
\le\Bigl[\sum_{i\in A}\sum_{j\in B}\frac{C_{ij}^2}{(\frac14+t_i)(\frac14+t_j)}\Bigr]
\Bigl[\sum_{j\in B}\Bigl(\frac14+t_j\Bigr)y_j^2\Bigr],
\]
方括号第一项即 $L(Z)$，故 $L(Z)\le1$ 蕴含 $C_0^\top D_A^{-1}C_0\preceq D_B$。
\end{proof}

\eqref{eq:schurS} 在双轴取等帧处\textbf{取等}（$L(Z^\ast)=1$，四项各 $\frac{1/16}{1/4}=\frac14$），
所以它是一个 sharp 的充分条件，不是随手放缩。但它在整体上\textbf{不成立}：

\begin{center}
\begin{tabular}{l r r r}
\toprule
帧族 & 有效帧数 & $\sup\lambda_{\max}(B)$ & $\sup L$（$>1$ 的帧数）\\
\midrule
Haar 随机 $\mathrm{St}(5,2)$ & $400\,000$ & $1.243115922865$ & $1.103649143330$（$527$）\\
双轴（含非等分 $t$） & $187\,614$ & $1.249993448627$ & $1.117916903090$（$1056$）\\
一原子近零（近退化） & $200\,000$ & $1.199539238870$ & $0.970453199530$（$0$）\\
近等角簇 & $150\,000$ & $1.236725364068$ & $0.994022138633$（$0$）\\
\midrule
合计 & $937\,614$ & $\mathbf{1.249993448627}$（违例 $\mathbf0$） & $1.117916903090$（违例 $1583$，$0.169\%$）\\
\bottomrule
\end{tabular}
\end{center}

\begin{remark}[诚实的解读]\label{rmk:schur}
(i) \eqref{eq:schurS} 把一条矩阵不等式换成一条\textbf{标量}不等式，并在 $99.831\%$ 的帧上成立；
逻辑一致性检查（「$L\le1$ 却 $\lambda_{\max}>\frac54$」的帧数）为 $\mathbf0$，说明推导本身没有错。
(ii) 但 $\sup L=1.1179>\!1$，所以它\textbf{不是}证明。失效集中在双轴族（$0.56\%$）与 Haar 族（$0.13\%$），
近退化与近等角两族则\textbf{完全}被覆盖——这提示失效区的几何是「两条直线上各有多个原子」。
(iii) Cauchy--Schwarz 那步的等号条件是 $y_j\propto C_{ij}/(\frac14+t_j)$，比例常数依赖于 $i$，
而 $A$ 部有两个 $i$，所以 C--S \emph{不可能}对两个 $i$ 同时取等——这正是 $\sup L>1$ 却仍有
$\lambda_{\max}\le\frac54$ 的原因，也指出了改进方向：不要用标量 C--S，
而要直接处理 $2\times2$ 的 Schur 条件 $D_A-C_0D_B^{-1}C_0^\top\succeq0$（它只需要 $2$ 个行列式条件，
且等价于原命题，见引理~\ref{lem:biaxparam} 的 \eqref{eq:biaxred}）。
\end{remark}

\subsection{本轮被定量否证的候选上界}\label{sec:refuted5}

下面每一条都给出了\textbf{具体反例帧/点}与\textbf{超出量}，不是「试了没成功」。
它们的共同点是：都能在取等帧 $Z^\ast$（或 $p^\ast$）处达到 $\frac54$，
但在别处超出——即取等点是\emph{非泛型}的，任何「一致放缩」都会被别的构型顶穿。

\begin{center}\small
\begin{tabular}{p{0.30\textwidth} p{0.30\textwidth} p{0.30\textwidth}}
\toprule
候选上界 & 反例处的值 & 失效机制\\
\midrule
无权 Collatz 行和 $\max_i\sum_jB_{ij}$ & $1.366025=\frac{1+\sqrt3}2$（帧 $t=(\frac12,1,\frac16,\frac16,\frac16)$，超出 $0.1160$；该帧真值 $\frac{4+\sqrt{10}}6=1.193712943361$） & 行和把「同部元为 $0$」与对角的信息一并平均掉\\
加权 Collatz，权 $h_i=1/\sqrt{t_i}$：$\max_i(1-t_i+t_iW_i)$，$W_i=\sum_{j\in\text{对部}}(-\cos\Delta_{ij})_+$ & $2.000000$（同一帧，Parseval 缺陷恰为 $0$，超出 $0.75$；真值 $1.193712943361$）
 & 该帧的真 $\lambda_{\max}$ 反而很低，说明 Collatz 权选错；它在 $Z^\ast$ 处\emph{取等}，故不能通过换权整体修复\\
块界 $\Psi=\frac{\lambda_A+\lambda_B}2+\sqrt{\bigl(\frac{\lambda_A-\lambda_B}2\bigr)^2+\sigma_{\max}(C)^2}$ & $1.500000$（帧 $t=(0,\frac12,\frac34,\frac34,0)$，缺陷 $2.2\times10^{-16}$，超出 $0.25$；真值恰为 $1.000000$） & 丢掉对角与跨块的耦合；$\sigma_{\max}(C)^2\le(\frac14+\min_At)(\frac14+\min_Bt)$ 不成立\\
无指示器版 Schur 条件 $L^+=\sum\langle z_i,z_j\rangle^2/[(\frac14+t_i)(\frac14+t_j)]\le1$ & $1.063770$（等角五框架 $t_i\equiv\frac25$，超出 $0.063770$；同一帧带指示器版 $L=0.991445\le1$、真值 $1.160503415378$） & 丢掉 $(\cdot)_+$ 的正交结构；带指示器版 $\sup L=1.117916903$ 仍 $>1$\\
秩二块对角优占 $\mathbf{B7}(\theta)=\Lam_3(D+sH(\theta))$，$H(\theta)=\theta e_Ae_A^\top+\theta^{-1}e_Be_B^\top$，$\theta\in[\frac1{12},12]$ & $\sup_\theta=1.999999999908$；在 $p^\ast$ 处 $\min_\theta\mathbf{B7}=1.750000$（$\theta=1$），均匀点处 $=1.579796$ & 用 $H\succeq e_Ae_B^\top+e_Be_A^\top$ 的 AM--GM 丢掉太多；连取等点都盖不住\\
配套下界型 $\mathbf{B4}(\theta)=1-(\lambda_4{+}\lambda_5)(D-sH(\theta))$（Ky Fan 最小原理） & $\sup_\theta=1.586391152793$；$p^\ast$ 处 $\min_\theta=1.500000$，均匀点处 $=1.579796$ & 同上\\
秩一优占 $\mathbf{B3}=\Lam_3(\operatorname{diag}(p)+\sqrt{ab}\,gg^\top)$，$g=(e_A{+}e_B)/\sqrt2$ & $\sup=1.499999999884$，$50.54\%$ 的点 $>\frac54$ & 丢掉 $R$ 的负特征值 $-\sqrt{ab}$；与 $\mathbf{B1}$ 同值，说明负部信息在此形式下完全损失\\
Ky Fan 界 $\mathbf{B1}=1-s_2+\sqrt{ab}$ & $\sup=1.500000$，$52.56\%$ 的点 $>\frac54$ & 在 $p^\ast$ 与均匀点两处取等，故本身已是最优的 Ky Fan 型界\\
迹--二次型界 $\mathbf{B2}$ & $\sup=1.579795897113$，$76.49\%$ 的点 $>\frac54$ & 只用 $\sum\lambda_i^2$，丢掉秩二结构\\
$\mathbf{B1}$ 与 $\mathbf{B2}$ 的联合 $\min(\mathbf{B1},\mathbf{B2})$ & $260486/400000=65.1\%$ 的点上 $>\frac54$ & 两条界的失效区大面积重叠\\
双轴族上 $\varphi(t)=\lambda_{\max}(B(t))$ 的凹性 & $360\,000$ 对中 $57\,149$ 对（$15.9\%$）违反 & $\lambda_{\max}$ 对矩阵是凸的，与逐条目凹接不起来\\
\bottomrule
\end{tabular}
\end{center}

\begin{remark}[从否证清单读出的方法论结论]\label{rmk:lesson5}
$\frac54$ 这个常数由\textbf{两个}互不相容的机制同时饱和：
$p^\ast=(\frac14,\frac14,\frac14,0,\frac14)$ 处的边界退化（一个原子权为零，问题降到 $N=4$，
即定理~\ref{thm:ladder}(ii) 的 $\frac54$），以及 $Z^\ast$ 处的双轴正交配对（$z_0=z_1=-z_2=-z_4$、$z_3\perp$）。
任何只用\emph{谱}信息（$\tr M$、$\sum\lambda^2$、Ky Fan 型不等式）或只用\emph{逐条目}信息
（Collatz、行和、$\ell_1$、Schur 标量化）的界，都只能抓住其中一个机制，因而在另一个机制的构型上被顶穿。
这解释了为什么这一族方法\textbf{系统性地}失败，也说明证明必须\emph{同时}使用
「秩二 $+$ 惯性」（命题~\ref{prop:k23rank2}）与「倍角--矩参数化」\eqref{eq:moment2} 两套结构。
\end{remark}

\subsection{本轮突破：完全二部掩码的全局闭合}\label{sec:bipglobal}

第~\ref{sec:k23}--\ref{sec:refuted5} 节的全部工作都把阶数钉死在 $N=5$：要么在一个 $4$ 维单纯形上研究一个
五次方程的两个最小根（命题~\ref{prop:k23poly}），要么等价地在 $\mathrm{St}(5,2)$ 上研究一个 $5\times5$
矩阵的最大特征值（命题~\ref{prop:framered}）。本节把方向\textbf{反过来}：不固定 $N$，而是对\emph{一切}
完全二部掩码 $I+A(K_{a,b})$、\emph{一切}阶数 $N=a+b\ge4$、\emph{一切}权 $p\in\Delta_{N-1}$，
用四步初等论证直接证明 $\Lam_3\le\frac54$。四步依次是：秩二分解与 Sylvester 惯性律、
Weyl 型逐特征值下界、尾和记账、以及长期方程之根的一条核心引理。整个证明\textbf{不含搜索、
不含半定规划、不含任何逐阶算法}，也不使用第~\ref{sec:frame}--\ref{sec:schur} 节的 Grassmannian 归约，
因此它把 \eqref{eq:frame5} 从主线的关键路径上摘掉了（注~\ref{rmk:bipvsframe}）。
更重要的是，这是一条对\emph{所有} $N$ 同时成立的\textbf{全局}陈述，而不是「算到某一阶为止」的有限核验；
它与第~\ref{sec:kappa} 节 $\kappa_3\le\rho_0$ 的表述在方向上是一致的：先把一大族掩码\emph{一次性}关掉，
再数还剩几个例外。

\begin{definition}[本节记号]\label{def:bip}
设 $a,b\ge1$、$N=a+b\ge4$，$\mathcal A=\{1,\dots,a\}$、$\mathcal B=\{a+1,\dots,N\}$，
$\Sigma=I_N+A(K_{a,b})$（即 $\Sigma_{ij}=1$ 当且仅当 $i=j$ 或 $i,j$ 分属两部）。
对 $p\in\Delta_{N-1}$ 记 $x=\sqrt p$（逐元素开方）、$M=M_p=D_x\Sigma D_x$，
其特征值降序排列为 $\lambda_1\ge\lambda_2\ge\cdots\ge\lambda_N$；记 $q^\downarrow$ 为 $p$ 的降序重排，
$\alpha=\sum_{i\in\mathcal A}p_i$、$\beta=\sum_{j\in\mathcal B}p_j$（$\alpha+\beta=1$），并定义两个关键量
\begin{equation}\label{eq:sig4}
\sigma_4(p):=\sum_{j=5}^{N}q^\downarrow_j\quad\bigl(\text{最小的 }N-4\text{ 个权之和；}N=4\text{ 时取 }0\bigr),
\qquad s:=-\lambda_N\ \ge\ 0 .
\end{equation}
（$s\ge0$ 由引理~\ref{lem:bipweyl} 保证；$N=5$ 时 $\sigma_4(p)=p_{\min}$。）
因为 $\tr M=\sum_ip_i=1$，所以
\begin{equation}\label{eq:tail}
\Lam_3(M)=\lambda_1+\lambda_2+\lambda_3=1-\sum_{i\ge4}\lambda_i ,
\end{equation}
即\textbf{上界问题完全等价于尾和的下界问题}。本节的全部内容就是给 $\sum_{i\ge4}\lambda_i$ 一个足够好的下界。
\end{definition}

\begin{lemma}[秩二分解与惯性]\label{lem:biprank}
$M=\operatorname{diag}(p)+R$，其中 $R_{ij}=x_ix_j$（$i,j$ 分属两部）、否则 $R_{ij}=0$。
令 $\mathbf a:=x$、$S:=\operatorname{diag}(\underbrace{1,\dots,1}_{a},\underbrace{-1,\dots,-1}_{b})$。则
\begin{equation}\label{eq:rank2}
R=\tfrac12\bigl(\mathbf a\mathbf a^\top-(S\mathbf a)(S\mathbf a)^\top\bigr),\qquad
v^\top Rv=\tfrac12\Bigl[(\mathbf a^\top v)^2-\bigl((S\mathbf a)^\top v\bigr)^2\Bigr]\quad(\forall v\in\R^N),
\end{equation}
故 $R$ 在 $\operatorname{span}\{\mathbf a,S\mathbf a\}$ 上的符号差为 $(1,1)$、在其正交补上恒为 $0$，
即 $\operatorname{neg}(R)=1$、$\rank R\le2$。于是由 Sylvester 惯性律
\[
\operatorname{neg}(M)\ \le\ \operatorname{neg}(\operatorname{diag}p)+\operatorname{neg}(R)\ =\ 0+1\ =\ 1 .
\]
\end{lemma}

\begin{proof}
写 $\varepsilon_i=(S)_{ii}\in\{+1,-1\}$，则 $\frac12\bigl(a_ia_j-(S\mathbf a)_i(S\mathbf a)_j\bigr)
=\frac12x_ix_j(1-\varepsilon_i\varepsilon_j)$，同部时 $\varepsilon_i\varepsilon_j=1$ 给出 $0$、
跨部时给出 $x_ix_j$，与 $R$ 的定义逐条目相同，这就是 \eqref{eq:rank2} 的第一个等式；
第二个等式是把 $v^\top(\cdot)v$ 展开。由此 $R$ 的像空间就是 $\operatorname{span}\{\mathbf a,S\mathbf a\}$，
且在该二维空间上二次型是 $\frac12(u^2-w^2)$ 型（取 $u=\mathbf a^\top v$、$w=(S\mathbf a)^\top v$，
二者在 $v$ 跑遍该空间时可独立取值），符号差 $(1,1)$。
惯性律的一步：设 $Y$ 有 $\nu=\operatorname{neg}(Y)$，取 $V=(\text{$Y$ 的负惯性子空间})^\perp$，
则 $\dim V=N-\nu$ 且 $Y|_V\succeq0$，于是 $(X+Y)|_V$ 的负惯性指数 $\le X|_V$ 的 $\le\operatorname{neg}(X)$，
故 $\operatorname{neg}(X+Y)\le\operatorname{neg}(X)+\nu$。
（这正是命题~\ref{prop:k23rank2} 在 $N=5$ 时用过的一步，此处对一般 $a,b$ 重述，因为下面要反复用它。）
\end{proof}

\begin{lemma}[Weyl 型逐特征值下界与根分布]\label{lem:bipweyl}
设 $e_{\mathcal A}=x_{\mathcal A}/\sqrt\alpha$、$e_{\mathcal B}=x_{\mathcal B}/\sqrt\beta$（单位、正交，支撑不交），
$e_\pm=(e_{\mathcal A}\pm e_{\mathcal B})/\sqrt2$。则
\begin{equation}\label{eq:rank2b}
R=\sqrt{\alpha\beta}\,\bigl(e_+e_+^\top-e_-e_-^\top\bigr)\ \succeq\ -\sqrt{\alpha\beta}\;e_-e_-^\top ,
\end{equation}
即 $M\succeq\operatorname{diag}(p)-\sqrt{\alpha\beta}\,e_-e_-^\top$（一个\textbf{秩一}负项）。
由 Weyl 不等式 $\lambda_{i+j-N}(X+Y)\ge\lambda_i(X)+\lambda_j(Y)$ 取 $j=N-1$
（此时 $\lambda_{N-1}(-\sqrt{\alpha\beta}e_-e_-^\top)=0$）得
\begin{equation}\label{eq:weyl}
\lambda_i(M)\ \ge\ q^\downarrow_{i+1}\qquad(i=1,2,\dots,N-1).
\end{equation}
特别地 $\lambda_{N-1}\ge q^\downarrow_N=p_{\min}\ge0$；与引理~\ref{lem:biprank} 的 $\operatorname{neg}(M)\le1$ 合并，
\begin{equation}\label{eq:rootdist}
\lambda_1,\dots,\lambda_{N-1}\ \ge\ p_{\min}\ \ge\ 0,\qquad \lambda_N=-s,\quad s\ge0,
\end{equation}
即 $M$ 的谱里\textbf{至多一个}负数、且它一定是最小的那个。
\end{lemma}

\begin{lemma}[秩二长期方程与单调性]\label{lem:bipsec}
记 $\mathcal A^+=\{i\in\mathcal A:p_i>0\}$、$\mathcal B^+=\{j\in\mathcal B:p_j>0\}$，
\begin{equation}\label{eq:AB}
A(\lambda)=\sum_{i\in\mathcal A^+}\frac{p_i}{p_i-\lambda},\qquad
B(\lambda)=\sum_{j\in\mathcal B^+}\frac{p_j}{p_j-\lambda},\qquad h(\lambda)=A(\lambda)B(\lambda).
\end{equation}
则（作为多项式恒等式）
\begin{equation}\label{eq:secular}
\det(\lambda I-M)=\prod_{i=1}^{N}(\lambda-p_i)\cdot\bigl[1-h(\lambda)\bigr],
\end{equation}
$h$ 在 $(-\infty,0)$ 上\textbf{严格递增}，$h(-\infty)=0$，$h(0^-)=|\mathcal A^+|\cdot|\mathcal B^+|$。
因此：$s>0$ 时必有 $h(-s)=1$；且对任意 $\sigma>0$，
\begin{equation}\label{eq:mono}
h(-\sigma)\le1\quad\Longleftrightarrow\quad \sigma\ge s .
\end{equation}
\end{lemma}

\begin{proof}
把 $M$ 按两部写成 $\begin{pmatrix}D_{\mathcal A}&x_{\mathcal A}x_{\mathcal B}^\top\\
x_{\mathcal B}x_{\mathcal A}^\top&D_{\mathcal B}\end{pmatrix}$（$D_{\mathcal A}=\operatorname{diag}(p_i)_{i\in\mathcal A}$），
用 Schur 补：
\[
\det(\lambda I-M)=\det(\lambda I-D_{\mathcal A})\det(\lambda I-D_{\mathcal B})
\Bigl[1-x_{\mathcal A}^\top(\lambda I-D_{\mathcal A})^{-1}x_{\mathcal A}\cdot
x_{\mathcal B}^\top(\lambda I-D_{\mathcal B})^{-1}x_{\mathcal B}\Bigr],
\]
而 $x_{\mathcal A}^\top(\lambda I-D_{\mathcal A})^{-1}x_{\mathcal A}=\sum_{i\in\mathcal A}\frac{p_i}{\lambda-p_i}=-A(\lambda)$
（$p_i=0$ 的项恒为 $0$，故可限制在 $\mathcal A^+$ 上），$B$ 同理，即得 \eqref{eq:secular}
（先在 $\lambda\notin\{p_i\}$ 上成立，两边都是多项式，故恒等）。
单调性：$\frac{\dd}{\dd\lambda}\frac{p_i}{p_i-\lambda}=\frac{p_i}{(p_i-\lambda)^2}>0$（$p_i>0$），
故 $A,B$ 都在 $(-\infty,0)$ 上严格递增且取正值，$h=AB$ 亦然；$h(-\infty)=0$、
$h(0^-)=|\mathcal A^+||\mathcal B^+|$。最后，\eqref{eq:rootdist} 说明 $-s$ 是 $M$ 的唯一可能负特征值，
而 $\prod_i(\lambda-p_i)$ 在 $\lambda<0$ 上无根（$p_i\ge0$），故 $s>0$ 时 $-s$ 只能是 $1-h$ 的根，
即 $h(-s)=1$；严格递增立刻给出 \eqref{eq:mono}。
\end{proof}
\begin{lemma}[核心引理]\label{lem:core4}
对定义~\ref{def:bip} 中的 $M=M_p$ 与 $s=-\lambda_N$，恒有
\begin{equation}\label{eq:core}
\boxed{\ s\ \le\ \tfrac14+\sigma_4(p)\ }
\end{equation}
等号成立当且仅当 $\sigma_4(p)=0$、$p$ 的四个非零分量都等于 $\frac14$、且每部各含其中两个。
\end{lemma}

\begin{proof}
记 $m=\sigma_4(p)$、$\sigma=\frac14+m>0$、$\phi(z)=\dfrac{z}{z+\sigma}$（$z\ge0$）。
设 $S$ 为 $p$ 的\emph{四个最大}分量的指标集、$T$ 为其余 $N-4$ 个（于是 $\sum_{i\in T}p_i=m$），
并记 $r=|\mathcal A\cap S|$（$0\le r\le4$）、$x=\sum_{i\in\mathcal A\cap S}p_i$、
$t=\sum_{i\in\mathcal A\cap T}p_i$、$y=\sum_{j\in\mathcal B\cap S}p_j=1-m-x$。
目标是证 $h(-\sigma)\le1$，然后由 \eqref{eq:mono} 收口。注意
\begin{equation}\label{eq:hUV}
h(-\sigma)=A(-\sigma)B(-\sigma)=\Bigl(\sum_{i\in\mathcal A^+}\phi(p_i)\Bigr)
\Bigl(\sum_{j\in\mathcal B^+}\phi(p_j)\Bigr)=:U\cdot V .
\end{equation}

\smallskip\noindent\textbf{第一步（组内 Jensen，写成纯恒等式）。}
$\phi$ 是凹的：$\phi''(z)=-2\sigma/(z+\sigma)^3<0$。但为了不依赖凹性的一般理论，
我们把二元 Jensen 缺陷写成显式非负分子：
\begin{equation}\label{eq:Delta}
\Delta(u,v):=2\phi\Bigl(\frac{u+v}{2}\Bigr)-\phi(u)-\phi(v)
=\frac{\sigma(u-v)^2}{(u+\sigma)(v+\sigma)(u+v+2\sigma)}\ \ge\ 0 .
\end{equation}
四元情形由三个 $\Delta$ 链式给出（这是恒等式，不是估计）：
\[
4\phi(\bar z)-\sum_{i=1}^{4}\phi(z_i)=\Delta(z_1,z_2)+\Delta(z_3,z_4)
+2\,\Delta\Bigl(\frac{z_1+z_2}{2},\frac{z_3+z_4}{2}\Bigr)\ \ge\ 0,\qquad \bar z=\tfrac14\textstyle\sum z_i ;
\]
三元情形补一个均值原子 $\bar z$（补进去后四元的均值不变）再用四元式；一元平凡。
于是对任何 $|G|\in\{1,2,3,4\}$ 的组 $G$，
\begin{equation}\label{eq:gjensen}
\sum_{i\in G}\phi(p_i)\ \le\ |G|\,\phi\Bigl(\frac{\sum_{i\in G}p_i}{|G|}\Bigr)
\ =\ g_{|G|}\Bigl(\sum_{i\in G}p_i\Bigr),\qquad g_c(z):=\frac{cz}{z+c\sigma}.
\end{equation}

\smallskip\noindent\textbf{第二步（小组用 Jensen，大组用线性放缩）。}
恒等式
\begin{equation}\label{eq:philin}
\phi(z)-\frac{z}{\sigma}=-\frac{z^2}{\sigma(z+\sigma)}\ \le\ 0
\end{equation}
对\emph{任意多个}原子求和都成立，故 $T$ 中那 $N-4$ 个原子（个数可以任意大）只需
$\sum_{i\in\mathcal A\cap T}\phi(p_i)\le t/\sigma$、$\sum_{j\in\mathcal B\cap T}\phi(p_j)\le(m-t)/\sigma$；
而 $S$ 中每部的原子数 $\le4$，可用 \eqref{eq:gjensen}。合起来：
\begin{equation}\label{eq:UVbound}
U\ \le\ g_r(x)+\frac{t}{\sigma},\qquad V\ \le\ g_{4-r}(y)+\frac{m-t}{\sigma}.
\end{equation}
这里 $r=0$ 或 $4$ 时对应的 $g$ 项按 $0$ 理解。

\smallskip\noindent\textbf{第三步（Titu / Cauchy--Schwarz）。}
把 $g_c$ 改写成 $g_c(z)=c-\dfrac{c^2\sigma}{z+c\sigma}$，则
\[
g_r(x)+g_{4-r}(y)=4-\sigma\Bigl[\frac{r^2}{x+r\sigma}+\frac{(4-r)^2}{y+(4-r)\sigma}\Bigr]
\ \le\ 4-\sigma\cdot\frac{(r+4-r)^2}{x+y+4\sigma}=4-\frac{16\sigma}{x+y+4\sigma},
\]
其中不等式是 Titu 引理 $\sum\frac{u_i^2}{v_i}\ge\frac{(\sum u_i)^2}{\sum v_i}$。
五个分支 $r=0,1,2,3,4$ 的缺口分子（通分后的分子，分母在 $x,y,\sigma>0$ 上恒正）分别为
\begin{equation}\label{eq:titugap}
16x,\quad (3x-y)^2,\quad 4(x-y)^2,\quad (x-3y)^2,\quad 16y,
\end{equation}
全部非负——这五式已用 \texttt{sympy} 逐个展开并判定为完全平方或正系数多项式。
代入 $x+y=1-m$ 与 $\sigma=\frac14+m$，$x+y+4\sigma=1-m+1+4m=2+3m$，再由 \eqref{eq:UVbound}
（注意 $\frac t\sigma+\frac{m-t}{\sigma}=\frac m\sigma$）得
\begin{equation}\label{eq:UVsum}
U+V\ \le\ 4-\frac{16\sigma}{2+3m}+\frac{m}{\sigma}\ =:\ \Xi(m),
\qquad
2-\Xi(m)=\frac{2m(1+14m)}{(2+3m)(1+4m)}\ \ge\ 0 .
\end{equation}
分母 $(2+3m)(1+4m)$ 的系数全正、分子 $2m(1+14m)$ 在 $m\ge0$ 上非负，故
\begin{equation}\label{eq:UVle2}
U+V\le2,\qquad\text{等号}\iff m=\sigma_4(p)=0 .
\end{equation}

\smallskip\noindent\textbf{第四步（AM--GM 与收口）。}
$UV=\bigl(\frac{U+V}{2}\bigr)^2-\bigl(\frac{U-V}{2}\bigr)^2\le\bigl(\frac{U+V}{2}\bigr)^2\le1$，
即 $h(-\sigma)\le1$；由 \eqref{eq:mono} 得 $\sigma\ge s$，也就是 \eqref{eq:core}。

\smallskip\noindent\textbf{等号条件。}要 \eqref{eq:core} 取等，必须 $h(-\sigma)=1$，
于是 \eqref{eq:UVle2} 取等、$m=0$（$p$ 至多四个非零分量）且 $\sigma=\frac14$；
AM--GM 取等要求 $U=V$；Titu 取等要求 $\frac{x}{r}=\frac{y}{4-r}$
（$r\in\{1,2,3\}$；$r\in\{0,4\}$ 时由 \eqref{eq:titugap} 需 $x=0$ 或 $y=0$，与 $x+y=1-m=1$ 矛盾，故排除）。
记公共值为 $w$，则 $x=rw$、$y=(4-r)w$、$1=x+y=4w$，即 $w=\frac14$；
Jensen 取等要求每组内的权相等，故四个非零权全是 $\frac14$。
最后 $U=g_r(x)=\frac{r\cdot r/4}{r/4+r/4}=\frac r2$、$V=\frac{4-r}{2}$，$U=V$ 给出 $r=2$。
反向：$a,b\ge2$、每部各取两个原子置权 $\frac14$（其余为 $0$）时，
$M$ 的支撑部分是 $\frac14\bigl(I+A(C_4)\bigr)$，谱为 $\{\frac34,\frac14,\frac14,-\frac14\}$，
其余特征值全为 $0$，故 $s=\frac14=\frac14+\sigma_4$，等号成立。
\end{proof}

\begin{theorem}[完全二部掩码的全局闭合]\label{thm:bipglobal}
设 $a,b\ge1$、$N=a+b\ge4$、$\Sigma=I_N+A(K_{a,b})$、$p\in\Delta_{N-1}$。则
\begin{equation}\label{eq:bipmain}
\Lam_3\bigl(D_{\sqrt p}\,\Sigma\,D_{\sqrt p}\bigr)\ \le\ \frac54 ,
\end{equation}
且常数 $\frac54$ 与 $a,b,N,p$ \textbf{全无关}。取等当且仅当 $a,b\ge2$、$p$ 恰有四个非零分量、
每部各两个且都等于 $\frac14$。特别地
\[
F\bigl(\Sigma_{K_{a,b}}\bigr)=\frac54\quad\textbf{精确}\ (a,b\ge2),\qquad
F\bigl(\Sigma_{K_{a,b}}\bigr)\le\frac54\quad(a,b\ge1),
\]
而 $\frac54<\rho_0$（余量 $\rho_0-\frac54=0.031568197$），故\textbf{一切完全二部掩码、一切阶数}
都无条件满足主线所需的界。这是\textbf{已证}结论，不含数值成分。
\end{theorem}

\begin{proof}
由 \eqref{eq:rootdist}，$\lambda_4,\dots,\lambda_{N-1}\ge0$；由 \eqref{eq:weyl}，
$\lambda_i\ge q^\downarrow_{i+1}$（$i=4,\dots,N-1$）。于是尾和
\begin{equation}\label{eq:tailbd}
\sum_{i\ge4}\lambda_i\ \ge\ \sum_{j=5}^{N}q^\downarrow_j\ +\ \lambda_N
\ =\ \sigma_4(p)-s\ \overset{\eqref{eq:core}}{\ge}\ \sigma_4(p)-\Bigl(\frac14+\sigma_4(p)\Bigr)=-\frac14 ,
\end{equation}
代入 \eqref{eq:tail} 得 $\Lam_3\le1+\frac14=\frac54$。
取等要求 \eqref{eq:tailbd} 的两处不等号同时取等：一处是引理~\ref{lem:core4} 的等号条件
（四个非零权均为 $\frac14$、每部各两个、$\sigma_4=0$），另一处是 $\lambda_i=q^\downarrow_{i+1}$（$i\ge4$）；
后者在前者的构型下自动成立（此时 $\lambda_4=\dots=\lambda_{N-1}=0=q^\downarrow_5=\dots=q^\downarrow_N$）。
下界方向：把权集中到每部各两个原子上，得到的支撑块是 $\frac14(I+A(C_4))$，
$\Lam_3=\frac34+\frac14+\frac14=\frac54$（与定理~\ref{thm:ladder}(ii) 的 $\kappa_3(4)=\frac54$ 是同一件事），
故 $a,b\ge2$ 时 $F=\frac54$ 精确。
\end{proof}

\begin{corollary}[$N=5$ 层的无条件闭合]\label{cor:n5uncond}
$\kappa_3(5)\le\rho_0$ \textbf{无条件成立}，且 $F(\Sigma_{K_{2,3}})=\frac54$ \textbf{精确}。
\end{corollary}

\begin{proof}
$K_{2,3}$、$K_{1,4}$、$K_{2,2}$ 都是 $N=5$ 的完全二部掩码，定理~\ref{thm:bipglobal} 给出 $F\le\frac54$；
下界由「把 $p_5$ 置零则 $\Sigma_{K_{2,3}}$ 退化为 $I+A(C_4)$」得到（注~\ref{rmk:faces}），
故 $F(\Sigma_{K_{2,3}})=\frac54<\rho_0$。于是定理~\ref{thm:n5frame} 的逻辑链中
最后一环由「未证」变成「已证」：
\[
\begin{gathered}
\kappa_3(5)\le\max_{34\ \text{类}}F(\Sigma)\ \textbf{[已证]}\quad\wedge\quad
\text{其余 }33\text{ 类 }F\le\Omega\le\rho_0\ \textbf{[已证，精确有理证书]}\\
\wedge\quad F(\Sigma_{K_{2,3}})=\tfrac54<\rho_0\ \textbf{[已证]} .
\end{gathered}
\]
数值上逐类 $\max F=1.253197265=\frac{9+4\sqrt6}{15}=\kappa_3(5)$，比 $\rho_0$ 低 $2.84\times10^{-2}$。
\end{proof}

\begin{remark}[与 \eqref{eq:frame5}、$(\ast_5)$ 的关系]\label{rmk:bipvsframe}
两条路线的难度不在一个量级上，值得说清楚。
(i) \eqref{eq:frame5} 是「固定 $N=5$、把权变量换成 $\mathrm{St}(5,2)$ 上的框架、
再证一个 $5\times5$ 矩阵的最大特征值 $\le\frac54$」；它的取等点 $Z^\ast$ 是\emph{非退化}的
（谱 $\{\frac54,\frac34,\frac34,\frac14,0\}$ 五个值都不为零），第~\ref{sec:refuted5} 节的十条候选界
全部被顶穿，正说明「解释这个取等点」有多难。
(ii) 定理~\ref{thm:bipglobal} 不换变量，而是直接在 $M_p$ 上走「秩二分解 $\to$ Weyl 逐特征值下界
$\to$ 尾和记账 $\to$ 长期方程单调性」四步，把一个 $N$ 维非凸极值问题压缩成\emph{一个}标量不等式
\eqref{eq:core}，而后者只用 \eqref{eq:Delta}、\eqref{eq:philin}、Titu 与 AM--GM 四条初等不等式。
关键在于取等构型（四个 $\frac14$、其余为零）使 $m=\sigma_4=0$，
于是「$T$ 部原子」整块消失、Titu 与 AM--GM 同时取等——取等是被\emph{解释}的，不是被硬凑出来的。
这也回答了注~\ref{rmk:lesson5} 的方法论诊断：那条诊断对\emph{一致放缩}成立，
而定理~\ref{thm:bipglobal} 不是放缩，它保留了完整的秩二结构与尾和记账。
(iii) 因此 \eqref{eq:frame5} 与旧的 $(\ast_5)$ 一起\textbf{退出关键路径}、降级为独立兴趣。
请注意方向：本定理\emph{不}蕴含 \eqref{eq:frame5}（后者是关于 $B(Z)$ 而不是关于 $M_p$ 的陈述，
而且已由命题~\ref{prop:biaxsharp} 知其常数 sharp），它只是使 \eqref{eq:frame5} 变得\textbf{不再必要}。
第~\ref{sec:frame}--\ref{sec:schur} 节的内容全部保留，作为 Grassmannian 恒等式（命题~\ref{prop:grassid}）
的应用示范与同一常数的第二种独立表述；定理~\ref{thm:n5frame} 保留原状并加注为「已被超越」。
\end{remark}

\begin{lemma}[完全多部掩码的秩 $k$ 长期方程与惯性]\label{lem:mpsec}
设 $\Sigma=I_N+A(K_{n_1,\dots,n_k})$（$k$ 部、$\sum_{t}n_t=N$、各部指标集 $P_1,\dots,P_k$），
$M=D_{\sqrt p}\Sigma D_{\sqrt p}$，$x^{(t)}=x\cdot\mathbf 1_{P_t}$、$c_t=\|x^{(t)}\|=\sqrt{\alpha_t}$
（$\alpha_t=\sum_{i\in P_t}p_i$），$A_t(\lambda)=\sum_{i\in P_t,p_i>0}\frac{p_i}{p_i-\lambda}$。则
\begin{equation}\label{eq:mpsec}
\det(\lambda I-M)=\prod_{i=1}^{N}(\lambda-p_i)\cdot\prod_{t=1}^{k}\bigl(1-A_t(\lambda)\bigr)
\cdot\Bigl[\,1+\sum_{t=1}^{k}\frac{A_t(\lambda)}{1-A_t(\lambda)}\Bigr],
\end{equation}
等价地（在 $\prod_t(1-A_t)\ne0$ 处）
\begin{equation}\label{eq:mpsec2}
\sum_{t=1}^{k}\frac{1}{1-A_t(\lambda)}\ =\ k-1 ;
\end{equation}
$k=2$ 时 \eqref{eq:mpsec2} 化归为 $A_1A_2=1$，即引理~\ref{lem:bipsec}。
此外 $R=\sum_{s\ne t}x^{(s)}(x^{(t)})^\top$ 在 $\operatorname{span}\{x^{(1)},\dots,x^{(k)}\}$ 上的限制
正合同于 $D_c(J_k-I_k)D_c$（$c=(c_1,\dots,c_k)$、$J_k$ 全一矩阵），
而 $\operatorname{spec}(J_k-I_k)=\{k-1,-1\}$，正对角缩放不改变惯性，故
\begin{equation}\label{eq:mpinertia}
\operatorname{inertia}(R)=(1,\,k-1),\qquad \operatorname{neg}(M)\ \le\ k-1 .
\end{equation}
\end{lemma}

\begin{proof}
因为 $R=xx^\top-\sum_tx^{(t)}(x^{(t)})^\top=\sum_{s\ne t}x^{(s)}(x^{(t)})^\top$，
令 $F=[f_1|\cdots|f_k]$（$f_t=x^{(t)}/c_t$，列正交归一），则 $R=F\,B\,F^\top$、$B=D_c(J_k-I_k)D_c$。
秩 $k$ 更新的行列式引理给
\[
\det(\lambda I-M)=\det(\lambda I-D_p)\cdot\det\bigl(I_k-B\,F^\top(\lambda I-D_p)^{-1}F\bigr),
\]
而 $F^\top(\lambda I-D_p)^{-1}F$ 是\emph{对角}的：
$\bigl(F^\top(\lambda I-D_p)^{-1}F\bigr)_{tt}=\frac1{c_t^2}\sum_{i\in P_t}\frac{p_i}{\lambda-p_i}
=-\frac{A_t(\lambda)}{c_t^2}$（非对角元为 $0$，因 $f_s\perp f_t$）。
于是 $BF^\top(\lambda I-D_p)^{-1}F=-D_c(J_k-I_k)\operatorname{diag}(A_t/c_t)$，
其 $(s,t)$ 元（$s\ne t$）为 $-c_sA_t/c_t$、对角元为 $0$，故
\[
\det\bigl(I_k-BF^\top(\cdot)F\bigr)=\det\Bigl[\bigl(I_k-\operatorname{diag}(A_t)\bigr)+\mathbf c\,\mathbf w^\top\Bigr]
=\prod_t(1-A_t)\cdot\Bigl[1+\sum_t\frac{A_t/c_t}{1-A_t}\cdot\frac{c_t}{1}\Bigr],
\]
其中用了 $\det(D+uv^\top)=\det D\,(1+v^\top D^{-1}u)$、$\mathbf w=(A_t/c_t)$，
即得 \eqref{eq:mpsec}；再由 $\frac1{1-A_t}=1+\frac{A_t}{1-A_t}$ 求和即得 \eqref{eq:mpsec2}。
惯性部分：$R$ 在 $\operatorname{span}\{x^{(t)}\}$ 上的 Gram 表示是 $B_{st}=c_sc_t(1-\delta_{st})=D_c(J_k-I_k)D_c$，
在其正交补上为 $0$，故 $\operatorname{inertia}(R)=\operatorname{inertia}(J_k-I_k)=(1,k-1)$；
配引理~\ref{lem:biprank} 证明中的那条惯性律即得 $\operatorname{neg}(M)\le k-1$。
\end{proof}

\begin{remark}[$k\ge3$ 时的确切障碍在哪里（诚实标注）]\label{rmk:mpopen}
上面的四步论证在 $k\ge3$ 时\textbf{断在第二步}，断点非常明确：
$\operatorname{neg}(R)=k-1$ 意味着 $R$ 不再是「秩一负项 $+$ 半正项」，
Weyl 的指标平移从 $1$ 变成 $k-1$，只能得到 $\lambda_i(M)\ge q^\downarrow_{i+k-1}$（$i\le N-k+1$），
尾和界随之退化为 $\sum_{i\ge4}\lambda_i\ge\sigma_{4+k-2}(p)-$（$k-1$ 个负根之和），
而核心引理那种「一个标量 $s$」的记账不再可用。这不是技术细节，而是这条证法的\textbf{结构性}边界。
数值状态（证书 G 段，每类 $15$ 万组随机权 $+$ 多起点上升，\textbf{数值}，非证明）：
\[
\begin{array}{l|llll}
\text{掩码} & K_{1,2,2}(N{=}5) & K_{1,1,2,2}(N{=}6) & K_{2,2,2}(N{=}6) & K_{1,2,3}(N{=}6)\\
\hline
\sup\Lam_3 & 1.2530756675 & 1.2526703644 & 1.2523919861 & 1.2500767598
\end{array}
\]
全部落在 $\bigl[\frac54,\ \frac{9+4\sqrt6}{15}\bigr]$ 内且\textbf{严格低于} $\rho_0$（余量 $\ge2.85\times10^{-2}$）。
其中 $K_{1,2,2}$ 的数值上确界与 $\kappa_3(5)=\frac{9+4\sqrt6}{15}=1.2531972647$ 只差 $1.2\times10^{-4}$，
提示 $\kappa_3(5)$ 的取等类可能正是它（$C_{4v}$-Ansatz 的对称类），但这是\textbf{数值}观察，未证。
另外两条否定性信息同样值得记录：这一族上 $\max\tr\bigl((-M)_+\bigr)=0.3726336456$（$K_{2,2,3}$，$N=7$），
远大于 $\tau=0.2815681969$，所以被命题~\ref{prop:A2no} 否证掉的 A2 型负质量凸化在这里\textbf{同样无效}；
而 $\min_p\sum_{i\ge4}\lambda_i=-0.2524$（$K_{1,2,2}$），比定理~\ref{thm:bipglobal} 所需的 $-\frac14$ 只差 $2.4\times10^{-3}$，
说明 $k\ge3$ 的真值确实\emph{略高于} $\frac54$、任何「$\frac54$ 型」的界都不可能对 $k\ge3$ 成立。
\end{remark}

\begin{remark}[$\min(a,b)=1$ 的族：定理不 sharp，真值 $\frac65$]\label{rmk:k1b}
$K_{1,3}$ 配 $p=\bigl(\frac25,\frac15,\frac15,\frac15\bigr)$ 给出\textbf{精确}谱
$\bigl\{\frac45,\frac15,\frac15,-\frac15\bigr\}$（\texttt{sympy} 在 $\Q$ 上算出特征多项式
$\lambda^4-\lambda^3+\dots$ 并因式分解），于是 $\Lam_3=\frac45+\frac15+\frac15=\frac65$；
$K_{1,b}$（$b\ge3$）都能容纳这个见证，故 $F(\Sigma_{K_{1,b}})\ge\frac65$。
这与注~\ref{rmk:faces} 的「删 $A$ 部原子得 $I+A(K_{1,3})$、$F=\frac65$」一致。
上界 $\frac65$ 目前\textbf{未证}（定理~\ref{thm:bipglobal} 只给 $\frac54$，松 $0.05$）。
失败原因可以精确指出：把 $\sigma=\frac14+m$ 换成 $\sigma=\frac15+m$ 后，
第三步在 $m=0$ 处给出的界是 $4-\frac{16\sigma}{1+4\sigma}=4-\frac{16/5}{9/5}=\frac{20}9>2$，
第四步的 AM--GM 就再也压不到 $UV\le1$（在见证点上 $U+V=\frac23+\frac32=\frac{13}6\le\frac{20}9$ 恰好成立、
$UV=\frac23\cdot\frac32=1$ 恰好取等，但这只是\emph{该点}的巧合，不是\emph{一致}的不等式）。
所以 $\frac65$ 需要一条与 $\frac54$ 不同的论证；这也是为什么定理~\ref{thm:bipglobal}
的常数 $\frac54$ 是这条证法能给出的\textbf{最优}常数。
\end{remark}

\begin{proposition}[$N=6$ 层剩下的 $19$ 个未闭合类]\label{prop:tab19}
$N=6$ 的 $156$ 个同构类中，已闭合 $137$ 类：$134$ 类有 $\Omega\le\rho_0$ 的精确有理证书（定理~\ref{thm:ncert6}）、
冠军类由定理~\ref{thm:champglobal} 覆盖（$F=\Omega=\rho_0$，闭式证书）、
$K_{2,4}$ 与 $K_{3,3}$ 由定理~\ref{thm:bipglobal} 得 $F=\frac54$ \textbf{精确}
（这两类此前只有 $\Omega=1.318182$、$1.333333>\rho_0$，故一直悬着）。
$\Omega>\rho_0$ 的例外类从 $22$ 降到 $20$，其中未闭合的是下表 $19$ 类。
表内 $F$ 列为\textbf{数值}（多起点上升，$\min_ip_i\le10^{-14}$，故极大化点都在边界上），
$\Omega$ 列为已证的 SDP 上界；结构列指出该类是否为完全多部、是否含诱导 $C_5$。
\end{proposition}

\begin{center}\small
\begin{longtable}{r l l r r l c}
\toprule
\# & graph6 & 度序列 & $\Omega$ & 数值 $\max F$ & 完全多部 & 含 $C_5$\\
\midrule
\endfirsthead
\toprule
\# & graph6 & 度序列 & $\Omega$ & 数值 $\max F$ & 完全多部 & 含 $C_5$\\
\midrule
\endhead
1 & \texttt{EznW} & $4,4,4,4,4,4$ & $1.333333$ & $1.253197265$ & $K_{2,2,2}$ & -- \\
2 & \texttt{EtTg} & $3,3,3,3,3,3$ & $1.333333$ & $1.251351067$ & -- & -- \\
3 & \texttt{Ep\textasciitilde{}o} & $4,4,4,4,3,3$ & $1.310920$ & $1.250000000$ & -- & -- \\
4 & \texttt{ErXw} & $4,4,4,3,3,2$ & $1.309612$ & $1.250000000$ & -- & -- \\
5 & \texttt{EhdW} & $3,3,3,3,2,2$ & $1.308890$ & $1.247213595$ & -- & 是 \\
6 & \texttt{EzNG} & $4,4,3,3,3,3$ & $1.306671$ & $1.250000000$ & -- & -- \\
7 & \texttt{ElMg} & $4,3,3,3,3,2$ & $1.306513$ & $1.251351067$ & -- & -- \\
8 & \texttt{EheO} & $3,3,2,2,2,2$ & $1.301608$ & $1.247213595$ & -- & 是 \\
9 & \texttt{EhUg} & $3,3,3,3,2,2$ & $1.300000$ & $1.250000000$ & -- & -- \\
10 & \texttt{E\}\^{}G} & $5,4,4,3,3,3$ & $1.298975$ & $1.253197265$ & $K_{3,2,1}$ & -- \\
11 & \texttt{Ele\_} & $4,3,3,2,2,2$ & $1.295778$ & $1.251351067$ & -- & -- \\
12 & \texttt{Exe\_} & $4,4,2,2,2,2$ & $1.295208$ & $1.251351067$ & -- & -- \\
13 & \texttt{EYOw} & $3,3,3,2,2,1$ & $1.294118$ & $1.250000000$ & -- & -- \\
14 & \texttt{E\^{}NG} & $4,4,3,3,3,3$ & $1.294118$ & $1.250000000$ & -- & -- \\
15 & \texttt{EO\textasciitilde{}o} & $4,4,3,3,2,2$ & $1.294118$ & $1.250000000$ & -- & -- \\
16 & \texttt{ErW?} & $3,3,2,2,2,0$ & $1.294118$ & $1.250000000$ & -- & -- \\
17 & \texttt{EXSg} & $4,3,2,2,2,1$ & $1.294118$ & $1.250000000$ & -- & -- \\
18 & \texttt{ElfO} & $4,4,3,3,2,2$ & $1.293216$ & $1.251351067$ & -- & -- \\
19 & \texttt{El\{G} & $5,3,3,3,3,1$ & $1.288675$ & $1.216091742$ & -- & -- \\
\bottomrule
\end{longtable}
\end{center}
\begin{remark}[从表~\ref{prop:tab19} 读出的四条结构信息]\label{rmk:tab19read}
\textbf{(i) 冠军型结构不是这 $19$ 类的共同特征。}$19$ 类里只有 $2$ 类是完全多部
（$K_{2,2,2}$=\texttt{EznW}、$K_{1,2,3}$=\texttt{E\}\^{}G}），只有 $2$ 类含诱导 $C_5$
（\texttt{EhdW}、\texttt{EheO}，$\Omega$ 分别 $1.308890$、$1.301608$），其余 $15$ 类两者都不是。
所以第~\ref{sec:exh7} 节「$F=\rho_0$ 的掩码全部含诱导 $C_5$」那条观察在这里\textbf{帮不上忙}；
这 $19$ 类必须逐个（或按新的分类）处理。

\smallskip\textbf{(ii) 数值 $\max F$ 只取五个值，其中三个有闭式。}
$\frac{9+4\sqrt6}{15}=\kappa_3(5)=1.2531972647$（$2$ 类：\texttt{EznW}、\texttt{E\}\^{}G}）；
$\frac54$（$9$ 类）；$\frac{4+\sqrt5}{5}=1.2472135955$（$2$ 类：\texttt{EhdW}、\texttt{EheO}，
恰是那两个含诱导 $C_5$ 的类）。另外两个值 $1.251351067$（$5$ 类）与 $1.216091742$（$1$ 类：\texttt{El\{G}）
在 $\bigl\{(a+b\sqrt c)/d:\ 0\le b\le39,\ |a|\le119,\ 1\le d\le60,\ c\le37\text{ 无平方因子}\bigr\}$
中\textbf{没有}匹配（穷举搜索，闭式待识别）。最大值\emph{恰好}等于 $\kappa_3(5)$，
这与「$N=6$ 的数值极大化点退化到 $5$ 个有效原子」一致（$19/19$ 类的 $\min_ip_i\le10^{-14}$）。

\smallskip\textbf{(iii) 一个把 $19$ 类一次性归约掉的观察。}若某类 $\Sigma$ 的 $F$-极大化点 $p$ 有零分量，
则 $M_p$ 有全零行列，$\Lam_3(M_p)$ 等于对应 $5$ 阶主子块的值，而该主子块是某个 $5$ 阶掩码在 $\Delta_4$ 上的取值，
故 $F(\Sigma)\le\max_{\Sigma'\in\mathcal S_5}F(\Sigma')\le\rho_0$——末一步正是推论~\ref{cor:n5uncond}
（$33$ 类有理证书 $+$ $\Sigma_{K_{2,3}}$ 的 $\frac54$）。所以：
\begin{equation}\label{eq:atom6}
\text{「$\Sigma$ 的每个 $F$-极大化点都有零分量」}\quad\Longrightarrow\quad F(\Sigma)\le\rho_0 .
\end{equation}
\eqref{eq:atom6} 对全部 $19$ 类\textbf{数值成立}（$\min_ip_i\le10^{-14}$），但\textbf{未证}。
必须注意它\emph{不能}推广到全部 $6$ 阶掩码：冠军类 $\Sigma_0$ 就是反例（其 $F$-极大化点是正权的六原子团吹起、
$F(\Sigma_0)=\rho_0$，定理~\ref{thm:champglobal}），这也正是冠军需要单独一份闭式证书的原因。

\smallskip\textbf{(iv) 比 (iii) 更弱、因此更可能证得动的一条。}引理~\ref{lem:kappamax} 与推论~\ref{cor:flipsc} 给出
\[
\kappa_3(6)=\max\Bigl\{\kappa_3(5),\ \sup_{\Sigma\in\mathcal S_6^{\mathrm{int}}}F(\Sigma)\Bigr\},
\]
其中 $\mathcal S_6^{\mathrm{int}}$ 只要求「\emph{存在}某个正权自洽见证」，不要求它是 $F$-极大化点。
由推论~\ref{cor:n5uncond}，$\kappa_3(5)\le\rho_0$ 已无条件；$\Omega\le\rho_0$ 的 $134$ 类自动满足 $F\le\Omega\le\rho_0$；
冠军类 $F=\rho_0$ 已证。于是
\begin{equation}\label{eq:s6int}
\boxed{\ \kappa_3(6)\le\rho_0\ }\qquad\Longleftarrow\qquad
\text{表~\ref{prop:tab19} 的 }19\text{ 类}\ \cap\ \mathcal S_6^{\mathrm{int}}=\varnothing
\end{equation}
——即\textbf{不需要}对它们逐个求 $F$ 的上界，只需要证明「不存在正权自洽见证」。
而引理~\ref{lem:bdryrig} 已经把这件事推进了一大步：$19$ 类中\textbf{$17$ 类没有支配点}（最大度 $\le4$），
对这 $17$ 类，任何自洽见证都必须是正权的，故「$\Sigma$ 完全不是腔掩码」\emph{蕴含}「$\Sigma\notin\mathcal S_6^{\mathrm{int}}$」
（\textbf{反方向不成立}：$K_{2,2,2}$ 是腔掩码却与 $F$ 无关，见命题~\ref{prop:chamberdead}），即引理~\ref{lem:chamber} 的纯几何陈述
\begin{equation}\label{eq:chamber6}
\Sigma\ \ne\ 1\{Q\ge0\}\qquad\text{对一切 }Q\in\mathcal P_3(6),
\end{equation}
一个 $\mathrm{Gr}(3,6)$（$9$ 维）上的问题，与 $F$ 的极大化点无关。

\smallskip
\noindent\fbox{\parbox{0.95\linewidth}{\small\textbf{第 12 轮更新（否证）}：
上面这条路线已经\textbf{死掉}。命题~\ref{prop:chamberdead} 给出一个精确有理证书
$Q=\frac12(I_6-M)\in\mathcal P_3(6)$，使 $1\{Q\ge0\}=I_6+A(K_{2,2,2})$，
即 \eqref{eq:chamber6} 对 $\Sigma=I+A(K_{2,2,2})$ 为假；数值扫描进一步显示
$19$ 个未闭合类\textbf{全部}是腔掩码（符号违背量 $\le5\times10^{-19}$）。
因此 \eqref{eq:chamber6} 不能作为 \eqref{eq:s6int} 的后件使用，
必须换成第~\ref{sec:chamberdead} 小节末尾指出的\emph{耦合}条件
$\Sigma=1\{\Pi_3(D_p^{1/2}\Sigma D_p^{1/2})\ge0\}$。详见第~\ref{sec:rowvar} 节。}}
剩下 $2$ 类有支配点（\texttt{E\}\^{}G}$=K_{1,2,3}$ 与 \texttt{El\{G}），
引理~\ref{lem:bdryrig} 只能排除其\emph{边界}自洽见证，需要单独处理；
其中 $K_{1,2,3}$ 是完全多部，引理~\ref{lem:mpsec} 的秩 $k$ 长期方程正是为它准备的。
\end{remark}


\subsection{装配}\label{sec:assemble6}

\begin{lemma}[$\kappa_3(N)$ 的极大化分解]\label{lem:kappamax}
$\kappa_3(N)=\max\bigl\{\kappa_3(N-1),\ \sup\{F(\Sigma):\Sigma\in\mathcal S_N^{\mathrm{int}}\}\bigr\}$，
其中 $\mathcal S_N^{\mathrm{int}}$ 是「存在\emph{正权}自洽见证」的掩码类。
\end{lemma}

\begin{theorem}[$N=5$ 与 $N=6$ 层的当前状态]\label{thm:n6closed}
\begin{enumerate}[label=(\roman*),leftmargin=2em]
\item $\kappa_3(3)\le\rho_0$、$\kappa_3(4)\le\rho_0$：\textbf{已证}（推论~\ref{cor:k34uncond}）。
\item $\kappa_3(5)\le\rho_0$：\textbf{已证，无条件}（推论~\ref{cor:n5uncond}）。
第六版之前这里是「在 \eqref{eq:frame5} 之下已证」（定理~\ref{thm:n5frame}），
现在最后一环由定理~\ref{thm:bipglobal} 提供：$F(\Sigma_{K_{2,3}})=\frac54$ \textbf{精确}，余量 $\rho_0-\frac54=0.031568197$。
逻辑链的每一环都是已证的：
\[
\begin{gathered}
\kappa_3(5)\le\max_{34\ \text{类}}F(\Sigma)\ \textbf{[已证]}\quad\wedge\quad
\text{其余 }33\text{ 类 }F\le\Omega\le\rho_0\ \textbf{[已证，精确有理证书]}\\
\wedge\quad F(\Sigma_{K_{2,3}})=\tfrac54<\rho_0\ \textbf{[已证]} .
\end{gathered}
\]
数值状态不变：逐类数值 $\max F=1.253197265=\frac{9+4\sqrt6}{15}=\kappa_3(5)$，比 $\rho_0$ 低 $2.84\times10^{-2}$。
\item $\kappa_3(6)\le\rho_0$：\textbf{未证}，但缺口从 $21$ 类缩到 $19$ 类，而且\emph{性质}变了（见 (iv)）。
$156$ 类中已闭合 $137$ 类：$134$ 类有 $\Omega\le\rho_0$ 的精确有理证书（定理~\ref{thm:ncert6}）、
冠军类由定理~\ref{thm:champglobal} 覆盖（$F=\Omega=\rho_0$，闭式证书）、
$K_{2,4}$ 与 $K_{3,3}$ 由定理~\ref{thm:bipglobal} 得 $F=\frac54$ \textbf{精确}
（这两类此前只有 $\Omega=1.318182$、$1.333333>\rho_0$，一直悬着）。
$\Omega>\rho_0$ 的例外类 $22\to20$，未闭合的 $19$ 类见命题~\ref{prop:tab19} 的完整表
（graph6 码、度序列、$\Omega$、数值 $\max F$、是否完全多部、是否含诱导 $C_5$）。
翻转链路线已被否证（命题~\ref{prop:kappa6red}、注~\ref{rmk:flipdead}）。
这 $19$ 类的数值 $\max F$ 均 $\le\frac{9+4\sqrt6}{15}=1.253197265<\rho_0$（余量 $2.84\times10^{-2}$），
且每类的数值极大化点都在单纯形边界上（$\min_ip_i\le10^{-14}$）。
\item \textbf{主线现在只剩一个缺口，而它不再是「求上界」。}

引理~\ref{lem:kappamax} 与推论~\ref{cor:flipsc} 给出
$\kappa_3(6)=\max\{\kappa_3(5),\,\sup_{\mathcal S_6^{\mathrm{int}}}F\}$，
而 (ii) 已使 $\kappa_3(5)\le\rho_0$ 无条件成立，$134$ 类由 $\Omega\le\rho_0$ 自动合格，冠军类已证 $F=\rho_0$，故
\[
\kappa_3(6)\le\rho_0\quad\Longleftarrow\quad
\text{命题~\ref{prop:tab19} 的 }19\text{ 类}\cap\mathcal S_6^{\mathrm{int}}=\varnothing .
\]
其中 $17$ 类没有支配点，引理~\ref{lem:bdryrig} 已把这条\emph{等价}化为引理~\ref{lem:chamber} 的纯几何陈述
$\Sigma\ne1\{Q\ge0\}$（对一切 $Q\in\mathcal P_3(6)$，$\mathrm{Gr}(3,6)$ 是 $9$ 维）；
剩 $2$ 类有支配点（$K_{1,2,3}$、\texttt{El\{G}），需要单独处理，
而 $K_{1,2,3}$ 正是引理~\ref{lem:mpsec} 的秩 $k$ 长期方程的目标。
$N=5$ 层与\emph{一切}完全二部掩码（任意阶数）都已无条件关闭；
\eqref{eq:frame5} 与 $(\ast_5)$ 退为独立兴趣（注~\ref{rmk:bipvsframe}）。
\end{enumerate}
\end{theorem}

\subsection{更新后的优先队列（第六版，\textbf{已被第~\ref{sec:pq7} 小节取代}）}\label{sec:pq6}

\textbf{阅读提示（第 12 轮）}：本小节保留原文以记录思路演化，但\textbf{不再有效}。
下面的 P1 已被命题~\ref{prop:chamberdead} 否证（$19/19$ 类都是腔掩码），
P6 已被引理~\ref{lem:slope1} 吸收为与主猜想逐字等价。现行队列见第~\ref{sec:pq7} 小节。

按「期望收益 $\times$ 成功概率 $/$ 代价」排序，已同步写入 \texttt{研究方法优先级协议.md}。
本轮最大的变化是：\textbf{P1 不再是「证一条上界不等式」，而是「证一条几何不存在性陈述」}，
因为推论~\ref{cor:n5uncond} 已经把 $\kappa_3(5)$ 关掉、定理~\ref{thm:bipglobal} 已经把全部完全二部掩码关掉。

\begin{itemize}[leftmargin=1.6em,itemsep=3pt]
\item \textbf{P1（队首，\textcolor{red}{第 12 轮已否证并弹出}）—— $N=6$ 的 $17$ 个无支配点类：证 \eqref{eq:chamber6}，即 $\Sigma\ne1\{Q\ge0\}$ 对一切 $Q\in\mathcal P_3(6)$。\ \textcolor{red}{【$19/19$ 类都是腔掩码，命题~\ref{prop:chamberdead}】}}
一旦成立，引理~\ref{lem:bdryrig} 给出 $\Sigma\notin\mathcal S_6^{\mathrm{int}}$，
再由 (iv) 的归约 $\kappa_3(6)\le\rho_0$ \textbf{立即}得证。
\emph{为什么现在排第一}：(a) 它不含 $F$、不含权、不含极大化，是一个关于秩三投影正部符号模式的\emph{纯}陈述；
(b) $\mathrm{Gr}(3,6)$ 只有 $9$ 维，且 $1\{Q\ge0\}$ 只能是 $156$ 个符号模式之一，
故等价于「$17$ 个开集 $U_\Sigma=\{Q:1\{Q\ge0\}=\Sigma\}$ 全为空」；
(c) 引理~\ref{lem:chamber} 已给出等价刻画：$\Sigma$ 是某个 Parseval 框架的\emph{非锐角图}
（$\langle z_i,z_j\rangle\le0$ 当且仅当 $ij\in\Sigma$ 的非边），
于是可以用\emph{钝角/非锐角框架}的现成理论（球面码、Gale 变换、$2$-design）来排除；
(d) 具体入口：对每个 $\Sigma$，其补图 $\bar\Sigma$ 必须是某 $6$ 个向量（$\sqrt3S^2$ 上、Parseval）的\emph{非锐角图}，
而非锐角框架在 $\R^3$ 中至多 $6$ 个向量时有强分类结果——先查文献（Danzer--Grünbaum 型钝角集、
Ball 的正部符号模式），把可移植定理写成「移植候选」再动手。
\item \textbf{P2 —— $2$ 个有支配点的类：$K_{1,2,3}$（\texttt{E\}\^{}G}）与 \texttt{El\{G}。}
引理~\ref{lem:bdryrig} 对它们只能排除\emph{边界}自洽见证，内部见证未被排除。
入口：支配点 $v$（度 $5$）意味着 $\Sigma=K_1\vee\Sigma'$，于是 $M$ 有 $\begin{pmatrix}p_v&\ast\\ \ast&\ast\end{pmatrix}$ 的结构，
可以先做一步\emph{Schur 消元}把问题降到 $5$ 阶（$\Sigma'$ 是 $5$ 阶掩码，$N=5$ 层已完全闭合），
剩下的只是支配点权 $p_v$ 这一个参数；引理~\ref{lem:mpsec} 的秩 $k$ 长期方程对 $K_{1,2,3}$ 直接可用（$k=3$）。
\item \textbf{P3 —— $k\ge3$ 的尾和界（引理~\ref{lem:mpsec} 的第二步推广）。}
障碍已精确定位（注~\ref{rmk:mpopen}）：$\operatorname{neg}(R)=k-1$ 使 Weyl 指标平移 $k-1$ 位，
尾和记账里出现 $k-1$ 个负根而非一个。
入口：不要对 $k-1$ 个负根分别记账，而是用\emph{秩 $k$ 长期方程} \eqref{eq:mpsec2} 直接研究
$\sum_{t}\frac1{1-A_t}=k-1$ 的根之和（Viète 型），目标是 $\sum_{i\ge4}\lambda_i\ge-\frac14+\varepsilon_k$，
其中 $\varepsilon_k>0$ 由数值给出（$k=3$ 时 $\min\sum_{i\ge4}\lambda_i=-0.2524$，故 $\varepsilon_3$ 至多 $2.4\times10^{-3}$，很小）。
\textbf{诚实提示}：注~\ref{rmk:mpopen} 已说明 $k\ge3$ 的真值\emph{略高于} $\frac54$，
所以这一条不可能靠「$\frac54$ 型」的界完成，必须真的改进尾和。
\item \textbf{P4 —— $N\ge7$ 的有效原子化约}（注~\ref{rmk:ladderinfo}(2)、\eqref{eq:atom6} 的推广）。
数值支持很强（$N=8$ 时冠军末两位权为 $3\times10^{-6}$、$2\times10^{-7}$；$19/19$ 个 $N=6$ 类 $\min_ip_i\le10^{-14}$），
但\textbf{没有}证明思路，而且 \eqref{eq:atom6} 对冠军类是\emph{假}的，
所以正确形式只能是「若 $\Sigma$ 的极大化点是正权的，则 $\Sigma$ 必含诱导 $C_5$ 且自洽」（定理~\ref{thm:selfcons} 的方向）。
\item \textbf{P5 —— $\min(a,b)=1$ 族的真值 $\frac65$}（注~\ref{rmk:k1b}）。
下界精确（$K_{1,3}$ 的有理见证，谱 $\{\frac45,\frac15,\frac15,-\frac15\}$），上界未证；
定理~\ref{thm:bipglobal} 的证法在此松 $0.05$，且注~\ref{rmk:k1b} 已精确指出失效点
（$\sigma=\frac15$ 时第三步给 $\frac{20}9>2$，AM--GM 断掉）。副线，收益是「一族掩码的精确值」。
\item \textbf{P6 —— 联合不等式 $\alpha+\ell\le2\rho_0$}（注~\ref{rmk:rigiditylimit} 的「第四条不等式」）。
命题~\ref{prop:A2no} 之后这是唯一还活着的凸化型入口；不含秩三结构，优先级最低。
\item \textbf{P7 —— 重写定理~\ref{thm:reduce} 的规范化证明}（注~\ref{rmk:normcaution}）。纯技术性，不影响主线。
\end{itemize}

\begin{remark}[两条已经死掉、不要再试的路线]\label{rmk:deadroutes}
(i) \textbf{一致放缩型上界}：第~\ref{sec:refuted5} 节定量否证了十条，注~\ref{rmk:lesson5} 给出了机制性解释。
(ii) \textbf{翻转链}：命题~\ref{prop:kappa6red} 与注~\ref{rmk:flipdead}，$647$ 个翻转像的 $\max\Omega=1.310920>\rho_0$。
(iii) \textbf{A2 型负质量凸化}：命题~\ref{prop:A2no} 的 $7$ 原子精确反例；
而且在完全多部族上同样无效（注~\ref{rmk:mpopen}：$\max\tr((-M)_+)=0.3726>\tau=0.2816$）。
(iv) \textbf{逐阶算法核验}：定理~\ref{thm:bipglobal} 是「对一切 $N$ 同时成立」的陈述，
这正是本研究协议要求的「数学优先、算法只用于认证」；
$N\ge7$ 的符号模式数以 $\binom{N}{2}$ 的划分个数增长，逐阶路线在原则上不可能收工。
\end{remark}


\section{本轮新结果（第 12 轮）：正部的行方差机制、天花板 $L(N)$，与 $Z_5$-等变弧的完全分类}\label{sec:rowvar}

本节是第 12 轮工作的完整记录。它有一条\textbf{坏消息}（旧队首 P1 的「腔掩码不存在性」路线被一个精确有理证书否证，
第~\ref{sec:chamberdead} 小节），和一条\textbf{好消息}（我们找到了一个全新的不等式机制——「行方差机制」——
它同时把 Mohar 的 $\frac{1+\sqrt3}2$ 与 \STW\ 的 $\Phi=\frac{3+\sqrt5}4$ 重导为\emph{同一个}闭式 $L(N)$ 在
$N\to\infty$ 与 $N=6$ 处的取值，并且在冠军处\emph{四重取等}；第~\ref{sec:rowid}--\ref{sec:symlayer} 小节）。
坏消息的意义在于：它说明「纯几何的存在性/不存在性」这条路上已经没有任何可关闭的类；
好消息的意义在于：它把主线缺口从「一条越来越难的不等式」换成了「一个可以逐条加约束的有限维松弛」，
并且\textbf{把缺口定量地压缩到 $3.2\times10^{-3}$}（在 $Z_5$-对称层内）。

\subsection{旧队首 P1 被否证：$K_{2,2,2}$ 是腔掩码（精确有理证书）}\label{sec:chamberdead}

注~\ref{rmk:tab19read}(iv) 把「$\Sigma\notin\mathcal S_6^{\mathrm{int}}$」\emph{等价}化为
\eqref{eq:chamber6}（$\Sigma\ne1\{Q\ge0\}$ 对一切 $Q\in\mathcal P_3(6)$）。这个等价化是\textbf{错的}。

\begin{proposition}[腔掩码否证，已证（$\Q$ 上精确）]\label{prop:chamberdead}
设 $M$ 为 $K_6$ 中某个完美匹配的置换矩阵（三对不相邻的 $2$-循环），令
\[
Q=\tfrac12\,(I_6-M)\ \in\ \Q^{6\times6}.
\]
则 $Q^\top=Q$、$Q^2=Q$、$\tr Q=3$（故 $Q\in\mathcal P_3(6)$，特征值为 $1,1,1,0,0,0$），
$q_{ii}=\frac12>0$，匹配对上 $q_{ij}=-\frac12<0$，其余位置 $q_{ij}=0$。于是
\[
1\{Q\ge0\}=J_6-M=I_6+A(K_{2,2,2}),
\]
即 \texttt{EznW}$\,=K_{2,2,2}$ \textbf{是}腔掩码。因此 \eqref{eq:chamber6} 对 $\Sigma=I+A(K_{2,2,2})$ 为\textbf{假}。
\end{proposition}

\begin{proof}
$M$ 是对合（$M^2=I_6$）且 $\tr M=0$，特征值为 $+1$（三重）与 $-1$（三重）。
故 $Q=\frac12(I_6-M)$ 的特征值为 $\frac12(1-1)=0$（三重）与 $\frac12(1+1)=1$（三重），
$Q^2=\frac14(I-2M+M^2)=\frac14(2I-2M)=Q$，$\tr Q=3$。符号模式由构造直接读出；
$K_6$ 去掉一个完美匹配正是完全三部图 $K_{2,2,2}$。全部计算在 $\Q$ 上进行，无任何数值成分。$\square$
\end{proof}

\begin{remark}[这条否证的真实后果]\label{rmk:chamberdead}
(i) 我们把 $N=6$ 的 $19$ 个未闭合类逐个做了数值腔搜索，结果是
\textbf{$19/19$ 都是腔掩码}（每个类都找到 $Q\in\mathcal P_3(6)$ 使 $1\{Q\ge0\}=\Sigma$，
符号违背量 $\le5\times10^{-19}$，即精确到机器精度）。所以
\textbf{「证明不是腔掩码」这条几何路线关闭的类数是 $0$}，旧 P1 应当弹出。

(ii) 引理~\ref{lem:bdryrig} 本身没有错，错的是对它的读法：对无支配点的类，
「$\Sigma$ 是腔掩码」只是「$\Sigma\in\mathcal S_6^{\mathrm{int}}$」的\textbf{必要条件}，不是充分条件。
真正缺的约束是\emph{耦合}条件
\[
\Sigma=1\bigl\{\Pi_3(D_p^{1/2}\Sigma D_p^{1/2})\ge0\bigr\},
\]
即掩码必须是\textbf{它自己}（带某个正权 $p$）的前三特征投影的符号模式。
命题~\ref{prop:chamberdead} 的 $Q=\frac12(I_6-M)$ 恰好说明了二者之差：
它的 $1\{Q\ge0\}=I+A(K_{2,2,2})$，但 $Q_+=\frac12I_6$ 的 Perron 向量是退化的、
$\rho(Q_+)=\frac12$，与 $F(I+A(K_{2,2,2}))=\kappa_3(5)=\frac{9+4\sqrt6}{15}$ 毫无关系。
「是腔掩码」完全不涉及自洽性，这就是它太弱的原因。

(iii) 因此注~\ref{rmk:tab19read}(iv) 中的「等价」必须降级为「必要条件」；
\eqref{eq:s6int} 仍然是正确的蕴含（$\Leftarrow$ 方向不受影响），但它的\emph{后件}不再能用
引理~\ref{lem:chamber} 的纯几何陈述来建立。$N=6$ 的 $19$ 类回到「必须使用自洽耦合」的状态。
\end{remark}

\subsection{本节记号与三条新恒等式}\label{sec:rowid}

\textbf{记号（本节专用，与 \S\ref{sec:nogo} 的框架行记号分离）。}
设 $Q=(q_{ij})\in\mathcal P_3(N)$，即 $Q^2=Q=Q^\top$、$\tr Q=3$（故 $0\preceq Q\preceq I$），
$x\ge0$、$\|x\|=1$、$Q_+x=\rho x$（$\rho=\rho(Q_+)=\lambda_{\max}(Q_+)$）。
沿用命题~\ref{prop:perronrig} 的 $\alpha=x^\top Qx$、$\beta=x^\top Q_-x$、$d=Q_-x$、$\delta=\|d\|^2$、$\ell=x^\top|Q|x$。
新增：
\[
g_i:=q_{ii}=c_i^2\quad(\text{故 }0\le g_i\le1,\ \textstyle\sum_ig_i=3),\qquad
\pi_i:=x_i^2\quad(\textstyle\sum_i\pi_i=1),
\]
\[
O:=\{(i,j):i\ne j,\ q_{ij}<0\},\qquad A:=\{(i,j):i\ne j,\ q_{ij}>0\},
\]
\[
a_i:=\!\!\sum_{j\ne i}\!\bigl((q_{ij})_+\bigr)^2,\quad
r_i:=\!\!\sum_{j\ne i}\!\bigl((q_{ij})_-\bigr)^2,\quad
v_i:=\!\!\sum_{(i,j)\in A}\!\!\pi_j,\quad
w_i:=\!\!\sum_{(i,j)\in O}\!\!\pi_j,
\]
于是 $v_i+w_i+\pi_i=1$（$\pi_i$ 自己那一项既不属于 $A$ 也不属于 $O$）。再记
\[
s_i:=\sum_{j\ne i}(q_{ij})_+x_j\quad(\text{正部行和}),\qquad
S_O:=\sum_ir_i,\qquad W_O:=\sum_i\pi_iw_i,\qquad R:=\sum_ir_iw_i .
\]

\begin{lemma}[三条新恒等式，已证]\label{lem:rowid}
在上述记号下，对\emph{任意} $Q\in\mathcal P_3(N)$ 与\emph{任意}满足 $Q_+x=\rho x$ 的 $x\ge0$、$\|x\|=1$：
\begin{enumerate}[label=(\roman*),leftmargin=2.2em]
\item \textbf{行能量恒等式}：$\displaystyle a_i+r_i=\sum_{j\ne i}q_{ij}^2=g_i(1-g_i)$；
特别地 $\sum_{ij}q_{ij}^2=3$、$\sum_i(a_i+r_i)=3-\sigma$，其中 $\sigma:=\sum_ig_i^2$。
\item \textbf{Perron 行方程}：$s_i=(\rho-g_i)x_i$；等价地 $\rho=g_i+s_i/x_i$（$x_i>0$ 时）。
\item \textbf{Hadamard 平方与三次迹}：$Q\circ Q\succeq0$、$(Q\circ Q)\one=g$、$\one^\top(Q\circ Q)\one=3$，
故 $\lambda_{\max}(Q\circ Q)\le\max_ig_i\le1$；又 $\tr Q^3=\tr Q=3$ 给出\textbf{三角形恒等式}
\begin{equation}\label{eq:tri}
\sum_{i<j<k}q_{ij}q_{jk}q_{ki}\ =\ \tfrac16\Bigl(3-3\sigma+2\sum_ig_i^3\Bigr).
\end{equation}
\end{enumerate}
\end{lemma}

\begin{proof}
(i) $(Q^2)_{ii}=\sum_jq_{ij}^2=Q_{ii}=g_i$，而 $q_{ii}^2=g_i^2$，故 $\sum_{j\ne i}q_{ij}^2=g_i-g_i^2$；
又 $q_{ij}^2=((q_{ij})_+)^2+((q_{ij})_-)^2$ 逐项成立。对 $i$ 求和并用 $\sum g_i=3$ 得第二句。

(ii) $Q_+x=\rho x$ 的第 $i$ 行是 $(Q_+)_{ii}x_i+\sum_{j\ne i}(q_{ij})_+x_j=\rho x_i$；
因 $g_i\ge0$ 有 $(Q_+)_{ii}=g_i$，故 $s_i=(\rho-g_i)x_i$。

(iii) $Q\succeq0\Rightarrow Q\circ Q\succeq0$（Schur 乘积定理）。
$(Q\circ Q)\one$ 的第 $i$ 分量就是 $\sum_jq_{ij}^2=g_i$，即 (i)；
再与 $\one$ 作内积得 $\one^\top(Q\circ Q)\one=\sum_ig_i=3$。
非负矩阵的谱半径介于最小行和与最大行和之间，故 $\lambda_{\max}(Q\circ Q)\le\max_ig_i$。
最后，$\tr Q^3=\sum_{ijk}q_{ij}q_{jk}q_{ki}$ 按「下标重合个数」分三类：
三个都相同给 $\sum_ig_i^3$；恰有两个相同给 $3\sum_{i\ne j}g_iq_{ij}^2=3\sum_ig_i\cdot g_i(1-g_i)=3\sigma-3\sum_ig_i^3$（用 (i)）；
三个互不相同给 $6\sum_{i<j<k}q_{ij}q_{jk}q_{ki}$。三者之和 $=\tr Q^3=\tr Q=3$，整理即 \eqref{eq:tri}。$\square$
\end{proof}

\begin{remark}\label{rmk:rowidread}
(i) 引理~\ref{lem:rowid}(i) 是\textbf{秩三投影的全部「逐行」内容}：它说每一行的能量 $g_i(1-g_i)$
被正部与负部按 $a_i:r_i$ 精确瓜分。第~\ref{sec:stwanat} 节的 \STW\ 解剖只用到了它的求和版
$\sum_{ij}q_{ij}^2=3$（即 Mohar 界里的那个 $\sqrt3$）；\textbf{逐行}使用它是本轮的新想法。

(ii) (ii) 把「$\rho$ 有多大」直接翻译成「每一行的正部行和有多大」，
于是对正部行和做一次 Cauchy--Schwarz 就得到 $\rho$ 的\emph{逐行}上界——这是第~\ref{sec:relax} 小节的发动机。

(iii) \eqref{eq:tri} 是一条\textbf{尚未使用}的强约束：它说全部三角形乘积之和由对角线单独决定。
冠军处 $\sigma=1.533964$、$\sum g_i^3=0.796424$，故右端 $=-0.0015073$（一个很小的负数），
这说明钝角三角形与锐角三角形几乎完全抵消。把它翻译成符号模式的组合约束是 P1$'$ 的入口之一。
\end{remark}


\subsection{四条 Cauchy--Schwarz 不等式，与冠军的「四重取等」}\label{sec:fourcs}

\begin{proposition}[行方差不等式组，已证]\label{prop:fourcs}
记号同第~\ref{sec:rowid} 小节。则下列四条\textbf{无条件}成立（不需要自洽性）：
\begin{align}
\text{\textbf{[f]}}\quad &(\rho-g_i)^2\,\pi_i\ \le\ a_iv_i
&&\text{（每一行，正部行和的 CS）}\label{eq:f}\\
\text{\textbf{[g]}}\quad &\beta^2\ \le\ S_OW_O
&&\text{（负部的整体 CS）}\label{eq:g}\\
\text{\textbf{[h1]}}\quad &\beta\ \le\ \sum_i\sqrt{\pi_ir_iw_i}
&&\text{（负部的逐行 CS）}\label{eq:h1}\\
\text{\textbf{[h2]}}\quad &\alpha(1-\alpha)+\beta^2\ \le\ R
&&\text{（命题~\ref{prop:perronrig}(c) $+$ 逐行 CS）}\label{eq:h2}
\end{align}
\end{proposition}

\begin{proof}
\eqref{eq:f}：由引理~\ref{lem:rowid}(ii) 与 Cauchy--Schwarz，
$s_i^2=\Bigl(\sum_{(i,j)\in A}(q_{ij})_+x_j\Bigr)^2\le
\Bigl(\sum_{(i,j)\in A}\bigl((q_{ij})_+\bigr)^2\Bigr)\Bigl(\sum_{(i,j)\in A}x_j^2\Bigr)=a_iv_i$，
而 $s_i=(\rho-g_i)x_i$。

\eqref{eq:g}：$\beta=\sum_{(i,j)\in O}(-q_{ij})x_ix_j\le
\sqrt{\bigl(\sum_Oq_{ij}^2\bigr)\bigl(\sum_Ox_i^2x_j^2\bigr)}$，
其中 $\sum_Oq_{ij}^2=\sum_ir_i=S_O$（每个非对角元被有序对数一次），$\sum_O\pi_i\pi_j=\sum_i\pi_iw_i=W_O$。

\eqref{eq:h1}：$d_i=(Q_-x)_i=\sum_{(i,j)\in O}(-q_{ij})x_j\le\sqrt{r_iw_i}$，
故 $\beta=\langle x,d\rangle=\sum_ix_id_i\le\sum_i\sqrt{\pi_ir_iw_i}$。

\eqref{eq:h2}：$\delta=\|d\|^2=\sum_id_i^2\le\sum_ir_iw_i=R$，代入命题~\ref{prop:perronrig}(c) 的
$\delta=\beta^2+\alpha(1-\alpha)$。$\square$
\end{proof}

\begin{remark}[冠军是这四条的\textbf{四重取等}点]\label{rmk:quadtight}
对冠军 $Q^*\in\mathcal P_3(6)$（$K_1\vee C_5$ 核心、权 $w^*=(a^*,b^*\times5)$、定理~\ref{thm:champ}），
数值到 $10^{-16}$：
\begin{center}\small
\begin{tabular}{lrr}
\toprule
量 & 值 & 说明\\
\midrule
$g_{\rm hub},g_{\rm cone}$ & $0.331767104061,\ 0.533646579188$ & 和为 $3$\\
$\pi_{\rm hub},\pi_{\rm cone}$ & $0.197271565265,\ 0.160545686947$ & 和为 $1$\\
$a_{\rm hub},a_{\rm cone}$ & $0.221697692724,\ 0.176698138437$ & \\
$r_{\rm hub},r_{\rm cone}$ & $0,\ 0.072169769272$ & 枢纽行无钝角\\
$v_{\rm hub},v_{\rm cone}$ & $0.802728434735,\ 0.518362939159$ & \\
$w_{\rm hub},w_{\rm cone}$ & $0,\ 0.321091373894$ & \\
\midrule
\eqref{eq:f} 枢纽行 & $0.177963041865=0.177963041865$ & \textbf{取等}（残差 $2.8\times10^{-17}$）\\
\eqref{eq:f} 锥顶点行 & $0.089807129438\le0.091593766384$ & 缺口 $1.787\times10^{-3}$\\
\eqref{eq:g} & $\beta^2=0.093008492800568=S_OW_O$ & \textbf{取等}（残差 $4.2\times10^{-17}$）\\
\eqref{eq:h1} & $\beta=0.304972937817=\sum\sqrt{\pi_ir_iw_i}$ & \textbf{取等}\\
\eqref{eq:h2} & $\alpha(1-\alpha)+\beta^2=0.115865451846=R$ & \textbf{取等}\\
\bottomrule
\end{tabular}
\end{center}
取等的原因很清楚：$Z_5$-对称使枢纽行的 $5$ 个正元全相等、$x$ 在锥顶点上也全相等，
于是 \eqref{eq:f} 的 CS 自动取等；钝角图是五角星、$10$ 个钝角元全相等，
于是 \eqref{eq:g}\eqref{eq:h1}\eqref{eq:h2} 也全部取等。
\textbf{唯一不取等的是 $5$ 个锥顶点行的 \eqref{eq:f}}（缺口 $1.787\times10^{-3}$），
原因同样清楚：锥顶点行有 $1$ 个枢纽正元 $0.210570$ 与 $2$ 个近邻正元 $0.257253$，两者不等，
故 CS 的「成比例」条件失效。\textbf{这 $1.787\times10^{-3}$ 就是全部主线的来源}
（第~\ref{sec:diag} 小节把这句话精确化）。
\end{remark}

\subsection{行方差松弛 REL 与天花板 $L(N)$：Mohar 与 \STW\ 常数的统一重导}\label{sec:relax}

把命题~\ref{prop:fourcs} 与引理~\ref{lem:rowid}(i) 里的 $a_i=g_i(1-g_i)-r_i$ 合起来，
就得到一个\textbf{只含 $(g,\pi,r,w)$ 的有限维松弛}：$Q$ 本身、符号图案、秩三条件统统被丢掉，
只留下「行能量如何瓜分」与「钝角质量如何分布」。

\begin{definition}[行方差松弛 REL]\label{def:rel}
对给定的 $N$，REL$(N)$ 是如下最大化问题的最优值：在
$g\in[0,1]^N$、$\sum g_i=3$；$\pi\ge0$、$\sum\pi_i=1$；$0\le r_i\le g_i(1-g_i)$；$0\le w_i\le1-\pi_i$
上，最大化 $\rho=\alpha+\beta$，约束为 $0\le\alpha\le1$、$\beta\ge0$，以及
\begin{equation}\label{eq:rel}
(\rho-g_i)^2\pi_i\le\bigl(g_i(1-g_i)-r_i\bigr)(1-\pi_i-w_i)\ \ \forall i,\qquad
\beta^2\le S_OW_O,\qquad \beta\le\sum_i\sqrt{\pi_ir_iw_i},\qquad
\alpha(1-\alpha)+\beta^2\le R .
\end{equation}
\end{definition}

由命题~\ref{prop:fourcs} 与引理~\ref{lem:rowid}(i)，任何真实构型都给出 REL 的一个可行点，故
\begin{equation}\label{eq:relvalid}
\boxed{\ \kappa_3(N)\ \le\ \mathrm{REL}(N)\ }\qquad(\text{无条件}).
\end{equation}
内层（给定 $(g,\pi,r,w)$ 后对 $(\alpha,\beta)$ 的）最大化是\textbf{一维且闭式}的：
$\beta$ 的上限是 $\beta_{\max}=\min\{\sqrt{S_OW_O},\ \sum_i\sqrt{\pi_ir_iw_i},\ \sqrt R\}$，
且 $\rho(\beta)=\beta+\min\{1,\ \frac12(1+\sqrt{(1-4(R-\beta^2))_+}),\ \min_i(g_i+\sqrt{a_iv_i/\pi_i})-\beta\}$
在 $[0,\beta_{\max}]$ 上单调增，故内层最优值就是
\begin{equation}\label{eq:inner}
\lambda^*(g,\pi,r,w)=\min\Bigl\{\,1+\beta_{\max},\ \ \min_i\Bigl(g_i+\sqrt{\tfrac{a_i(1-\pi_i-w_i)}{\pi_i}}\Bigr)\Bigr\}
\quad(\text{当 }\beta_{\max}\text{ 使 }\alpha=1\text{ 可行时}).
\end{equation}
\textbf{冠军处 $\lambda^*=\rho_0$ 精确成立}（数值到 $10^{-16}$；因为 \eqref{eq:f} 的枢纽行取等、
且 $\beta_{\max}=\beta$）。也就是说 REL 在冠军身上\textbf{一点都不松}——松的是别处。

\begin{theorem}[均匀层的天花板 $L(N)$，已证]\label{thm:LN}
设 $Q\in\mathcal P_3(N)$、$Q_+x=\rho x$，且\textbf{双均匀}：$g_i\equiv\frac3N$、$\pi_i\equiv\frac1N$。
记
\[
\kappa_N:=\sqrt{\frac{3(N-1)}{N-3}}\qquad(N\ge4),\qquad
L(N):=\frac{2+\kappa_N}{1+\kappa_N}.
\]
则 $\rho\le L(N)$，且等号成立的必要条件是钝角度
$n^*:=N\omega^*=\dfrac{N\kappa_N}{3(1+\kappa_N)}$（其中 $\omega^*=\dfrac{\kappa_N}{3(1+\kappa_N)}$）。
具体地
\[
L(3)=1,\quad L(4)=\frac54,\quad L(5)=\frac{4+\sqrt6}5=1.289897949,\quad
L(6)=\frac{3+\sqrt5}4=\Phi=1.309016994,
\]
\[
L(7)=\frac{5+3\sqrt2}7=1.320377241,\quad
L(8)=\frac{11+\sqrt{105}}{16}=1.327934423,\quad
L(12)=\frac{5+\sqrt{33}}8=1.343070331,
\]
$L(N)$ 关于 $N$ 严格递增，且 $\displaystyle\lim_{N\to\infty}L(N)=\frac{2+\sqrt3}{1+\sqrt3}=\frac{1+\sqrt3}2=1.366025404$。
\end{theorem}

\begin{remark}[这条定理为什么值得单独立一条]\label{rmk:LNread}
(i) \textbf{Mohar 界与 \STW\ 界是同一个闭式的两端。}
$\lim_{N\to\infty}L(N)=\frac{1+\sqrt3}2$ 正是 Mohar 的常数；$L(6)=\frac{3+\sqrt5}4=\Phi$ 正是 \STW\ 的常数。
两者在我们的框架里\textbf{不是两条独立的不等式}，而是「行方差机制」在 $N=\infty$ 与 $N=6$ 处的取值。
这解释了第~\ref{sec:stwanat} 节留下的疑问：为什么 \STW\ 的 $\Phi$ 偏偏是 $\frac{3+\sqrt5}4$
——因为 $\kappa_6=\sqrt5$，而 $\sqrt5$ 来自 $3(N-1)/(N-3)$ 在 $N=6$ 处的值，即「行能量 $3$」与「$N-3=\dim\ker Q$」。

(ii) $L(4)=\frac54=\kappa_3(4)$（报告已知的精确值），而且\textbf{取等}：
$N=4$ 时 $n^*=\frac{4\cdot3}{3\cdot4}=1$ 是整数，极值构型正是定理~\ref{thm:bipglobal} 的
$K_{2,2}$ 见证 $Q=\frac34I_4+\frac14A(K_{2,2})-\frac14M=I_4-vv^\top$，$v=(\frac12,\frac12,-\frac12,-\frac12)$（谱 $\{1,1,1,0\}$，见注~\ref{rmk:n4exact} 的精确计算）。
所以这个机制\textbf{不是永远松的}：在 $N=4$ 它精确。

(iii) 均匀层只是 REL 的一个子类。我们对 REL 的\emph{自由}（$4N$ 维）上确界做了全局搜索
（差分进化，$170$--$350$ 代，Sobol 初始化，多起点），结果与 $L(N)$ 逐位吻合：
\begin{center}\small
\begin{tabular}{c|cccccc}
\toprule
$N$ & 5 & 6 & 7 & 8 & 10 & 12\\
\midrule
REL 自由搜索 & $1.2898973$ & $1.3090164$ & $1.3203769$ & $1.3279311$ & $1.3373855$ & $1.3430693$\\
$L(N)$ 闭式 & $1.2898979$ & $1.3090170$ & $1.3203772$ & $1.3279344$ & $1.3373864$ & $1.3430703$\\
差 & $-6.5\!\times\!10^{-7}$ & $-5.6\!\times\!10^{-7}$ & $-3.1\!\times\!10^{-7}$ & $-3.4\!\times\!10^{-6}$ & $-9.0\!\times\!10^{-7}$ & $-1.0\!\times\!10^{-6}$\\
\bottomrule
\end{tabular}
\end{center}
（\textbf{数值}标签：这一行的吻合度强烈提示 $\mathrm{REL}(N)=L(N)$，但我们\textbf{只对均匀子类给出了证明}。
若 $\mathrm{REL}(N)=L(N)$ 对一切 $N$ 成立，则 $\kappa_3(N)\le L(N)$ 成为无条件定理，
它在 $N\le6$ 优于或等于 \STW，在 $N\ge7$ 弱于 \STW。）
\end{remark}

\begin{proof}[定理~\ref{thm:LN} 的证明]
双均匀时令 $\varrho:=r_i$（各行的负部能量，由对称性可设相同）、$\omega:=w_i$、$c:=N\varrho$。则
$a_i=\frac3N(1-\frac3N)-\varrho$，$v_i=1-\frac1N-\omega$，
\[
S_O=c,\qquad W_O=\sum_i\pi_iw_i=\omega,\qquad R=c\omega,\qquad
\sum_i\sqrt{\pi_ir_iw_i}=N\sqrt{\tfrac1N\varrho\omega}=\sqrt{c\omega},
\]
即 $\beta_{\max}=\min\{\sqrt{c\omega},\sqrt{c\omega},\sqrt{c\omega}\}=\sqrt{c\omega}$（三个上限\emph{退化成一个}）。
于是 \eqref{eq:inner} 给出 $\rho\le\min\{A,B\}$，其中
\[
A:=\frac3N+\sqrt{\Bigl(3-\frac9N-c\Bigr)\Bigl(1-\frac1N-\omega\Bigr)},\qquad
B:=1+\sqrt{c\omega}.
\]
$A$ 关于 $c$ 严格减、$B$ 关于 $c$ 严格增（$\omega$ 固定时同理反向），故 $\min\{A,B\}$ 的上确界在
$A=B$ 处取得（边界 $c=0$ 或 $\omega=0$ 给 $B=1$，边界 $c=3-\frac9N$ 给 $A=\frac3N$，都更小）。
设 $A=B=L$，记 $K=3-\frac9N=\frac{3(N-3)}N$、$M=1-\frac1N=\frac{N-1}N$，则
\[
(K-c)(M-\omega)=\Bigl(L-\frac3N\Bigr)^2,\qquad c\omega=(L-1)^2 .
\]
把第二式代入第一式并展开，$L^2$ 相消，得
\[
\frac{4(N-3)}N=\frac{3(N-3)}N\,\omega+\frac{N-1}N\,c+\frac{2(N-3)}N\,L,
\]
其中用到 $KM+1-\frac9{N^2}=\frac{3(N-3)(N-1)+N^2-9}{N^2}=\frac{4N^2-12N}{N^2}=\frac{4(N-3)}N$。两边乘 $\frac{N}{2(N-3)}$：
\begin{equation}\label{eq:Lstar}
L=2-\frac32\,\omega-\frac{N-1}{2(N-3)}\,c .
\end{equation}
对 \eqref{eq:Lstar} 的后两项用 AM--GM，并用 $c\omega=(L-1)^2$：
\[
\frac32\omega+\frac{N-1}{2(N-3)}c\ \ge\ 2\sqrt{\frac{3(N-1)}{4(N-3)}\,c\omega}
\ =\ \kappa_N(L-1)\qquad(L\ge1),
\]
代入 \eqref{eq:Lstar} 得 $2-L\ge\kappa_N(L-1)$，即 $L\le\frac{2+\kappa_N}{1+\kappa_N}=L(N)$。
AM--GM 取等当且仅当 $\frac32\omega=\frac{N-1}{2(N-3)}c$，解出
$\omega^*=\frac{\kappa_N}{3(1+\kappa_N)}$、$c^*=\frac{(N-3)\kappa_N}{(N-1)(1+\kappa_N)}$
（直接代入验证 $c^*\omega^*=(L(N)-1)^2$ 成立，故该点确实可行）。
单调性与 $N\to\infty$ 极限由 $\kappa_N$ 关于 $N$ 严格减、$\kappa_N\to\sqrt3$ 得到；
$L(3)$ 单独说明：$N=3$ 时 $g_i\equiv1$，故 $a_i+r_i=g_i(1-g_i)=0$，$Q=I_3$，$\rho=1$。$\square$
\end{proof}


\subsection{分数度障碍：$L(N)$ 对 $N\ge5$ 不可达；整数化阶梯}\label{sec:fracdeg}

\begin{proposition}[分数度障碍与例外阶数集 $\mathcal E$，已证]\label{prop:fracdeg}
在双均匀层中 $w_i=\sum_{j\in O(i)}\pi_j=n_i/N$，其中 $n_i:=|O(i)|\in\{0,1,\dots,N-1\}$ 是钝角图的度。
定理~\ref{thm:LN} 的取等\textbf{强制}两件事：
\begin{enumerate}[label=(\roman*),leftmargin=2.6em]
\item 逐行的 Cauchy--Schwarz 与 AM--GM 全部取等，故 $r_i\equiv\varrho^*$、$w_i\equiv\omega^*$，
从而钝角图 $A_O$ 是\emph{正则}的，度数必须是
\begin{equation}\label{eq:nstar}
n^*:=N\omega^*=\frac{N\kappa_N}{3(1+\kappa_N)};
\end{equation}
\item 四条 CS 同时取等，故 $Q$ 是\emph{两值}的：
$Q=\frac3NI_N+m^+(J-I-A_O)-m^-A_O$，$m^\pm\ge0$（推导见第~\ref{sec:twovalue} 小节开头）。
\end{enumerate}
精确值：
\[
n^*(4)=1,\quad n^*(5)=2-\frac{\sqrt6}3,\quad n^*(6)=\frac{5-\sqrt5}2,\quad
n^*(7)=3-\sqrt2,\quad n^*(8)=\frac72-\frac{\sqrt{105}}6,\quad n^*(9)=\mathbf{2},
\]
\[
n^*(10)=\frac92-\frac{\sqrt{21}}2,\quad n^*(12)=\frac{11}2-\frac{\sqrt{33}}2,\quad
n^*(28)=\mathbf{6},\quad n^*(99)=\mathbf{21}.
\]
于是 $n^*$ 必须是\textbf{整数}。而 $n^*\in\Z$ 当且仅当 $\kappa_N\in\Q$，
即 $\frac{3(N-1)}{N-3}=\frac{p^2}{q^2}$（既约）。写成 $N-3=kq^2$ 后方程化为 $k(p^2-3q^2)=6$，
故 $k\in\{1,2,3,6\}$，其中 $k=2,3$ 无解（模 $3$ 矛盾），只余两支 Pell 型方程：
\begin{equation}\label{eq:pell}
\mathcal E=\underbrace{\{\,3+q^2:\ q^2-3p'^2=-2\,\}}_{\text{A 支}}\ \cup\
\underbrace{\{\,3+6q^2:\ p^2-3q^2=1\,\}}_{\text{B 支}}
=\{4,9,28,99,364,1353,5044,18819,70228,262089,\dots\}.
\end{equation}
$\mathcal E$ 是\textbf{无穷集}（两支都是标准 Pell 方程）。
\emph{结论}：对 $N\ge5$ 且 $N\notin\mathcal E$，$L(N)$ 在双均匀层内不可达；
$\mathcal E$ 中的例外阶数需要另行处理（定理~\ref{thm:unattain}）。
\end{proposition}

\begin{proof}
(i) 双均匀层里 \eqref{eq:f} 给 $\rho\le\frac3N+\sqrt{N a_iv_i}$，$a_i=\frac3N(1-\frac3N)-r_i$、$v_i=1-\frac1N-w_i$；
对 $i$ 取最小值后，用 Cauchy--Schwarz 与 AM--GM 两次：
\[
\min_i\sqrt{a_iv_i}\le\frac1N\sum_i\sqrt{a_iv_i}\le\frac1N\sqrt{\Bigl(\sum_ia_i\Bigr)\Bigl(\sum_iv_i\Bigr)}
=\frac1N\sqrt{\bigl(3-{\textstyle\frac9N}-c\bigr)N\bigl(1-{\textstyle1N}-\omega\bigr)},
\]
于是 $\rho\le\frac3N+\sqrt{(3-\frac9N-c)(1-\frac1N-\omega)}=A$，与定理~\ref{thm:LN} 的证明中的 $A$ 逐字相同
——所以\textbf{不需要}预先假设 $r_i,w_i$ 均匀。取等要求两次放缩都取等：
$\min=$ 均值 $\Rightarrow a_iv_i$ 全部相等；CS 取等 $\Rightarrow a_i\propto v_i$；两者合起来给 $r_i\equiv\varrho^*$、$w_i\equiv\omega^*$。
再由 $w_i=n_i/N$ 得 $A_O$ 正则、度 $n^*=N\omega^*$，而 $\omega^*$ 由 AM--GM 取等条件唯一确定为
$\frac{\kappa_N}{3(1+\kappa_N)}$，即 \eqref{eq:nstar}。(ii) 见第~\ref{sec:twovalue} 小节。

$\frac{3(N-1)}{N-3}=\frac{p^2}{q^2}$（$\gcd(p,q)=1$）$\Rightarrow p^2(N-3)=q^2(3N-3)=q^2(3(N-3)+6)$。
令 $N-3=kq^2$（$k\in\Z_{>0}$）得 $k(p^2-3q^2)=6$，故 $k\mid6$。
$k=2$：$p^2-3q^2=3\Rightarrow3\mid p\Rightarrow3p'^2-q^2=1\Rightarrow q^2\equiv-1\pmod3$，无解；
$k=3$：$p^2-3q^2=2\Rightarrow p^2\equiv2\pmod3$，无解。
$k=1$：$p^2-3q^2=6$，模 $3$ 给 $3\mid p$，写 $p=3p'$ 得 $q^2-3p'^2=-2$，且 $N=3+q^2$（A 支）；
$k=6$：$p^2-3q^2=1$（Pell），$N=3+6q^2$（B 支）。
两支的基本解分别是 $(q,p')=(1,1)$、$(p,q)=(2,1)$，递推 $(q,p')\mapsto(2q+3p',q+2p')$、
$(p,q)\mapsto(2p+3q,p+2q)$ 生成全部解，故 $\mathcal E$ 无穷。$n^*(4)=1$、$n^*(9)=2$、$n^*(28)=6$、$n^*(99)=21$
由 \eqref{eq:nstar} 直接代入；$N=5,6,7,8,10,12$ 处 $\kappa_N$ 无理故 $n^*$ 无理。$\square$
\end{proof}

\begin{theorem}[例外阶数上 $L(N)$ 仍不可达：强正则图障碍，已证（$9\le N\le7.1\times10^8$）]\label{thm:unattain}
设 $N\in\mathcal E$、$N\ge9$，$k:=n^*$。若 $L(N)$ 在双均匀层内可达，则钝角图 $A_O$ 必是
$k$-正则且\textbf{恰有三个不同特征值}的图，从而只有两种可能：
\begin{enumerate}[label=(\alph*),leftmargin=2.2em]
\item $A_O$ 连通，于是是强正则图 $\mathrm{SRG}(N,k,\lambda,\mu)$，$\mu\ge1$；
\item $A_O$ 不连通，于是 $A_O=mK_{k+1}$，$m(k+1)=N$。
\end{enumerate}
两种可能都被排除：
\begin{itemize}
\item[(b)] $(k+1)\mid N$ 在 $\mathcal E$ 中\textbf{只}发生于 $N=9$（$A_O=3K_3$）与 $N=28$（$A_O=4K_7$）
（以及可达的 $N=4$，$A_O=2K_2$）。此时 $A_O$ 有特征值 $-1$、重数 $N-m\ge6$，
对应的 $Q$-特征值是 $\frac3N+m^-$，它必须属于 $\{0,1\}$：
$m^-=0$ 给 $\frac3N\notin\{0,1\}$；$m^-=1-\frac3N$ 给 $N-m\ge6>3$ 个特征值等于 $1$，与 $\rank Q=3$ 矛盾。
\item[(a)] 由 $Q=\frac3NI+m^+(J-I-A_O)-m^-A_O$，在 $\one^\perp$ 上 $Q$ 的作用是 $\beta I-uA_O$，
$\beta=\frac3N-m^+$、$u=m^++m^->0$。$Q$ 的特征值只能是 $0$ 或 $1$，而 $\theta\mapsto\beta-u\theta$ 是\emph{单射}，
故 $A_O$ 在 $\one^\perp$ 上至多两个特征值，且「映射到 $1$」的那个特征值重数 $\le3$。
用强正则图的重数公式
$f,g=\frac12\bigl[(N-1)\mp\frac{2k+(N-1)(\lambda-\mu)}{\sqrt{\Delta}}\bigr]$，$\Delta=(\lambda-\mu)^2+4(k-\mu)$，
以及参数关系 $\mu(N-k-1)=k(k-\lambda-1)$，对每个 $N\in\mathcal E$ 穷举 $\lambda\in[0,k)$：
\[
\begin{array}{c|cccccccccc}
N&99&364&1353&5044&18819&70228&262089&978124&3650403&13623484\\
k&21&77&286&1066&3977&14841&55386&206702&771421&2878981\\
\hline
N&50843529&189750628&708158979&&&&&&&\\
k&10744502&40099026&149651601&&&&&&&
\end{array}
\]
\textbf{没有一个} $(\lambda,\mu)$ 使 $\Delta$ 成为完全平方（故 $f,g$ 必为无理数，不可能是重数）。
情形 (a) 因此不可能。
\end{itemize}
\emph{结论}：$L(N)$ 对 $\mathcal E$ 中一切 $9\le N\le7.1\times10^8$ 都不可达。
合并命题~\ref{prop:fracdeg}：\textbf{$L(N)$ 对一切 $5\le N\le7.1\times10^8$ 不可达，$N=4$ 是唯一取等的阶数。}
\end{theorem}

\begin{remark}[对第六版一条陈述的自我修正]\label{rmk:fracdegfix}
本条命题的第一版写的是「$n^*$ 对一切 $N\ge5$ 无理，故 $L(N)$ 对一切 $N\ge5$ 不可达」。
\textbf{前半句是错的}：$n^*(9)=2$、$n^*(28)=6$、$n^*(99)=21$ 都是整数。
错误是在证书脚本 \texttt{证书\_行方差松弛与天花板.py} 的第 (ii) 项被发现的——
把「无理」的判定从 \texttt{sympy.is\_rational}（对根式返回 \texttt{None}）改成
「$\frac{3(N-1)}{N-3}$ 是否为 $\Q$ 中的平方」的纯整数检验后，立刻扫出 $N=9$。
这是本报告「所有『对一切 $N$』型陈述都必须附带可执行检验」这条纪律的又一次兑现。

\medskip
\noindent\textbf{诊断性读法（这条为什么是本轮最重要的结构性结论）。}
它说明 $\Phi-\rho_0=0.027449$ 这个缺口\emph{不是}因为 Cauchy--Schwarz 太粗：
四条 CS 在冠军处已经同时取等（注~\ref{rmk:quadtight}），把松弛做到「双均匀层内最优」也只能到 $\Phi$。
缺口的真正来源是两条\emph{离散}约束——钝角度必须是整数、钝角图的谱必须是整数——
它们在连续松弛 REL 里完全看不见。这也解释了为什么 \STW\ 的常数偏偏是 $\Phi=L(6)$：
\STW\ 的证明同样没有用到度与谱的整性，它停在的正是行方差机制的连续天花板。
$N=4$ 之所以精确，是因为 $n^*(4)=1$ 与 $A_O=2K_2$ 同时满足了这两条离散约束（注~\ref{rmk:n4exact}）。
修正后的结论\emph{更强}：例外集 $\mathcal E$ 虽然无穷，但它稀疏到（前 $13$ 项就跨过 $7\times10^8$），
而且在 $\mathcal E$ 上还有第二道、性质完全不同的障碍（强正则图的重数整性）。
换句话说，「行方差机制在 $N\ge5$ 必然松」这个结论\textbf{成立}，但它的原因不止一个：
$N\notin\mathcal E$ 时是度的整性，$N\in\mathcal E$ 时是谱的整性。
\end{remark}

\begin{proposition}[$N=6$ 的整数化阶梯，已证（精确闭式）]\label{prop:ladder6}
在双均匀层（$g_i\equiv\frac12$、$\pi_i\equiv\frac16$）且钝角图正则、度为 $n$ 时，
REL 给出 $\rho\le\Lambda_n$，其中
\[
\Lambda_0=1,\quad \Lambda_1=\frac{13}{10},\quad \Lambda_2=\frac{4+\sqrt6}5=1.289897949,\quad
\Lambda_3=\frac7{10}+\frac{\sqrt6}5=1.189897949,\quad \Lambda_4=1,\quad \Lambda_5=\frac12 .
\]
取等要求 $\varrho^*=\frac{9}{100}$（$n=1$）、$\frac7{50}-\frac{\sqrt6}{25}$（$n=2$）、$\frac{11}{100}-\frac{\sqrt6}{25}$（$n=3$）。
故双均匀层在 $N=6$ 满足 $\rho\le\frac{13}{10}=1.300000$，比 $L(6)=\Phi=1.309017$ 改进了 $9.0\times10^{-3}$。
\end{proposition}

\begin{proof}
$\omega=n/6$ 固定后，$A=\frac12+\sqrt{(\frac32-c)\frac{5-n}6}$ 关于 $c$ 严格减、$B=1+\sqrt{cn/6}$ 严格增，
故 $\max\min\{A,B\}$ 在 $A=B$ 处取得；这是关于 $\sqrt\varrho$ 的一元二次方程，
$n=1$：$20\varrho+4\sqrt\varrho-3=0\Rightarrow\sqrt\varrho=\frac3{10}$；
$n=2$：$5\varrho+\sqrt2\sqrt\varrho-\frac12=0$；$n=3$：$5\varrho+\sqrt3\sqrt\varrho-\frac14=0$，
（一般地 $5\varrho+\sqrt{n\varrho}=(4-n)/4$，其中 $\varrho=c/6$；
第三式旧版误作 $4.5\varrho+\sqrt3\sqrt\varrho-\frac18=0$，本版已更正，
$\varrho^\ast$ 与 $\Lam_3$ 的数值原本就是对的。）逐个求解即得上述闭式。$n=0$ 给 $B=1$；$n=4,5$ 时 $A$ 与 $B$ 在可行域内不相交，
上确界在 $c\to0$（$\Lambda_4=1$）或 $A$ 的端点（$\Lambda_5=\frac12$）处取得。$\square$
\end{proof}

\subsection{两值刚性：正则图案的两值秩三投影在 $N=6$ 只有 $\frac12(I_6-M)$}\label{sec:twovalue}

若 REL 的 \eqref{eq:f} 在每一行都取等，则该行的全部\emph{锐角}元相等（因为 CS 取等要求
$(q_{ij})_+\propto x_j$，而双均匀层里 $x_j$ 是常数）；同理 \eqref{eq:g}\eqref{eq:h1} 取等要求全部\emph{钝角}元相等。
配合图案的连通性（锐角图或钝角图连通时，$q_{ij}=q_{ji}$ 把各行的常数粘在一起），
极值点必须是\textbf{两值}的：$q_{ij}=m^+$（锐）或 $-m^-$（钝）。

\begin{proposition}[两值情形的精确谱，已证（符号计算）]\label{prop:twovalue}
令 $Q=\frac12I_6+m^+(J-I-A_O)-m^-A_O$，$A_O$ 为 $6$ 点正则图案。则：
\begin{enumerate}[label=(\alph*),leftmargin=2.2em]
\item $n=1$（完美匹配 $M$）：$\mathrm{spec}(Q)=\{\frac12-m^-+4m^+,\ \frac12-m^--2m^+,\ (\frac12+m^-)^{[3]}\}$。
三重特征值 $\frac12+m^-$ 只能取 $1$（取 $0$ 会要求 $m^-=-\frac12<0$），故 $m^-=\frac12$；
余下两个必须为 $0$，故 $m^+=0$。于是\textbf{唯一}可能是 $Q=\frac12(I_6-M)$，$Q_+=\frac12I_6$，$\rho=\frac12$。
\item $n=2$（$C_6$ 或 $2C_3$）：\textbf{无解}。$C_6$ 用 Fourier 基得
$\mathrm{spec}=\{\frac12+3m^+-2m^-,\ \frac12+m^++2m^-,\ (\frac12-2m^+-m^-)^{[2]},\ (\frac12+m^-)^{[2]}\}$，
逐个讨论「三个 $1$ 三个 $0$」的所有分支均导出 $m^\pm<0$ 或矛盾；$2C_3$ 有四重特征值 $\frac12+m^-$，
重数 $4>3$，无论取 $0$ 或 $1$ 都无法凑出 $(1,1,1,0,0,0)$。
\item $n=3$（$K_{3,3}$ 或棱柱图）：\textbf{无解}。两者都含一个四重特征值（$K_{3,3}$：$\frac12-m^+$；棱柱：$\frac12-m^+$），同上排除。
\item $n=0$ 与 $n=5$：\textbf{无解}（$n=0$ 给五重特征值 $\frac12-m^+$；$n=5$ 给五重 $\frac12+m^-$）。
\end{enumerate}
因此\textbf{双均匀层内不存在达到 $\Lambda_1=\frac{13}{10}$ 或 $\Lambda_2=\frac{4+\sqrt6}5$ 的秩三投影}。
\end{proposition}

\begin{remark}[$N=4$ 处同一套计算\textbf{成功}了]\label{rmk:n4exact}
$N=4$ 的双均匀层是 $g_i\equiv\frac34$，正则图案只有完美匹配（$n=1=n^*$）。
取 $Q=\frac34I_4+\frac14A(K_{2,2})-\frac14M$（即 $Q=I_4-vv^\top$，$v=(\frac12,\frac12,-\frac12,-\frac12)$）：
\[
\mathrm{spec}(Q)=\Bigl\{\underbrace{\tfrac34+\tfrac14}_{=1}\ (\text{二重：反对配对子空间}),\ \
\tfrac34+2\cdot\tfrac14-\tfrac14=1,\ \ \tfrac34-2\cdot\tfrac14-\tfrac14=0\Bigr\}=\{1,1,1,0\},
\]
确是秩三投影；$Q_+=\frac34I_4+\frac14A(K_{2,2})$ 是正则矩阵，故 $x=\frac12\one$ 是 Perron 向量，
$\rho=\frac34+\frac14\cdot2=\frac54=L(4)$。\textbf{行方差机制在 $N=4$ 精确取到 $\kappa_3(4)=\frac54$}，
这是定理~\ref{thm:bipglobal} 的 $a=b=2$ 情形的一个\emph{独立}新证（证法完全不同：
那里用的是秩二长期方程与根分布，这里用的是行能量瓜分与 AM--GM）。
\end{remark}

\subsection{$Z_5$-对称层：八个图案中六个被直接排除，冠军图案缺口 $3.2\times10^{-3}$}\label{sec:symlayer}

冠军的对称群是 $Z_5$（枢纽 $+$ 正五边形锥）。$Z_5$-不变的钝角图案由两个数据决定：
$b\in\{0,1\}$（枢纽--锥顶点是否为钝角）与 $m\subseteq\{1,2\}$（锥顶点之间哪些圆周距离类为钝角），共 $8$ 个。
在这一层内，REL 只剩 $4$ 个自由参数（$g_{\rm hub},\pi_{\rm hub},r_{\rm hub},r_{\rm cone}$；
$g_{\rm cone}=\frac{3-g_{\rm hub}}5$、$\pi_{\rm cone}=\frac{1-\pi_{\rm hub}}5$ 由两个求和约束定死，
而 $w_{\rm hub}=5b\,\pi_{\rm cone}$、$w_{\rm cone}=b\,\pi_{\rm hub}+2|m|\,\pi_{\rm cone}$ 由图案\textbf{整性地}决定）。

\begin{center}\small
\begin{tabular}{clrl}
\toprule
$b$ & $m$ & REL 天花板 & 与 $\rho_0$ 之差\\
\midrule
0 & $\varnothing$ & $1.000000000000$ & $-0.281568$ \textbf{（排除）}\\
1 & $\varnothing$ & $1.000000000000$（$g_{\rm hub}\to1$ 的极限） & $-0.281568$ \textbf{（排除）}\\
1 & $\{1\}$ & $1.000000000000$（同上） & $-0.281568$ \textbf{（排除）}\\
1 & $\{2\}$ & $1.000000000000$（同上） & $-0.281568$ \textbf{（排除）}\\
0 & $\{1,2\}$ & $1.243046694070$ & $-0.038522$ \textbf{（排除）}\\
1 & $\{1,2\}$ & $0.500000024835$ & $-0.781568$ \textbf{（排除）}\\
0 & $\{1\}$ & $1.284780995477$ & $+0.003213$\\
0 & $\{2\}$ & $1.284780995477$ & $+0.003213$ \textbf{（冠军所在类）}\\
\bottomrule
\end{tabular}
\end{center}

（\textbf{数值}标签：每格由差分进化全局搜索得到，$1200$ 代、种群 $24$、Sobol 初始化、内层用 \eqref{eq:inner} 闭式；
$m=\{1\}$ 与 $m=\{2\}$ 给出同一数值，因为 $C_5$ 的两个圆周距离类互为同构（五边形 $\leftrightarrow$ 五角星），
而冠军的钝角图恰是五角星。$b=1$ 的三格是 $g_{\rm hub}\to1$ 处的极限值，故写成 $1.000000$。）

\begin{remark}[这张表读出的三件事]\label{rmk:symread}
(i) \textbf{$8$ 个对称图案中 $6$ 个被 REL 单独否证}（天花板严格低于 $\rho_0$，最小余量 $0.0385$）。
这是本项目第一次用\emph{同一条}解析不等式族把一批图案整类排除，而不是逐个跑 $F$ 的优化器。
特别地，$b=1$（枢纽与锥顶点成钝角）整类被排除，$m=\{1,2\}$（锥顶点之间全钝）整类被排除：
冠军的「枢纽全锐 $+$ 五角星钝」是\textbf{唯一存活}的对称图案。

(ii) 存活类的天花板 $1.284780995477$ 与 $\rho_0$ 只差 $3.213\times10^{-3}$，
其极大化点是 $(g_{\rm hub},\pi_{\rm hub},r_{\rm hub},r_{\rm cone})=(0.339344,0.200520,0,0.074883)$，
与冠军的 $(0.331767,0.197272,0,0.072170)$ \textbf{很近但不相等}。
也就是说：\textbf{主线缺口已经从 $2.7\times10^{-2}$（$\Phi-\rho_0$）压缩到 $3.2\times10^{-3}$}，
而且剩下的缺口在一个 $4$ 维、图案已定死的显式优化问题里。

(iii) 对照第~\ref{sec:conefamily} 小节（$Z_5$-锥族定理）：那里的解析计算给出\emph{同一层}的精确上确界就是 $\rho_0$。
两者之差 $3.2\times10^{-3}$ 正是 REL 丢掉的秩三信息的价值——
\textbf{这就是下一步要买回来的东西}（P1$'$）。
\end{remark}

\subsection{诊断：REL 缺的到底是什么，以及新的优先队列}\label{sec:diag}

把本轮所有上界排成一条阶梯（$N=6$）：
\begin{equation}\label{eq:ladder}
\underbrace{1.366025}_{\text{Mohar }=\lim L(N)}\ >\
\underbrace{1.309017}_{\Phi=L(6)=\text{REL 自由}}\ >\
\underbrace{1.300000}_{\frac{13}{10}=\text{REL}+\text{度的整性}}\ >\
\underbrace{1.284781}_{\text{REL}+Z_5\text{对称}+\text{整性}}\ >\
\underbrace{1.281568}_{\rho_0=\text{目标}} .
\end{equation}
REL 丢掉的约束，按「买回来的难度 $\times$ 预计收益」排序如下：
\begin{enumerate}[label=(\Alph*),leftmargin=2.2em]
\item \textbf{图案的整性与对称性}（$w_i$ 必须是某个对称 $0$--$1$ 图案的 $\pi$-子集和；
$r_i$ 与 $r_j$ 通过共享的钝角元 $q_{ij}$ 耦合，因为 $q_{ij}^2$ 同时计入 $r_i$ 与 $r_j$）。
收益：$1.309017\to1.300000$（已兑现）。继续兑现的入口是把「子集和」写成对 $(\pi,w)$ 的显式不等式。
\item \textbf{秩三}（\eqref{eq:tri} 的三角形恒等式、四点 Gram 行列式 $\det(q_{ij})_{i,j\in S}=0$（$|S|=4$）、
以及等角线在 $\R^3$ 中至多 $6$ 条这一经典事实）。收益：预计是 $1.300\to\rho_0$ 的主要部分，
因为命题~\ref{prop:twovalue} 已经显示「两值 $+$ 正则图案 $+$ 秩三」在 $N=6$ 几乎无解。
\item \textbf{Perron 向量的二阶信息}：我们只用了行方程 $s_i=(\rho-g_i)x_i$。
$Q_+$ 的\emph{全部}特征值（不只是最大的）与 $Q_+$ 的不可约性还没用。
\item \textbf{自洽耦合} $\Sigma=1\{\Pi_3(D_p^{1/2}\Sigma D_p^{1/2})\ge0\}$：REL 完全没用，
但注~\ref{rmk:chamberdead} 说明它才是 $N=6$ 那 $19$ 类的真正约束。
\end{enumerate}

\subsection{冠军的闭式导出：$Z_5$-等变秩三投影恰好是一条代数弧}\label{sec:conefamily}

到目前为止，$\rho_0=\frac{9+\sqrt5+2\sqrt{35}}{18}$ 在报告里的地位是「一个被构造出来、并用高精度数值验证过的下界见证」：
定理~\ref{thm:champ} 给出了它的构造，证书脚本核验了它的每一位数字，但\textbf{没有人解释它为什么是这几个根式}。
本小节把这件事变成一个定理。结论是：

\begin{quote}\small
在 $N=6$ 的\emph{全部} $Z_5$-等变秩三投影上（这是一个一维代数弧，不是离散集合），
$\lambda_{\max}(Q_+)$ 是一个显式的一元根式函数 $\rho(s)$，它在弧内部有唯一临界点 $s^*$，
而 $\rho(s^*)$ \textbf{精确等于} $\rho_0$。冠军不是数值搜索的偶然产物，而是一个二次方程的根。
\end{quote}

更意外的是副产品：定理~\ref{thm:zonal} 的见证——正二十面体等角线系（ETF）——
\textbf{就在同一条弧上}，位于 $s=\frac12$；它的值是 $\rho(\frac12)=\frac{5+\sqrt5+\sqrt{30}}{10}=1.271329355255$，
正是第~\ref{sec:kappa} 节数值扫描里那个始终排在冠军之后的「亚军」。
也就是说，本报告里两个最重要的六阶对象（$\Phi$ 的见证与 $\rho_0$ 的见证）是同一条代数弧上的两点，
而 $\rho_0$ 是这条弧的最高点、ETF 是弧的中点。这个统一是我们之前完全没有预料到的。

\subsubsection*{等变分解与分类}

把六个点记作 $\{h\}\cup\{1,2,3,4,5\}$，$Z_5$ 平凡地作用在 $h$ 上、以 $j\mapsto j+1$ 循环作用在锥点上。
一个 $Z_5$-等变的对称矩阵由五个实参数决定：
\[
g_h:=q_{hh},\quad g_c:=q_{jj},\quad A:=q_{hj},\quad d_1:=q_{jk}\ (|j-k|\equiv1),\quad
d_2:=q_{jk}\ (|j-k|\equiv2).
\]
$\R^6$ 分解为三个不可约子空间：平凡分量 $T=\mathrm{span}\{e_h,\ \one_c/\sqrt5\}$（维 $2$），
以及两个二维分量 $\psi_1$（特征标 $\chi_1,\chi_4$）、$\psi_2$（特征标 $\chi_2,\chi_3$）。
等变矩阵在这三个分量上分别作用为
\begin{equation}\label{eq:isoT}
Q\big|_T=\begin{pmatrix} g_h & \sqrt5\,A\\[2pt] \sqrt5\,A & g_c+2d_1+2d_2\end{pmatrix},\qquad
\mu_1=g_c+\tfrac{\sqrt5-1}2d_1-\tfrac{\sqrt5+1}2d_2,\qquad
\mu_2=g_c-\tfrac{\sqrt5+1}2d_1+\tfrac{\sqrt5-1}2d_2 .
\end{equation}

\begin{lemma}[$Z_5$-等变秩三投影的完全分类，已证]\label{lem:z5class}
$Q\in\mathcal P_3(6)$ 是 $Z_5$-等变的，当且仅当存在 $s\in[0,1]$ 与 $\varepsilon\in\{+1,-1\}$，使得
\begin{equation}\label{eq:arc}
g_h=1-s,\qquad g_c=\frac{s+2}5,\qquad A=\varepsilon\sqrt{\frac{s(1-s)}5},\qquad
\{d_1,d_2\}=\Bigl\{\frac{2s-1+\sqrt5}{10},\ \frac{2s-1-\sqrt5}{10}\Bigr\},
\end{equation}
其中 $\{d_1,d_2\}$ 是\emph{无序}对：$d_1>d_2$ 与 $d_1<d_2$ 两种分配都出现，各占一支。
每个这样的 $Q$ 的特征多项式恒等于 $L^3(L-1)^3$（对 $s$ 是恒等式，不是逐点数值事实）。
\end{lemma}

\begin{proof}
$\mathrm{spec}(Q)=\mathrm{spec}(Q|_T)\cup\{\mu_1,\mu_1\}\cup\{\mu_2,\mu_2\}$，而投影的特征值只能是 $0$ 或 $1$，
故 $\mu_1,\mu_2\in\{0,1\}$，且 $Q|_T$ 的特征值也在 $\{0,1\}$ 中。秩条件
$3=\rank(Q|_T)+2\mu_1+2\mu_2$ 只有唯一解：$\rank(Q|_T)=1$，且 $\mu_1,\mu_2$ 恰有一个为 $1$、另一个为 $0$
（$\rank(Q|_T)\le2$ 排除 $0$ 或 $2$ 之外的情形：$0+2+2=4$、$2+0+0=2$、$2+2+0=4$ 都不等于 $3$）。

\textbf{支 I}（$\mu_2=0,\mu_1=1$）。两式相减给 $\mu_1-\mu_2=\sqrt5(d_1-d_2)=1$，即 $d_1-d_2=1/\sqrt5$；
再把 $d_2=d_1-1/\sqrt5$ 代入 $\mu_1=1$：
\[
g_c+d_1\Bigl(\tfrac{\sqrt5-1}2-\tfrac{\sqrt5+1}2\Bigr)+\tfrac{\sqrt5+1}{2\sqrt5}=1
\ \Longrightarrow\ d_1=g_c-\tfrac{5-\sqrt5}{10},\quad d_2=g_c-\tfrac{5+\sqrt5}{10}.
\]
$\rank(Q|_T)=1$ 意味着 $\tr(Q|_T)=1$ 与 $\det(Q|_T)=0$。前者给 $g_h+5g_c=3$（这正是 $\tr Q=3$），
后者给 $5A^2=g_h(g_c+2d_1+2d_2)=(3-5g_c)(5g_c-2)$。令 $s:=5g_c-2$，
则 $g_c=\frac{s+2}5$、$g_h=1-s$、$5A^2=s(1-s)$，代回 $d_1,d_2$ 得 \eqref{eq:arc} 中 $d_1>d_2$ 的那一支；
$A^2\ge0$ 恰好等价于 $s\in[0,1]$。

\textbf{支 II}（$\mu_1=0,\mu_2=1$）。完全平行地得到 $d_2-d_1=1/\sqrt5$，
即 $d_1=g_c-\frac{5+\sqrt5}{10}$、$d_2=g_c-\frac{5-\sqrt5}{10}$，而 $g_h,g_c,A$ 与支 I 逐字相同。
这就是 \eqref{eq:arc} 中 $d_1<d_2$ 的那一支。

最后，特征多项式 $L^3(L-1)^3$ 的恒等式由符号计算直接验证（两支各自展开后所有 $s$ 的幂次相消），
证书脚本 \texttt{证书\_行方差松弛与天花板.py} 的第 (iv) 项给出 \texttt{sympy.factor} 的输出。$\square$
\end{proof}

\begin{remark}[两支在正部意义下是同一条弧]\label{rmk:twobranch}
支 I 满足 $d_1>0>d_2$（对一切 $s\in[0,1]$：$2s-1-\sqrt5\le2-1-\sqrt5<0$），
支 II 满足 $d_2>0>d_1$。于是 $Q_+$ 的做法在两支上都是「保留那个正的锥--锥值、把负的那个抹成 $0$」。
锥点上的置换 $j\mapsto 2j\pmod 5$（固定 $h$）把「距离 $1$」对换成「距离 $2$」对，
因此它是一个置换矩阵 $P$，使 $Q_+^{\mathrm{II}}(s)=P\,Q_+^{\mathrm{I}}(s)\,P^\top$。
\textbf{两支的 $\lambda_{\max}(Q_+)$ 逐点相等}，所以「$\lambda_{\max}(Q_+)$ 的等变弧」实际上只有一条，
参数 $s\in[0,1]$，$\varepsilon$ 与支的选取都不影响它。这条观察把后面所有计算减半。
\end{remark}

\subsubsection*{弧上的 Perron 特征值}

以下固定支 I（由注~\ref{rmk:twobranch} 无损失）。抹去 $d_2$ 之后，$Q_+$ 仍是 $Z_5$-等变的，
它在 $\psi_1,\psi_2$ 上的标量是
\begin{equation}\label{eq:psip}
\mu_1(Q_+)=g_c+\tfrac{\sqrt5-1}2d_1=\frac{7-\sqrt5+(1+\sqrt5)s}{10},\qquad
\mu_2(Q_+)=g_c-\tfrac{\sqrt5+1}2d_1=\frac{2-(\sqrt5-1)s}{10},
\end{equation}
而在平凡分量上是
\[
T_+(s)=\begin{pmatrix}1-s & \sqrt{s(1-s)}\\[2pt] \sqrt{s(1-s)} & \dfrac{3s+1+\sqrt5}{5}\end{pmatrix}
\qquad(\text{右上元}=\sqrt5A=\sqrt{s(1-s)}).
\]
由 \eqref{eq:psip}，$\mu_1(Q_+)\le\frac{7-\sqrt5+1+\sqrt5}{10}=\frac8{10}<1$ 且
\[
T_+(s)_{22}-\mu_1(Q_+)=\frac{(5-\sqrt5)s+3\sqrt5-5}{10}\ \ge\ \frac{3\sqrt5-5}{10}=0.170820>0 ,
\]
故 $\lambda_{\max}(Q_+(s))=\lambda_{\max}(T_+(s))$ 对\emph{一切} $s\in[0,1]$ 成立（不需要任何数值验证）。

\begin{theorem}[弧上的精确 Perron 谱与 $\rho_0$ 的闭式导出，已证]\label{thm:conefam}
定义 $\rho(s):=\lambda_{\max}\bigl(Q_+(s)\bigr)$，$s\in[0,1]$。则
\begin{equation}\label{eq:rhos}
\rho(s)=\frac{6+\sqrt5-2s+\sqrt{P(s)}}{10},\qquad
P(s):=21-8\sqrt5+(36+16\sqrt5)s-36s^2,
\end{equation}
并且：
\begin{enumerate}[label=(\alph*),leftmargin=2.2em]
\item $\rho'(s)=0$ 在 $[0,1]$ 内\textbf{恰有一个}解
$s^*=\dfrac{9+4\sqrt5-\sqrt{35}}{18}=0.668232895938863$；
另一个根 $\frac{9+4\sqrt5+\sqrt{35}}{18}=1.325575094$ 落在弧外。
\item $\rho(s^*)=\dfrac{9+\sqrt5+2\sqrt{35}}{18}=\rho_0$，\textbf{精确相等，无任何数值误差}。
\item 端点值 $\rho(0)=1$、$\rho(1)=\frac{4+\sqrt5}5=1.247213595$，都严格小于 $\rho_0$。
\item 故 $\displaystyle\max_{s\in[0,1]}\lambda_{\max}(Q_+(s))=\rho_0$，且最大化点唯一。
\item Perron 向量在弧的最高点处是 $x=(x_h,x_c,x_c,x_c,x_c,x_c)$，其平方权重\textbf{精确}为
\begin{equation}\label{eq:perronw}
\pi_h=x_h^2=\frac{91-25\sqrt7}{126}=0.197271565264962,\qquad
\pi_c=x_c^2=\frac{7+5\sqrt7}{126}=0.160545686947008 .
\end{equation}
\item $s=\frac12$ 处（支 II）给出 $Q=\frac12I_6+\frac1{2\sqrt5}S$，其中 $S$ 是 $6$ 阶会议矩阵
（$S^2=5I_6$，对角为 $0$、非对角为 $\pm1$），即\textbf{正二十面体等角线系}；
其值为 $\rho(\frac12)=\frac{5+\sqrt5+\sqrt{30}}{10}=1.271329355255145$。
\end{enumerate}
\end{theorem}

\begin{proof}
(a) 由 $2\times2$ 公式，$\rho(s)=\frac12\bigl(1-s+\frac{3s+1+\sqrt5}5\bigr)+\frac12\sqrt{D(s)}$，
$D(s)=\bigl(1-s-\frac{3s+1+\sqrt5}5\bigr)^2+4s(1-s)$。通分化简得
$D(s)=P(s)/25$（$25D(s)$ 展开后与 $P(s)$ 逐项相同），于是 \eqref{eq:rhos} 成立。
求导：$\rho'(s)=-\frac15+\frac{P'(s)}{20\sqrt{P(s)}}$，$P'(s)=36+16\sqrt5-72s$，故 $\rho'=0$ 等价于
$P'(s)=4\sqrt{P(s)}$；两边平方（在 $[0,1]$ 上 $P'(s)\ge36+16\sqrt5-72=35.78>0$，故平方不引入伪根）得
$P'(s)^2=16P(s)$，展开相消后是
\[
18s^2-(18+8\sqrt5)s+(7+4\sqrt5)=0,
\]
判别式 $(18+8\sqrt5)^2-72(7+4\sqrt5)=140$，故 $s=\frac{18+8\sqrt5\pm2\sqrt{35}}{36}=\frac{9+4\sqrt5\pm\sqrt{35}}{18}$，
其中只有取 $-$ 的那个落在 $[0,1]$。

(b) 在临界点上 $\sqrt{P(s^*)}=P'(s^*)/4=9+4\sqrt5-18s^*$，代回 \eqref{eq:rhos}：
\[
\rho(s^*)=\frac{6+\sqrt5-2s^*+9+4\sqrt5-18s^*}{10}=\frac{15+5\sqrt5}{10}-2s^*
=\frac{3+\sqrt5}2-\frac{9+4\sqrt5-\sqrt{35}}9
=\frac{9(3+\sqrt5)-2(9+4\sqrt5-\sqrt{35})}{18}=\frac{9+\sqrt5+2\sqrt{35}}{18}.
\]

(c) $P(0)=21-8\sqrt5=(4-\sqrt5)^2$、$P(1)=21+8\sqrt5=(4+\sqrt5)^2$，
故 $\rho(0)=\frac{6+\sqrt5+4-\sqrt5}{10}=1$、$\rho(1)=\frac{4+\sqrt5}{10}+\frac{4+\sqrt5}{10}=\frac{4+\sqrt5}5$。
两者都小于 $\rho_0=1.281568$。

(d) 由 (a)(c) 与 $\rho$ 在 $[0,1]$ 上连续、在内部只有一个临界点。

(e) $T_+(s^*)$ 的 Perron 向量满足 $(g_h-\rho)x_h+\sqrt5A\,x_c'=0$（这里 $x_c'$ 是 $\one_c/\sqrt5$ 方向的分量），
即 $x_h/x_c=\sqrt5\cdot\frac{\sqrt5A}{\rho-g_h}=\frac{5A}{\rho(s^*)-g_h}$。
用 $\rho(s^*)=\frac{3+\sqrt5}2-2s^*$、$g_h=1-s^*$、$5A=\sqrt{5s^*(1-s^*)}$ 化简得
$x_h/x_c=1.1084925$，再由 $x_h^2+5x_c^2=1$ 得 $x_c^2=\frac1{(x_h/x_c)^2+5}$；
符号计算给出 \eqref{eq:perronw} 的两个根式（证书脚本第 (v) 项逐位核验
$\frac{91-25\sqrt7}{126}$、$\frac{7+5\sqrt7}{126}$ 与数值 Perron 向量之差为 $0$）。

(f) $s=\frac12$ 代入 \eqref{eq:arc}：$g_h=g_c=\frac12$、$|A|=\frac1{2\sqrt5}$、
$\{|d_1|,|d_2|\}=\{\frac1{2\sqrt5}\}$，即所有非对角元的绝对值都是 $\frac1{2\sqrt5}$，
$Q=\frac12I_6+\frac1{2\sqrt5}S$，$S:=2\sqrt5(Q-\frac12I_6)$ 是 $\pm1$ 对称矩阵、零对角；
符号计算给 $S^2=5I_6$，故 $S$ 是 $6$ 阶会议矩阵，$Q$ 是 $6$ 点在 $\R^3$ 中的等角紧标架（正二十面体的 $6$ 条对角轴）。
$\rho(\frac12)=\frac{6+\sqrt5-1+\sqrt{30}}{10}$，其中 $P(\frac12)=21-8\sqrt5+18+8\sqrt5-9=30$。$\square$
\end{proof}

\begin{remark}[这条定理在主线上意味着什么]\label{rmk:coneread}
(i) \textbf{$\rho_0$ 从「见证」升级为「极值」。}
定理~\ref{thm:lower} 说 $\rho_0$ 可达；本小节说在整个 $Z_5$-等变层内它\emph{只有}这一个极大值点，
并且这个极大值是二次方程 $18s^2-(18+8\sqrt5)s+(7+4\sqrt5)=0$ 的根的显式函数。
根式 $\sqrt{35}$ 的来源至此清楚了：它是判别式 $140=4\cdot35$ 的平方根，
而 $140$ 来自「弧的参数化」与「$2\times2$ Perron 公式」两次代换的合成，没有任何神秘之处。

(ii) \textbf{它同时解释了 $\Phi$ 的见证与 $\rho_0$ 的见证为什么长得像。}
定理~\ref{thm:zonal} 说 $\Phi$ 是「逐点主函数 $+$ 逐通道线性计费」这一\emph{方案族}的天花板，见证是 ETF；
本小节说 ETF 是等变弧的中点，而弧的最高点比它高 $\rho_0-\frac{5+\sqrt5+\sqrt{30}}{10}=9.9\times10^{-3}$。
换句话说，\STW\ 的 $\Phi$ 与我们的 $\rho_0$ 之间的差距，在等变层内\emph{不是}两个孤立常数的差距，
而是同一条弧上「中点」与「顶点」的差距。

(iii) \textbf{它与 REL 松弛的对照是量化的。}
第~\ref{sec:symlayer} 小节的 $Z_5$-对称数值扫描给出五角星图案类的 REL 天花板 $1.284780995477$，
比 $\rho_0$ 高 $3.2128\times10^{-3}$。本小节证明真正的等变秩三最大值是 $\rho_0$。
所以\textbf{那 $3.2\times10^{-3}$ 完全来自 REL 丢掉的秩三信息}，而不是来自 Cauchy--Schwarz 的粗糙：
REL 保留的是「逐行能量 $a_i+r_i=g_i(1-g_i)$」这一\emph{二次}信息，
丢掉的是引理~\ref{lem:rowid}(iii) 的三角形恒等式 \eqref{eq:tri} 与四点 Gram 行列式。
这正是第~\ref{sec:diag} 小节排在 (B) 位的缺失成分，也是新队首 P1 的靶心。

(iv) \textbf{它给出一条可执行的改进路线。}
弧是一条\emph{一维}对象，秩三约束在它上面是恒等式；要压缩那 $3.2\times10^{-3}$，
必须允许 $Q$ 离开等变层（否则已经做完了）。
最小的下一步是把对称性从 $Z_5$ 降到「hub $+$ 锥点上的一个二面角子群」，
此时等变参数从 $1$ 维变成 $3$--$4$ 维，秩三约束（\eqref{eq:tri} 与 Gram 行列式）第一次成为\emph{非平凡}的方程，
可以期望把 REL 的天花板往下压。这是 P1 的具体形态。
\end{remark}

\subsection{垂直二标架分解：把秩三拆成秩一 $+$ 秩二}\label{sec:perp2}

冠军的 $x$ 不是 $Q$ 的特征向量（$\alpha=0.9766\ne1$），所以 $Qx$ 与 $x$ 之间有一个夹角。
把这件事做成一个分解，就能把 $\mathcal P_3$ 上的问题改写成 $\mathcal P_2$ 上的问题——后者的几何要简单得多
（$\R^3$ 中的 $2$-标架就是一个单位向量加一个正交方向）。

\begin{lemma}[垂直二标架分解，已证]\label{lem:perp2}
设 $Q\in\mathcal P_3(N)$、$x\ge0$、$\|x\|=1$、$\alpha=x^\top Qx>0$。令
\[
p:=\frac{Qx}{\sqrt\alpha},\qquad B:=Q-pp^\top,\qquad q:=p-\sqrt\alpha\,x .
\]
则：
\begin{enumerate}[label=(\roman*),leftmargin=2.2em]
\item $\|p\|=1$，且 $B\in\mathcal P_2(N)$（$B^2=B=B^\top$、$\tr B=2$、$B\succeq0$）；
\item $Bp=0$、$Bx=0$；
\item $\|q\|^2=1-\alpha$，且 $q\perp x$；
\item $Q=pp^\top+B$，从而 $\Lam_3$ 的 Ky Fan 对偶量分裂为「秩一部分 $pp^\top$」与「垂直部分 $B$」之和。
\end{enumerate}
\end{lemma}

\begin{proof}
$Q$ 是投影，故 $Qp=Q^2x/\sqrt\alpha=Qx/\sqrt\alpha=p$，$\|p\|^2=x^\top Q^2x/\alpha=\alpha/\alpha=1$。
于是 $B^2=Q^2-Qpp^\top-pp^\top Q+pp^\top pp^\top=Q-pp^\top-pp^\top+pp^\top=B$，
$\tr B=3-1=2$，$B=Q-pp^\top\succeq0$（因为 $p$ 在 $Q$ 的像空间内且 $\|p\|=1$，$pp^\top\preceq Q$）。
$Bp=Qp-p=0$；$Bx=Qx-p(p^\top x)=\sqrt\alpha p-\sqrt\alpha p=0$（$p^\top x=x^\top Qx/\sqrt\alpha=\sqrt\alpha$）。
$\|q\|^2=1-2\sqrt\alpha\,p^\top x+\alpha=1-2\alpha+\alpha=1-\alpha$，$q^\top x=p^\top x-\sqrt\alpha=0$。$\square$
\end{proof}

\begin{remark}[冠军处的分裂，精确/数值]\label{rmk:perp2read}
冠军处 $\|q\|^2=1-\alpha=0.023404740944563$，
恰好等于命题~\ref{prop:decouple} 里「投影能量损失」那一项——这不是巧合：
$1-\alpha=\|q\|^2$ 就是 $x$ 偏离 $Q$ 的像空间的平方距离。
垂直部分的能量 $E_\perp:=\tfrac12x^\top|B|x$ 满足 $x^\top B_+x=-x^\top B_-x=E_\perp$
（因为 $Bx=0$ 使 $x^\top Bx=0$），冠军处 $E_\perp=0.519536756454917$（数值）。
这个分解的价值在于：$B\in\mathcal P_2(N)$ 且 $Bx=0$，
于是 $\kappa_3$ 可以改写成
\[
\kappa_3=\sup\Bigl\{x^\top\bigl(pp^\top+B\bigr)_+x:\ B\in\mathcal P_2(N),\ Bp=Bx=0,\ \|p\|=1,\ x\ge0,\ \|x\|=1\Bigr\},
\]
其中秩三条件被\emph{完全}吸收进了 $B\in\mathcal P_2$ 与两个正交约束里。
$\mathcal P_2(N)$ 是 $\R^3$ 中的 $2$-标架空间，它的秩二结构（例如四点 Gram 行列式为零）
比秩三情形的对应约束要\textbf{显式}得多。这是 P1 路线的第二种可能形态。
\end{remark}

\subsection{$\Phi(Q)=\lambda_{\max}(Q_+)$ 是凸的：谱盒归约定理}\label{sec:convex}

这是本轮一个短小但\textbf{方向性很强}的定理。它说明「秩三」这个离散约束在求上确界时可以\emph{无损}丢掉，
同时也说明为什么丢掉之后仍然做不出 $\rho_0$。

\begin{theorem}[凸性与谱盒归约，已证]\label{thm:convex}
记 $\mathcal S_N:=\{Q\in\R^{N\times N}_{\rm sym}:0\preceq Q\preceq I,\ \tr Q=3\}$（谱盒），
$\Phi(Q):=\lambda_{\max}(Q_+)$。则
\begin{enumerate}[label=(\alph*),leftmargin=2.2em]
\item $\Phi$ 在\emph{全体}对称矩阵上是凸函数；
\item $\operatorname{conv}\bigl(\mathcal P_3(N)\bigr)=\mathcal S_N$，且 $\mathcal S_N$ 的极点恰是 $\mathcal P_3(N)$；
\item 因此 $\kappa_3(N)=\max_{Q\in\mathcal P_3(N)}\Phi(Q)=\max_{Q\in\mathcal S_N}\Phi(Q)$。
\end{enumerate}
\end{theorem}

\begin{proof}
(a) 因为 $Q_+\ge0$ 逐点非负，对任意单位向量 $y$ 有 $y^\top Q_+y\le|y|^\top Q_+|y|$，故
\[
\Phi(Q)=\max_{\substack{y\ge0\\ \|y\|=1}}y^\top Q_+y=\max_{\substack{y\ge0\\ \|y\|=1}}\sum_{i,j}(q_{ij})_+\,y_iy_j .
\]
固定 $y\ge0$ 时，每一项 $(q_{ij})_+y_iy_j$ 是 $q_{ij}$ 的函数 $t\mapsto c\,(t)_+$ 与 $c=y_iy_j\ge0$ 的复合，
而 $c\,(t)_+$ 在 $c\ge0$ 时是凸函数（它是 $\max(ct,0)$）。有限个凸函数之和仍凸，
故内层函数对 $Q$ 凸；$\Phi$ 是一族凸函数的逐点上确界，因此凸。

(b) $\mathcal P_3(N)$ 是对角阵 $\mathrm{diag}(1,1,1,0,\dots,0)$ 的\emph{酉轨道}。
酉轨道的凸包由 Schur--Horn 定理刻画为 $\{Q:\mathrm{spec}(Q)\prec(1,1,1,0,\dots,0)\}$，
而这正是 $\mathcal S_N$（$0\preceq Q\preceq I$ 且 $\tr Q=3$ 等价于特征值落在 $[0,1]$ 且和为 $3$，
即被 $(1,1,1,0,\dots)$ 优超）。凸紧集的极点：$Q\in\mathcal S_N$ 是极点当且仅当它不能写成两个不同点的中点，
这在谱语言里等价于每个特征值都是 $0$ 或 $1$，即 $Q\in\mathcal P_3(N)$。

(c) 凸函数在凸紧集上的最大值在极点取得，故 $\max_{\mathcal S_N}\Phi=\max_{\mathcal P_3(N)}\Phi$。$\square$
\end{proof}

\begin{remark}[这条定理的两副面孔]\label{rmk:convexread}
(i) \textbf{好的一面}：$\kappa_3(N)$ 可以在谱盒 $\mathcal S_N$ 上计算，秩三约束是\emph{免费}的。
这为「用半定/矩方法给 $\kappa_3$ 配上界」提供了合法性：任何在 $\mathcal S_N$ 上成立的 $\Phi$ 的上界，
自动就是 $\kappa_3(N)$ 的上界，不需要额外论证极值点是投影。
第~\ref{sec:relax} 小节的 REL 松弛与第~\ref{sec:conefamily} 小节的等变弧都落在这个框架里。

(ii) \textbf{坏的一面}：$\Phi$ 是凸的，所以「在 $\mathcal S_N$ 上最大化 $\Phi$」是一个\emph{凸最大化}问题，
不是半定规划；任何把它变成可解问题的做法都必须用一个\emph{凹的或线性的}函数从上方控制 $\Phi$，
而这一步必然损失。两个已经量化过的例子：
\begin{itemize}
\item $\Phi(Q)\le\frac12\bigl(1+\lambda_{\max}(|Q|)\bigr)$（因为 $Q_+=\frac12(Q+|Q|)$）。
右端在 $\mathcal S_N$ 上可以取到 $\Phi=\frac{3+\sqrt5}4$——见证正是 ETF（$\lambda_{\max}(|Q|)=\varphi$）；
再用 $\lambda_{\max}(|Q|)\le\||Q|\|_F=\|Q\|_F=\sqrt3$ 就退化成 Mohar 的 $\frac{1+\sqrt3}2$。
\item 完全丢掉自洽性 $A=1\{Q>0\}$ 之后，命题~\ref{prop:octnogo} 给出八面体测试集上的值 $4/3>\rho_0$。
\end{itemize}
所以凸化路线的\textbf{最好成绩恰好是 \STW\ 的 $\Phi$}——这与定理~\ref{thm:zonal} 的结论完全一致，
而且是同一个现象的两个视角：$\Phi$ 是「线性/凹控制」的天花板，$\rho_0$ 需要非凸信息。
\end{remark}

\subsection{余秩一层：$\kappa_2(3)=\frac87$ 与 $\kappa_3(4)=\frac54$ 的自包含证明}\label{sec:corank1}

$N$ 阶的秩 $N-1$ 投影恰好是 $\{I_N-vv^\top:\|v\|=1\}$，即「余秩一层」。
定理~\ref{thm:convex} 在这一层把 $\kappa_{N-1}(N)$ 化成一个\textbf{球面上的}极值问题，
而球面问题的符号分类是可以穷尽的。这一小节给出两个精确常数。

\begin{proposition}[余秩一层的闭式，已证]\label{prop:corank1}
设 $v\in\R^N$、$\|v\|=1$、$Q=I_N-vv^\top$（于是 $Q\in\mathcal P_{N-1}(N)$，即 $Q$ 在余秩一层）。
\begin{enumerate}[label=(\alph*),leftmargin=2.2em]
\item 若 $v$ 的非零坐标有 $m$ 个正、$1$ 个负（$N\ge m+1$），记 $v=(c\one_m,-d,0,\dots,0)$、$mc^2+d^2=1$。
此时 $Q_+$ 在剩下的 $N-m-1$ 个坐标方向上是孤立的对角元 $1$，而在前 $m+1$ 个方向上给出
\[
\lambda_{\max}(Q_+)=\max\Bigl\{1,\ \tfrac12\bigl(1+(m-1)c^2\bigr)+\tfrac12\sqrt{(d^2-c^2)^2+4mc^2d^2}\Bigr\},
\]
其中第二项在 $c^2=\frac2{3m+1}$、$d^2=\frac{m+1}{3m+1}$ 处取最大值
\begin{equation}\label{eq:fm1}
f_{m,1}=\frac{4m}{3m+1}\qquad(m\ge2\ \text{时}\ f_{m,1}>1).
\end{equation}
\item 若 $v$ 有 $m$ 个正、$m$ 个负且两组内部分别相等（$N=2m$，$v=(c\one_m,-c\one_m)$、$c^2=\frac1{2m}$），则
\begin{equation}\label{eq:fmm}
\lambda_{\max}(Q_+)=f_{m,m}=1+\frac{m-1}{2m}=\frac{3m-1}{2m}.
\end{equation}
\end{enumerate}
\end{proposition}

\begin{proof}
(a) $Q_+$ 的非对角正元恰是「正坐标 $\times$ 负坐标」的 $m$ 个交叉位置，值都是 $cd$；对角是 $1-c^2$（$m$ 个）与 $1-d^2$。
由 $m$ 个正坐标的对称性，Perron 向量形如 $(y,y,\dots,y,z)$，于是问题化到二维子空间
$\mathrm{span}\{\one_m/\sqrt m,\ e_{m+1}\}$ 上的矩阵
$\begin{pmatrix}1-c^2&\sqrt m\,cd\\ \sqrt m\,cd&1-d^2\end{pmatrix}$，其特征值即所给公式。
令 $z=c^2$、$d^2=1-mz$，被开方式成为 $1+(2m-2)z+(-3m^2+2m+1)z^2$，
求导并整理得临界点方程 $z\bigl[4m(3m+1)z-8m\bigr]=0$，非零根 $z=\frac2{3m+1}$。
代入后开方式\emph{恰好}等于 $1$（因为 $(m-1)^2+8m(m+1)=(3m+1)^2$），故
$\lambda_{\max}=\frac12\cdot\frac{5m-1}{3m+1}+\frac12=\frac{4m}{3m+1}$。端点 $z=0$ 与 $z=1/m$ 都给 $1<\frac{4m}{3m+1}$。

(b) $Q_+$ 是「对角 $1-\frac1{2m}$ $+$ $\frac1{2m}K_{m,m}$ 的邻接」，
故 $\lambda_{\max}=1-\frac1{2m}+\frac1{2m}\cdot m=1+\frac{m-1}{2m}$。$\square$
\end{proof}

\begin{corollary}[$\kappa_2(3)=\frac87$，已证]\label{cor:k23}
$\kappa_2(3)=\max_{\|v\|=1}\lambda_{\max}\bigl((I_3-vv^\top)_+\bigr)=\frac87$，
极大化点 $v=\bigl(\sqrt{2/7},\sqrt{2/7},-\sqrt{3/7}\bigr)$，
此时 $(I_3-vv^\top)_+$ 的特征多项式\textbf{精确}为 $\frac1{343}(7L-8)(7L-5)(7L-1)$。
\end{corollary}

\begin{proof}
$\mathcal P_2(3)=\{I_3-vv^\top:\|v\|=1\}$（$3$ 维中秩二投影就是补一条直线），
故不需要定理~\ref{thm:convex}，直接对 $v$ 的符号模式分类（三类）：
\emph{三个坐标同号，或某个坐标为零且其余两个同号}时 $Q_+$ 的非对角元全部被抹成 $0$，于是它是对角阵，
$\lambda_{\max}(Q_+)=\max_i(1-v_i^2)\le1<\frac87$；\emph{恰有一个坐标为零、另两个异号}时是 (a) 的 $m=1$ 情形，$f_{1,1}=1$；
\emph{$2$ 正 $1$ 负}（或 $1$ 正 $2$ 负，由 $v\mapsto-v$ 对称）时正是命题~\ref{prop:corank1}(a) 的 $m=2$ 情形，
最大值 $\frac{4\cdot2}{7}=\frac87$，$c^2=\frac27$、$d^2=\frac37$。
特征多项式由符号计算直接展开。$\square$
\end{proof}

\begin{remark}\label{rmk:k23read}
(i) 推论~\ref{cor:k23} 是 $k=2$ 情形常数 $c_2^\Sigma=\frac87$ 的一条\textbf{新的、三行的}证明；
已知证明要走完全二部吹起的谱计算。这条副线值得单独写成一个小节（见第~\ref{sec:pq7} 小节的 P4）。
(ii) $m=3$ 时 \eqref{eq:fm1} 给 $\frac65$，$m=2$ 时 \eqref{eq:fmm} 给 $\frac54$；
两者都在 $N=4$ 的余秩一层内，而 $\frac54>\frac65$，所以 $\kappa_3(4)$ 的余秩一见证是\emph{平衡}的那一个：
$v=(\frac12,\frac12,-\frac12,-\frac12)$，$\lambda_{\max}=\frac54$。这与注~\ref{rmk:n4exact} 的
$Q=I_4-vv^\top$ 是\textbf{同一个}矩阵——于是 $\kappa_3(4)=\frac54$ 现在有了两条独立证明：
行方差机制（定理~\ref{thm:LN} 在 $N=4$ 取等）与余秩一层穷举。
(iii) $f_{m,m}=\frac{3m-1}{2m}\to\frac32$、$f_{m,1}=\frac{4m}{3m+1}\to\frac43$：
余秩一层在 $N$ 大时给出的常数\emph{超过} $\rho_0$，
这提醒我们 $Q=I_N-vv^\top$ 的秩是 $N-1$ 而不是 $3$，它只对 $\kappa_{N-1}(N)$ 有意义，
不能拿来反驳 $\kappa_3\le\rho_0$。
\end{remark}

\subsection{斜率恒等式 $\rho=\frac{\alpha+\ell}2$、$\alpha=1$ 层与通道线性规划}\label{sec:slope1}

\begin{lemma}[斜率恒等式，已证]\label{lem:slope1}
在命题~\ref{prop:perronrig} 的记号下（$\alpha=x^\top Qx$、$\beta=x^\top Q_-x$、$\ell=x^\top|Q|x$、$\rho=\alpha+\beta$），
恒有
\begin{equation}\label{eq:slope1}
\ell=\alpha+2\beta,\qquad\text{从而}\qquad \rho=\frac{\alpha+\ell}2 .
\end{equation}
于是主猜想 $\kappa_3\le\rho_0$ \textbf{逐字等价}于「第四条不等式」$\alpha+\ell\le2\rho_0$，
并且冠军处取等（$\ell=1.586541134688899=2\rho_0-\alpha$，精确到 $10^{-15}$）。
\end{lemma}

\begin{proof}
$|Q|=Q_++Q_-=Q+2Q_-$（逐点），故 $\ell=\alpha+2\beta$；$\rho=\alpha+\beta$ 是命题~\ref{prop:perronrig} 的 (b)。
两式相加除以 $2$ 得 \eqref{eq:slope1} 的第二个式子。$\square$
\end{proof}

\begin{remark}\label{rmk:slope1read}
(i) 引理~\ref{lem:slope1} 把注~\ref{rmk:rigiditylimit} 里那条「缺失的第四不等式」的\emph{形状}钉死了：
它必须是 $\alpha+\ell$ 的线性界，斜率\textbf{只能}是 $1$（任何 $\alpha+c\ell\le C$ 形式在 $c\ne1$ 时
都无法被冠军饱和，因为冠军同时给出 $\rho$ 与 $\alpha,\ell$ 的值）。
(ii) 但它也说明这条路\textbf{没有提供新信息}：$\alpha+\ell\le2\rho_0$ 与 $\rho\le\rho_0$ 是同一句话。
真正的杠杆在于 $\ell=x^\top|Q|x$ 是一个\emph{绝对值}量，
而绝对值量有独立的工具（例如 $|Q|$ 的谱、$\|Q\|_1$、通道分解）。
第~\ref{sec:nogo} 节的定理~\ref{thm:zonal} 正是把「对 $|Q|$ 做逐点主函数」这一整族方案封顶在 $\Phi$；
所以剩下的可能性只有\textbf{耦合}通道（同时用到 $x$ 与 $Q$ 的联合结构）。
\end{remark}

\textbf{$\alpha=1$ 层（数值）。}若 $\alpha=1$，则 $Qx=x$，$x$ 落在 $Q$ 的像空间里，$\rho=1+\beta$。
我们在这一层上做了全局搜索（差分进化 $+$ 多起点，$N=6$），得到 $\max\beta=0.258695<\tau=\rho_0-1=0.281568$，
即 $\alpha=1$ 层给出 $\rho\le1.2587<\rho_0$。这说明\textbf{冠军不在 $\alpha=1$ 层的边界上}，
「$\alpha$ 尽量大」与「$\beta$ 尽量大」之间存在真实的权衡；zonal 见证在这一层不可行。

\textbf{通道线性规划（数值）。}把 $\ell=x^\top|Q|x$ 按「通道」（$q_{ij}$ 的符号与 $x_ix_j$ 的大小组合）拆成
$8701$ 个构型，对每个构型写出线性约束并用 HiGHS 求解，得到：
度 $2$ 通道的最优值 $z=2.7320505=1+\sqrt3$（与 Mohar 常数一致），
度 $4,6,8,10$ 通道相对 $2\rho_0=2.563136$ 的缺口分别是 $+7.6\times10^{-3}$、$+5.2\times10^{-3}$、
$+4.0\times10^{-5}$、$+1.1\times10^{-5}$。
结论：\emph{单独}的 $x$-加权通道族并\textbf{不}被 $\Phi$ 封顶（度 $2$ 通道能到 $1+\sqrt3>2\Phi$），
真正封顶的是\textbf{联合}通道界。这把第~\ref{sec:nogo} 节的否定性结论精确化了一层：
损失不在单个通道，而在通道之间的耦合。

\subsection{若干新的精确恒等式，以及一条被否证的充分性}\label{sec:miscid}

\textbf{(1) $\|Q_-\|_F^2=\rho_0\tau$，精确。}冠军处
\[
\|Q_-\|_F^2=10d_2^2=\frac{16+5\sqrt7}{81}=0.360848846362012=\rho_0(\rho_0-1)=\rho_0\tau ,
\]
符号计算核验差为 $0$。这个值同时等于第~\ref{sec:fourcs} 小节的 $S_O=\sum_ir_i$（不等式 [g] 在冠军处取等），
所以「负部的 Frobenius 能量」$=\rho_0\tau$ 是一个\emph{恒等式级别}的事实，不是近似。

\textbf{(2) $Q_+$ 的完整谱，精确。}由 \eqref{eq:psip} 与定理~\ref{thm:conefam}，冠军的 $Q_+$ 的谱是
\[
\mathrm{spec}(Q_+)=\Bigl\{\rho_0,\ \mu_1^{[2]},\ \mu_2^{[2]},\ \lambda_-\Bigr\}
=\{1.281568197,\ 0.692637910^{[2]},\ 0.117401872^{[2]},\ 0.098352240\},
\]
\[
\mu_1=\tfrac{155-5\sqrt5-5\sqrt7-\sqrt{35}}{180},\quad
\mu_2=\tfrac{25-5\sqrt5+5\sqrt7-\sqrt{35}}{180},\quad
\lambda_-=\tfrac12+\tfrac{\sqrt5}{18}-\tfrac{4\sqrt{35}}{45}
\]
（与推论~\ref{cor:champclosed} 完全一致；此前版本此处误排为 $\{\rho_0,\mu_1^{[2]},\mu_2,0,0\}$，
把 $\mu_2$ 写成单重并多出两个零，已于第 17 轮更正）。
$Q_-$（即 $|Q|$ 的负部支撑）的谱是
$\{0.379920437,\ 0.117401872^{[2]},\ 0,\ -0.307366^{[2]}\}=2g^\ast|s_2|\cdot\mathrm{spec}(A(C_5))$，
其中 $0.117401872$ 与 $Q_+$ 的 $\mu_2$ \emph{相同}——这是 $Z_5$ 对称性的直接后果
（同一个 $\psi_2$ 分量在 $Q$ 中是 $0$，在 $Q_+$ 与 $Q_-$ 中各留一份）。
注意 $\nu_-(Q_-)=2$ 而 $\nu_-(Q_+)=0$，与 Weyl 型控制 $\nu_-(Q_+)\le\nu_-(Q_-)$ 一致。

\textbf{(3) $|Q|$-路线的一行式与 Mohar 常数的第二次导出。}
$Q_+=\frac12(Q+|Q|)$ 给
\[
\lambda_{\max}(Q_+)\le\tfrac12\bigl(\lambda_{\max}(Q)+\lambda_{\max}(|Q|)\bigr)\le\tfrac12\bigl(1+\|Q\|_F\bigr)=\tfrac{1+\sqrt3}2,
\]
最后一步用 $\lambda_{\max}(|Q|)\le\||Q|\|_F=\|Q\|_F=\sqrt3$。这就是 Mohar 的 $\frac{1+\sqrt3}2$，
与定理~\ref{thm:LN} 的 $\lim_{N\to\infty}L(N)$ \textbf{完全相同}——
「行方差机制的 $N\to\infty$ 极限」与「一次 Cauchy--Schwarz」给同一个常数，
这不是偶然：定理~\ref{thm:LN} 的证明在双均匀层里退化成一次 CS。
冠军处 $\lambda_{\max}(|Q|)=1.602533590389977<\varphi$，故这条路线给 $1.301267$，
而 ETF 处 $\lambda_{\max}(|Q|)=\frac12+\frac{\sqrt5}2=\varphi$ 恰好给 $\frac{1+\varphi}2=\Phi$。
\textbf{\STW\ 的 $\Phi$ 就是 $|Q|$-路线在 ETF 上的取值。}

\textbf{(4) \STW\ 命题 6.1 的一行重导。}
对图 $G$，令 $M=\Sigma\circ Q$（$\Sigma=1\{Q>0\}$ 为掩码，$Q$ 为 $\Lam_3$ 的 Ky Fan 对偶投影），
则 $\Lam_3(G)=\tr(M\cdot A)$ 型的计算给出 $\Lam_3(G)\le n\cdot\lambda_{\max}(Q_+)-3$，
所以 $\kappa_3:=\sup_N\max_{Q\in\mathcal P_3(N)}\lambda_{\max}(Q_+)$ 一旦 $\le\rho_0$ 就得到主定理。
这条重导比 \STW\ 原文短得多，且把「$-3$」的来源讲清楚了（$\tr Q=3$）。

\textbf{(5) KKT 对偶变量（已验证）。}极值问题的 KKT 条件给出一个对偶矩阵 $\Lambda$，
冠军处 $\tr\Lambda=\lambda=\rho_0$ \emph{精确}，$\mathrm{spec}(\Lambda)=\{0.762031,\ 0.259768^{[2]}\}$。
这条恒等式说明冠军的 KKT 乘子没有退化，局部刚性（定理~\ref{thm:champ} 附近的 Hess 认证）是\emph{严格}的局部极大，
不是鞍点。

\textbf{(6) 一条被否证的充分性（否证也是交付物）。}
第~\ref{sec:fourcs} 小节的不等式 [g] $\beta\le\sqrt{S_OW_O}$ 在冠军处取等，
因此很自然地会猜想「$[g]$ 取等 $+$ 秩三 $\Rightarrow$ 冠军」。这是\textbf{假的}：
在 $N=4$ 上，$[g]$ 单独（配上秩三与 Perron 条件）允许 $\rho=1.4186>\rho_0$。
所以 [g] 的取等只是冠军的\emph{必要}特征，不是刻画；四条 CS 同时取等（注~\ref{rmk:quadtight}）才是。
这也再次说明：把不等式一条条加进 REL 时，必须\textbf{同时}检查它们能否被非冠军构型饱和。

\subsection{更新后的方法优先队列（第七版）}\label{sec:pq7}

第~\ref{sec:pq6} 小节的队列（第六版）在本轮之后需要大改：旧队首 P1 被命题~\ref{prop:chamberdead} 否证，
旧 P6 被引理~\ref{lem:slope1} 吸收。新队列如下（完整版见 \texttt{研究方法优先级协议.md} 第七版）：

\begin{center}\small
\begin{tabular}{p{1.1cm}p{3.1cm}p{8.6cm}p{1.6cm}}
\toprule
位次 & 名称 & 内容与判据 & 现状\\
\midrule
\textbf{P1} & 全局：秩三约束 $\times$ REL &
把引理~\ref{lem:rowid}(iii) 的三角形恒等式 \eqref{eq:tri}、四点 Gram 行列式 $\det G_4\ge0$、
以及 $\R^3$ 中等角线 $\le6$ 的约束加进第~\ref{sec:relax} 小节的 REL，
在 $Z_5$-对称五角星图案类内闭合 $3.2128\times10^{-3}$ 的缺口。
判据：REL$^+$ 的天花板降到 $\rho_0$，或找到一个 $>\rho_0$ 的合法构型（后者同样是重大结果）。 & 主攻\\
P2 & 整数化锐化（脱离 $Z_5$） &
把命题~\ref{prop:ladder6} 的整数化阶梯推广到 $N=6$ 的一般图案（不正则的钝角图），
特别是 $n=1$ 子层的精确上确界。判据：双均匀层的上界降到 $\le1.29$。 & 排队\\
P3 & 证明 $\mathrm{REL}(N)=L(N)$ &
定理~\ref{thm:LN} 只对双均匀层成立；REL 的自由上确界在 $N=5,6,7,8,10,12$ 数值上等于 $L(N)$（误差 $\le4\times10^{-6}$）。
若对一切 $N$ 证明相等，则 $\kappa_3(N)\le L(N)$ 成为无条件定理（$N\le6$ 优于或等于 \STW）。 & 排队\\
P4 & 副线：$\kappa_2=\frac87$ 对一切 $N$ &
推论~\ref{cor:k23} 已给 $N=3$。若能把余秩一层 $+$ 定理~\ref{thm:convex} 的组合推到一般 $N$，
就得到 $c_2^\Sigma=\frac87$ 的新证明。判据：$\mathcal P_2(N)$ 上的完整符号分类。 & 副线\\
P5 & 旧队列降级项 &
$k\ge3$ 断点定位、原子化、$\min(a,b)=1$ 族的 $\frac65$、范数改写：全部降为「有余力再做」。 & 降级\\
\midrule
弹出 & 旧 P1（腔掩码不存在性） & 命题~\ref{prop:chamberdead} 否证：$K_{2,2,2}$ 是腔掩码，$19/19$ 未闭合类都是腔掩码。 & 已否证\\
吸收 & 旧 P6（联合 $\alpha+\ell$） & 引理~\ref{lem:slope1} 说明它与主猜想逐字等价，无独立内容。 & 已吸收\\
\bottomrule
\end{tabular}
\end{center}

\textbf{P1 的三种具体形态}（按我们的偏好排序，都是数学方法而非搜索）：
\begin{enumerate}[label=(\alph*),leftmargin=2.2em]
\item \textbf{三角形恒等式路线}。由 \eqref{eq:tri}，$\sum_{i<j<k}q_{ij}q_{jk}q_{ki}$ 被对角线单独决定。
在 $Z_5$-对称层里这个和可以显式算出（只有三种三角形：$(h,j,j{+}1)$、$(h,j,j{+}2)$、$(j,j{+}1,j{+}2)$），
于是 \eqref{eq:tri} 变成 $(g_h,g_c,d_1,d_2,A)$ 的一个\emph{方程}。第~\ref{sec:conefamily} 小节说明这个方程在等变层里是恒等式，
所以必须\textbf{先降低对称性}（见 (c)）它才会咬人。
\item \textbf{四点 Gram 路线}。$Q=VV^\top$、$V\in\mathrm{St}(N,3)$，故任意四行的 Gram 行列式 $\det G_4\ge0$，
且在 $N\ge4$ 时四点必线性相关、$\det G_4=0$ 是\emph{等式}。冠军的六个向量在 $\R^3$ 中，
所有 $\binom64=15$ 个四点行列式都为零，这是 $15$ 个方程——远比 REL 的四条不等式强。
把它们写成 $(g,\pi,r,w)$ 语言的代理形式是 P1 的核心技术工作。
\item \textbf{降对称路线}。把 $Z_5$ 降为 $Z_2$（hub $+$ 锥点上的一个对合），
等变参数从 $1$ 维升到 $4$ 维，秩三约束第一次非平凡。
此时第~\ref{sec:conefamily} 小节的技巧（分类 $\mu$ 特征值）仍然适用，只是分量的重数变了。
\end{enumerate}

\section{本轮新结果（第 13--14 轮）：调和恒等式、帧测度重构、二阶变分公式与「六原子归约」}\label{sec:harm}

\subsection{为什么要换语言}\label{sec:harm:motiv}

前十一轮的工作有一个共同的形状：固定 $N$，把 $\mathcal P_3(N)$ 按符号图案分类，逐类认证。
这条路在 $N=5$ 上走通了（定理~\ref{thm:n5frame}），在 $N=6$ 上把 $22$ 类压到 $20$ 类，
但每增加一个 $N$，类数按 $2^{\binom N2}$ 增长，而注~\ref{rmk:ladderinfo}(2) 已经告诉我们
$N\ge7$ 的最大值点是\emph{退化}的——冠军加几个几乎零权的原子。
也就是说，逐 $N$ 穷举在原理上就在做重复劳动。本轮我们换了两次语言，
目的都是让 $N$ 从问题里消失：

\begin{itemize}[leftmargin=2em]
\item 第 13 轮：把 $Q=VV^\top$ 拆成「质量 $g_i$ $+$ 方向 $u_i\in S^2$」，
得到一组精确恒等式，其中最重要的是\textbf{调和恒等式} $\rho H=3$（定理~\ref{thm:harmonic}）。
它把主猜想 $\kappa_3\le\rho_0$ 改写成一个\emph{下界}问题：调和泛函 $H\ge 3/\rho_0$。
同时 $N$ 消失了：$\kappa_3$ 变成球面上一类「紧框架测度」的二次型上确界（定理~\ref{thm:nfree}）。
\item 第 14 轮：既然约束 $\sum g_i=3$、$\sum g_iu_iu_i^\top=I_3$ 对质量 $g$ 是\emph{线性}的，
可行集就是一个弱$^*$紧凸集，而它的极端点只有 $\le6$ 个原子（命题~\ref{prop:extreme}，Carathéodory）。
于是问题变成：目标 $\rho$ 沿这个凸集的直线是不是拟凸的？
定理~\ref{thm:secondvar} 给出二阶变分的\emph{精确闭式}，
它把拟凸性化归为 $Q_+$ 的谱符号这一件可计算的事情；
定理~\ref{thm:atomreduce} 给出条件性的「六原子归约定理」，
把全局猜想化归为 $N=6$ 的有限问题。
\end{itemize}

另外，定理~\ref{thm:Fn} 把冠军放进一个自然的对称族 $F_n$（枢纽 $+$ 正 $n$ 边形）里逐 $n$ 精确求解，
结论是：这个族在 $n=5$（即 $N=6$）处取到唯一峰值，值恰好是 $\rho_0$，
而且 $\rho_0$ 第一次有了一个不依赖数值搜索的解析来源。

\subsection{记号与字典}\label{sec:harm:dict}

设 $Q=VV^\top\in\mathcal P_3(N)$，$V\in\mathrm{St}(N,3)$，$Q_+$ 为逐元正部，
$\rho=\lambda_{\max}(Q_+)$，$x\ge0$、$|x|=1$ 为 $Q_+$ 的 Perron 向量。
记 $v_i\in\R^3$ 为 $V$ 的第 $i$ 行，$g_i=|v_i|^2=Q_{ii}$，$u_i=v_i/\sqrt{g_i}\in S^2$。
因为 $V^\top V=I_3$，
\begin{equation}\label{eq:frame}
\sum_{i=1}^N g_i=\tr Q=3,\qquad
\sum_{i=1}^N g_iu_iu_i^\top=\sum_{i=1}^N v_iv_i^\top=V^\top V=I_3 ,
\end{equation}
第二条即\textbf{紧框架条件}，等价于
\begin{equation}\label{eq:framez}
\sum_{i=1}^N\langle v_i,z\rangle^2=1\qquad\forall z\in S^2 .
\end{equation}
再定义
\[
c_i=\sqrt{g_i}\,x_i,\quad
G_+=\bigl(\langle u_i,u_j\rangle_+\bigr)_{ij},\quad
p=G_+c,\quad
\omega_i=\frac{c_i}{p_i},\quad
H=\sum_i\omega_i,\quad
\Psi=\sum_ic_i ,
\]
\[
\alpha=x^\top Qx=\Bigl|\sum_ic_iu_i\Bigr|^2,\quad
\ell=x^\top|Q|x=\sum_{ij}c_ic_j|\langle u_i,u_j\rangle|,\quad
\varphi(z)=\sum_ic_i|\langle u_i,z\rangle| .
\]

\begin{lemma}[字典]\label{lem:dict}
在上述记号下：
\begin{enumerate}[label=(\roman*),leftmargin=2.2em]
\item $\rho x_i=\sqrt{g_i}\,p_i$，从而 $\rho=c^\top G_+c=\sum_ic_ip_i$；
\item $g_i=\rho\,\omega_i$（\textbf{Perron$\leftrightarrow$框架字典}）；
\item $\rho H=3$，即 $\rho=3/H$；
\item $\sum_i\omega_iu_iu_i^\top=I_3/\rho$（\textbf{调和权的各向同性}）；
\item $2\rho=\alpha+\ell$，且 $\alpha\le1$；
\item $\varphi(z)\le1$ 对一切 $z\in S^2$（\textbf{带位势不等式}），
从而 $p_i\le\varphi(u_i)\le1$、$\ell\le\Psi\le\sqrt3$。
\end{enumerate}
\end{lemma}

\begin{proof}
(i) $\rho x_i=(Q_+x)_i=\sum_j\bigl(\sqrt{g_ig_j}\langle u_i,u_j\rangle\bigr)_+x_j
=\sqrt{g_i}\sum_j\langle u_i,u_j\rangle_+\sqrt{g_j}x_j=\sqrt{g_i}p_i$。
两边乘 $x_i$ 求和：$\rho=\sum_i\sqrt{g_i}x_ip_i=\sum_ic_ip_i=c^\top G_+c$。

(ii) 由 (i)，$p_i=\rho x_i/\sqrt{g_i}$，故 $\omega_i=c_i/p_i=\sqrt{g_i}x_i\cdot\sqrt{g_i}/(\rho x_i)=g_i/\rho$。

(iii) $H=\sum_i\omega_i=\sum_ig_i/\rho=3/\rho$，用了 \eqref{eq:frame}。

(iv) $\sum_i\omega_iu_iu_i^\top=\frac1\rho\sum_ig_iu_iu_i^\top=I_3/\rho$。

(v) $(t)_+=\frac12(t+|t|)$ 逐元成立，故 $\rho=x^\top Q_+x=\frac12(\alpha+\ell)$；
又 $Q\preceq I$（秩三正交投影），故 $\alpha=x^\top Qx\le|x|^2=1$；
$\alpha=|V^\top x|^2=|\sum_ix_iv_i|^2=|\sum_ic_iu_i|^2$。

(vi) 由 \eqref{eq:framez} 与 Cauchy--Schwarz，
\[
\varphi(z)^2=\Bigl(\sum_jx_j\cdot\sqrt{g_j}|\langle u_j,z\rangle|\Bigr)^2
\le\Bigl(\sum_jx_j^2\Bigr)\Bigl(\sum_jg_j\langle u_j,z\rangle^2\Bigr)=1\cdot|z|^2=1 .
\]
于是 $p_i=\sum_jc_j\langle u_j,u_i\rangle_+\le\sum_jc_j|\langle u_j,u_i\rangle|=\varphi(u_i)\le1$，
$\ell=\sum_ic_i\varphi(u_i)\le\sum_ic_i=\Psi$，
最后 $\Psi=\sum_i\sqrt{g_i}x_i\le\sqrt{\sum_ig_i}\sqrt{\sum_ix_i^2}=\sqrt3$。
\end{proof}

引理~\ref{lem:dict} 的六条都是几行的计算，但它们把原来分散的量全部串起来了：
$g$（质量）、$x$（Perron）、$u$（方向）、$p$（位势）、$\omega$（调和权）
五个对象之间只有\emph{一个}自由度——字典 (ii) 说，知道 $\rho$ 与 $\omega$ 就知道 $g$。
冠军处的数值（脚本 \texttt{\_n6.py}，全部残差 $\le4.4\times10^{-16}$）：
\[
\begin{aligned}
&g=(0.331767104061,\ 0.533646579188^{[5]}),\quad
c=(0.255828489269,\ 0.292702334536^{[5]}),\\
&p=(0.988228343580,\ 0.702933397722^{[5]}),\quad
\omega=(0.258875887269,\ 0.416401234433^{[5]}),\\
&H=2.340882059435,\quad \Psi=1.719340161947,\quad
\alpha=0.976595259055,\quad \ell=1.586541134689 .
\end{aligned}
\]

\subsection{调和恒等式与主猜想的新形式}\label{sec:harm:H}

\begin{theorem}[调和恒等式]\label{thm:harmonic}
对任意 $Q\in\mathcal P_3(N)$，
\[
\Bigl(\sum_ic_ip_i\Bigr)\Bigl(\sum_i\frac{c_i}{p_i}\Bigr)=3,\qquad\text{即}\qquad \rho\cdot H=3,\quad \rho=\frac3H .
\]
特别地
\[
\frac3{\rho_0}=\frac{54}{9+\sqrt5+2\sqrt{35}}
=\frac{15\sqrt7+9\sqrt{35}-18\sqrt5-48}{2}=2.340882059434595\ldots
\]
\end{theorem}

\begin{proof}
这就是引理~\ref{lem:dict}(i)$+$ (iii)：第一个因子是 $\rho$，第二个是 $H=3/\rho$。
闭式由 $\rho_0=(9+\sqrt5+2\sqrt{35})/18$ 有理化得到（\texttt{sympy} 化简差为 $0$）。
\end{proof}

这条恒等式的价值不在它难（它是 Cauchy--Schwarz 取等缺陷的另一种写法，
见定理~\ref{thm:psi}），而在于它\textbf{把上界问题翻成下界问题}：

\begin{conjecture}[主猜想的新形式]\label{conj:Hform}
对一切可行的 $(u,g)$（即满足 \eqref{eq:frame} 且 $x$ 为 $Q_+$ 的 Perron 向量），
\[
H=\sum_{i=1}^N\frac{c_i}{(G_+c)_i}\ \ge\ \frac3{\rho_0}=2.340882059434595 .
\]
等价地，在「调和权各向同性」$\sum_i\omega_iu_iu_i^\top=(H/3)I_3$
（引理~\ref{lem:dict}(iv) 的改写）这一约束下极小化 $H$。
冠军处取等：$\rho_0\cdot2.340882059434594=3$（残差 $4.4\times10^{-16}$）。
\end{conjecture}

\begin{remark}[为什么这个形式可能更好]\label{rmk:Hform}
(i) 约束 $\sum_i\omega_iu_iu_i^\top=(H/3)I_3$ 是\emph{齐次}的：把 $\omega$ 整体缩放不改变它，
所以 $H$ 的极小化只跟 $\omega$ 的\emph{比例}和方向 $u$ 有关，尺度由 $\rho=3/H$ 自动确定。
(ii) $\omega_i=c_i/p_i$ 是「质量除以位势」，物理上像静电容量：
$H$ 是总容量，而 $\rho=3/H$ 说\textbf{容量越大、谱半径越小}。
(iii) 它给出一个立刻可用的\emph{反伪证判据}：任何松弛如果产出的极值点不满足 $\rho H=3$，
就是伪的。第~\ref{sec:harm:fail} 小节里 REL2 的 $1.314006$ 正是这样被识破的。
\end{remark}


\subsection{$\Psi\le\sqrt3$ 的精确缺陷式，与 Mohar 界的两个新证明}\label{sec:harm:mohar}

\begin{theorem}[$\Psi$ 的上界与精确缺陷]\label{thm:psi}
$\Psi\le\sqrt3$，且缺陷有\emph{精确}表达式
\begin{equation}\label{eq:psidef}
3-\Psi^2=\sum_{i<j}\frac{c_ic_j(p_i-p_j)^2}{p_ip_j}\ \ge\ 0 .
\end{equation}
因此 $\Psi=\sqrt3$ 当且仅当 $p\equiv$ 常数；而 $p\equiv$ 常数与 $\rho=\frac12(\alpha+\ell)$、
$\ell\le\Psi$ 联立会推出 Mohar 界取等，这在秩三情形不可能（定理~\ref{thm:projmohar}(b)）。
冠军处 $3-\Psi^2=0.043869407516$，与 \eqref{eq:psidef} 右端之差 $1.9\times10^{-16}$。
\end{theorem}

\begin{proof}
对 $a_i=\sqrt{c_ip_i}$、$b_i=\sqrt{c_i/p_i}$ 用 Cauchy--Schwarz：
$\sum a_ib_i=\Psi$、$\sum a_i^2=\sum c_ip_i=\rho$、$\sum b_i^2=\sum c_i/p_i=H$，
而标准的 CS 缺陷式给
$\bigl(\sum a_i^2\bigr)\bigl(\sum b_i^2\bigr)-\bigl(\sum a_ib_i\bigr)^2=\sum_{i<j}(a_ib_j-a_jb_i)^2$，
右端 $=\sum_{i<j}c_ic_j\bigl(\sqrt{p_i/p_j}-\sqrt{p_j/p_i}\bigr)^2=\sum_{i<j}c_ic_j(p_i-p_j)^2/(p_ip_j)$。
左端由定理~\ref{thm:harmonic} 等于 $3-\Psi^2$。取等当且仅当 $a\propto b$，即 $p\equiv$ 常数。

若 $p\equiv p_0$，则由引理~\ref{lem:dict}(i) 与 $x=\sqrt g\,^{-1}c$：
$\rho x_i=\sqrt{g_i}p_0$，故 $x_i\propto\sqrt{g_i}$；结合 $|x|=1$ 与 $\sum g_i=3$ 得
$x_i=\sqrt{g_i}/\sqrt3$、$c_i=g_i/\sqrt3$，于是 $\Psi=\sum g_i/\sqrt3=\sqrt3$。
此时引理~\ref{lem:dict}(vi) 的两次 Cauchy--Schwarz 都必须取等：
$\varphi(u_i)=1$ 对所有 $i$ 成立，其取等条件是
$x_j/\sqrt{g_j}\propto|\langle u_j,u_i\rangle|$（对每个 $i$、一切 $j$），
而 $x_j/\sqrt{g_j}=1/\sqrt3$ 为常数，故 $|\langle u_i,u_j\rangle|=s$（$i\ne j$）为常数，
且 $g_i\equiv3/N$（等权紧框架）。这样的等角紧框架满足 $s^2=(N-3)/(3(N-1))$；
另一方面 $\ell=c^2(N+s^2N(N-1))=3/N+3s^2(N-1)/N$，令其等于 $\sqrt3$ 得
$s^2=(\sqrt3N-3)/(3(N-1))$，两式相减给 $N=\sqrt3N$，即 $N=0$，矛盾。
\end{proof}

\begin{corollary}[Mohar 界，三行新证]\label{cor:mohar3}
$\rho\le\frac12(1+\sqrt3)=1.366025403784439$。
\end{corollary}

\begin{proof}
$\rho=\frac12(\alpha+\ell)\le\frac12(1+\Psi)\le\frac12(1+\sqrt3)$，
用了引理~\ref{lem:dict}(v)、(vi)。$\square$
\end{proof}

第二个证明完全不同，而且顺带给出「永不取等」的最短理由。它只用 $Q$ 是\emph{投影}这一事实。

\begin{theorem}[投影子表述下的 Mohar 界]\label{thm:projmohar}
设 $Q\in\mathcal P_3(N)$。则：
\begin{enumerate}[label=(\alph*),leftmargin=2.2em]
\item \textbf{行能量恒等式}：$\sum_jQ_{ij}^2=(Q^2)_{ii}=Q_{ii}=g_i$，从而 $\bigl\||Q|\bigr\|_F^2=\sum_{ij}Q_{ij}^2=\tr Q^2=3$；
\item $\rho\le\frac12\bigl(1+\lambda_{\max}(|Q|)\bigr)\le\frac12(1+\sqrt3)$，且\textbf{第二个不等号永不取等}；
\item \textbf{Perron 方程的对称形式}：$2\rho x=Qx+|Q|x$，且 $(|Q|x)_i\le\sqrt{g_i}$；
\item 由此 $2\rho x_i\le\sqrt{g_i}(1+\sqrt\alpha)$，故 $\rho\le\frac{\sqrt3}2(1+\sqrt\alpha)$。
\end{enumerate}
\end{theorem}

\begin{proof}
(a) $Q^2=Q$（秩三正交投影）。(b) $Q_+=\frac12(Q+|Q|)$ 逐元，故
$\rho=\lambda_{\max}(Q_+)\le\frac12\lambda_{\max}(Q)+\frac12\lambda_{\max}(|Q|)
\le\frac12\cdot1+\frac12\bigl\||Q|\bigr\|_F=\frac12(1+\sqrt3)$。
若取等，则 $\lambda_{\max}(|Q|)=\||Q|\|_F$，即 $|Q|$ 的其余特征值全为零，$|Q|=aa^\top$ 为秩一；
比较对角元得 $a_i^2=g_i$，于是 $|\langle v_i,v_j\rangle|=\sqrt{g_ig_j}=|v_i||v_j|$，
所有 $v_i$ 两两平行，$\rank Q=1\ne3$，矛盾。

(c) 逐元 $(Q_+)_{ij}=\frac12(Q_{ij}+|Q_{ij}|)$ 作用在 $x$ 上即得。
再由 (a) 与 Cauchy--Schwarz，$(|Q|x)_i=\sum_j|Q_{ij}|x_j\le\sqrt{\sum_jQ_{ij}^2}\sqrt{\sum_jx_j^2}=\sqrt{g_i}$。

(d) $(Qx)_i=\langle v_i,w\rangle$，其中 $w=V^\top x$、$|w|^2=\alpha$，故
$(Qx)_i\le|v_i||w|=\sqrt{g_i\alpha}$；代入 (c) 得
$2\rho x_i\le\sqrt{g_i}(1+\sqrt\alpha)$。平方求和并用 $\sum g_i=3$：
$4\rho^2\le3(1+\sqrt\alpha)^2$。
\end{proof}

\begin{remark}[哪一步松、松多少]\label{rmk:projmoharslack}
冠军处 $\lambda_{\max}(|Q|)=1.602533590390$，于是 (b) 的第一段给
$\rho\le\frac12(1+1.602534)=1.301266795195$——\textbf{已经优于} $\Phi=1.309017$，
但仍比 $\rho_0$ 大 $1.97\times10^{-2}$。我们一开始以为「把 $\lambda_{\max}(|Q|)$ 直接压到
$2\rho_0-1=1.563136393744$ 就行了」，但冠军本身就违反它（$1.602534>1.563136$），
所以这条路\textbf{被否证}：$\lambda_{\max}(|Q|)$ 不是正确的中间量，
必须回到「同一个 $x$ 下的 $\alpha$ 与 $\ell$ 的\emph{联合}约束」。
(d) 更弱：$\alpha=0.976595$ 时它给 $\rho\le\frac{\sqrt3}2(1+0.988228)=1.721$，
远不如 Mohar，只作为 (c) 的一个副产品记录在此。
(c) 中 $(|Q|x)_i\le\sqrt{g_i}$ 在冠军处的最大余量是 $-6.78\times10^{-3}$，
即这一条其实相当紧，但把 $N$ 条不等式直接平方相加会丢掉它们之间的相位关系。
\end{remark}

\subsection{方差族与 zonotope 读法}\label{sec:harm:var}

\begin{theorem}[方差恒等式]\label{thm:variance}
\[
\langle p^2\rangle_\omega=\frac{\rho^2}3,\qquad
\operatorname{Var}_\omega(p)=\frac{\rho^2(3-\Psi^2)}9,\qquad
\sum_{ij}\omega_i\omega_j\langle u_i,u_j\rangle^2=\frac{H^2}3,\qquad
\sum_ig_ip_i^2=\rho^2 .
\]
冠军处依次为 $0.547472347745$、$0.008005762509$、$1.826576272061$、$1.642417043234$
（与闭式之差 $\le3.3\times10^{-16}$）。
\end{theorem}

\begin{proof}
$\sum_i\omega_ip_i^2=\sum_ic_ip_i=\rho$，除以 $H=3/\rho$ 得第一条；
$\langle p\rangle_\omega=\Psi/H=\Psi\rho/3$，代入即第二条。
第三条：$\sum_{ij}\omega_i\omega_j\langle u_i,u_j\rangle^2
=\tr\bigl(\sum_i\omega_iu_iu_i^\top\bigr)^2=\tr(I_3/\rho)^2=3/\rho^2=H^2/3$，
用了引理~\ref{lem:dict}(iv)。第四条：$\sum_ig_ip_i^2=\rho\sum_i\omega_ip_i^2=\rho^2$。
\end{proof}

注意 $\operatorname{Var}_\omega(p)$ 与 $3-\Psi^2$ 成正比：\textbf{$p$ 越接近常数，$\Psi$ 越接近 $\sqrt3$}，
而定理~\ref{thm:psi} 已证这不可能。所以方差恒等式其实是「Mohar 界不取等」的定量版本：
冠军处 $\operatorname{Var}_\omega(p)=8.0\times10^{-3}\ne0$，$p$ 的两个值
$0.988228$（枢纽）与 $0.702933$（五边形）差得很远。

\begin{proposition}[zonotope 读法]\label{prop:zonotope}
令 $Z=\sum_ic_i[-u_i,u_i]\subset\R^3$（三维 zonotope）。则引理~\ref{lem:dict}(vi) 恰好是
$Z\subseteq B^3$；支撑函数 $h_Z(z)=\varphi(z)$，平均宽度 $w(Z)=\frac2\pi\cdot\pi\Psi$ 正比于 $\Psi\le\sqrt3$，
且 $\ell=\sum_ic_ih_Z(u_i)$、$\rho=\frac12\bigl(|\sum_ic_iu_i|^2+\sum_ic_ih_Z(u_i)\bigr)$。
\end{proposition}

这条改写没有立即产出新不等式（我们试过用 $Z\subseteq B^3$ 的体积/宽度不等式，
都被 $Z$ 的高度非均匀性抵消），但它把问题变成了一个\emph{凸体}的语言：
「在单位球内的 zonotope 中，最大化 $\frac12(|$对角线$|^2+\sum c_ih_Z(u_i))$」，
并且生成向量 $u_i$ 还带着质量约束 \eqref{eq:frame}。记在这里，作为可能的后续接口。

\subsection{本轮两次未闭合的尝试（诚实记录）}\label{sec:harm:fail}

\begin{proposition}[挫折界，未闭合]\label{prop:frustration}
对任意符号染色 $\varepsilon\in\{\pm1\}^N$，令
$\mathrm{FRU}_\varepsilon=\sum_{i<j:\ \varepsilon_i\varepsilon_j\langle u_i,u_j\rangle<0}c_ic_j|\langle u_i,u_j\rangle|$
（被违反对的权重），$\mathrm{FRU}=\min_\varepsilon\mathrm{FRU}_\varepsilon$。则
$\ell\le1+4\,\mathrm{FRU}$，从而 $\rho\le1+2\,\mathrm{FRU}$。
冠军处 $\mathrm{FRU}=0.152486468908=\beta/2$（$\beta=x^\top Q_-x=0.304972937817$）（全 $+$ 染色最优，被违反的恰是五角星的五条边），
故该界给 $\rho\le1.304972937817$；要闭合需要 $\mathrm{FRU}\le\tau/2=0.140784098436$，
差 $1.17\times10^{-2}$，\textbf{未闭合}。
\end{proposition}

\begin{proof}
$\ell-\alpha_\varepsilon=\sum_{ij}c_ic_j(|\langle u_i,u_j\rangle|-\varepsilon_i\varepsilon_j\langle u_i,u_j\rangle)
=4\,\mathrm{FRU}_\varepsilon$（被违反对每对贡献 $2|\cdot|$，双向计两次），
而 $\alpha_\varepsilon=|\sum c_i\varepsilon_iu_i|^2\le1$（引理~\ref{lem:dict}(v) 对 $\varepsilon_i v_i$ 同样成立）。
\end{proof}

第二条尝试是把带位势不等式逐点量化：由 $\sum_ig_i\langle u_i,z\rangle^2=1$ 与 CS 缺陷可得
$p_i^2\le1-\rho O_i$，其中 $O_i=\sum_{j:\langle u_i,u_j\rangle<0}\omega_j\langle u_i,u_j\rangle^2$
是「钝角方向上的调和占用」；求和得 $\rho^2(1+\sum_i\omega_iO_i)\le3$。
冠军处左端 $=2.003265889596$，右端 $=3$，余量 $0.997$——\textbf{太粗，没有信息量}。
我们记录它是因为它的形状是对的（占用越界、$\rho$ 越小），
但系数 $\rho^2$ 前面的 $1$ 太弱，需要的是把 $\alpha$ 与 $\ell$ 的\emph{同一个} $x$ 关联起来，
而不是分别估计。

第三条是数值松弛 REL2（只保留 $\Psi\le\sqrt3$、$\sum c^2\le1$、$\varphi\le1$ 三族约束，
丢掉各向同性）：它在二十面体 ETF 方向上给出 $1.314006274436$，
质量 $(0.408211,0.264768^{[5]})$。这个值\textbf{是伪的}：
在该点上 $\rho H=3.078911\ne3$，直接违反定理~\ref{thm:harmonic}。
换句话说，调和恒等式是一个 REL2 看不见的约束，它一眼就识别出了松弛产物。
（把 $H\ge3/\rho$ 加回去得到 REL4；REL4 的差分进化在 $N=4$ 只搜到 $1.249796<\frac54$，
低于已知精确值，说明\emph{该松弛本身的数值求解不可靠}，我们据此放弃了这条数值路线，
不再把它当作证据。）


\subsection{$N$ 消失：帧测度重构}\label{sec:harm:nfree}

\begin{definition}[帧测度]\label{def:framemeas}
令 $\mathcal M$ 为所有形如 $\mu=\sum_{i=1}^Ng_i\delta_{u_i}$（$N$ 任意，$g_i>0$，$u_i\in S^2$）
且满足
\[
\mu(S^2)=3,\qquad \int_{S^2}uu^\top\dd\mu(u)=I_3
\]
的测度全体。对 $\mu\in\mathcal M$，令 $T_\mu$ 为 $L^2(\mu)$ 上以 $(\langle u,u'\rangle)_+$ 为核的积分算子，
$\rho(\mu)=\lambda_{\max}(T_\mu)$（核非负，故由 Perron--Frobenius 这就是谱半径，且特征函数可取非负）。
\end{definition}

\begin{theorem}[$N$-自由重构]\label{thm:nfree}
\[
\kappa_3=\sup_{\mu\in\mathcal M}\rho(\mu)
=\sup\Bigl\{\iint(\langle u,u'\rangle)_+h(u)h(u')\dd\mu\dd\mu'\ :\ \mu\in\mathcal M,\ h\ge0,\ \int h^2\dd\mu=1\Bigr\}.
\]
\end{theorem}

\begin{proof}
给定 $V$，取 $g_i=Q_{ii}$、$u_i=v_i/\sqrt g_i$，则 $\mu=\sum g_i\delta_{u_i}\in\mathcal M$
正是 \eqref{eq:frame}；取 $h_i=x_i/\sqrt{g_i}$，则 $\int h^2\dd\mu=\sum_ix_i^2=1$，且
$\iint(\langle\rangle)_+hh'\dd\mu\dd\mu'=\sum_{ij}(\langle u_i,u_j\rangle)_+g_ig_jh_ih_j
=\sum_{ij}(\langle u_i,u_j\rangle)_+c_ic_j=\rho$（引理~\ref{lem:dict}(i)）。
反向：给定 $(\mu,h)$，令 $x_i=\sqrt{g_i}h_i$（则 $|x|=1$）、$V$ 的行为 $\sqrt{g_i}u_i^\top$；
帧条件给 $V^\top V=I_3$，故 $Q=VV^\top\in\mathcal P_3(N)$，
而 $x\ge0$ 的 Rayleigh 商正是上式，由 Perron--Frobenius 不超过 $\lambda_{\max}(Q_+)$。
\end{proof}

于是 $\kappa_3$ 是一个球面上的 Motzkin--Straus 型二次型的上确界，
「图」被连续核 $(\langle u,u'\rangle)_+$ 取代，$N$ 只是 $\mu$ 的原子个数。

\begin{proposition}[$\mathcal M$ 的凸性与极端点]\label{prop:extreme}
$\mathcal M$ 在弱$^*$拓扑下的闭包 $\overline{\mathcal M}$ 是凸紧集，且其\textbf{极端点由至多 $6$ 个原子的测度构成}。
\end{proposition}

\begin{proof}
约束对 $\mu$ 是线性的，故 $\overline{\mathcal M}$ 凸；总质量固定为 $3$、$S^2$ 紧，故弱$^*$紧。
矩映射 $\phi(u)=(1,uu^\top)\in\R^7$ 的像落在 $5$ 维仿射子空间
$\{t=1,\ \tr=1\}$ 中，目标 $(3,I_3)$ 也在其中。
若极端点 $\mu$ 有 $\ge7$ 个原子，则这 $\ge7$ 个 $\phi(u_i)$ 在 $5$ 维仿射空间中仿射相关，
存在 $a\ne0$ 使 $\sum_ia_i=0$、$\sum_ia_iu_iu_i^\top=0$；
于是 $\mu\pm t\sum a_i\delta_{u_i}\in\overline{\mathcal M}$（$|t|$ 小），与极端性矛盾。
\end{proof}

命题~\ref{prop:extreme} 本身还不够：$\rho(\mu)$ 对 $\mu$ 既非凸也非凹
（$Q_+$ 的元是 $\sqrt{g_ig_j}$，对 $g$ 是\emph{凹}的，而 $\lambda_{\max}$ 是凸的，两者复合方向不明）。
所以需要下面的精确二阶计算。

\subsection{矩核直线上的二阶变分公式}\label{sec:harm:secvar}

固定方向 $u_1,\dots,u_N$，令 $K_{ij}=(\langle u_i,u_j\rangle)_+$、$b_i=\sqrt{g_i}$，则
$Q_+(g)=\bigl(b_ib_jK_{ij}\bigr)_{ij}$。设 $M\in\R^{7\times N}$ 为矩矩阵
（行 $1,u_1^2,u_2^2,u_3^2,u_1u_2,u_1u_3,u_2u_3$，$\rank M=6$）。
若 $N\ge7$，则 $\ker M\ne\{0\}$；取 $a\in\ker M$，
$g(t)=g+ta$ 保持矩约束，最大区间 $[t_-,t_+]$ 的端点处某个 $g_i$ 归零（原子被「杀掉」）。

\begin{theorem}[二阶变分公式]\label{thm:secondvar}
设 $\rho(t)=\lambda_{\max}(Q_+(g(t)))$ 在 $t$ 处为单特征值，$x(t)>0$、$|x(t)|=1$ 为对应 Perron 向量，
$(\rho_k(t),y_k(t))_{k\ge2}$ 为其余特征对。记 $\hat a_i=a_i/g_i$、$q=\hat a\odot x$、
$d_k=\langle q,y_k\rangle$。则
\begin{equation}\label{eq:firstvar}
\rho'(t)=\rho\sum_i\frac{a_i}{g_i}x_i^2=\rho\,\langle q,x\rangle,
\end{equation}
\begin{equation}\label{eq:secvar}
\boxed{\ \rho''(t)=2\rho\sum_{k\ge2}\frac{\rho_k}{\rho-\rho_k}\,d_k^{\,2}\ }
\qquad\text{等价地}\qquad
\rho''(t)=2\rho\,a^\top\!\Bigl(\sum_{k\ge2}\frac{\rho_k}{\rho-\rho_k}w_kw_k^\top\Bigr)a ,
\end{equation}
其中 $w_k=(x_iy_{ki}/g_i)_{i=1}^N$（故 $d_k=\langle a,w_k\rangle$）。
\end{theorem}

\begin{proof}
记 $A(t)=Q_+(g(t))$，$A_{ij}=b_ib_jK_{ij}$。由 $b_i'=a_i/(2b_i)$、$b_i''=-a_i^2/(4b_i^3)$，
\[
A'_{ij}=\tfrac12(\hat a_i+\hat a_j)A_{ij},\qquad
A''_{ij}=\Bigl(\tfrac12\hat a_i\hat a_j-\tfrac14\hat a_i^2-\tfrac14\hat a_j^2\Bigr)A_{ij}.
\]
包络定理给 $\rho'=x^\top A'x=\sum_i\hat a_ix_i(Ax)_i=\rho\sum_i\hat a_ix_i^2$，即 \eqref{eq:firstvar}。
二阶项用标准公式 $\rho''=x^\top A''x+2\sum_{k\ge2}\frac{|\langle y_k,A'x\rangle|^2}{\rho-\rho_k}$。
第一段：
\[
x^\top A''x=\tfrac12q^\top Aq-\tfrac\rho2|q|^2
=\tfrac12\Bigl(\rho\langle q,x\rangle^2+\sum_{k\ge2}\rho_kd_k^2\Bigr)-\tfrac\rho2\Bigl(\langle q,x\rangle^2+\sum_{k\ge2}d_k^2\Bigr)
=\tfrac12\sum_{k\ge2}(\rho_k-\rho)d_k^2 ,
\]
其中用了 $q=\langle q,x\rangle x+\sum_{k\ge2}d_ky_k$。
第二段：$\langle y_k,A'x\rangle=\frac12\bigl[\rho\langle \hat a\odot y_k,x\rangle+\langle \hat a\odot x,Ay_k\rangle\bigr]
=\frac12(\rho+\rho_k)d_k$。合并：
\[
\rho''=\tfrac12\sum_{k\ge2}d_k^2\Bigl[(\rho_k-\rho)+\frac{(\rho+\rho_k)^2}{\rho-\rho_k}\Bigr]
=\tfrac12\sum_{k\ge2}\frac{d_k^2}{\rho-\rho_k}\Bigl[(\rho+\rho_k)^2-(\rho-\rho_k)^2\Bigr]
=2\rho\sum_{k\ge2}\frac{\rho_kd_k^2}{\rho-\rho_k}. \quad\square
\]
\end{proof}

公式 \eqref{eq:secvar} 的结构值得停下来看一眼：
$\sqrt{g_ig_j}$ 的凹性贡献（$x^\top A''x\le0$）与特征向量旋转的贡献
（$2\sum\frac{|\langle y_k,A'x\rangle|^2}{\rho-\rho_k}\ge0$）\textbf{恰好合并}成一个单项和，
每一项的符号只由 $\rho_k$ 的符号决定。这把一个看起来无望的凸性问题变成了一个\emph{谱符号}问题。

\begin{corollary}[正部半定 $\Rightarrow$ 沿直线凸]\label{cor:psdconvex}
若在区间 $[t_-,t_+]$ 上恒有 $Q_+(g(t))\succeq0$（即 $\rho_k\ge0$ 对一切 $k$），
则 $\rho$ 在该区间上是凸函数，从而最大值在端点取得。
\end{corollary}

\begin{corollary}[一阶条件 $=$ 自洽场]\label{cor:scf}
$\rho'(0)=0$ 对一切 $a\in\ker M$ 成立，当且仅当向量 $(x_i^2/g_i)_i$ 属于 $M$ 的行空间，
即存在 $\lambda_0\in\R$ 与对称 $\Lambda\in\R^{3\times3}$ 使
\begin{equation}\label{eq:scfnew}
\frac{x_i^2}{g_i}=\lambda_0+u_i^\top\Lambda u_i\qquad(i=1,\dots,N),
\end{equation}
即 Perron 向量的平方除以质量是一个\emph{二次球面多项式}。
这与第~\ref{sec:n6global} 节的自洽场方程是同一条件的两种写法，
但这里是作为「沿矩核直线的一阶最优性」直接导出的，不涉及拉格朗日乘子的显式求解。
\end{corollary}

\begin{corollary}[内部极值的必要条件：负特征值个数]\label{cor:nuneg}
设 $\mu\in\mathcal M$ 有 $N\ge7$ 个原子且 $g$ 是 $\rho$ 在仿射矩集
$\{g':Mg'=(3,1,1,1,0,0,0)\}$ 内的局部最大值点（即一切 $g_i>0$ 的内部极值）。则
\[
\nu_-(Q_+)\ \ge\ N-6 ,
\]
其中 $\nu_-(\cdot)$ 计负特征值个数（含重数）。
\end{corollary}

\begin{proof}
内部局部最大要求 \eqref{eq:secvar} 中的二次型 $\Sigma=\sum_{k\ge2}\frac{\rho_k}{\rho-\rho_k}w_kw_k^\top$
在 $\ker M$（维数 $N-6$）上半负定，故 $N-6\le\nu_-(\Sigma|_{\ker M})\le\nu_-(\Sigma)$。
写 $W=[w_2\ \cdots\ w_N]=D_g^{-1}\operatorname{diag}(x)Y$，$Y=[y_2\ \cdots\ y_N]$ 列正交归一，
则 $\Sigma=D_g^{-1}\operatorname{diag}(x)\,YC Y^\top\operatorname{diag}(x)D_g^{-1}$，
$C=\operatorname{diag}\bigl(\frac{\rho_k}{\rho-\rho_k}\bigr)_{k\ge2}$；
由合同变换的惯性律 $\nu_-(\Sigma)\le\nu_-(C)=\#\{k\ge2:\rho_k<0\}=\nu_-(Q_+)$。
\end{proof}

\subsection{六原子归约：条件定理、猜想与数值证据}\label{sec:harm:reduce}

\begin{theorem}[条件性归约]\label{thm:atomreduce}
设 $\mu\in\mathcal M$ 有 $N\ge7$ 个原子，$a\in\ker M$，$[t_-,t_+]$ 为最大可行区间。
若 $Q_+(g(t))\succeq0$ 对一切 $t\in[t_-,t_+]$，则
$\rho(\mu)\le\max\{\rho(g(t_-)),\rho(g(t_+))\}$，
且两个端点各自至少有一个原子质量为 $0$（即原子数 $\le N-1$）。
\end{theorem}

\begin{proof}
推论~\ref{cor:psdconvex} 给凸性，凸函数在闭区间上的最大值在端点取得。
端点处 $g_i(t_\pm)=0$ 对某个 $i$ 成立，该原子对 $Q_+$ 的贡献整行整列为零，删去后仍是帧测度。
\end{proof}

反复使用定理~\ref{thm:atomreduce}，只要每一步都能找到一条「全程 $Q_+\succeq0$」的矩核直线，
就能把任意构型化约到 $\le6$ 个原子而不减小 $\rho$。于是：

\begin{conjecture}[六原子归约]\label{conj:sixatom}
\[
\kappa_3=\max_{3\le N\le6}\kappa_3(N)=\rho_0 .
\]
\end{conjecture}

\begin{remark}[为什么这是本轮最重要的一步]\label{rmk:reducepayoff}
注~\ref{rmk:ladderinfo}(2) 早就指出「只需要证明 $N\ge7$ 的投影可化约到 $\le6$ 个有效原子」，
但没有给出任何方法。猜想~\ref{conj:sixatom} 把这件事变成了一个\emph{可计算的判据}：
沿矩核直线检查 $Q_+$ 是否保持半定。一旦成立，
主猜想 $\kappa_3\le\rho_0$ 就\textbf{完全化归为} $\kappa_3(6)\le\rho_0$，
即第~\ref{sec:n6global} 节剩下的那 $20$ 个符号类——一个有限问题。
这是本项目第一次把「$N$ 任意」这个真正的困难从主线里移除。
\end{remark}

\textbf{数值证据}（脚本 \texttt{\_n2.py}、\texttt{\_n3.py}，输出 \texttt{\_n2\_out.txt}、\texttt{\_n3\_out.txt}）：
\begin{enumerate}[label=(\alph*),leftmargin=2.2em]
\item \textbf{二阶变分公式}：$N=7,8,9$ 各取随机帧测度与矩核方向，
\eqref{eq:firstvar}\eqref{eq:secvar} 与中心差分之差 $\le3.9\times10^{-8}$（即差分精度本身），共 $45$ 组。
\item \textbf{端点最大}：$N=7,8,9,10,12,14$ 共 $96$ 条完整矩核线段，
$\rho$ 的最大值\textbf{$96/96$ 都在端点取得}（拟凸性），其中 $90$ 条全程凸、
$6$ 条在某处轻微凹（最差二阶差分 $-4.5\times10^{-7}$）但最大值仍在端点。
\item \textbf{杀原子链单调}：从 $N=8,10,12$ 的随机帧测度出发，每步选端点较优者，
$6$ 条链\textbf{全部单调不减}，例如 $N=12$：$0.841622\to0.996704\to1.057454\to1.115438\to1.138190\to1.143178\to1.143313$。
\item \textbf{冠军邻域}（最危险的测试）：给冠军的 $6$ 个原子加 $K=1,2,3,4,6$ 个随机方向原子，
用最小二乘 $+$ 仿射投影把矩约束修回精确，得到 $N=7,8,9,10,12$ 的近极值构型；
其 $\rho$ 分别为 $1.281544230485$、$1.281520032516$、$1.281493080203$、$1.281424228047$、$1.281424252153$，
\textbf{全部严格小于} $\rho_0$；对它们执行杀原子归约，
每一步单调上升，末值与 $\rho_0$ 之差 $\le2.2\times10^{-15}$，
即归约\emph{精确地}把冠军恢复出来。
\end{enumerate}

\begin{remark}[必要条件 $\nu_-\ge N-6$ 有多强]\label{rmk:nunegstat}
推论~\ref{cor:nuneg} 说明：任何 $N\ge7$ 的内部局部最大值点必须有相当多的负特征值。
对 $500$ 个随机秩三投影统计（\texttt{\_n5.py}）：
\[
\begin{array}{c|cccccccc}
N&6&7&8&9&10&12&16&20\\ \hline
\nu_-(Q_+)\ \text{均值}&0.42&0.83&1.22&1.62&2.17&3.12&5.23&7.24\\
\nu_-(Q_+)\ \text{最大}&2&2&3&3&4&5&7&9\\
\#\{\nu_-\ge N-6\}/500&500&376&148&25&5&0&0&0\\
\max\rho\ \text{(随机)}&1.238020&1.225943&1.219509&1.215660&1.202749&1.205608&1.174544&1.138442
\end{array}
\]
即 $N\ge12$ 时随机构型\emph{从不}满足内部极值的必要条件，
且随机构型的 $\rho$ 随 $N$ 单调下降。这与注~\ref{rmk:ladderinfo}(2) 的退化性观察一致：
「多原子」不但不帮忙，反而系统性地压低 $\rho$。
冠军本身 $\nu_-(Q_+)=0$（$\lambda_{\min}(Q_+)=0.098352>0$），
落在推论~\ref{cor:psdconvex} 的凸性一侧——这是归约纲领能启动的原因。
\end{remark}


\subsection{族 $F_n$：枢纽 $+$ 正 $n$ 边形的精确解}\label{sec:harm:Fn}

第~\ref{sec:ladder} 小节在 $C_{4v}$ Ansatz 下解出了 $\kappa_3(5)$。
本轮我们把这个 Ansatz 推广成整个族，并\textbf{对每个 $n$ 精确求解}。
结果不只是漂亮：它给出了 $\rho_0$ 的第一个不依赖数值搜索的解析来源。

\begin{definition}[族 $F_n$]\label{def:Fn}
$n\ge3$。取枢纽 $u_*=e_3$ 与正 $n$ 边形
$u_k=(r\cos\frac{2\pi k}n,\ r\sin\frac{2\pi k}n,\ h)$，$k=0,\dots,n-1$；
质量为 $g_*$（枢纽）与 $g$（$n$ 个多边形顶点）。
帧条件 \eqref{eq:frame} 等价于
\[
\frac n2gr^2=1,\qquad ngh^2+g_*=1,\qquad g_*=3-ng ,
\]
故
\[
r^2=\frac2{ng},\qquad h^2=1-\frac2{ng},\qquad g\in\Bigl[\frac2n,\frac3n\Bigr],
\]
只有一个自由参数 $g$。Gram 元为 $\langle u_*,u_k\rangle=h$、
$\langle u_k,u_{k+m}\rangle=r^2\cos\frac{2\pi m}n+h^2=1-\frac4{ng}\sin^2\frac{\pi m}n=:s_m$。
\end{definition}

记 $\gamma_m=\frac4n\sin^2\frac{\pi m}n$，则 $gs_m=g-\gamma_m$，
于是多边形块是 circulant 矩阵，其元素为 $c_m=(g-\gamma_m)_+$（$m=1,\dots,n-1$），对角为 $g$。
枢纽与所有多边形顶点的耦合相等（$w=\sqrt{g_*g}\,h$），故全 $1$ 向量是多边形块的特征向量，
而枢纽只与它耦合。Fourier 分解立即给出：

\begin{theorem}[$F_n$ 的精确谱]\label{thm:Fn}
\[
\mu_l=g+\sum_{m=1}^{n-1}c_m\cos\frac{2\pi lm}n\quad(l=0,\dots,n-1),\qquad
\lambda_\pm=\frac{\mu_0+g_*\pm\sqrt{(\mu_0-g_*)^2+4g_*(ng-2)}}2 ,
\]
且
\[
\mathrm{spec}(Q_+(F_n))=\{\lambda_+,\lambda_-\}\cup\{\mu_l:l=1,\dots,n-1\},\qquad
\rho(F_n(g))=\max\Bigl\{\lambda_+(g),\ \max_{l\ge1}\mu_l(g)\Bigr\}.
\]
其中 $4g_*(ng-2)$ 来自 $(\sqrt n\,w)^2=ng_*gh^2=g_*(ng-2)$。
\end{theorem}

在每一个符号图案不变的区间上，$\lambda_+(g)=\frac12[S(g)+\sqrt{D(g)}]$ 是两个多项式的组合，
驻点方程 $4S'^2D=D'^2$ 是\emph{二次}的，故 $g^*$ 与 $\rho(F_n)$ 都有闭式。
\texttt{sympy} 逐 $n$ 求解的结果（\texttt{\_n4.py}，输出 \texttt{\_n4\_out.txt}）：

\[
\begin{array}{c|c|c|c|c}
n&N=n+1&\text{激活的 }m&g^*&\max_g\rho(F_n)\\ \hline
3&4&\varnothing&\frac45&\frac65=1.200000000000\\
4&5&\{1,3\}&\frac{24-\sqrt6}{30}&\frac{9+4\sqrt6}{15}=1.253197264742\\
5&6&\{1,4\}&\frac{45+4\sqrt5-\sqrt{35}}{90}&\frac{9+\sqrt5+2\sqrt{35}}{18}=\rho_0=1.281568196872\\
6&7&\{1,5\}&\frac{30-\sqrt{11}}{63}&\frac{13+4\sqrt{11}}{21}=1.250785674353
\end{array}
\]
驻点多项式分别为 $-1920g^2+3456g-1536$、$-960g^2+1536g-608$、
$-5760g^2+(5760+512\sqrt5)g-(1472+256\sqrt5)$、$-18144g^2+17280g-4064$；
$n=5$ 时 $\lambda_+(g^*)-\rho_0$ 经 \texttt{simplify} 化简为 \textbf{0}（不是数值近似）。
$n=7,\dots,14$ 的激活集随 $g$ 分段变化，我们只做数值最大化：
\[
1.240186898264,\ 1.260043241072,\ 1.254191504574,\ 1.238325394688,\
1.252993053595,\ 1.253419542550,\ 1.245607856493,\ 1.249638553443 ,
\]
全部 $\le1.260044<\rho_0$。$n\to\infty$ 的极限（多边形变成连续圆环，总质量 $M=ng\in[2,3]$）为
\[
\max_M\ \tfrac12\Bigl[I(M)+3-M+\sqrt{(I(M)-3+M)^2+4(3-M)(M-2)}\Bigr]=1.248748900844\quad(M=2.635),
\]
其中 $I(M)=\int_0^1(M-4\sin^2\pi s)_+\dd s$；同样低于 $\rho_0$（差 $3.28\times10^{-2}$）。

\begin{corollary}[族的峰值唯一在 $n=5$]\label{cor:Fnpeak}
$\max_{n\ge3}\max_g\rho(F_n)=\rho_0$，且唯一的最大化对是 $(n,g)=(5,g^*)$，即 $N=6$。
$n=3,4,5,6$ 为精确闭式，$n\ge7$ 与 $n\to\infty$ 为数值。
\end{corollary}

\begin{corollary}[冠军的闭式参数]\label{cor:champclosed}
冠军 $=$ $F_5$ 在 $g^*=\frac{45+4\sqrt5-\sqrt{35}}{90}=0.533646579188$ 处，且
\[
g_*=\frac{9-4\sqrt5+\sqrt{35}}{18}=0.331767104061,\qquad
h^2=1-\frac2{5g^*}=0.250440243412,\qquad
r^2=\frac2{5g^*}=0.749559756588,
\]
\[
\langle u_*,u_k\rangle=h=0.500440049768,
\]
\[
\text{五边形边}\ s_1=0.482066946497,\qquad \text{五角星}\ s_2=-0.355966337967,
\]
\[
\mathrm{spec}(Q_+)=\Bigl\{\rho_0,\ \mu_1^{[2]},\ \mu_2^{[2]},\ \lambda_-\Bigr\},\quad
\mu_1=\frac{155-5\sqrt5-5\sqrt7-\sqrt{35}}{180}=0.692637909856,
\]
\[
\mu_2=\frac{25-5\sqrt5+5\sqrt7-\sqrt{35}}{180}=0.117401871582,\qquad
\lambda_-=\tfrac12+\tfrac{\sqrt5}{18}-\tfrac{4\sqrt{35}}{45}=0.098352240252 .
\]
与第~\ref{sec:ladder} 节数值冠军的差：$g$ 为 $8.9\times10^{-16}$，特征值为 $2.2\times10^{-16}$。
\end{corollary}

\begin{remark}[族读出的三条信息]\label{rmk:Fnread}
(i) \textbf{$\rho_0$ 的来历}。$\rho_0$ 含 $\sqrt5$（来自正五边形的 $\cos72^\circ$）与
$\sqrt{35}$（来自驻点方程的判别式 $\sqrt{(5760+512\sqrt5)^2-4\cdot5760(1472+256\sqrt5)}$ 化简），
这解释了为什么这个常数「长得不像」$\Phi=\frac{1+\sqrt5}4\cdot$ 之类的量：
它是\emph{五重对称 $+$ 二次驻点}的产物，而不是某个松弛的天花板。
(ii) \textbf{$n$ 的方向}。$\rho(F_n)$ 在 $n=5$ 处达到峰值，$n=3$ 偏低
（$F_3$ 的掩码是星形 $K_{1,3}$，只有 $1.2$，而 $\kappa_3(4)=\frac54$ 由 $K_{2,2}$ 取得，
即 $N=4$ 时最优掩码\emph{不在}这个族里），$n\ge6$ 又因为
$\gamma_m=\frac4n\sin^2\frac{\pi m}n$ 变密、正部掩掉的边变多而下降。
「$6$ 个原子」因此不是偶然：它是这个自然族的全局峰。
(iii) \textbf{与归约纲领的接口}。推论~\ref{cor:Fnpeak} 加上猜想~\ref{conj:sixatom}
会把主猜想化归为「$N=6$ 的最优掩码是 $I+A(K_1\vee C_5)$」这一条有限断言；
反过来，若 $N=6$ 的 $20$ 个未闭合符号类中有一类超过 $\rho_0$，
它必然\emph{不在} $F_n$ 族里（族内已被精确排除），这本身就是一条结构信息。
\end{remark}

\subsection{数值：$\kappa_3(N)=\rho_0$ 推进到 $N=20$}\label{sec:harm:num}

第~\ref{sec:ladder} 节的 $\kappa_3(N)=\rho_0$ 只验证到 $N=14$（且 $13,14$ 只有一套求解器）。
本轮用帧测度参数化（方向 $+$ 质量，矩约束以罚函数施加）做独立的差分进化搜索
（\texttt{\_m2.py}，输出 \texttt{\_m2\_out.txt}）：
\[
\begin{array}{c|ccccccccccc}
N&4&5&6&7&8&9&10&12&14&16&20\\ \hline
\kappa_3(N)\ge&1.250000000000&1.253197264742&\rho_0&\rho_0&\rho_0&\rho_0&\rho_0&\rho_0&\rho_0&\rho_0&\rho_0\\
\text{与}\rho_0\text{之差}&-3.16\times10^{-2}&-2.84\times10^{-2}&\le2\times10^{-14}&\le6\times10^{-14}&\le2\times10^{-13}&\le2\times10^{-13}&\le2\times10^{-13}&\le3\times10^{-13}&\le5\times10^{-13}&\le7\times10^{-13}&\le6\times10^{-13}
\end{array}
\]
$N=4,5$ 的搜索值与精确闭式 $\frac54$、$\frac{9+4\sqrt6}{15}$ 吻合到 $10^{-13}$
（$\kappa_3(5)$ 的闭式由 \texttt{nsimplify} 独立重新发现，见 \texttt{\_m6\_out.txt}）。
$N\ge6$ 一律停在 $\rho_0$，从未超过——与猜想~\ref{conj:sixatom} 一致，
也与注~\ref{rmk:ladderinfo}(2) 的退化性一致。
\textbf{标签：数值}（$N\ge6$ 只给下界 $\kappa_3(N)\ge\rho_0$，上界仍需第~\ref{sec:n6global} 节的分类）。

\begin{remark}[$F_5$ 就是第 12 轮的 $Z_5$-等变弧]\label{rmk:F5arc}
把 $F_n$ 的参数换成 $s=ng-2\in[0,1]$（即「轴向质量」），则 $F_5$ 的 $g^*=\frac{45+4\sqrt5-\sqrt{35}}{90}$
对应 $s^*=5g^*-2=\frac{9+4\sqrt5-\sqrt{35}}{18}$，与引理~\ref{lem:z5class}、定理~\ref{thm:conefam}
的弧参数与临界点\textbf{完全一致}。所以本轮在 $n=5$ 上没有重复劳动：
贡献是把那条弧\emph{推广成整个族}并解出 $n=3,4,6$（以及 $n\ge7$ 的数值与 $n\to\infty$ 的极限），
从而第一次说明\textbf{为什么是五边形}——它是族内的全局峰，而不是某条人为选取的弧的顶点。
\end{remark}

\begin{remark}[一次诊断实验：强制 $Q_+$ 奇异并不能压低 $\rho$]\label{rmk:rhosing}
为了检验推论~\ref{cor:nuneg} 的补集一侧（$\nu_-(Q_+)\ge1$ 的构型能有多大 $\rho$），
我们在 $\mathrm{St}(7,3)$ 上做了带罚项 $\max(0,\lambda_{\min}(Q_+))$ 的差分进化（\texttt{\_n5.py}，$470$ 秒）。
结果是 $\rho=1.281568196116$（与 $\rho_0$ 差 $-7.6\times10^{-10}$）、$\lambda_{\min}(Q_+)=-1.8\times10^{-11}$，
质量 $g=(0.081784,0.533618,0.533650,0.331751,0.451863,0.533680,0.533656)$：
其中 $0.081784+0.451863=0.533647=g^*$，即\textbf{冠军的一个五边形原子被劈成两个}，
$\lambda_{\min}=0$ 正来自这一对近重合原子。
所以「$Q_+$ 奇异」这一侧的最大值仍是 $\rho_0$、且由\emph{退化}构型达到，
它既不能否证也不能支持猜想~\ref{conj:sixatom}，只是说明该实验\textbf{没有信息量}，
必须改成「$\nu_-(Q_+)\ge1$ 且各原子质量有下界」才有意义（下一轮 P1(b)）。
\end{remark}

\subsection{本轮状态、诚实标签与方法优先队列}\label{sec:harm:status}

\textbf{已证（符号推导 $+$ 数值证书）}：
引理~\ref{lem:dict}（字典六条）；定理~\ref{thm:harmonic}（$\rho H=3$ 与 $3/\rho_0$ 的闭式）；
定理~\ref{thm:psi}（$\Psi\le\sqrt3$ 与精确缺陷式、Mohar 永不取等）；
推论~\ref{cor:mohar3} 与定理~\ref{thm:projmohar}（Mohar 的两个新证明）；
定理~\ref{thm:variance}（方差恒等式）；命题~\ref{prop:zonotope}；
命题~\ref{prop:frustration}（挫折界，界本身已证、但不足以闭合）；
定理~\ref{thm:nfree}（$N$-自由重构）；命题~\ref{prop:extreme}（Carathéodory，$\le6$ 原子极端点）；
定理~\ref{thm:secondvar}（二阶变分公式）；推论~\ref{cor:psdconvex}、\ref{cor:scf}、\ref{cor:nuneg}；
定理~\ref{thm:atomreduce}（条件性六原子归约）；
定理~\ref{thm:Fn} 与推论~\ref{cor:Fnpeak}、\ref{cor:champclosed}（族 $F_n$ 的精确解，$n=3,4,5,6$ 闭式）。

\textbf{数值（不作成果）}：$\kappa_3(N)=\rho_0$ 对 $N=6,\dots,20$；
$n\ge7$ 的 $\rho(F_n)$；$n\to\infty$ 极限 $1.248748900844$；
$96$ 条矩核线段的拟凸性；$18$ 条杀原子链的单调性；$\nu_-$ 统计表。

\textbf{否证 / 未闭合（同样是成果）}：
$\lambda_{\max}(|Q|)\le2\rho_0-1$ 在冠军处为假（注~\ref{rmk:projmoharslack}）；
挫折界差 $1.17\times10^{-2}$（命题~\ref{prop:frustration}）；
$p_i^2\le1-\rho O_i$ 求和后余量 $0.997$，无信息量；
REL2 的 $1.314006274436$ 是伪值（$\rho H=3.078911\ne3$）；
REL4 的数值求解不可靠（$N=4$ 只得 $1.249796<\frac54$），已放弃。

\textbf{剩下的唯一缺口}：猜想~\ref{conj:sixatom} 的\emph{无条件}版本，
即证明「每个 $N\ge7$ 的构型都存在一条全程 $Q_+\succeq0$ 的矩核直线」，
或在 $\nu_-(Q_+)\ge N-6$ 的构型上直接证明 $\rho\le\rho_0$。
这是下一轮的队首。


\section{本轮新结果（第 15--17 轮）：Stiefel 刚性、能量分解与全局二分}\label{sec:r1517}

本节记录第 15 至第 17 轮的工作。三轮之间有一条清晰的主线：
第~\ref{sec:harm} 节把问题翻译成了「帧测度 $+$ 方向」的语言，并得到了一个\emph{条件性}的六原子归约
（定理~\ref{thm:atomreduce}）。本轮做三件事：
第一，把条件性归约所依赖的临界点方程组\textbf{完整解出}（\S\ref{sec:r17:erratum}--\S\ref{sec:r17:kkt}），
其间修正了一处会导致整套刚性论证被误判为失效的\textbf{代码错误}；
第二，找到主不等式的一个\textbf{最简等价形式}（\S\ref{sec:r17:tau}），它把「$\rho\le\rho_0$」压缩成
一个单一的标量关系 $\tau=\beta-(1-\alpha)$，并配上一条精确的 Frobenius 恒等式（\S\ref{sec:r17:frob}）；
第三，也是本轮的主定理：给出一个\textbf{无条件的全局二分定理}（定理~\ref{thm:dichotomy}），
它把「$N\ge7$ 能否成为极大元」这件事压缩成两个各自良定的子问题，
并且在数值上把这两个分支都\textbf{精确}地实现出来了（\S\ref{sec:r17:num}）。

同时本轮有一条\textbf{否证}：上一轮准备的「第四不等式」候选 $\alpha+\sqrt{A_1A_2}\le\rho_0$
虽然在冠军处尖锐取等，但\textbf{不是全局真命题}（命题~\ref{prop:fourthfalse}）。
按既定规则，否证与正定理同等记录。

\subsection{勘误：Stiefel 梯度公式的一次误报}\label{sec:r17:erratum}

第 15 轮曾报告「解析 Stiefel 梯度与数值差分相差 $2.03$ 倍（随机构型最大相对误差 $10$）」，
并据此把整条 $h_iW_i=\rho\Lam u_i$ 的刚性论证标为可疑。本轮查明：\textbf{这是脚本错误，不是数学错误}。

\begin{proposition}[正部的环境梯度]\label{prop:gradformula}
固定 $g>0$，令 $Q=VV^\top$、$V=U\operatorname{diag}(\sqrt g)$、$Q_+=K\circ Q$，
其中 $K_{ij}=\mathbf 1\{\langle u_i,u_j\rangle>0\}$（$K_{ii}=1$）。
设 $\lambda_{\max}(Q_+)$ 单重、Perron 向量 $x>0$、$\|x\|=1$，并记 $h_i=x_i/\sqrt{g_i}$、$c_i=g_ih_i=\sqrt{g_i}\,x_i$。
只要没有 $\langle u_i,u_j\rangle=0$，$\rho(g,U)=\lambda_{\max}(Q_+)$ 在 $(u_i)$ 方向可微，且
\begin{equation}\label{eq:gradu}
\frac{\partial\rho}{\partial u_i}=2c_iW_i,\qquad
W_i:=\sum_{j\in A(i)}c_ju_j,\qquad A(i):=\{j:\langle u_i,u_j\rangle>0\}\ \ (\text{含 }j=i).
\end{equation}
\end{proposition}

\begin{proof}
单重特征值的包络定理给 $\mathrm d\rho=x^\top(\mathrm dQ_+)\,x$。
在符号模式不变处 $\mathrm d(Q_+)_{ij}=K_{ij}\,\mathrm dQ_{ij}$，故
$\mathrm d\rho=\sum_{ij}x_ix_jK_{ij}\,\mathrm dQ_{ij}=2\bigl\langle (xx^\top\!\circ K)V,\mathrm dV\bigr\rangle$。
而 $\bigl((xx^\top\!\circ K)V\bigr)_i=x_i\sum_{j\in A(i)}x_j\sqrt{g_j}\,u_j=x_iW_i$，
且 $\mathrm dv_i=\sqrt{g_i}\,\mathrm du_i$，代入即得 $\mathrm d\rho=2\sum_ic_i\langle W_i,\mathrm du_i\rangle$。
\end{proof}

\textbf{错误的根源}。脚本用 \texttt{np.maximum(K,0)} 去构造 $W_i$，
即把\emph{截断余弦值} $(\langle u_i,u_j\rangle)_+$ 当成了\emph{指示函数} $\mathbf 1\{\langle u_i,u_j\rangle>0\}$。
这两个矩阵\textbf{只在再与 $K$ 缩并一次时才相同}：
$\sum_j(\langle u_i,u_j\rangle)_+c_j=\sum_{j\in A(i)}c_j\langle u_i,u_j\rangle=\langle u_i,W_i\rangle=p_i$，
而 $\sum_j(\langle u_i,u_j\rangle)_+c_ju_j\ne W_i$。
第~\ref{sec:harm} 节里所有用到 $p_i$ 的恒等式因此\textbf{不受影响}；受影响的只有把 $W_i$ 当向量使用的三处
（梯度、KKT 拟合、椭球面恒等式）。

\begin{center}\small
\begin{tabular}{lcc}
\hline
公式 & 冠军处 $\max_i\lvert$数值$-$解析$\rvert$ & 随机构型 \\
\hline
\eqref{eq:gradu}（指示掩码，正确） & $4.65\times10^{-10}$ & 全部一致 \\
截断余弦当掩码（错误） & $1.87\times10^{-1}$ & 相对误差最大 $10$ \\
\hline
\end{tabular}
\end{center}

\noindent 数值列是环境坐标中心差分（步长 $10^{-6}$）与 \eqref{eq:gradu} 的逐分量比较；
$4.65\times10^{-10}$ 正是该步长的截断误差量级。\textbf{结论：命题~\ref{prop:gradformula} 成立，
第 15 轮的刚性结果全部有效。}

\subsection{Stiefel KKT、椭球面恒等式与枢纽刚性}\label{sec:r17:kkt}

把 \eqref{eq:gradu} 投影到 Stiefel 流形 $\{V:V^\top V=I_3\}$ 的法空间 $\{VS:S=S^\top\}$ 上，
就得到临界点的完整方程组。

\begin{theorem}[Stiefel KKT 与迹一刚性]\label{thm:stiefelkkt}
设 $(U,g)$ 是 $\rho$ 在帧测度簇 $\mathcal M$（定义~\ref{def:framemeas}）上的临界点，
$\lambda_{\max}(Q_+)$ 单重且符号模式非退化。令 $S:=2V^\top(xx^\top\!\circ K)V$。则
\begin{equation}\label{eq:kkt}
2h_iW_i=Su_i\quad(\forall i),\qquad \tr S=2\rho .
\end{equation}
特别地，取 $\Lam:=S/(2\rho)$ 有
\begin{equation}\label{eq:trace1}
h_iW_i=\rho\Lam u_i,\qquad \tr\Lam=1,\qquad h_i^2=u_i^\top\!\Lam u_i,\qquad
\langle u_i,W_i\rangle=p_i=\rho h_i .
\end{equation}
\end{theorem}

\begin{proof}
法空间投影：临界 $\Leftrightarrow$ 存在对称 $S$ 使 $\nabla_V\rho=VS$（逐行 $\nabla_{v_i}\rho=(Sv_i)^\top$）。
由命题~\ref{prop:gradformula} 的证明，$\nabla_V\rho=2(xx^\top\!\circ K)V$，故 \eqref{eq:kkt} 的前一式；
左乘 $V^\top$ 并用 $V^\top V=I_3$ 得 $S=2V^\top(xx^\top\!\circ K)V$。
迹：$\tr S=2\sum_i v_i^\top(x_iW_i)=2\sum_ix_i\sqrt{g_i}\langle u_i,W_i\rangle=2\sum_ic_ip_i=2x^\top Q_+x=2\rho$。
把 $2h_iW_i=Su_i=2\rho\Lam u_i$ 与 $u_i$ 作内积，并用 $\langle u_i,W_i\rangle=p_i$，
得 $2h_ip_i=2\rho\,u_i^\top\!\Lam u_i$；代入 $p_i=\rho h_i$（由 \eqref{eq:kkt} 与 $\tr\Lam=1$ 亦可独立导出）
即得 $h_i^2=u_i^\top\!\Lam u_i$。
\end{proof}

\begin{remark}[迹一归一化的几何意义]\label{rmk:trace1}
$h_i^2=u_i^\top\!\Lam u_i$ 与 $\tr\Lam=1$ 合起来说：
\textbf{Perron 向量的「相对振幅」$h_i$ 恰好是原子方向 $u_i$ 在一个固定迹一椭球 $\Lam$ 上的支撑函数}。
冠军处 $\Lam$ 的谱是 $\{0.594608575866,\ 0.202695712067^{[2]},\ \}$，
即一个\emph{旋转椭球}，长半轴 $1/\sqrt{0.2027}$ 沿枢纽方向；
$\min_{S^2}u^\top\!\Lam u=0.2027>0$，故松弛「$\theta\ge0$ 在整个 $S^2$ 上」在冠军处是\textbf{紧}的。
\end{remark}

\begin{corollary}[椭球面恒等式]\label{cor:ellipsoid}
记 $m:=\sum_ic_iu_i=V^\top x$（于是 $\alpha=|m|^2$），
$S_i:=\sum_{j\in O(i)}c_ju_j$（钝部部分和，$O(i)=\{j:\langle u_i,u_j\rangle<0\}$），
则 $W_i=m-S_i$，且在定理~\ref{thm:stiefelkkt} 的条件下
\begin{equation}\label{eq:ellipsoid}
(m-S_i)^\top\Lam^{-1}(m-S_i)=\rho^2\qquad\text{对一切 }i\text{ 成立},
\end{equation}
等价地 $S_i^\top\Lam^{-1}S_i-2\langle z,S_i\rangle=-D$ 与 $i$ 无关，其中 $z=\Lam^{-1}m$、$D=m^\top\Lam^{-1}m-\rho^2$。
\end{corollary}

\begin{proof}
由 \eqref{eq:trace1}，$u_i=(h_i/\rho)\Lam^{-1}(m-S_i)$；代入 $h_i^2=u_i^\top\!\Lam u_i$ 得
$h_i^2=(h_i^2/\rho^2)(m-S_i)^\top\Lam^{-1}(m-S_i)$，约去 $h_i^2>0$ 即得。
\end{proof}

\begin{remark}[几何读法]\label{rmk:ellipsoidread}
\eqref{eq:ellipsoid} 说：\textbf{所有钝部部分和 $S_i$ 落在以 $m$ 为中心、以 $\Lam^{-1}$ 定形、
半径 $\rho$ 的同一张椭球面上}。这是一个把「$N$ 个原子的组合信息」压缩成「一张二次曲面」的强约束：
$N$ 越大，落在同一张椭球面上的点越多。这是 \S\ref{sec:r17:dichotomy} 全局二分的几何原型。
\end{remark}

\begin{corollary}[枢纽无关式与精确 Cauchy--Schwarz 缺陷]\label{cor:hubfree}
在定理~\ref{thm:stiefelkkt} 的条件下：
\begin{enumerate}\itemsep1pt
\item 若存在\textbf{枢纽}（某 $i_0$ 使 $O(i_0)=\varnothing$，即 $S_{i_0}=0$），则
\begin{equation}\label{eq:hubfree}
\rho^2=m^\top\Lam^{-1}m,\qquad
\tr(\Lam^{-1})-\rho^2=\sum_ig_i\bigl|\Lam^{-1/2}(u_i-h_im)\bigr|^2\ \ge0 ;
\end{equation}
\item $m=(\rho/h_{i_0})\Lam u_{i_0}$，从而 $u_{i_0}$ 是 $\Lam$ 的特征向量、$\lambda_{\max}(\Lam)=h_{i_0}^2$，且
$\alpha=\rho^2h_{i_0}^2$。
\end{enumerate}
\end{corollary}

\begin{proof}
(1) $S_{i_0}=0$ 代入 \eqref{eq:ellipsoid}。第二式是把 $\tr(\Lam^{-1})=\sum_ig_i\,u_i^\top\Lam^{-1}u_i$
（帧条件 $\sum g_iu_iu_i^\top=I_3$）展开并配方。
(2) 由 \eqref{eq:trace1} 在 $i=i_0$：$h_{i_0}m=\rho\Lam u_{i_0}$。
再取 $u_{i_0}$ 与两边内积：$h_{i_0}\langle u_{i_0},m\rangle=\rho h_{i_0}^2$，
即 $\langle u_{i_0},m\rangle=\rho h_{i_0}=p_{i_0}$；结合 $\Lam u_{i_0}=(h_{i_0}/\rho)m$ 知
$m$ 与 $u_{i_0}$ 平行且 $\Lam u_{i_0}=h_{i_0}^2u_{i_0}$，于是 $\alpha=|m|^2=\rho^2h_{i_0}^2$。
\end{proof}

冠军处的数值（全部到机器精度）：
\begin{center}\small
\begin{tabular}{ll}
\hline
检验 & 残差 \\
\hline
$h_iW_i=\rho\Lam u_i$（\eqref{eq:trace1}） & $4.44\times10^{-16}$ \\
$\tr\Lam-1$ & $0$（$\tr\Lam=1.000000000000000$） \\
$h_i^2=u_i^\top\!\Lam u_i$（无常数项，$\lambda_0=0$） & $2.2\times10^{-16}$ \\
$\rho^2-m^\top\Lam^{-1}m$ & $6.66\times10^{-16}$ \\
$\tr(\Lam^{-1})-\rho^2-\sum g_i|\Lam^{-1/2}(u_i-h_im)|^2$ & $<10^{-11}$ \\
$m-(\rho/h_0)\Lam u_0$ & $4.45\times10^{-16}$ \\
$\kappa_{\rm ang}-1$（$u_0$ 是 $\Lam$ 特征向量） & $0$ \\
$\alpha-\rho^2h_0^2$ & $<10^{-12}$ \\
$m$ 与 $u_0$ 的夹角余弦 & $1.000000000000000$ \\
\hline
\end{tabular}
\end{center}


\subsection{主不等式的最简等价形式：$\tau=\beta-(1-\alpha)$}\label{sec:r17:tau}

第~\ref{sec:harm} 节把问题写成 $\rho=\alpha+\beta\le\rho_0$，其中 $\alpha=x^\top Qx=|m|^2\le1$、
$\beta=x^\top Q_-x\ge0$。本轮把它进一步压缩：\textbf{需要控制的其实只是一个标量差}。

\begin{lemma}[$\tau$-形式]\label{lem:tauform}
记 $\tau:=\rho-1$、$\tau_0:=\rho_0-1=\frac{\sqrt5+2\sqrt{35}-9}{18}=0.281568196872168$。则
\begin{equation}\label{eq:tauform}
Q_-x=\tau x+(I-Q)x,\qquad\text{从而}\qquad
\boxed{\ \tau=\beta-(1-\alpha)=x^\top Q_-x-\|(I-Q)x\|^2\ } .
\end{equation}
于是\textbf{主猜想 $\kappa_3\le\rho_0$ 等价于：对一切 $N$ 与一切 Perron 对 $(Q_+,x)$，}
\begin{equation}\label{eq:maincrux}
x^\top Q_-x-\|(I-Q)x\|^2\ \le\ \tau_0 .
\end{equation}
\end{lemma}

\begin{proof}
$Q_+=Q+Q_-$ 且 $Q_+x=\rho x$ 给 $Q_-x=\rho x-Qx=(\rho-1)x+(I-Q)x$。
与 $x$ 作内积：$\beta=\tau+\|(I-Q)x\|^2$（用 $(I-Q)x\perp x$），
而 $\|(I-Q)x\|^2=x^\top(I-Q)x=1-\alpha$（$Q$ 是秩三正交投影，$Q^2=Q$）。
\end{proof}

\begin{remark}[缺口在哪里]\label{rmk:taugap}
\eqref{eq:maincrux} 的左端是两项之差：
\emph{钝部能量} $x^\top Q_-x=\beta$ 减去 \emph{Perron 向量离开列空间的距离} $\|(I-Q)x\|^2=1-\alpha$。
冠军处 $\beta=0.304972937817$、$1-\alpha=0.023404740945$，差恰为 $0.281568196872=\tau_0$
（数值核对到 $1.22\times10^{-15}$）——\textbf{两项都不小，是它们的差在取等}。
这解释了为什么过去所有「单独控制 $\beta$」的尝试（例如 $\beta\le\lambda_{\max}(Q_-)$）都注定松弛：
冠军处 $\lambda_{\max}(Q_-)=0.379920437>\tau_0$，
单独控制 $\beta$ 最好也只能给 $\rho\le1.37992$，比 Mohar 的 $1.36603$ 还差。
\textbf{必须让 $(I-Q)x$ 这一项参与抵消。}
\end{remark}

\begin{corollary}[$1+\lambda_{\max}(Q_-)$ 界的一行新证]\label{cor:oneneg}
$\rho\le1+\lambda_{\max}(Q_-)$。
\end{corollary}

\begin{proof}
由 \eqref{eq:tauform}，$\tau=\beta-\|(I-Q)x\|^2\le\beta=x^\top Q_-x\le\lambda_{\max}(Q_-)$。
\end{proof}

\begin{lemma}[$Q_-x$ 的精确范数]\label{lem:normqmx}
$\|Q_-x\|^2=(1-\alpha)(1+2\tau)+\tau^2=\beta^2+\alpha(1-\alpha)$，
且 $(I-Q)Q_-x=\rho(I-Q)x$、$QQ_-x=\tau Qx$。
\end{lemma}

\begin{proof}
第一式：把 \eqref{eq:tauform} 取范数平方，用 $(I-Q)x\perp x$。
第二式：$(I-Q)$ 作用在 \eqref{eq:tauform} 两端（$(I-Q)Q=0$）。
第三式：$Q$ 作用在两端（$Q(I-Q)=0$、$Qx=Qx$）。
末一等式 $\|Q_-x\|^2=\beta^2+\alpha(1-\alpha)$ 就是第 15 轮的 $\sigma$-恒等式
$\sum_ig_i\sigma_i^2=\beta^2+\alpha(1-\alpha)$（因 $(Q_-x)_i=\sqrt{g_i}\,\sigma_i$），
两者在此互相印证。
\end{proof}

\subsection{精确 Frobenius 恒等式与一行式谱地板}\label{sec:r17:frob}

\begin{lemma}[正部与负部的能量守恒]\label{lem:frob}
记 $e:=\|Q_-\|_F^2$。则
\begin{equation}\label{eq:frob}
\langle Q,Q_-\rangle=-e,\qquad \|Q_+\|_F^2=3-e,\qquad \sum_k\lambda_k(Q_+)^2=3-e,\qquad \sum_k\lambda_k(Q_+)=3 .
\end{equation}
\end{lemma}

\begin{proof}
$(Q_-)_{ij}>0$ 仅当 $Q_{ij}<0$，此时 $Q_{ij}=-(Q_-)_{ij}$，故
$\langle Q,Q_-\rangle=\sum_{ij}Q_{ij}(Q_-)_{ij}=-\sum(Q_-)_{ij}^2=-e$（\emph{精确}等式，非数值）。
$Q_+=Q+Q_-$ 给 $\|Q_+\|_F^2=\|Q\|_F^2+2\langle Q,Q_-\rangle+\|Q_-\|_F^2=3-2e+e=3-e$，
其中 $\|Q\|_F^2=\tr Q^2=\tr Q=3$（$Q$ 是秩三正交投影，即 Welch 界取等）。
迹：$\tr Q_+=\tfrac12(\tr Q+\tr|Q|)=\tfrac12(3+3)=3$，因 $\tr|Q|=\sum|Q_{ii}|=\sum g_i=3$。
\end{proof}

\begin{remark}[这条恒等式的价值]\label{rmk:frobval}
\eqref{eq:frob} 把「正部的谱」与「负部的能量」用一个\textbf{等式}锁在一起：
$\rho$ 想大，$e$ 就必须小；$e$ 小，则 $\|Q_+\|_F^2=3-e$ 大，正部的谱就必须更「铺开」。
冠军处 $e=\rho_0\tau_0=\frac{16+5\sqrt7}{81}=0.360848846362012$，
$\|Q_+\|_F^2=2.639151153638$，与 $\sum\lambda_k^2$ 直接计算吻合到 $1.3\times10^{-15}$。
\end{remark}

\begin{corollary}[一行式谱地板]\label{cor:floor2}
$n_+(Q_+)\ge\dfrac{9}{3-e}$（$n_+$ $=$ 正特征值个数）。
特别地，若 $e\le\rho_0\tau_0$ 则 $n_+(Q_+)\ge4$。
\end{corollary}

\begin{proof}
$3-e=\sum\lambda_k^2\ge(\sum\lambda_k)^2/n_+=9/n_+$。第二个断言：$9/(3-0.360848846)=3.410\Rightarrow n_+\ge4$。
\end{proof}

\begin{remark}
这与第 15 轮的谱地板 $n_+\ge1+(3-\rho)^2/(3-\rho^2)$（在 $\rho=\rho_0$ 处同为 $4$）结论一致，
但证明只有一行，且\textbf{不需要} $\rho$ 的取值，只需要 $e$。两条地板互为对偶。
\end{remark}

\begin{proposition}[条件性重排上界，优于 \STW]\label{prop:majbound}
若 $Q_+\succeq0$ 且 $n_+(Q_+)\le4$，则
\begin{equation}\label{eq:majbound}
\rho\ \le\ \frac{3+\sqrt{9-12e}}{4}\ \le\ \frac34+\frac{\sqrt{9-12\cdot0}}4=\frac32 ,
\end{equation}
且在冠军的 $e=\rho_0\tau_0$ 处右端 $=1.290244$，\textbf{严格小于 \STW\ 的 $\Phi=(3+\sqrt5)/4=1.309017$}。
\end{proposition}

\begin{proof}
设 $Q_+$ 的正特征值为 $\rho=\lambda_1>\lambda_2,\lambda_3,\lambda_4\ge0$（其余为 $0$）。
则 $\sum_{k=2}^4\lambda_k=3-\rho$、$\sum_{k=2}^4\lambda_k^2=3-e-\rho^2$，
Cauchy--Schwarz 给 $3-e-\rho^2\ge(3-\rho)^2/3$，整理得 $4\rho^2-6\rho+3e\le0$，解出 \eqref{eq:majbound}。
\end{proof}

\begin{remark}[诚实标注]\label{rmk:majcond}
命题~\ref{prop:majbound} 是\textbf{条件性}的：冠军自己 $n_+=6$，不满足 $n_+\le4$，
所以 \eqref{eq:majbound} 在冠军处\textbf{不适用}（把 $n_+=6$ 代入同样的重排论证只得 $1.474$）。
它的意义是：\textbf{如果}能把 $n_+$ 压到 $4$（推论~\ref{cor:floor2} 说 $4$ 是下界，所以 $n_+=4$ 是唯一可能的极值层），
那么 $\rho\le1.290244$ 就已经\textbf{无条件地优于} \STW。这把「改进 \STW」与「压低 $n_+$」两件事
精确地联系起来了，也给出了下一轮的一条明确路线（见 \S\ref{sec:r1517:status} 的 P2）。
\end{remark}


\subsection{第 16 轮：能量分解与双 Cauchy--Schwarz}\label{sec:r16:energy}

本小节的对象是\textbf{每个原子}的能量分配。对原子 $i$ 令 $O(i)=\{j:\langle u_i,u_j\rangle<0\}$、
$A(i)=\{j:\langle u_i,u_j\rangle>0\}$（含 $j=i$），并定义
\begin{align}
\eta_i&:=\sum_{j\in O(i)}g_j\langle u_i,u_j\rangle^2 &&\text{（$i$ 的\emph{钝部}方向能量）}\notag\\
\zeta_i&:=\sum_{j\in A(i),\,j\ne i}g_j\langle u_i,u_j\rangle^2 &&\text{（$i$ 的\emph{锐部}离对角能量）}\notag\\
\kappa_i&:=\sum_{j\in O(i)}g_jh_j^2=\sum_{j\in O(i)}x_j^2=1-\sum_{j\in A(i)}x_j^2 &&\text{（钝邻域的 Perron 质量）}\label{eq:energydef}\\
\sigma_i&:=\sum_{j\in O(i)}c_j|\langle u_i,u_j\rangle|,\qquad
p_i:=\langle u_i,W_i\rangle=\rho h_i,\qquad
A_1:=\sum_ic_i\kappa_i,\quad A_2:=\sum_ic_i\eta_i .\notag
\end{align}
\eqref{eq:energydef} 中 $\kappa_i$ 的第二个等号用 $g_jh_j^2=x_j^2$，第三个用帧条件
$\sum_jg_j\langle u_i,u_j\rangle^2=u_i^\top(\sum g_ju_ju_j^\top)u_i=1$ 与 $A(i)\cup O(i)$ 是全集。

\begin{lemma}[能量分解]\label{lem:framesplit}
$g_i+\eta_i+\zeta_i=1$ 对一切 $i$ 精确成立。特别地 $\|Q_-\|_F^2=e=\sum_ig_i\eta_i$。
\end{lemma}

\begin{proof}
帧条件按 $j=i$、$j\in O(i)$、$j\in A(i)\setminus\{i\}$ 三段拆分，$j=i$ 段贡献 $g_i\cdot1$。
第二个断言：$e=\sum_{ij}(Q_-)_{ij}^2=\sum_{ij}g_ig_j\langle u_i,u_j\rangle^2\mathbf 1\{ij\text{ 钝}\}=\sum_ig_i\eta_i$。
\end{proof}

\begin{lemma}[第一 Cauchy--Schwarz，冠军处取等]\label{lem:cs1}
$\sigma_i^2\le\kappa_i\eta_i$，等号 $\Leftrightarrow$ 在 $O(i)$ 上 $h_j$ 与 $|\langle u_i,u_j\rangle|$ 成比例。
\end{lemma}

\begin{proof}
$\sigma_i=\sum_{j\in O(i)}(\sqrt{g_j}h_j)(\sqrt{g_j}|\langle u_i,u_j\rangle)$，用 Cauchy--Schwarz。
\end{proof}

\begin{lemma}[第二 Cauchy--Schwarz]\label{lem:cs2}
$p_i^2\le(1-\kappa_i)(1-\eta_i)$，且
$\varphi(u_i):=\sum_jc_j|\langle u_i,u_j\rangle|\le\sqrt{(1-\kappa_i)(1-\eta_i)}+\sqrt{\kappa_i\eta_i}\le1$，
右端第二个不等号取等 $\Leftrightarrow\kappa_i=\eta_i$。
\end{lemma}

\begin{proof}
$p_i=\langle u_i,W_i\rangle$，把 $W_i$ 按 $A(i)$ 拆分并用引理~\ref{lem:framesplit}；
最后一句是 $\sqrt{(1-a)(1-b)}+\sqrt{ab}\le1$（$a,b\in[0,1]$，两边平方后化简为 $(\sqrt{a(1-b)}-\sqrt{b(1-a)})^2\ge0$）。
\end{proof}

\begin{theorem}[乘积界 \textrm{thm:betaAA}]\label{thm:betaAA}
对一切帧测度与一切 $N$：
\begin{equation}\label{eq:betaAA}
\beta^2\ \le\ A_1A_2 ,
\end{equation}
其中 $\beta=x^\top Q_-x$。冠军处取等：$\beta^2-A_1A_2=+1.39\times10^{-17}$。
\end{theorem}

\begin{proof}
两步 Cauchy--Schwarz：
\[
\beta=\sum_ic_i\sigma_i\overset{\text{引理}\ref{lem:cs1}}{\le}\sum_ic_i\sqrt{\kappa_i\eta_i}
\le\sqrt{\Bigl(\sum_ic_i\kappa_i\Bigr)\Bigl(\sum_ic_i\eta_i\Bigr)}=\sqrt{A_1A_2}.
\]
$c_i>0$ 保证第二步是真正的 Cauchy--Schwarz。取等数值核对：
$\beta=0.304972937817$、$A_1=0.469920973690$、$A_2=0.197923689318$。
\end{proof}

\begin{proposition}[取等刻画]\label{prop:betaAAeq}
\eqref{eq:betaAA} 取等 $\Leftrightarrow$ 同时满足
\begin{enumerate}\itemsep1pt
\item[(E1)] 每条钝边 $\{i,j\}$ 上 $h_i=h_j$；
\item[(E2)] $(h_i+h_j)^2/|\langle u_i,u_j\rangle|$ 在全部钝边上是同一常数。
\end{enumerate}
\end{proposition}

\begin{proof}
第一步取等要求在每个 $O(i)$ 上 $h_j\propto|\langle u_i,u_j\rangle|$；
对两条共享端点的钝边比较比例常数即得 (E1)，回代即得 (E2)。
\end{proof}

冠军的钝边集是五芒星的五条边 $\{(1,3),(1,4),(2,4),(2,5),(3,5)\}$；
(E1) 的最大偏离 $1.22\times10^{-15}$、(E2) 的比例散布 $3.28\times10^{-15}$。
在 $300$ 个\textbf{人为破坏} (E2) 的随机构型中，$299$ 个给出严格不等式
（余下 $1$ 个是钝图为空的退化情形，此时 $\beta=A_1=A_2=0$，两边同为 $0$）。

\begin{remark}[为什么这条定理不够]\label{rmk:betaAAlimit}
\eqref{eq:betaAA} 在冠军处\textbf{尖锐}，所以它本身没有损失。损失发生在下一步：
要用它闭合主猜想，还需要 $\alpha+\sqrt{A_1A_2}\le\rho_0$，而这一条\textbf{是假的}
（命题~\ref{prop:fourthfalse}）。这是本轮最重要的一次「否证」，见下节。
\end{remark}

\subsection{否证：全局「第四不等式」$\alpha+\sqrt{A_1A_2}\le\rho_0$}\label{sec:r16:refute}

\begin{proposition}[否证]\label{prop:fourthfalse}
不等式
\begin{equation}\label{eq:fourth}
\alpha+\sqrt{A_1A_2}\ \le\ \rho_0
\end{equation}
\textbf{不是}对所有帧测度成立的真命题。存在 $N=7$ 的构型使左端 $=1.287835562520>\rho_0=1.281568196872$，
超出 $6.27\times10^{-3}$。
\end{proposition}

\begin{proof}[数值证据]
在 $360$ 个随机构型（$N=4,\dots,14$，由 \texttt{random\_frame\_measure} 通过 LP 构造，保证帧条件精确成立）
上扫描，最大值为上式。脚本 \texttt{\_o2.py} 第 (E) 段，输出文件 \texttt{\_o2\_out.txt}。
\end{proof}

\begin{remark}[否证的含义，以及它\textbf{不}否证什么]\label{rmk:fourthmean}
\begin{enumerate}\itemsep2pt
\item \eqref{eq:fourth} 在冠军处\textbf{取等}（残差 $-4.44\times10^{-16}$），
所以它是一条「尖锐的候选」；被否证的是它的\textbf{全局有效性}，不是它的尖锐性。
\item 它\textbf{不}否证主猜想 $\kappa_3\le\rho_0$：反例构型的真实 $\rho=\alpha+\beta$ 仍远小于 $\rho_0$，
失效的只是 $\beta\le\sqrt{A_1A_2}$ 这一步的\textbf{松弛量}。
\item 因此 \textbf{thm:betaAA} 这条路线只可能在\emph{临界点集合}上完成：
在满足定理~\ref{thm:stiefelkkt} 全部方程的构型上，$\kappa_i,\eta_i$ 不再自由
（它们被 $\Lam$、$h$、钝图三者共同锁定），\eqref{eq:fourth} 有可能在那里成立。
这把问题从「$N$ 维不等式」缩小成「$6+1$ 个参数的刚性方程组上的不等式」，
是\textbf{一次真实的降维}，但尚未完成。
\item 反例出现在 $N=7$，正好是 $\dim\ker M\ge1$ 的第一个阶数。这提示
\eqref{eq:fourth} 的失效与「矩核方向上 $\rho$ 不取极值」是同一现象的两面，
应当与定理~\ref{thm:dichotomy} 合并使用。
\end{enumerate}
\end{remark}

\textbf{仍然全局成立}（$360$ 个构型 $0$ 违背，最大超出如下）：
\begin{center}\small
\begin{tabular}{lcc}
\hline
不等式 / 恒等式 & 最大超出 & 冠军处 \\
\hline
$g_i+\eta_i+\zeta_i=1$（引理~\ref{lem:framesplit}） & $3.11\times10^{-15}$ & 恒等式 \\
$\sigma_i^2\le\kappa_i\eta_i$（引理~\ref{lem:cs1}） & $+1.39\times10^{-17}$ & \textbf{取等} \\
$p_i^2\le(1-\kappa_i)(1-\eta_i)$（引理~\ref{lem:cs2}） & $+8.40\times10^{-16}$ & 余量 $2.34\times10^{-2}$ \\
$\varphi(u_i)\le\sqrt{(1-\kappa_i)(1-\eta_i)}+\sqrt{\kappa_i\eta_i}$ & $0$ & 余量 $1.18\times10^{-2}$ \\
$\sqrt{(1-\kappa_i)(1-\eta_i)}+\sqrt{\kappa_i\eta_i}\le1$ & $0$ & 严格 \\
$\beta^2\le A_1A_2$（定理~\ref{thm:betaAA}） & $0$ & \textbf{取等} \\
$\rho^2\le\sum_ig_i(1-\kappa_i)(1-\eta_i)$ & $+3.33\times10^{-16}$ & 余量 $2.56\times10^{-1}$（约 $15\%$ 松） \\
\hline
\end{tabular}
\end{center}


\subsection{新全局定理：矩核二阶变分的符号二分}\label{sec:r17:dichotomy}

这是本轮的主定理，也是对「先攻全局」这一指令的直接回应。
它的目标是：\textbf{不依赖任何未证猜想}，说明 $N\ge7$ 的构型为什么不能成为极大元，
并把残留的困难精确地分割成两个各自良定的分支。

\paragraph{设定}
固定方向组 $U=(u_1,\dots,u_N)\subset S^2$。帧测度簇
$\mathcal G_U=\{g\in\R^N_{>0}:\sum_ig_iu_iu_i^\top=I_3\}$
是一个开凸集内的仿射切片，其方向空间是
$\ker M$，其中 $M$ 是 $7\times N$ 的\emph{矩矩阵}，
行为 $\bigl(1,\ u_{i,0}^2,\ u_{i,1}^2,\ u_{i,2}^2,\ u_{i,0}u_{i,1},\ u_{i,0}u_{i,2},\ u_{i,1}u_{i,2}\bigr)_{i=1}^N$。
因为 $u_{i,0}^2+u_{i,1}^2+u_{i,2}^2\equiv1$，$\rank M\le6$，故
\begin{equation}\label{eq:kerdim}
\dim\ker M=N-\rank M\ \ge\ N-6 ,
\end{equation}
\textbf{这正是「$N\ge7$ 才有矩核方向」的来源}，也是 $N=6$ 成为冠军的结构性原因：
$N=6$ 时一般 $\dim\ker M=0$，质量 $g$ 被方向 $U$ 完全锁死，没有内部扰动可用。

令 $\mathcal K:=\{\,q=D\,a:a\in\ker M\,\}$，$D=\operatorname{diag}(x_i/g_i)$，$\dim\mathcal K=\dim\ker M$。

\begin{lemma}[一阶条件是无条件的]\label{lem:firstorder}
若 $g$ 是 $\rho(\cdot,U)$ 在 $\mathcal G_U$ 上的局部极值点且 $\lambda_{\max}(Q_+)$ 单重，则
$\langle a,h^2\rangle=0$ 对一切 $a\in\ker M$ 成立，即 $h^2=(h_i^2)_i\in\operatorname{rowspace}(M)$：
存在对称 $\Lam$ 与常数 $\lambda_0$ 使 $h_i^2=u_i^\top\!\Lam u_i+\lambda_0$。
\end{lemma}

\begin{proof}
$g+ta\in\mathcal G_U$ 对小 $|t|$ 成立（$g>0$），故 $\rho'(0)=0$。
由定理~\ref{thm:secondvar} 的一阶部分 $\rho'(0)=\rho\sum_i(a_i/g_i)x_i^2=\rho\langle a,h^2\rangle$。
\end{proof}

\begin{theorem}[二阶变分的谱形式，复核]\label{thm:secondvar2}
在引理~\ref{lem:firstorder} 的条件下，沿 $a\in\ker M$：
\begin{equation}\label{eq:secondvar2}
\rho''(0)=2\rho\,q^\top R\,q,\qquad q_i=\frac{a_ix_i}{g_i},\qquad
R:=\sum_{k\ge2}\frac{\rho_k}{\rho-\rho_k}\,w_kw_k^\top ,
\end{equation}
其中 $\rho=\rho_1>\rho_2\ge\cdots$ 与 $w_k$ 是 $Q_+$ 的特征值与单位特征向量。
\end{theorem}

\begin{proof}[证明要点（本轮补全了此前缺失的抵消步骤）]
$Q(t)_{ij}=\sqrt{(g_i+ta_i)(g_j+ta_j)}\,\langle u_i,u_j\rangle$，故
$Q'(0)=\tfrac12(\hat AQ+Q\hat A)$、$Q''(0)_{ij}=-Q_{ij}(\hat a_i-\hat a_j)^2/4$，$\hat a_i=a_i/g_i$。
符号掩码在 $t=0$ 附近不变，故 $\mathrm d(Q_+)_{ij}=K_{ij}\mathrm dQ_{ij}$，
且 $K\circ Q'(0)$ 满足 $x^\top(K\circ Q')w_k=\tfrac12(\rho+\rho_k)(q^\top w_k)$。
标准二阶摄动给
$\rho''=2\sum_{k\ge2}\frac{|x^\top(K\circ Q')w_k|^2}{\rho-\rho_k}+x^\top(K\circ Q'')x$，
其中 $x^\top(K\circ Q'')x=-\tfrac12\bigl[\rho|q|^2-\sum_k\rho_k(q^\top w_k)^2\bigr]$。
由引理~\ref{lem:firstorder}，$q^\top x=\langle a,h^2\rangle=0$，故 $|q|^2=\sum_{k\ge2}(q^\top w_k)^2$；
代入后 $\rho^2$ 与 $\rho_k^2$ 的项\textbf{精确抵消}，只剩 $2\rho\rho_k$ 项，得 \eqref{eq:secondvar2}。
数值复核：$N=7,\dots,10$ 共 $12$ 个构型、每个取 $4$ 条矩核方向，
$\max|\rho''_{\rm num}-\eqref{eq:secondvar2}|=5.86\times10^{-8}$（步长 $2\times10^{-4}$ 的中心二阶差分截断误差量级）。
\end{proof}

\begin{theorem}[全局二分定理（本轮主定理）]\label{thm:dichotomy}
设 $N\ge7$，$g_i>0$ 对一切 $i$，$\lambda_{\max}(Q_+)$ 单重，且 $(U,g)$ 是 $\rho$ 在 $\mathcal G_U$ 上的局部极大点。则
\begin{equation}\label{eq:Rneg}
R\big|_{\mathcal K}\preceq0 ,
\end{equation}
从而：
\begin{enumerate}\itemsep2pt
\item[(i)] $n_+(Q_+)\le7$；
\item[(ii)] \textbf{分支 I}：$\nu_-(Q_+)\ge1$；\quad\textbf{或}\quad
\textbf{分支 II}：$Q_+\succeq0$ 且 $\mathcal K\subseteq\ker Q_+$（此时 $n_+(Q_+)\le6$）；
\item[(iii)] 特别地，\textbf{若 $Q_+\succ0$ 则必有 $N\le6$}。
\end{enumerate}
\end{theorem}

\begin{proof}
局部极大给 $\rho''(0)\le0$ 对一切 $a\in\ker M$，即 \eqref{eq:Rneg}。
把 $\R^N$ 按 $Q_+$ 的谱分解为
$\mathcal P=\operatorname{span}\{w_k:\rho_k>0,k\ge2\}$、
$\mathcal Z=\ker Q_+$、$\mathcal N=\operatorname{span}\{w_k:\rho_k<0\}$，
则 $R$ 在 $\mathcal P$ 上 $\succ0$、在 $\mathcal N$ 上 $\prec0$、在 $\mathcal Z\oplus\R x$ 上 $=0$。

(i) 若 $0\ne q\in\mathcal K\cap\mathcal P$ 则 $q^\top Rq>0$，与 \eqref{eq:Rneg} 矛盾。故
$\dim\mathcal K+\dim\mathcal P\le N$，即 $(N-6)+(n_+-1)\le N$，得 $n_+\le7$。

(ii) 若 $\nu_-=0$ 则 $\mathcal N=\{0\}$，于是 $q^\top Rq=0\Leftrightarrow q\in\mathcal Z\oplus\R x$；
又 $q\in\mathcal K$ 已蕴含 $q\perp x$，故 \eqref{eq:Rneg}（此时 $R\succeq0$）强制 $R|_{\mathcal K}\equiv0$，
即 $\mathcal K\subseteq\mathcal Z=\ker Q_+$。此时 $\dim\ker Q_+=N-n_+\ge\dim\mathcal K=N-6$，得 $n_+\le6$。

(iii) $Q_+\succ0$ 时 $\mathcal Z=\{0\}$、$\mathcal N=\{0\}$，(ii) 给 $\mathcal K=\{0\}$，
故 $\dim\ker M=0$，与 \eqref{eq:kerdim} 及 $N\ge7$ 矛盾。
\end{proof}

\begin{remark}[与旧猜想的关系]\label{rmk:dichvsold}
猜想~\ref{conj:sixatom} 要求「每个 $N\ge7$ 的构型都存在一条全程 $Q_+\succeq0$ 的矩核直线」。
定理~\ref{thm:dichotomy} 把这句话\textbf{削弱并精确化}为：只需要排除分支 I 与分支 II 两种情形，
而这两个分支都是\textbf{可以在有限维上陈述和检验}的条件，
不再需要对「整条直线上的 $Q_+$ 半定性」这种全局性质作断言。这是实质性的进展。
\end{remark}

\begin{corollary}[分支 II 的球面二次曲线刚性]\label{cor:branch2rigid}
在分支 II 中，令 $B=Q_+\operatorname{diag}(x_j/g_j)$，则 $B_{ij}=\sqrt{g_i}\,h_j(\langle u_i,u_j\rangle)_+$，
且 $\mathcal K\subseteq\ker Q_+$ 等价于 $B\,a=0$ 对一切 $a\in\ker M$，即 $B$ 的\textbf{每一行属于} $\operatorname{rowspace}(M)$。
写成几何形式：对每个 $i$ 存在常数 $\mu_i$ 与对称 $3\times3$ 矩阵 $\Gamma_i$，使
\begin{equation}\label{eq:conic}
h_j\,\bigl(\langle u_i,u_j\rangle\bigr)_+\ =\ \mu_i+u_j^\top\Gamma_iu_j\qquad(\forall j),
\qquad\text{且}\qquad 3\mu_i+\tr\Gamma_i=\rho h_i .
\end{equation}
于是：二次函数 $u\mapsto\mu_i+u^\top\Gamma_iu$ 在球面上的零集（一条\textbf{球面二次曲线}）
必须\textbf{恰好包含} $u_i$ 的全部钝伙伴；而在锐伙伴处它必须取值 $h_j\langle u_i,u_j\rangle>0$。
\end{corollary}

\begin{proof}
$B\,a=0\ \forall a\in\ker M\Leftrightarrow$ 行空间条件。
第二式：\eqref{eq:conic} 两端乘 $g_j$ 求和，左端 $\sum_jg_jh_j(\langle u_i,u_j\rangle)_+=p_i=\rho h_i$，
右端 $3\mu_i+\tr\Gamma_i$（帧条件）。
\end{proof}

\begin{remark}[分支 II 就是「重合分裂」退化]\label{rmk:coincident}
若两个原子\textbf{重合}（$u_i=u_j$），则 $a=e_i-e_j\in\ker M$，
且 $\rho$ 沿之\textbf{恒等不变}——因为 $\rho$ 只依赖于测度 $\mu=\sum_ig_i\delta_{u_i}$
（第 15 轮的合并/分裂等谱性：合并保 $\rho$ 精确、分裂添一个零特征值）。
于是 $R|_{\mathcal K}$ 在这些方向上恰为 $0$，\textbf{分支 II 被精确实现}。
反过来，\eqref{eq:conic} 说明分支 II 中 $h_j(\langle u_i,u_j\rangle)_+$ 作为 $u_j$ 的函数
必须被一个二次多项式在原子集上插值；而 $\max(\langle u_i,u\rangle,0)$ 在 $\langle u_i,u\rangle=0$ 处有\textbf{折点}，
$\sqrt{u^\top\Lam u}$ 又不是多项式——除非原子大量重合，否则这种插值刚性极难满足。
\textbf{这给出了「分支 II 迫使原子重合」这条待证命题的精确形式}，是下一轮的队首之一。
\end{remark}

\subsection{数值：$N\ge7$ 极大元的精确结构}\label{sec:r17:num}

用命题~\ref{prop:gradformula} 的解析梯度 $\nabla_V\rho=2(xx^\top\!\circ K)V$ 作 Stiefel 切向投影
$G\leftarrow G-V\,\mathrm{sym}(V^\top G)$ 加 QR 收缩，从随机初值上升。结果：

\begin{center}\small
\begin{tabular}{cccccl}
\hline
$N$ & $\max\rho-\rho_0$ & $\min_ig_i$ & $\nu_-(Q_+)$ & $n_+(Q_+)$ & 极大元的质量分组 \\
\hline
5 & $-3.157\times10^{-2}$ & $2.04\times10^{-1}$ & 0 & 4 & （$\kappa_3(5)=\frac{9+4\sqrt6}{15}$） \\
6 & $-2.693\times10^{-2}$ & $2.61\times10^{-1}$ & 0 & 6 & 局部吸引域（非全局） \\
7 & $-3.678\times10^{-12}$ & $5.91\times10^{-11}$ & 0 & 6 & 冠军 $+$ 一个零质量行 \\
8 & $+1.110\times10^{-15}$ & $1.52\times10^{-1}$ & 0 & 6 & 冠军，其中 $2$ 个原子各作重合二分裂 \\
9 & $-5.400\times10^{-10}$ & $8.68\times10^{-9}$ & 1 & 6 & 冠军 $+$ 零质量行 $+$ 重合分裂 \\
10 & $-1.345\times10^{-10}$ & $3.48\times10^{-10}$ & 0 & 6 & 同上 \\
12 & $+8.882\times10^{-16}$ & $1.07\times10^{-2}$ & 0 & 6 & 冠军，$6$ 个方向作 $(3,3,2,1,2,1)$ 重合分裂 \\
\hline
\end{tabular}
\end{center}

$N=12$ 的分组是\textbf{精确}的（不是数值巧合）：
\[
\{0,3,7\},\ \{1,5,8\},\ \{2,6\},\ \{4\},\ \{11\}\ \text{各自}\ \textstyle\sum g=0.533646579188=g^\ast,
\qquad \{9,10\}:\ \textstyle\sum g=0.331767104061=g_0^\ast .
\]
即 $12$ 个原子只占据 $6$ 个不同方向，正是冠军测度 $\mu^\ast=g_0^\ast\delta_{u_0}+g^\ast\sum_{k=0}^4\delta_{u_k}$。

\textbf{分支 II 在这个极大元上被机器精度级地验证}：
\[
\max_{a\in\ker M\text{ 基}}|Ba|=3.91\times10^{-13},\qquad
\operatorname{eig}\bigl(R|_{\mathcal K}\bigr)=\{-3.0\times10^{-15},\dots,+2.6\times10^{-15}\}\ (\text{全为 }0,\ \dim\mathcal K=6).
\]
$N=8$ 的情形同样给出 $4$ 个 $g^\ast$ 分组、一个二分裂 $\{4,6\}$（和 $=0.533646579186$）
与一个 $g=4\times10^{-12}$ 的零质量行。

\begin{remark}[数值结论的强度]\label{rmk:numstrength}
以上\textbf{不是}「跑算法找极值」，而是用定理~\ref{thm:dichotomy} 预言的结构去\emph{检验}算法输出：
预言是「$N\ge7$ 的极大元必在分支 II（$R|_{\mathcal K}=0$）且必为冠军测度的重合分裂」，
检验结果是\textbf{完全吻合}，$g$ 的分组和精确到 $10^{-12}$。
因此数值在这里的作用是\textbf{证伪探测器}：如果分支 II 只能由重合原子实现，
那么 $\kappa_3(N)=\rho_0$（$N\ge6$）就与 $\kappa_3(6)=\rho_0$ 是同一件事，
全局归约就归结为 $N=6$ 的 $20$ 个符号类（第~\ref{sec:n6global} 节）与推论~\ref{cor:branch2rigid} 的刚性。
\end{remark}


\subsection{对偶容量形式的修正与强对偶}\label{sec:r17:dual}

把 $\kappa_3$ 写成「方向组 $U$ 固定、只对质量 $g$ 优化」的形式，会得到一个干净的对偶。

\begin{lemma}[对偶容量]\label{lem:phid}
对固定的 $U=(u_i)\subset S^2$ 与 $c\in\R^N$，令
$\Phi_U(c):=\inf\bigl\{\sum_ic_i^2/g_i:\ g>0,\ \sum_ig_iu_iu_i^\top=I_3\bigr\}$。则
\begin{equation}\label{eq:phid}
\Phi_U(c)=\max_{M\,:\,u_i^\top Mu_i\ge0\ \forall i}\ \frac{\bigl(\sum_i|c_i|\sqrt{u_i^\top Mu_i}\bigr)^2}{\tr M},
\end{equation}
右端是\textbf{尺度不变}的（$M\to tM$ 不变），因而是紧集 $\{M\succeq0,\tr M=1\}$ 上的凹函数最大化。
\end{lemma}

\begin{proof}
对任意满足 $u_i^\top Mu_i\ge0$ 的 $M$，令 $\psi_i=\sqrt{u_i^\top Mu_i}$。
AM--GM 给 $c_i^2/g_i+g_i\psi_i^2\ge2|c_i|\psi_i$；对 $i$ 求和并用帧条件
$\sum_ig_i\psi_i^2=\tr\bigl(M\sum_ig_iu_iu_i^\top\bigr)=\tr M$，得
$\sum c_i^2/g_i\ge2\sum|c_i|\psi_i-\tr M$。
对固定的方向 $M_0$ 再优化标度 $M=tM_0$（$2\sqrt tA-tB$ 的最大值是 $A^2/B$）即得 \eqref{eq:phid} 的「$\ge$」；
「$\le$」由 KKT：最优 $g_i=|c_i|/\psi_i$ 使 AM--GM 逐点取等。
\end{proof}

\begin{corollary}[$\kappa_3$ 的容量表述]\label{cor:phikappa}
$\kappa_3=\sup_{U,\,c\ne0}\dfrac{c^\top K_+(U)c}{\Phi_U(c)}$，其中 $K_+(U)_{ij}=(\langle u_i,u_j\rangle)_+$。
故 $\kappa_3\le\rho_0$ 等价于：对一切 $U$ 与 $c$，$c^\top K_+(U)c\le\rho_0\,\Phi_U(c)$。
\end{corollary}

\begin{proof}
$\lambda_{\max}(D_g^{1/2}K_+D_g^{1/2})=\max\{c^\top K_+c:\sum_ic_i^2/g_i=1\}$，再对 $g\in\mathcal G_U$ 取上确界。
\end{proof}

冠军处的数值：最优 $M=\Lam$（正是定理~\ref{thm:stiefelkkt} 的迹一乘子），$\tr M=1$，
$\psi_i=\sqrt{u_i^\top\Lam u_i}=h_i$，$\Phi=1$，$c^\top K_+c=\rho_0$，
\textbf{零对偶间隙}（primal $=0.999999978$，dual $=1.000000000$，差 $2.2\times10^{-8}$；
由 $g_i=c_i/\sqrt{t\,u_i^\top Mu_i}$ 复原的质量与 primal 最优质量之差 $2.25\times10^{-8}$）。
在 \texttt{random\_frame\_measure} 给出的可行构型（$N=5,\dots,8$）上，随机方向 $c$ 的强对偶相对差 $\le1.3\times10^{-4}$。

\begin{remark}[一处必须记录的可行性陷阱]\label{rmk:phitrap}
若 $U$ 的原子使 $I_3$ \textbf{不}在 $\{\sum_ig_iu_iu_i^\top:g\ge0\}$ 的内部（例如 $6$ 个随机方向），
则 primal 不可行、$\Phi_U(c)=+\infty$，对偶 \eqref{eq:phid} 也发散到 $+\infty$
（脚本上一度出现 $1.6\times10^{17}$）。这\textbf{不是}强对偶的反例，而是可行域为空。
上一版脚本还把 primal 目标误写成 $\sum c_i^2g_i$（应为 $\sum c_i^2/g_i$），
造成 $\Phi_{\rm primal}=0.2503$ 的假违背；两处均已修正，修正后 \S(F) 全绿。
\end{remark}

\subsection{钝图的团数与 Motzkin--Straus}\label{sec:r17:ms}

\begin{lemma}[钝团数]\label{lem:obtuseclique}
钝图 $G_-=(\{1,\dots,N\},\{\{i,j\}:\langle u_i,u_j\rangle<0\})$ 满足 $\omega(G_-)\le4$，且 $4$ 可达。
\end{lemma}

\begin{proof}
$\R^3$ 中两两成钝角的单位向量至多 $4$ 个（若 $u_1,\dots,u_5$ 两两钝，
则 Gram 矩阵 $G$ 满足 $G\mathbf 1$ 的各分量 $=1+\sum_{j\ne i}\langle u_i,u_j\rangle$，
而 $\rank G\le3$ 使 $G$ 有一个零空间向量；标准论证给 $k\le d+1$）。
可达性：数值扫描中确实出现 $\omega(G_-)=4$ 的构型（$360$ 个随机构型的最大值恰为 $4$）。
\end{proof}

\begin{corollary}[Motzkin--Straus 型界]\label{cor:ms}
记 $\varpi:=\sum_{\{i,j\}\text{ 钝}}x_i^2x_j^2$（无序对求和的两倍，即 $\sum_{i\sim j}x_i^2x_j^2$）。则
$\varpi\le\frac12\bigl(1-\frac1{\omega(G_-)}\bigr)\le\frac38$，且 $\beta\le\sqrt{e\,\varpi}\le\sqrt{3e/8}$。
\end{corollary}

\begin{proof}
以 $y_i=x_i^2$（$\sum y_i=1$）代入 Motzkin--Straus 定理；
$\beta^2\le e\varpi$ 是把 $\beta=\sum_{ij}(Q_-)_{ij}x_ix_j$ 与 $\|Q_-\|_F^2=e$、$\varpi$ 作 Cauchy--Schwarz。
\end{proof}

数值：$360$ 个构型上 $\max\varpi=0.359148483447<3/8=0.375$，$\beta^2-e\varpi\le0$ 无一违背。

\begin{remark}[这条路线的定量天花板（诚实记录）]\label{rmk:mslimit}
把 $\beta\le\sqrt{3e/8}$、$\rho^2\le\|Q_+\|_F^2=3-e$（引理~\ref{lem:frob}）与 $\tau\le\beta$ 联立：
$\tau\le\sqrt{3(3-(1+\tau)^2)/8}$，解得 $\tau\le0.5146$、$\rho\le1.5146$，
\textbf{比 Mohar 的 $1.36603$ 还差}。原因见注~\ref{rmk:taugap}：
Motzkin--Straus 只控制 $\beta$，而\textbf{抵消项 $-(1-\alpha)$ 完全没有被用上}。
本小节因此记为「已证但不足」，并作为「必须让 $(I-Q)x$ 参与」这一判断的定量佐证。
\end{remark}

\subsection{本轮状态、诚实标签与方法优先队列}\label{sec:r1517:status}

\textbf{已证（符号推导 $+$ 数值证书）}：
命题~\ref{prop:gradformula}（正部的环境梯度，勘误后重新确立）；
定理~\ref{thm:stiefelkkt}（Stiefel KKT 与迹一刚性）；
推论~\ref{cor:ellipsoid}（椭球面恒等式）、\ref{cor:hubfree}（枢纽无关式 $\rho^2=m^\top\Lam^{-1}m$ 与精确 CS 缺陷）；
引理~\ref{lem:tauform}（$\tau$-形式）与推论~\ref{cor:oneneg}、引理~\ref{lem:normqmx}；
引理~\ref{lem:frob}（能量守恒 $\|Q_+\|_F^2=3-e$，精确）、推论~\ref{cor:floor2}（一行式谱地板）；
引理~\ref{lem:framesplit}、\ref{lem:cs1}、\ref{lem:cs2}、定理~\ref{thm:betaAA}、命题~\ref{prop:betaAAeq}；
引理~\ref{lem:firstorder}、定理~\ref{thm:secondvar2}、\textbf{定理~\ref{thm:dichotomy}（全局二分）}、
推论~\ref{cor:branch2rigid}（分支 II 的球面二次曲线刚性）；
引理~\ref{lem:phid} 与推论~\ref{cor:phikappa}（对偶容量、强对偶）；
引理~\ref{lem:obtuseclique} 与推论~\ref{cor:ms}（$\omega(G_-)\le4$、$\varpi\le3/8$、$\beta\le\sqrt{3e/8}$）。

\textbf{条件性（已证但需附加假设）}：
命题~\ref{prop:majbound}（$Q_+\succeq0$ 且 $n_+\le4$ 时 $\rho\le1.290244<\Phi$）；
推论~\ref{cor:ellipsoid} 的 $\Lam\succ0$ 要求（冠军处成立：谱 $\{0.5946,0.2027^{[2]}\}>0$）。

\textbf{否证（同样是成果）}：
命题~\ref{prop:fourthfalse}——全局第四不等式 $\alpha+\sqrt{A_1A_2}\le\rho_0$ 为\textbf{假}，
$N=7$ 反例超出 $6.27\times10^{-3}$；
注~\ref{rmk:mslimit}——Motzkin--Straus 路线的定量天花板 $1.5146$，不足；
注~\ref{rmk:majcond}——重排路线在 $n_+=6$（冠军的实际层）只给 $1.474$，不足。

\textbf{数值（不作成果）}：
$N=5,\dots,12$ 的 Stiefel 梯度上升极大元及其质量分组；
分支 II 在 $N=12$ 极大元上的精确实现（$R|_{\mathcal K}=0$ 到 $3\times10^{-15}$）；
$360$ 个随机构型上的能量不等式扫描。

\textbf{方法优先队列（第七版，第 15--17 轮；已被第八版取代，见第~\ref{sec:r18:status}~节）}
\begin{enumerate}\itemsep3pt
\item \textbf{P1（队首，全局）——排除分支 II}：证明推论~\ref{cor:branch2rigid} 的插值刚性
\eqref{eq:conic} 在原子两两不重合时不可能成立（$N\ge7$）。
这是一道\textbf{有限维球面插值}题：一个二次多项式要在 $N$ 个点上等于
$\sqrt{u^\top\Lam u}\cdot(\langle u_i,u\rangle)_+$，后者的折点结构是主要武器。
\item \textbf{P2（全局）——压低 $n_+$ 到 $4$}：推论~\ref{cor:floor2} 说 $n_+\ge4$，
定理~\ref{thm:dichotomy}(i) 说 $n_+\le7$，命题~\ref{prop:majbound} 说 $n_+=4$ 即可完成「优于 \STW」。
缺口是 $n_+\in\{5,6,7\}$ 三层（冠军在 $n_+=6$）。
\item \textbf{P3（局部）——分支 I}：在 $\nu_-(Q_+)\ge1$ 的构型上直接证 $\rho\le\rho_0$。
可用 $\nu_-(Q_+)\le\nu_-(Q_-)$ 与 $\omega(G_-)\le4$（引理~\ref{lem:obtuseclique}）。
\item \textbf{P4（$N=6$ 收尾）}：第~\ref{sec:n6global} 节余下的 $20$ 个符号类。
\item \textbf{P5（副方向，下界）}：$\Lambda_3(G)\ge\rho_0n-3$ 的可达到性——冠军是 hub$+F_5$，
对应的图族已由定理~\ref{thm:Fn} 给出闭式谱，下界一侧只差「无穷族的渐近」这一步。
\end{enumerate}

\textbf{本轮把「唯一缺口」重写为}：定理~\ref{thm:dichotomy} 的分支 I 与分支 II 二者之一必须被排除。
分支 II 的形式已完全代数化（\eqref{eq:conic}），分支 I 的形式是一个负惯性指数条件。
两者都不再需要「整条矩核直线上 $Q_+\succeq0$」这种难以处理的断言。




\section{第 18 轮：钝性稀释恒等式与 $\Lam$-参数化（消去特征值的代数重构）}\label{sec:r18}

\subsection{动机：换一个正交的方向}\label{sec:r18:motiv}

第~\ref{sec:r1517}~节的全局二分定理~\ref{thm:dichotomy}~把缺口压缩为「排除分支 I 或分支 II」，
两个分支都已经是有限维的。但这两个分支都仍然建立在同一个对象之上：
$Q_+$ 的\textbf{谱}——$\lambda_{\max}(Q_+)$ 与它的 Perron 向量 $x$。
谱对象天然带着两个技术麻烦：一是 $\lambda_{\max}$ 只在单重时才光滑，
二是 Perron 向量由一个 $N$ 维特征方程隐式定义，$N$ 还不固定。

本轮换一个与之\textbf{正交}的方向：不再研究 $Q_+$ 的谱，而是设法把 $\rho=\lambda_{\max}(Q_+)$
本身改写成一个\textbf{不含特征值、不含 Perron 向量}的代数泛函，
其自变量只有几何数据 $(g_i,u_i)$ 与一个 $3\times3$ 的半正定矩阵 $\Lam$。
结果（定理~\ref{thm:lamparam}）是：主猜想等价于一族\textbf{五个额外参数}上的代数不等式，
且这个代数泛函在 $\Lam$ 上是\textbf{一次正齐次}的，数值上还是\textbf{凹}的
（命题~\ref{prop:concave}）。这使内层极大化成为一个凸规划，从而原则上可以给出
\textbf{对偶证书}（推论~\ref{cor:dualcert}），而不是只能跑算法。

沿路我们得到两个精确恒等式（引理~\ref{lem:deltautau}、引理~\ref{lem:dilution}）、
一个 $u$-KKT 的闭式解（命题~\ref{prop:uktkt}）、一个初等的结构定理
（定理~\ref{thm:transitive}，它解释了 $\rho_0$ 中为什么出现 $\sqrt5$），
以及\textbf{两个否证}（命题~\ref{prop:refA}、命题~\ref{prop:refB}），
它们精确指出「哪一步松弛一定会崩」。

\subsection{记号}\label{sec:r18:not}

沿用 \STW\ Theorem 4.1 的坐标并做如下固定记法。设 $N\ge3$，$g_i>0$，$u_i\in S^2$，
\begin{equation}\label{eq:frame18}
\sum_{i=1}^N g_iu_iu_i^\top=I_3\qquad(\text{从而 } \textstyle\sum_ig_i=3),
\end{equation}
记 $v_i:=\sqrt{g_i}\,u_i\in\R^3$，则 $\sum_iv_iv_i^\top=I_3$，$Q=VV^\top$ 的条目为
$Q_{ij}=v_i^\top v_j=\sqrt{g_ig_j}\langle u_i,u_j\rangle$，且 $Q$ 是秩三正交投影。
写
\[
(t)_+=\max(t,0),\qquad (t)_-=\max(-t,0),\qquad
P_{ij}:=(\langle u_i,u_j\rangle)_-,\qquad K_{ij}:=(\langle u_i,u_j\rangle)_+ .
\]
$Q_+=(Q)_{ij}=(\sqrt{g_ig_j}\langle u_i,u_j\rangle)_+=\sqrt{g_ig_j}K_{ij}$，
$Q_-=(Q_-)_{ij}=\sqrt{g_ig_j}P_{ij}$，$Q=Q_+-Q_-$，$Q_+=\tfrac12(Q+|Q|)$。
设 $\rho=\lambda_{\max}(Q_+)$ 单重，$x\ge0$，$\|x\|=1$，$Q_+x=\rho x$，并令
\begin{equation}\label{eq:hdef18}
h_i:=\frac{x_i}{\sqrt{g_i}}>0\qquad\Bigl(\text{即 } x_i=\sqrt{g_i}h_i,\quad \sum_ig_ih_i^2=1\Bigr).
\end{equation}
于是 Perron 方程等价于
\begin{equation}\label{eq:perron18}
\rho h_i=\sum_jg_jh_jK_{ij}=\langle u_i,m\rangle+d_i ,
\end{equation}
其中我们引入两个核心量：
\begin{equation}\label{eq:md18}
m:=\sum_{i}g_ih_iu_i\in\R^3,\qquad d_i:=\sum_jg_jh_jP_{ij}\ge0 .
\end{equation}
$m$ 是\textbf{矩向量}，$d=(d_i)$ 是\textbf{钝场}。由第~\ref{sec:r1517}~节的 $\tau$-形式（引理~\ref{lem:tauform}）：
\[
\alpha:=\|m\|^2=x^\top Qx,\qquad \beta:=x^\top Q_-x=\sum_ig_ih_id_i,\qquad
\rho=\alpha+\beta,\qquad \tau:=\rho-1=\beta-(1-\alpha).
\]
最后再引入一个\textbf{纯几何}量（不含 $h$）：
\begin{equation}\label{eq:Didef}
D_i:=\sum_jg_jP_{ij}=\sum_jg_j(\langle u_i,u_j\rangle)_-\ \ge0 ,
\end{equation}
它是原子 $i$ 处的「钝质量」。注意 $d_i$ 与 $D_i$ 的关系是 $d_i=\sum_jg_jh_jP_{ij}$：
$d$ 是 $D$ 的 $h$-加权版本。相应的钝图记为 $G_-$：$i\sim j\Leftrightarrow\langle u_i,u_j\rangle<0$。

\subsection{第一个精确恒等式：钝场向量恰为 $\tau m$}\label{sec:r18:delta}

\begin{lemma}[钝场向量恒等式]\label{lem:deltautau}
设 $\delta:=\sum_ig_id_iu_i$。则
\begin{equation}\label{eq:deltatau}
\boxed{\ \delta=\tau\,m\ }
\end{equation}
等价地，
\begin{align}
QQ_-x&=\tau\,Qx=\tau m, \label{eq:QQmx}\\
(I-Q)Q_-x&=\rho\,(I-Q)x, \label{eq:orthQmx}\\
\|Q_-x\|^2&=\tau^2\alpha+\rho^2(1-\alpha)=(1-\alpha)(1+2\tau)+\tau^2, \label{eq:normQmx2}\\
\langle m,\delta\rangle&=\alpha\tau . \label{eq:mdelta}
\end{align}
\end{lemma}

\begin{proof}
$Q_-=Q_+-Q$，故 $Q_-x=\rho x-Qx$。用 $Q^2=Q$：
$QQ_-x=\rho Qx-Qx=\tau Qx$，$(I-Q)Q_-x=\rho(I-Q)x-(I-Q)Qx=\rho(I-Q)x$，
这就是 \eqref{eq:QQmx}\eqref{eq:orthQmx}。又 $Qx=V^\top Vx$ 而
$(Vx)_k=\sum_i\sqrt{g_i}x_iu_{i,k}=m$ 的第 $k$ 分量乘以 $\sqrt{g_i}$ 的约定——直接算：
$\sum_i\sqrt{g_i}(Qx)_iu_i=\sum_i\sqrt{g_i}\sum_j\sqrt{g_ig_j}\langle u_i,u_j\rangle x_ju_i
=\sum_j x_j\sqrt{g_j}\sum_ig_i\langle u_i,u_j\rangle u_i=\sum_jx_j\sqrt{g_j}u_j=m$，
另一方面 $\sum_i\sqrt{g_i}(Q_-x)_iu_i=\sum_ig_id_iu_i=\delta$，
而 $\sum_i\sqrt{g_i}(Q_+x)_iu_i=\rho m$；由 $Q_-x=Q_+x-Qx$ 得 $\delta=\rho m-m=\tau m$，即 \eqref{eq:deltatau}。
\eqref{eq:normQmx2}：$Q_-x$ 沿 $\operatorname{im}Q$ 与 $\ker Q$ 正交分解，
由 \eqref{eq:QQmx}\eqref{eq:orthQmx} 得 $\|Q_-x\|^2=\tau^2\|Qx\|^2+\rho^2\|(I-Q)x\|^2$，
而 $\|Qx\|^2=x^\top Qx=\alpha$、$\|(I-Q)x\|^2=1-\alpha$。
展开 $\tau^2\alpha+\rho^2(1-\alpha)$ 并用 $\rho=1+\tau$ 即得引理~\ref{lem:normqmx}~的旧形式，
故 \eqref{eq:normQmx2} 是它的\textbf{勾股细化}：它把 $\|Q_-x\|^2$ 分成了「行空间内」与「行空间外」两块。
\eqref{eq:mdelta}~是 \eqref{eq:deltatau}~与 $\|m\|^2=\alpha$ 的直接推论。
\end{proof}

\begin{remark}[读法]\label{rmk:deltaread}
\eqref{eq:deltatau}~的几何含义是：把每个原子的钝场 $d_i$ 以质量 $g_i$ 堆到方向 $u_i$ 上，
得到的向量\textbf{必定}与矩向量 $m$ 平行，比例系数恰好是超出量 $\tau=\rho-1$。
因此 $\tau$ 有了一个纯几何的读法：
\begin{equation}\label{eq:tauread}
\tau=\frac{\bigl\|\sum_ig_id_iu_i\bigr\|}{\bigl\|\sum_ig_ih_iu_i\bigr\|}\qquad(m\ne0),
\end{equation}
即「钝场向量长度 / 矩向量长度」。这把主猜想变成
\[
\Bigl\|\sum_ig_id_iu_i\Bigr\|\le\tau_0\Bigl\|\sum_ig_ih_iu_i\Bigr\|,
\]
两端都是 $\R^3$ 中的向量长度，而不再是矩阵的最大特征值。
数值验证：冠军处 \eqref{eq:deltatau}~的残差为 $6.1\times10^{-16}$，
在 $800$ 个随机帧构型上最大残差 $2.9\times10^{-14}$
（\texttt{证书\_Lam参数化与稀释恒等式.py}~(A) 段）。
\end{remark}

\subsection{第二个精确恒等式：钝性稀释}\label{sec:r18:dil}

\begin{lemma}[钝性稀释恒等式]\label{lem:dilution}
定义
\begin{equation}\label{eq:Dhdef}
\mathcal D(h):=\sum_ig_ih_i^2D_i=\sum_{i\sim j}g_ig_j|\langle u_i,u_j\rangle|\,\frac{h_i^2+h_j^2}{2},
\qquad
\mathcal E_-(h):=\sum_{i\sim j}g_ig_j|\langle u_i,u_j\rangle|\,(h_i-h_j)^2 .
\end{equation}
则对\textbf{一切}构型与一切 $h\ge0$ 有精确等式
\begin{equation}\label{eq:dilution}
\boxed{\ \mathcal D(h)=\beta+\tfrac12\,\mathcal E_-(h)\ }\qquad\text{即}\qquad
\beta=\mathcal D(h)-\tfrac12\mathcal E_-(h)\ge \mathcal D(h)-\tfrac12\mathcal E_-(h),
\end{equation}
且 $\mathcal E_-(h)\ge0$，$\mathcal E_-(h)=0\Leftrightarrow h$ 在钝图 $G_-$ 的每个连通支上取常值。
特别地 $\beta\le\mathcal D(h)$（\textbf{钝性稀释不等式}），等号 $\Leftrightarrow\mathcal E_-=0$。
\end{lemma}

\begin{proof}
对固定的无序钝对 $\{i,j\}$，$\frac{h_i^2+h_j^2}2-h_ih_j=\frac{(h_i-h_j)^2}2$；
把这一等式乘以 $g_ig_j|\langle u_i,u_j\rangle|$ 并对所有钝对求和，
左端给出 $\mathcal D(h)$（因为 $\sum_{j\sim i}g_j|\langle u_i,u_j\rangle|=D_i$），
中项给出 $\beta=\sum_{i\sim j}g_ig_jh_ih_j|\langle u_i,u_j\rangle|$，右端给出 $\frac12\mathcal E_-(h)$。
\end{proof}

\begin{corollary}[$\tau$ 的势--能量分解]\label{cor:tauenergy}
对一切构型，
\begin{equation}\label{eq:tauenergy}
\tau=\underbrace{\bigl[\mathcal D(h)-(1-\alpha)\bigr]}_{\text{势项 }\Phi}
\;-\;\underbrace{\tfrac12\mathcal E_-(h)}_{\text{钝 Dirichlet 能量}} .
\end{equation}
\end{corollary}

\begin{remark}[为什么这条分解是有信息的]\label{rmk:dilread}
\eqref{eq:tauenergy}~把主猜想写成
\begin{equation}\label{eq:star}
\Phi=\mathcal D(h)-(1-\alpha)\ \le\ \tau_0+\tfrac12\mathcal E_-(h),
\end{equation}
左边是「势」：$\mathcal D(h)$ 对 $h^2$ 是\textbf{线性}的，$-(1-\alpha)$ 是 $-\|m\|^2+1$；
右边是「能量」：$h$ 在钝图上的加权 Dirichlet 能量。
这是一个 Poincar\'e 型不等式的形状，而不是谱不等式的形状。
它的关键优点是：$\mathcal D(h)$ 只通过 $h^2$ 依赖 $h$，因此一旦用上
$g$-KKT 的 $h_i^2=u_i^\top\tilde\Lam u_i$（第~\ref{sec:r1517}~节的 R15-2），
$\mathcal D(h)$ 就变成 $\tr(\tilde\Lam\Xi)$，其中
\begin{equation}\label{eq:Xidef}
\Xi:=\sum_ig_iD_iu_iu_i^\top\succeq0
\end{equation}
是一个\textbf{只依赖几何}的 $3\times3$ 矩阵。于是势项完全落在 $3\times3$ 矩阵的世界里。
冠军处：$\operatorname{eig}(\Xi)=\{0.379920,\,0.379920,\,0.253875\}$，$\tr\Xi=1.013716$，
$\tr(\tilde\Lam\Xi)=0.304973=\beta$（因为 $\mathcal E_-=0$），$\Phi=\tau_0$ 精确取等。
\end{remark}

\begin{lemma}[$\mathcal E_-$ 的第二显式]\label{lem:Em2}
设 $h$ 是 Perron 向量（即 \eqref{eq:perron18}~成立）。则
\begin{equation}\label{eq:Em2}
\mathcal E_-(h)=\frac1{\rho^2}\sum_{i\sim j}g_ig_j|\langle u_i,u_j\rangle|\,
\zeta_{ij}^2,\qquad
\zeta_{ij}:=\langle u_i-u_j,m\rangle+d_i-d_j .
\end{equation}
\end{lemma}

\begin{proof}
由 \eqref{eq:perron18}，$\rho h_i=\langle u_i,m\rangle+d_i$，故 $\rho(h_i-h_j)=\zeta_{ij}$。
\end{proof}

\begin{remark}\label{rmk:Em2read}
\eqref{eq:Em2}~的重要性在于：右端\textbf{不再显含} $h$，只含几何量 $u_i,g_i$ 与两个
辅助量 $m,d$（而 $m,d$ 又由 $h$ 决定）。它把「需要多少 Dirichlet 能量来补偿」这个问题
改写成了一个\textbf{显式的非负二次型}的下界问题，是下一轮 P1 的具体攻击目标。
数值验证：\eqref{eq:Em2}~在冠军处残差 $<10^{-15}$，在 $1200$ 个随机构型上最大残差 $2.1\times10^{-15}$。
\end{remark}

\subsection{$u$-KKT 的闭式解}\label{sec:r18:ukkt}

第~\ref{sec:r1517}~节以「R15-2 / R15-3」的形式记录了两个在冠军处数值成立的临界点关系。
本轮把它们\textbf{解出来}：乘子不是「存在」，而是有闭式。

\begin{proposition}[$u$-KKT 的显式形式]\label{prop:uktkt}
设 $(U,g)$ 是 $\rho$ 在帧切片 \eqref{eq:frame18}~上的临界点，$g_i>0$，$\lambda_{\max}(Q_+)$ 单重。
令 $W_i:=\sum_{j\in S_i^+}g_jh_ju_j$（$S_i^+=\{j:\langle u_i,u_j\rangle>0\}$）。
设 Lagrange 函数为
$\mathcal L=\rho-\langle\Lam,\sum_ig_iu_iu_i^\top-I_3\rangle-\sum_i\frac{\nu_i}2(\|u_i\|^2-1)$。
则：
\begin{enumerate}\itemsep2pt
\item[(i)] \textbf{球面乘子恒为零}：$\nu_i\equiv0$；
\item[(ii)] $\tilde\Lam u_i=\dfrac{h_i}{\rho}W_i$，其中 $\tilde\Lam:=\frac2\rho\Lam$；
\item[(iii)] \textbf{闭式}：
\begin{equation}\label{eq:lamclosed}
\tilde\Lam=\frac1\rho\sum_{i}\sum_{j\in S_i^+}g_ig_jh_ih_j\,u_iu_j^\top
=\frac1\rho\sum_{(i,j)\ \text{锐}}g_ig_jh_ih_j\,u_iu_j^\top ;
\end{equation}
\item[(iv)] \eqref{eq:lamclosed}~\textbf{自动对称}，$\tr\tilde\Lam=1$ 是恒等式（不依赖临界性），
且 $h_i^2=u_i^\top\tilde\Lam u_i$。
\end{enumerate}
\end{proposition}

\begin{proof}
$\partial\rho/\partial u_i=\sum_jx_ix_j\sqrt{g_ig_j}\,\one[\langle u_i,u_j\rangle>0]\,u_j
=x_i\sqrt{g_i}W_i=g_ih_iW_i$，
$\partial\mathcal L/\partial u_i=g_ih_iW_i-2g_i\Lam u_i-\nu_iu_i=0$，故
$h_iW_i=2\Lam u_i+\hat\nu_iu_i=\rho\tilde\Lam u_i+\hat\nu_iu_i$（$\hat\nu_i=\nu_i/g_i$）。
与 $u_i$ 作内积：左端 $h_i\langle u_i,W_i\rangle=h_i\sum_{j\in S_i^+}g_jh_j\langle u_i,u_j\rangle
=h_i\sum_jg_jh_jK_{ij}=h_i\cdot\rho h_i$（用 \eqref{eq:perron18}），
右端 $\rho u_i^\top\tilde\Lam u_i+\hat\nu_i$。而 $g$-KKT 给
$\frac{\partial\rho}{\partial g_i}=\frac{\rho x_i^2}{2g_i}=\frac{\rho h_i^2}2=\langle u_i,\Lam u_i\rangle$，
即 $u_i^\top\tilde\Lam u_i=h_i^2$。代入得 $\rho h_i^2=\rho h_i^2+\hat\nu_i$，故 $\hat\nu_i=0$，(i)(ii) 成立。
(iii)：$\tilde\Lam=\tilde\Lam\cdot I_3=\sum_ig_i\tilde\Lam u_iu_i^\top=\frac1\rho\sum_ig_ih_iW_iu_i^\top$，展开 $W_i$ 即得。
(iv)：\eqref{eq:lamclosed}~的和式在 $(i,j)$ 上的系数 $g_ig_jh_ih_j$ 对称，而「锐」关系本身对称，
故转置后仍是同一个和；$\tr\tilde\Lam=\frac1\rho\sum_ig_ih_i\langle u_i,W_i\rangle
=\frac1\rho\sum_ig_ih_i\cdot\rho h_i=\sum_ig_ih_i^2=1$；
$u_i^\top\tilde\Lam u_i=\frac{h_i}{\rho}\langle u_i,W_i\rangle=h_i^2$。
\end{proof}

\begin{remark}\label{rmk:ukktread}
$\nu_i\equiv0$ 是一个此前未记录的干净事实：\textbf{球面约束在临界点上是「松」的}，
全部约束信息都来自帧条件。这一点在数值上同样成立：冠军处
$|h_iW_i-\rho\tilde\Lam u_i|$ 的最大残差为 $7.2\times10^{-16}$，
\eqref{eq:lamclosed}~与最小二乘拟合出的 $\tilde\Lam$ 相差 $5.0\times10^{-16}$，
$\operatorname{eig}(\tilde\Lam)=\{0.594609,\,0.202696,\,0.202696\}$，$\lambda_{\min}=0.2027>0$。
\end{remark}

\subsection{主结果：$\Lam$-参数化}\label{sec:r18:main}

\begin{definition}[$\Lam$-泛函]\label{def:lamfunc}
对固定的几何 $(g,U)$ 满足 \eqref{eq:frame18}，与任意 $\Lam\succeq0$，令
\begin{equation}\label{eq:hlam}
h_i(\Lam):=\sqrt{u_i^\top\Lam u_i}=\|\Lam^{1/2}u_i\|,\qquad
x_i(\Lam):=\sqrt{g_i}\,h_i(\Lambda),
\end{equation}
并定义
\begin{equation}\label{eq:rholam}
\rho(\Lam):=x(\Lam)^\top Q_+x(\Lam)
=\Bigl\|\sum_ig_ih_i(\Lam)u_i\Bigr\|^2+\sum_{i,j}g_ig_jh_i(\Lam)h_j(\Lam)(\langle u_i,u_j\rangle)_- .
\end{equation}
\end{definition}

\begin{theorem}[$\Lam$-参数化：消去特征值的等价重构]\label{thm:lamparam}
\begin{enumerate}\itemsep3pt
\item[(i)] $\|x(\Lam)\|^2=\sum_ig_ih_i(\Lam)^2=\tr\Lam$；
$\rho(t\Lam)=t\,\rho(\Lam)$（\textbf{一次正齐次}）；
且当 $\tr\Lam\le1$ 时 $\rho(\Lam)\le\lambda_{\max}(Q_+)$。
\item[(ii)] 在任一临界点上（命题~\ref{prop:uktkt}），$\rho(\tilde\Lam)=\rho=\lambda_{\max}(Q_+)$，且
\begin{equation}\label{eq:gradlam}
\nabla_\Lam\rho(\Lam)\big|_{\Lam=\tilde\Lam}=\rho\,I_3 ,
\end{equation}
即 $\tilde\Lam$ 是 $\rho(\cdot)$ 在超平面 $\{\tr=1\}$ 上的\textbf{驻点}。
更一般地，$\nabla_\Lam\rho(\Lam)=\sum_i\frac{g_i[\langle u_i,m(\Lam)\rangle+d_i(\Lam)]}{h_i(\Lam)}u_iu_i^\top$，
故 \eqref{eq:gradlam}~成立当且仅当 $h(\Lam)$ 满足 Perron 方程 \eqref{eq:perron18}。
\item[(iii)] \textbf{等价重构}：
\begin{equation}\label{eq:kappalam}
\kappa_3=\sup_{(g,U)}\ \max_{\substack{\Lam\succeq0\\ \tr\Lam=1}}\ \rho_{g,U}(\Lam) .
\end{equation}
因此 $\kappa_3\le\rho_0$（等价地 $\Lam_3(G)\le\rho_0n-3$ 对一切 $G$）由下面这一条
\textbf{不含特征值、不含 Perron 向量}的代数不等式推出：
\begin{equation}\label{eq:starmain}
\Bigl\|\sum_ig_ih_iu_i\Bigr\|^2+\sum_{i,j}g_ig_jh_ih_j(\langle u_i,u_j\rangle)_-\ \le\ \rho_0 ,
\qquad h_i=\sqrt{u_i^\top\Lam u_i},\ \Lam\succeq0,\ \tr\Lam=1 .
\end{equation}
\end{enumerate}
\end{theorem}

\begin{proof}
(i) $\sum_ig_ih_i(\Lam)^2=\sum_ig_iu_i^\top\Lam u_i=\tr(\Lam\sum_ig_iu_iu_i^\top)=\tr\Lam$。
$h(t\Lam)=\sqrt t\,h(\Lam)$，而 \eqref{eq:rholam}~的右端对 $h$ 是二次的，故一次正齐次。
若 $\tr\Lam=1$ 则 $\|x(\Lam)\|=1$，于是 $\rho(\Lam)=x(\Lam)^\top Q_+x(\Lam)\le\lambda_{\max}(Q_+)$；
$\tr\Lam<1$ 时用齐次性。
(ii) 在临界点，$h_i(\tilde\Lam)^2=u_i^\top\tilde\Lam u_i=h_i^2$，故 $h(\tilde\Lam)=h$（两者皆正），
从而 $\rho(\tilde\Lam)=x^\top Q_+x=\rho$。梯度：$\partial h_i/\partial\Lam=u_iu_i^\top/(2h_i)$，
$\partial\rho(\Lam)/\partial h_i=2g_i\langle u_i,m(\Lam)\rangle+2\sum_jg_ig_jh_j(\Lam)P_{ij}
=2g_i[\langle u_i,m(\Lam)\rangle+d_i(\Lam)]$，链式法则给出所断言的公式；
在 $\Lam=\tilde\Lam$ 处方括号为 $\rho h_i$（\eqref{eq:perron18}），于是
$\nabla_\Lam\rho=\rho\sum_ig_iu_iu_i^\top=\rho I_3$。
(iii) 「$\le$」：由 (i)，对每个几何 $\max_{\tr\Lam=1}\rho(\Lam)\le\lambda_{\max}(Q_+)$，取上确界即得。
「$\ge$」：在 $\rho$ 的极大点（$g_i>0$，$\lambda_{\max}$ 单重）由 (ii) 得
$\max_\Lam\rho(\Lam)\ge\rho(\tilde\Lam)=\rho$；对 $g_i=0$ 或重根的退化情形，
用第~\ref{sec:r1517}~节的合并/分裂等谱性化归到 $g_i>0$ 且 $\lambda_{\max}$ 单重的情形
（合并保 $\rho$ 精确，分裂添一个零特征值），再用极限论证。
\end{proof}

\begin{remark}[这消除了什么]\label{rmk:lamparamread}
\eqref{eq:starmain}~与原来的问题 $\lambda_{\max}((VV^\top)_+)\le\rho_0$ 相比：
\begin{itemize}\itemsep2pt
\item 不再有 $N\times N$ 矩阵的特征值，也不再有 Perron 向量（$h$ 由 $\Lam$ 显式给出）；
\item 新增的自由度只有 $\Lam$ 的 $5$ 个（半正定、迹一）；
\item 约束集 $\{(g,u,\Lam): g\ge0,\ \sum g_iu_iu_i^\top=I_3,\ \Lam\succeq0,\ \tr\Lam=1\}$
是\textbf{半代数}的；把 $h_i\ge0$、$h_i^2=u_i^\top\Lam u_i$ 显式引入后，
目标函数在固定符号模式（即固定钝图 $G_-$）下是\textbf{多项式}的。
这为矩-SOS 层级或显式对偶证书留出了接口。
\end{itemize}
诚实标注：(iii) 中「$\ge$」一侧的退化情形化归沿用第~\ref{sec:r1517}~节的等谱性结论，
其细节（$g_i\to0$ 的极限与重根的分歧）在本报告中仍按\textbf{条件性}处理。
\end{remark}

\begin{proposition}[凹性，数值；\textbf{已在第~\ref{sec:r19}~节升级为定理~\ref{thm:geo}}]\label{prop:concave}
对固定几何，$\rho(\Lam)$ 在 $\Lam$ 上是\textbf{凹}的：在冠军几何上，
$600$ 组随机 $(\Lam_1,\Lam_2,s)$ 的最小凹性余量
$\rho((1-s)\Lam_1+s\Lam_2)-[(1-s)\rho(\Lam_1)+s\rho(\Lam_2)]$ 为 $+1.99\times10^{-5}$；
在 $10$ 个随机几何（$N=5,\dots,11$）上全局最小余量为 $+1.57\times10^{-5}$，无一为负。
\end{proposition}

\begin{remark}[第 19 轮回填]\label{rmk:concavefill}
本命题现在是\textbf{定理}。写 $L_i=g_iu_i^\top\Lam u_i=x_i(\Lam)^2$，则
$\rho(\Lam)=\sum_{ij}(Q_+)_{ij}\sqrt{L_iL_j}$，而 $\sqrt{ab}=\min_{t>0}\frac{ta+b/t}2$
是仿射族的下确界，故凹性由「线性 $\to$ 几何平均 $\to$ 非负组合」三步解析得到，
见定理~\ref{thm:geo}。同一表示还给出闭式梯度（推论~\ref{cor:gradpsi}）、
精确对偶（定理~\ref{thm:exactdual}）与内层规划的全局极大点刻画（推论~\ref{cor:argmax}）。
\end{remark}

\begin{remark}\label{rmk:concave}
（第 19 轮更新：凹性\textbf{已}解析证明，见定理~\ref{thm:geo}。下面保留当时的判断；
其中「不能由 $\nabla^2\rho$ 读出」仍然正确，正确做法是把 $\rho$ 写成
线性函数的几何平均的非负组合。）
注意凹性\textbf{不能}由 $\nabla^2\rho$ 的符号直接读出：
\eqref{eq:gradlam}~只在 Perron 点成立，不能对 $\Lam$ 再求导。
$\rho(\Lam)$ 作为一次正齐次凹函数具有支撑超平面表示
$\rho(\Lam)=\inf\{\langle\Lam,S\rangle:\ S\in\mathcal S\}$，
其中 $\mathcal S=\{S\in\mathrm{Sym}_3:\ \langle\Lam,S\rangle\ge\rho(\Lam)\ \forall\Lam\succeq0\}$，
于是由 Sion 极小极大定理：
\begin{equation}\label{eq:dualcert}
\max_{\substack{\Lam\succeq0\\\tr\Lam=1}}\rho(\Lam)=\min_{S\in\mathcal S}\ \lambda_{\max}(S).
\end{equation}
\end{remark}

\begin{corollary}[对偶证书形式]\label{cor:dualcert}
（第 19 轮更新：\eqref{eq:dualcert}~已由定理~\ref{thm:exactdual}~无条件证明，
且其中的 $\min$ 可限制在显式对偶族 $\{S(T):T>0\}$ 上，最小值在
$t^*_{ij}=x_j/x_i$ 处取到，此时 $S(T^*)=\rho I_3$。本推论的结论因此无条件成立；
但它同时被命题~\ref{prop:P1dead}~判为\textbf{无用}——冠军处可行证书只有 $\rho_0I_3$ 一个。）
则要证某几何满足 $\max_\Lam\rho(\Lam)\le\rho_0$，
\textbf{只需给出一个} $S\in\mathrm{Sym}_3$ 使得
\begin{equation}\label{eq:certcond}
\lambda_{\max}(S)\le\rho_0\qquad\text{且}\qquad \rho(\Lam)\le\langle\Lam,S\rangle\quad\forall\Lam\succeq0 .
\end{equation}
在冠军几何上这个证书是\textbf{被强制的且是紧的}：由 $\rho(\tilde\Lam)=\rho_0$ 与 $\tr\tilde\Lam=1$，
\eqref{eq:certcond}~给出 $\rho_0\le\langle\tilde\Lam,S\rangle\le\lambda_{\max}(S)\tr\tilde\Lam\le\rho_0$，
故必须处处取等，即 $\lambda_{\max}(S)=\rho_0$ 且 $\tilde\Lam$ 的支撑落在 $S$ 的最高特征空间里。
特别地 $S=\rho_0I_3$ 是唯一的「各向同性」候选，而它给出的条件
$\rho(\Lam)\le\rho_0\tr\Lam$ 恰与 \eqref{eq:starmain}~等价。
\end{corollary}

\begin{remark}\label{rmk:certsearch}
因此冠军几何上不存在「比切平面更好的」各向同性证书；要前进必须利用 $\tilde\Lam$ 的
特征方向 $u_0$（枢纽方向，$\tilde\Lam$ 的最高特征向量）来构造\textbf{各向异性}的 $S$，
形如 $S=\rho_0u_0u_0^\top+\kappa(I_3-u_0u_0^\top)$，$\kappa\le\rho_0$。
这是下一轮 P1 的具体构造任务。
\end{remark}

\begin{proposition}[尖锐性与刚性，数值]\label{prop:r18sharp}
\begin{enumerate}\itemsep2pt
\item[(i)] 冠军几何上 $\max_{\Lam\succeq0,\tr\Lam=1}\rho(\Lam)=1.281568196872167$，
与 $\rho_0=1.281568196872168$ 之差为 $-4.4\times10^{-16}$；
$26$ 组多起点（L-BFGS-B 与 Nelder--Mead 各一次）的极大点全部收敛到
$\operatorname{eig}=\{0.59460858,0.20269572,0.20269570\}$，与 $\tilde\Lam$ 的距离 $1.1\times10^{-8}$。
\item[(ii)] 在 $25$ 个几何上（Stiefel 梯度上升极大元 $N=6,7,9,10,12$；
$16$ 个随机帧构型 $N=5,\dots,12$；冠军的三档微扰 $\varepsilon=10^{-3},10^{-2},4\times10^{-2}$）
全部满足 $\max_\Lam\rho(\Lam)\le\rho_0$，全局最大值 $1.281568196584$（出现在 $N=10$ 的极大元上，
即冠军加零质量行）。
\item[(iii)] 在临界点上 $\max_\Lam\rho(\Lam)=\rho$ 精确成立
（$N=6$：两边同为 $-3.02\times10^{-2}$；$N=7$：同为 $-6.55\times10^{-2}$；
$N=9,10,12$：一致到 $10^{-9}$），
印证定理~\ref{thm:lamparam}(iii)~的「$\ge$」一侧。
\end{enumerate}
\end{proposition}

\subsection{一个初等的结构定理：可传递符号模式}\label{sec:r18:trans}

\begin{theorem}[可传递符号模式定理]\label{thm:transitive}
设存在 $\sigma\in\{\pm1\}^N$ 使
\begin{equation}\label{eq:transig}
\sigma_i\sigma_j\langle u_i,u_j\rangle\ge0\qquad(\forall i,j).
\end{equation}
则 $\rho=\lambda_{\max}(Q_+)\le1$。
逆否命题：\textbf{$\rho>1$ 蕴含符号模式不可传递}，
等价地，钝图 $G_-$ \textbf{不是完全二部图}（在把 $\langle u_i,u_j\rangle=0$ 的对任意归部之后）。
\end{theorem}

\begin{proof}
$x^\top|Q|x=\sum_{ij}x_ix_j|v_i^\top v_j|
=\sum_{ij}x_ix_j\sigma_i\sigma_jv_i^\top v_j
=\bigl\|\sum_i\sigma_ix_iv_i\bigr\|^2
\le\lambda_{\max}\bigl(\sum_iv_iv_i^\top\bigr)\,\|\sigma\odot x\|^2=1\cdot1=1$，
其中用到 $\sum_iv_iv_i^\top=I_3$ 与 $x\ge0$、$\|x\|=1$。
另一方面 $\alpha=x^\top Qx=\|V^\top x\|^2\le\lambda_{\max}(VV^\top)\|x\|^2=1$。
由 $Q_+=\tfrac12(Q+|Q|)$（逐项）得 $\rho=\tfrac12\alpha+\tfrac12x^\top|Q|x\le1$。
\end{proof}

\begin{remark}[为什么这解释了 $\sqrt5$]\label{rmk:transread}
冠军的钝图是 $G_-=C_5\sqcup\{\text{枢纽}\}$：枢纽与所有原子成锐角（$D_0=0$），
五个五边形原子之间恰有「次近邻」为钝，构成一个 $5$-圈。
$C_5$ \textbf{不是}完全二部图，故定理~\ref{thm:transitive}~不适用，$\rho_0>1$ 得以发生；
而且 $C_5$ 是\textbf{最小的}非二部钝图（$C_3$ 在 $S^2$ 上不可能全部两两成钝角，
因为三个两两钝角的单位向量需要 $\sum g_iu_iu_i^\top$ 的秩结构与之相容——
引理~\ref{lem:obtuseclique}~已给出 $\omega(G_-)\le4$，而 $C_5$ 的团数为 $2$）。
这给出 $\rho_0=\frac{9+\sqrt5+2\sqrt{35}}{18}$ 中 $\sqrt5$ 的\textbf{结构解释}：
它来自 $A(C_5)$ 的谱 $\{2,\frac{\sqrt5-1}2,\frac{\sqrt5-1}2,-\frac{\sqrt5+1}2,-\frac{\sqrt5+1}2\}$，
第~\ref{sec:r1517}~节已验证 $\operatorname{spec}(Q_-)=2g^*|s_2|\cdot\operatorname{spec}(A(C_5))$。
数值验证：$662\to193$ 个严格可传递构型上 $\max\rho=1+1.3\times10^{-10}$
（超出量与帧残差 $6.9\times10^{-11}$ 同阶，即纯数值误差）；
$2208$ 个 $\rho>1$ 的随机构型中可传递者 $0$ 个。
\end{remark}

\subsection{两个否证：哪一步松弛一定会崩}\label{sec:r18:ref}

按本协议，否证也是成果。本轮我们试验了两条看起来最有希望的松弛，二者\textbf{都是假的}，
而且失败方式恰好指明了正确路线。

\begin{proposition}[否证 R18-A：自由 $\Lam$ 松弛失效]\label{prop:refA}
考虑把 \eqref{eq:tauenergy}~中的 Dirichlet 项丢掉、并让 $\Lam$ 在整个
$\{\Lam\succeq0,\tr\Lam=1\}$ 上自由变化的松弛
\begin{equation}\label{eq:phifree}
\Phi_{\mathrm{free}}(g,U):=\max_{\substack{\Lam\succeq0\\\tr\Lam=1}}
\Bigl[\tr(\Lam\Xi)+\Bigl\|\sum_ig_ih_i(\Lam)u_i\Bigr\|^2-1\Bigr].
\end{equation}
则\textbf{即使在冠军几何上}也有
\[
\Phi_{\mathrm{free}}=0.302149011933>\tau_0=0.281568196872\qquad(\text{超出 }+2.058\times10^{-2}),
\]
极大元是\textbf{秩一}的（$\operatorname{eig}=\{1,1.7\times10^{-14},-1.4\times10^{-17}\}$）。
因此「$\Lam$ 是 KKT 乘子」这一信息\textbf{不可省略}：
只有把 $\Lam$ 限制到满足 $u$-KKT（命题~\ref{prop:uktkt}(ii)）的那一点，
才能得到 $\Phi=\tau_0$。
\end{proposition}

\begin{proposition}[否证 R18-B：强稀释不等式失效]\label{prop:refB}
不等式
\begin{equation}\label{eq:strongdil}
\alpha+\mathcal D(h)\le\rho_0\qquad(h=h_{\mathrm{Perron}})
\end{equation}
是\textbf{假}的。一个 $N=4$ 的帧构型给出 $\alpha+\mathcal D(h)=1.462928395681$，
超出 $\rho_0$ 达 $+1.814\times10^{-1}$。但该构型的真实 $\rho=0.798522010308\ll\rho_0$，
其补偿项为 $\tfrac12\mathcal E_-(h)=0.664406385373$，
恰好把 $\alpha+\mathcal D(h)-\tfrac12\mathcal E_-$ 拉回 $0.798522=\rho$。
因此\textbf{钝 Dirichlet 能量 $\mathcal E_-$ 是不可或缺的补偿项}，
任何丢弃它的路线都注定失败——这与第~\ref{sec:r1517}~节「单独界 $\beta$ 的路线全部失败」
是同一现象的两个侧面：$\mathcal E_-$ 与 $-(1-\alpha)$ 是主猜想中仅有的两个「负项」。
\end{proposition}

\begin{remark}\label{rmk:refsafe}
两个反例的真实 $\rho$ 都远低于 $\rho_0$（分别为 $0.798522$ 与冠军几何上的 $\rho_0$ 本身），
所以被否证的是\textbf{路线}，主猜想~\ref{conj:A1}~未受任何影响。
它们的价值在于把「哪一项不能丢」变成了定理级的事实。
\end{remark}

\subsection{本轮状态、诚实标签与方法优先队列（第八版）}\label{sec:r18:status}

\textbf{已证（解析，含完整证明）}：
引理~\ref{lem:deltautau}（钝场向量 $=\tau m$ 及三个推论式）；
引理~\ref{lem:dilution}（钝性稀释恒等式，含等号刻画）；
推论~\ref{cor:tauenergy}（$\tau$ 的势--能量分解，精确）；
引理~\ref{lem:Em2}（$\mathcal E_-$ 的第二显式）；
命题~\ref{prop:uktkt}（$u$-KKT 闭式解、$\nu_i\equiv0$、$\tr\tilde\Lam=1$ 恒等）；
定理~\ref{thm:lamparam}(i)(ii)（$\Lam$-泛函的齐次性、$\rho(\tilde\Lam)=\rho$、$\nabla\rho=\rho I_3$）；
定理~\ref{thm:transitive}（可传递符号模式 $\Rightarrow\rho\le1$）。

\textbf{条件性}：定理~\ref{thm:lamparam}(iii)~的「$\ge$」一侧对退化情形（$g_i=0$、重根）
的化归；推论~\ref{cor:dualcert}~依赖命题~\ref{prop:concave}~的凹性。
\emph{（第 19 轮更新：凹性已解析证明，故推论~\ref{cor:dualcert}~无条件；
「$\ge$」一侧在 $x>0$ 时亦无条件，见推论~\ref{cor:argmax}。）}

\textbf{数值（不作成果）}：
命题~\ref{prop:concave}（凹性，冠军 $+$ $10$ 个随机几何）；
命题~\ref{prop:r18sharp}（尖锐性、刚性、$25$ 个几何的扫描）；
引理~\ref{lem:deltautau}\eqref{eq:dilution}\eqref{eq:Em2}~的随机验证
（$800/1500/1200$ 个构型，残差 $\le2.9\times10^{-14}$）；
定理~\ref{thm:transitive}~的 $193+2208$ 个构型检验。

\textbf{否证}：命题~\ref{prop:refA}（自由 $\Lam$ 松弛在冠军几何上超出 $\tau_0$ 达 $2.06\times10^{-2}$）；
命题~\ref{prop:refB}（强稀释不等式在 $N=4$ 上超出 $\rho_0$ 达 $1.81\times10^{-1}$）。

\textbf{方法优先队列（第八版；P2 已由定理~\ref{thm:geo}~完成，P1 已由命题~\ref{prop:P1dead}~否证并弹出，现行队列为第~\ref{sec:r19:status}~节的第九版）}
\begin{enumerate}\itemsep3pt
\item \textbf{P1（队首，全局）——构造对偶证书 $S$}：由推论~\ref{cor:dualcert}，
只需对每个几何给出 $S\in\mathrm{Sym}_3$ 满足 \eqref{eq:certcond}。
由注~\ref{rmk:certsearch}，冠军处证书必须是紧的且 $\tilde\Lam$ 落在 $S$ 的最高特征空间，
故候选族为 $S=\rho_0u_0u_0^\top+\kappa(I_3-u_0u_0^\top)$，$\kappa\le\rho_0$；
第一步是确定 $\kappa$ 的可行区间（一个 $3\times3$ 的 LMI 问题），
第二步是把 $\kappa$ 的表达式推广到一般几何（预期与 $\Xi$、$\tr\Xi$、$D_i$ 的分布有关）。
这是\textbf{纯数学}任务，不需要对每阶图跑算法。
\item \textbf{P2（全局）——解析证明 $\rho(\Lam)$ 的凹性}：
$\rho(\Lam)=\sum_k\lambda_k(Q_+)\langle w_k,x(\Lam)\rangle^2$，
而 $x_i(\Lam)=\sqrt{g_iu_i^\top\Lam u_i}$ 是线性函数的平方根（凹）。
可尝试的路径：(a) 逐项证明 $\langle w,x(\Lam)\rangle^2$ 在 $w\ge0$ 时凹；
(b) 用 \eqref{eq:dualcert}~的极小极大直接构造 $S$ 绕过凹性；
(c) 计算 $\nabla^2\rho$ 在 $\tilde\Lam$ 处的显式并证明其在该点负半定（局部凹性）。
\item \textbf{P3（全局）——$\mathcal E_-$ 的下界}：用引理~\ref{lem:Em2}~的显式
$\zeta_{ij}=\langle u_i-u_j,m\rangle+d_i-d_j$，在 \eqref{eq:star}~中给出
$\frac12\mathcal E_-\ge\Phi-\tau_0$ 的直接证明。这是把否证 R18-B 变正的路径。
\item \textbf{P4（局部）——第~\ref{sec:r1517}~节遗留的分支 I / 分支 II}：
分支 II 的排除（推论~\ref{cor:branch2rigid}~的插值刚性）；分支 I 经
$\nu_-(Q_+)\le\nu_-(Q_-)$ 与 $\omega(G_-)\le4$（引理~\ref{lem:obtuseclique}）。
注意定理~\ref{thm:transitive}~给出分支 I 的一个新入口：若 $\nu_-(Q_+)\ge1$ 则 $\rho$ 的
符号结构必然不可传递，$G_-$ 不是完全二部。
\item \textbf{P5（$N=6$ 收尾）}：第~\ref{sec:n6global}~节余下的 $20$ 个符号类。
\item \textbf{P6（副方向，下界）}：$\Lam_3(G)\ge\rho_0n-3$ 的可达到性，
冠军族已由定理~\ref{thm:Fn}~给出闭式谱，只差无穷族的渐近。
\end{enumerate}

\textbf{本轮把主线改写成}：定理~\ref{thm:lamparam}~的 \eqref{eq:starmain}。
这是一个\textbf{五个额外参数}上的代数不等式，目标函数在固定钝图符号模式后是多项式的，
约束集是半代数的，并且（数值上）对 $\Lam$ 凹。
与第~\ref{sec:r1517}~节的二分定理相比，它\textbf{完全不含}「$N\times N$ 矩阵的谱」这一对象，
因此更适合用凸对偶、SOS 或显式不等式来处理；而命题~\ref{prop:refA}\ref{prop:refB}
已经排除了两条最自然的错误走法。


\section{第 19 轮：精确对偶、对偶证书的不可能性，与 Bloch 截面分解}\label{sec:r19}

\subsection{本轮要交代的三件事}\label{sec:r19:motiv}

第~\ref{sec:r18}~节把主线改写成 \eqref{eq:starmain}，并留下三个未了项：
命题~\ref{prop:concave}~的凹性当时只有数值证据；
推论~\ref{cor:dualcert}~因此是条件性的；
第八版队列的队首 P1（构造各向异性对偶证书 $S$）还没有动手。
本轮的处理是：凹性给出\textbf{解析证明}（定理~\ref{thm:geo}），
并顺带得到闭式梯度（推论~\ref{cor:gradpsi}）；
在此基础上证明 $\Lam$-参数化的对偶是\textbf{精确紧}的（定理~\ref{thm:exactdual}），
于是 P1 不是一个「还没做成」的任务，而是一个\textbf{不可能做成}的任务
（命题~\ref{prop:P1dead}：冠军几何上唯一可行的证书就是 $\rho_0I_3$）；
最后给出一个新的结构分解（定理~\ref{thm:section}），
把主线从「一个维数无界的帧流形上的上确界」改写成
「一个\textbf{二维单纯形}上的一族子问题」，并且指出冠军恰好落在其中的
\textbf{单轴}子族上。

顺带记录两条被定量否证的放缩族（命题~\ref{prop:ceilweyl}\ref{prop:ceilharm}），
它们分别给出 $\varphi=1.618034$ 与 $1.585642$ 的天花板，
都\textbf{劣于}报告已有的 $\Phi=\frac{3+\sqrt5}4=1.309017$，故只作为死路记录，不作为上界成果。

\subsection{凹性的解析证明与闭式梯度}\label{sec:r19:conc}

凹性在第~\ref{sec:r18}~节只有数值证据。它的解析证明其实只需要一行恒等式，
而这一行恒等式同时给出对偶族的显式形式，是本轮一切结论的起点。

\begin{lemma}[几何平均的仿射下确界表示]\label{lem:geo}
对 $a,b\ge0$，
\begin{equation}\label{eq:geoam}
\sqrt{ab}=\min_{t>0}\frac{ta+b/t}{2},
\end{equation}
极小值在 $t=\sqrt{b/a}$ 处取到（$a=0$ 或 $b=0$ 时下确界为 $0$，在 $t\to\infty$ 或 $t\to0$ 时逼近）。
故 $(a,b)\mapsto\sqrt{ab}$ 是 $\R_+^2$ 上一族仿射函数的逐点下确界，从而是\textbf{凹}的。
\end{lemma}

\begin{proof}
AM--GM 给 $\sqrt{ab}\le\frac{ta+b/t}2$，取 $t=\sqrt{b/a}$ 时等号成立。
一族仿射函数的下确界是凹函数，这是凸分析的标准事实。
\end{proof}

\begin{theorem}[$\rho(\Lam)$ 的凹性与几何平均表示]\label{thm:geo}
令 $D_i:=g_iu_iu_i^\top$、$L_i(\Lam):=\langle\Lam,D_i\rangle=g_iu_i^\top\Lam u_i=x_i(\Lam)^2$。
则
\begin{equation}\label{eq:geo}
\rho(\Lam)=\sum_{i,j}(Q_+)_{ij}\sqrt{L_i(\Lam)L_j(\Lam)}
=\inf_{T>0}\ \langle\Lam,S(T)\rangle ,
\end{equation}
其中 $T=(t_{ij})$ 取遍正对称矩阵，
\begin{equation}\label{eq:ST}
S(T):=\frac12\sum_{i,j}(Q_+)_{ij}\bigl(t_{ij}D_i+t_{ij}^{-1}D_j\bigr)\in\mathrm{Sym}_3 ,
\end{equation}
且当 $x(\Lam)>0$ 时下确界是\textbf{最小值}，在 $t^*_{ij}=x_j(\Lam)/x_i(\Lam)$ 处取到。
特别地 $\rho$ 在 $\{\Lam\succeq0\}$ 上\textbf{凹}且一次正齐次。
\end{theorem}

\begin{proof}
由 \eqref{eq:hlam}，$x_i(\Lam)=\sqrt{g_i}h_i(\Lam)=\sqrt{g_iu_i^\top\Lam u_i}=\sqrt{L_i(\Lam)}$，
代入 \eqref{eq:rholam}~的 $\rho(\Lam)=x(\Lam)^\top Q_+x(\Lam)$ 即得 \eqref{eq:geo}~的第一个等号。
对每一对 $(i,j)$ 用引理~\ref{lem:geo}（权 $(Q_+)_{ij}\ge0$）：
\[
\sum_{ij}(Q_+)_{ij}\sqrt{L_iL_j}
=\inf_{T>0}\frac12\sum_{ij}(Q_+)_{ij}\bigl(t_{ij}L_i+t_{ij}^{-1}L_j\bigr)
=\inf_{T>0}\Bigl\langle\Lam,\ \frac12\sum_{ij}(Q_+)_{ij}(t_{ij}D_i+t_{ij}^{-1}D_j)\Bigr\rangle ,
\]
这里用了 $L_i=\langle\Lam,D_i\rangle$ 的线性性。$L_i$ 是 $\Lam$ 的线性函数，
$\sqrt{\cdot\,\cdot}$ 凹，非负系数求和保凹，故 $\rho$ 凹；
一次正齐次性已见于定理~\ref{thm:lamparam}(i)。
\end{proof}

\begin{remark}[凹性为什么此前看起来「不好证」]\label{rmk:concwhy}
第~\ref{sec:r18}~节的注~\ref{rmk:concave}~指出凹性不能由 $\nabla^2\rho$ 的符号读出，
因为 \eqref{eq:gradlam}~只在 Perron 点成立。定理~\ref{thm:geo}~说明正确的看法是：
不要去算二阶导，而是先把 $\rho$ 写成\textbf{线性函数的几何平均的非负组合}。
这个表示的代价是引入了 $N(N+1)/2$ 个辅助参数 $t_{ij}$，
但回报是 $\rho$ 变成了一个\textbf{下确界型}（因此凹、因此可对偶）的泛函。
数值复核：冠军几何上 \eqref{eq:geo}~两种写法的最大差为 $6.7\times10^{-16}$，
$400$ 组随机 $(\Lam_1,\Lam_2,s)$ 的凹性余量最小值 $+1.2\times10^{-4}$，无一为负。
\end{remark}

\begin{corollary}[闭式梯度恒半正定]\label{cor:gradpsi}
在 $x(\Lam)>0$ 处 $\rho$ 可微，且
\begin{equation}\label{eq:gradpsi}
\nabla_\Lam\rho(\Lam)=S(T^*(\Lam))=\sum_i\frac{(Q_+x(\Lam))_i}{x(\Lam)_i}\,D_i
=\sum_i\frac{g_i[\langle u_i,m(\Lam)\rangle+d_i(\Lam)]}{h_i(\Lam)}u_iu_i^\top\ \succeq\ 0 .
\end{equation}
即梯度\textbf{恒落在锥} $\operatorname{cone}\{D_1,\dots,D_N\}$ \textbf{中}，
并且满足 Euler 恒等式 $\langle\Lam,\nabla_\Lam\rho(\Lam)\rangle=\rho(\Lam)$。
\end{corollary}

\begin{proof}
由 \eqref{eq:geo}~与 Danskin 定理（下确界在 $T^*$ 处唯一取到，因 AM--GM 的取等点唯一），
$\nabla_\Lam\rho=S(T^*(\Lam))$。把 $t^*_{ij}=x_j/x_i$ 代入 \eqref{eq:ST}：
\[
S(T^*)=\frac12\sum_i D_i\frac1{x_i}\sum_j(Q_+)_{ij}x_j+\frac12\sum_j D_j\frac1{x_j}\sum_i(Q_+)_{ij}x_i
=\sum_i\frac{(Q_+x)_i}{x_i}D_i ,
\]
第二步用 $Q_+$ 对称。$D_i\succeq0$ 与 $(Q_+x)_i/x_i\ge0$ 给出半正定。
第二个等号：$(Q_+)_{ij}=\sqrt{g_ig_j}(\langle u_i,u_j\rangle)_+$ 与 $x_j=\sqrt{g_j}h_j$ 给出
$(Q_+x)_i/\sqrt{g_i}=\sum_jg_jh_j(\langle u_i,u_j\rangle)_+
=\langle u_i,m\rangle+d_i$（把 $(\langle u_i,u_j\rangle)_+$ 拆成
$\langle u_i,u_j\rangle+(\langle u_i,u_j\rangle)_-$），再除以 $x_i=\sqrt{g_i}h_i$。
Euler 恒等式是一次正齐次性的标准推论。
\end{proof}

\begin{remark}[这条梯度公式的用处]\label{rmk:gradpsi}
其一，它把 \eqref{eq:gradlam}~从「临界点上的巧合」提升为「处处成立的闭式」，
因此\textbf{不需要}特征值或 Perron 向量就能求梯度，数值上可直接驱动内层凸规划
（本轮证书用它做 L-BFGS，$15$ 个起点全部收敛到 $\rho$，超出量 $\le1.3\times10^{-15}$）。
其二，$\nabla\rho(\Lam)\succeq0$ 说明 $\rho$ 关于 Loewner 序\textbf{单调}：
$\Lam_1\succeq\Lam_2\Rightarrow\rho(\Lam_1)\ge\rho(\Lam_2)$。
其三，$\nabla\rho$ 恒在 $\operatorname{cone}\{D_i\}$ 中，而 $\sum_iD_i=I_3$，
故 $\tr\nabla_\Lam\rho(\Lam)=\sum_ig_i(Q_+x(\Lam))_i/x(\Lam)_i$，
这是一个只涉及帧数据的显式标量；它与 $\rho(\Lam)$ 的关系由 Euler 恒等式给出。
\end{remark}

\subsection{精确对偶：$\Lam$-参数化没有松弛间隙，最优证书恒为各向同性}\label{sec:r19:dual}

\begin{theorem}[精确对偶]\label{thm:exactdual}
设 $x>0$（即 $Q_+$ 不可约；一般情形按锐图的分块取最大值化归）。令
$t^*_{ij}:=x_j/x_i$，其中 $x$ 是 $Q_+$ 的 Perron 向量。则
\begin{equation}\label{eq:stariso}
S(T^*)=\rho\,I_3 ,
\end{equation}
并且
\begin{equation}\label{eq:dualtight}
\max_{\substack{\Lam\succeq0\\ \tr\Lam=1}}\rho(\Lam)
\ =\ \inf_{T>0}\lambda_{\max}\!\bigl(S(T)\bigr)
\ =\ \lambda_{\max}(Q_+)\ =\ \rho .
\end{equation}
等式 \eqref{eq:stariso}~对\textbf{每一个}几何都成立，不需要任何临界点或极值假设。
\end{theorem}

\begin{proof}
\eqref{eq:stariso}：如推论~\ref{cor:gradpsi}~的证明，$S(T^*)=\sum_i\frac{(Q_+x)_i}{x_i}D_i$。
$x$ 是 $Q_+$ 的 Perron 向量，故 $Q_+x=\rho x$，于是 $\frac{(Q_+x)_i}{x_i}=\rho$，
$S(T^*)=\rho\sum_iD_i=\rho I_3$（帧条件 \eqref{eq:frame18}）。
\eqref{eq:dualtight}：三段夹逼。
(a) 弱对偶：对一切 $T$ 与一切 $\Lam\succeq0$，由 \eqref{eq:geo}~有
$\rho(\Lam)\le\langle\Lam,S(T)\rangle\le\lambda_{\max}(S(T))\tr\Lam$，
故 $\sup_{\tr\Lam=1}\rho(\Lam)\le\inf_T\lambda_{\max}(S(T))$。
(b) 由 \eqref{eq:stariso}，$\inf_T\lambda_{\max}(S(T))\le\lambda_{\max}(S(T^*))=\rho$。
(c) 由定理~\ref{thm:lamparam}(i)，$\sup_{\tr\Lam=1}\rho(\Lam)\le\lambda_{\max}(Q_+)=\rho$。
三者合起来迫使全部取等。
\end{proof}

\begin{remark}[这条定理说的是「坏消息」]\label{rmk:dualtight}
\eqref{eq:dualtight}~说明 $\Lam$-参数化\textbf{不产生任何松弛间隙}：
它的对偶最优值恰好是 $\lambda_{\max}(Q_+)$，也就是原问题本身。
这是好事也是坏事。好事是：定理~\ref{thm:lamparam}(iii)~的等价重构是\textbf{精确}的，
不存在「凸化之后丢掉信息」的隐患，因此 \eqref{eq:starmain}~确实值得作为主线。
坏事是：对偶证书路线\textbf{不能}给出比 $\lambda_{\max}(Q_+)$ 更好的东西——
下一小节把这句话变成对 P1 的否证。
\end{remark}

\begin{corollary}[内层规划是凸的，$\tilde\Lam$ 是全局极大点]\label{cor:argmax}
设 $x>0$。则对一切 $\Lam\succeq0$，
\begin{equation}\label{eq:tangent}
\rho(\Lam)\ \le\ \rho\cdot\tr\Lam ,
\end{equation}
从而 $\max_{\Lam\succeq0,\tr\Lam=1}\rho(\Lam)=\rho(\tilde\Lam)=\rho$。
于是定理~\ref{thm:lamparam}(iii)~的「$\ge$」一侧在 $x>0$ 时\textbf{无条件}成立
（第~\ref{sec:r18}~节里它被列为条件性结论）。
极大点集恰为仿射面
\begin{equation}\label{eq:face}
\mathcal F=\{\Lam\succeq0:\ \langle\Lam,D_i\rangle=x_i^2,\ i=1,\dots,N\},
\end{equation}
它退化为单点 $\{\tilde\Lam\}$ 当且仅当 $\{D_i\}_{i=1}^N$ 在 $\mathrm{Sym}_3$ 中线性无关
（等价地，那个 $N\times6$ 矩阵的秩为 $6$）。
\end{corollary}

\begin{proof}
$\rho$ 凹（定理~\ref{thm:geo}）且在 $\tilde\Lam$ 处可微、$\nabla\rho(\tilde\Lam)=\rho I_3$
（定理~\ref{thm:lamparam}(ii)），故一阶条件给
$\rho(\Lam)\le\rho(\tilde\Lam)+\langle\rho I_3,\Lam-\tilde\Lam\rangle=\rho+\rho(\tr\Lam-1)$。
对 $\tr\Lam\le1$ 得 $\rho(\Lam)\le\rho$；一般 $\Lam$ 用一次正齐次性化到 $\tr\Lam=1$，
即得 \eqref{eq:tangent}。
取等条件：\eqref{eq:geo}~中等号要求对每一对 $(Q_+)_{ij}>0$ 都有
$t^*_{ij}=\sqrt{L_j/L_i}$，即 $L_i/L_j=x_i^2/x_j^2$；锐图连通故存在 $c>0$ 使 $L_i=cx_i^2$。
而 $\tr\Lam=\sum_iL_i=c\sum_ix_i^2=c$，于是 $\tr\Lam=1$ 时 $c=1$，$L_i=x_i^2$，即 $\Lam\in\mathcal F$。
反过来 $\Lam\in\mathcal F$ 显然取等。$\mathcal F$ 是 $\mathrm{Sym}_3$ 中一个仿射子空间与半正定锥之交，
该仿射子空间由 $\{D_i\}$ 的正交补决定，故为单点当且仅当 $\{D_i\}$ 张成 $\mathrm{Sym}_3$。
\end{proof}

\begin{remark}[数值复核，以及一个意外发现]\label{rmk:argmax}
$\|\nabla\rho(\tilde\Lam)-\rho I_3\|_{\max}=1.6\times10^{-15}$；
$\rho(\tilde\Lam)-\rho=1.3\times10^{-15}$；
以 \eqref{eq:gradpsi}~为梯度的 L-BFGS 从 $15$ 个起点出发，最大超出量 $1.3\times10^{-15}$，
从未越过 $\rho$。切向界 \eqref{eq:tangent}~在 $960$ 个随机 $\Lam$（含 $\tr\Lam\ne1$）上最大违背 $0$。
\textbf{意外发现}：极大点\textbf{不总是唯一}的。$N=6$ 的一个上升极大元处
$\operatorname{rank}\{D_i\}=4<6$，此时 $\mathcal F$ 是一个二维仿射面，
L-BFGS 收敛到的极大点与 $\tilde\Lam$ 相距 $1.2\times10^{-1}$ 而\textbf{取值完全相同}（差 $1.1\times10^{-15}$）；
$N=7,9,10$ 处 $\operatorname{rank}\{D_i\}=6$，距离分别为 $0$、$1.2\times10^{-5}$、$2.1\times10^{-10}$。
这印证了 \eqref{eq:face}~的秩判据：\textbf{值}是刚性的，\textbf{极大点}只在 $\{D_i\}$ 线性相关时才不刚性。
\end{remark}

\subsection{各向异性对偶证书不存在：P1 被否证}\label{sec:r19:p1dead}

第八版队列的队首 P1 是：按注~\ref{rmk:certsearch}，构造形如
$S=\rho_0u_0u_0^\top+\kappa(I_3-u_0u_0^\top)$ 的\textbf{各向异性}证书。
本小节证明这件事在冠军几何上\textbf{不可能}，因此 P1 必须弹出。
证明只需要一条关于 $\tilde\Lam$ 的精确谱信息。

\begin{lemma}[冠军的 $\tilde\Lam$ 是单轴的，且正定]\label{lem:lampos}
冠军几何在绕枢纽轴的二面体群 $D_5$ 作用下不变，故 $\tilde\Lam$ 与之交换，从而
\begin{equation}\label{eq:uniax}
\tilde\Lam=\lambda_\parallel\,u_{\rm hub}u_{\rm hub}^\top+\lambda_\perp\bigl(I_3-u_{\rm hub}u_{\rm hub}^\top\bigr),
\qquad
\lambda_\parallel=h_{\rm hub}^2=\frac{\pi_{\rm hub}}{g_0^*},\quad
\lambda_\perp=\frac{1-h_{\rm hub}^2}{2},
\end{equation}
其中 $\pi_{\rm hub}=x_{\rm hub}^2=\dfrac{91-25\sqrt7}{126}=0.197271565265$，
$g_0^*=\dfrac{9-4\sqrt5+\sqrt{35}}{18}=0.331767104061$，故
$\lambda_\parallel=0.594608575866$、$\lambda_\perp=0.202695712067$。
特别地 $\tilde\Lam\succ0$，且这两个不等式是\textbf{精确}的：
$0<h_{\rm hub}^2<1$ 等价于 $28+28\sqrt5<25\sqrt7+7\sqrt{35}$，两边平方得 $6272<6090+350\sqrt{35}$，显然。
\end{lemma}

\begin{proof}
$\tilde\Lam$ 由命题~\ref{prop:uktkt}~的闭式 \eqref{eq:lamclosed}~从 $(g,U)$ 规范地定义，
故与 $D_5$ 作用交换。$D_5$ 在 $\R^3$ 上的表示分解为「轴上一维平凡」$\oplus$「轴外二维不可约」
（$n=5\ge3$ 时二维表示不可约），由 Schur 引理，交换代数中的元素在这两个不可约分量上是标量，
即 \eqref{eq:uniax}。$\lambda_\parallel=u_{\rm hub}^\top\tilde\Lam u_{\rm hub}=h_{\rm hub}^2$
（命题~\ref{prop:uktkt}(iv) 的 $u_i^\top\tilde\Lam u_i=h_i^2$）；
$\tr\tilde\Lam=1$（同处已证）给 $\lambda_\parallel+2\lambda_\perp=1$。
$h_{\rm hub}^2=x_{\rm hub}^2/g_{\rm hub}=\pi_{\rm hub}/g_0^*$ 是定义。
\end{proof}

\begin{proposition}[各向异性证书不存在]\label{prop:P1dead}
在冠军几何上，满足 \eqref{eq:certcond}~的 $S\in\mathrm{Sym}_3$ \textbf{只有一个}：$S=\rho_0I_3$。
因此推论~\ref{cor:dualcert}~的证书路线\textbf{等价于}原始不等式 \eqref{eq:starmain}，
不携带任何新信息；注~\ref{rmk:certsearch}~提出的各向异性候选族是空的。
\end{proposition}

\begin{proof}
设 $S$ 满足 \eqref{eq:certcond}。则
\[
\rho_0=\rho(\tilde\Lam)\ \le\ \langle\tilde\Lam,S\rangle\ \le\ \lambda_{\max}(S)\,\tr\tilde\Lam
\ =\ \lambda_{\max}(S)\ \le\ \rho_0 ,
\]
第一式是定理~\ref{thm:lamparam}(ii)~的 $\rho(\tilde\Lam)=\rho=\rho_0$，
第三式是引理~\ref{lem:lampos}~的 $\tr\tilde\Lam=1$。故三处全部取等。
由 $\langle\tilde\Lam,S\rangle=\lambda_{\max}(S)\tr\tilde\Lam$ 知 $\tilde\Lam$ 的像含于 $S$ 的
$\lambda_{\max}(S)=\rho_0$ 特征空间；而引理~\ref{lem:lampos}~给 $\tilde\Lam\succ0$，
其像为全空间，故该特征空间 $=\R^3$，$\lambda_{\max}$ 的重数为 $3$，$S=\rho_0I_3$。
\end{proof}

\begin{remark}[否证的定量内容]\label{rmk:p1dead}
这条否证有两层定量含义。
\textbf{第一层}：$S$ 的可行集是\textbf{单点}，不是「很小」；
任何试图调节 $\kappa$ 的做法都会立刻破坏 $\rho(\Lam)\le\langle\Lam,S\rangle$。
数值上，$3000$ 次随机各向异性尝试（$\lambda_{\max}(S)\le\rho_0$）中最好的一个仍然违反
证书条件达 $-4.2\times10^{-1}$。
\textbf{第二层}（更本质）：由定理~\ref{thm:exactdual}~的 \eqref{eq:stariso}，
\textbf{每个}几何的 AM--GM 取等参数都是 $t^*_{ij}=x_j/x_i$，都给出 $S=\rho I_3$。
所以对偶族 $\{S(T)\}$ 里「好的那一个」永远是各向同性的；
各向异性只能来自\emph{非}取等的 $T$，而那样 $\lambda_{\max}(S(T))>\rho$，不可能 $\le\rho_0$
（在冠军处 $\rho=\rho_0$）。
\textbf{结论}：凸对偶这条路在冠军几何上已经走到头，
主线必须回到不等式 \eqref{eq:starmain}~本身，或者改变分解方式。
\end{remark}

\subsection{两条被定量否证的放缩族}\label{sec:r19:ceil}

按协议，被排除的路线必须留下定量天花板。本轮排除两条自然但无效的族。

\begin{proposition}[Weyl$+$Frobenius 自举的天花板]\label{prop:ceilweyl}
记 $Q=VV^\top$、$Q_\pm=(\pm Q)_+$（故 $Q_+=Q+Q_-$），
$\alpha:=x^\top Qx=\|V^\top x\|^2$、$\beta:=x^\top Q_-x\ge0$，
则 $\rho=x^\top Q_+x=\alpha+\beta$。只用下面四条数据
\begin{equation}\label{eq:bootdata}
\tr Q_+=\textstyle\sum_ig_i=3,\qquad
\|Q_+\|_F^2+\|Q_-\|_F^2=\|Q\|_F^2=\tr(Q^2)=3,\qquad
\alpha\le\lambda_{\max}(Q)=1,\qquad \beta\le\lambda_{\max}(Q_-)\le\|Q_-\|_F ,
\end{equation}
最多只能推出
\begin{equation}\label{eq:boot}
\rho\ \le\ \frac{1+\sqrt5}{2}=\varphi=1.618033988750 ;
\end{equation}
若再加上 $N=6$ 时的两条迹精细化
$\|Q_+\|_F^2\ge\rho^2+\frac{(3-\rho)^2}{5}$（因 $\tr Q_+=3$）与
$\|Q_-\|_F^2\ge\frac65\lambda_{\max}(Q_-)^2$（因 $\tr Q_-=0$），则最多推出 $\rho\le\frac32$。
两者都\textbf{劣于}报告已有的 $\Phi=\frac{3+\sqrt5}4=1.309017$，故此族死。
\end{proposition}

\begin{proof}
$\rho=\alpha+\beta$ 来自 $Q_+=Q+Q_-$ 的逐项分解。由 \eqref{eq:bootdata}，
$\rho\le1+\beta\le1+\|Q_-\|_F$ 且 $\rho^2\le\|Q_+\|_F^2=3-\|Q_-\|_F^2$，
消去 $\|Q_-\|_F$ 得 $(\rho-1)^2\le3-\rho^2$，即 $\rho^2-\rho-1\le0$，故 $\rho\le\varphi$。
$N=6$ 版本：$3\ge\rho^2+\frac{(3-\rho)^2}5+\frac65(\rho-1)^2$，
整理为 $2.4\rho^2-3.6\rho\le0$，即 $\rho\le\frac32$。
天花板是\textbf{可达}的（不是估计松紧的问题）：把 \eqref{eq:bootdata}~当约束做线性/二次规划，
$(\alpha,\beta)=(1,\varphi-1)$ 与 $\|Q_-\|_F^2=3-\rho^2$ 恰好相容，故 $\varphi$ 是该族的真上确界。
数值：$9$ 个阶数（$N=4,\dots,14$）的上升极大元全部满足 $(\rho-1)^2\le3-\rho^2$，
最大余量 $-1.28$；该不等式本身远未饱和，说明损失发生在 \eqref{eq:bootdata}~的第四条。
\end{proof}

\begin{proposition}[Legendre/优占族的天花板]\label{prop:ceilharm}
把核 $(t)_+$ 在 $[-1,1]$ 上按 Legendre 多项式展开，
\begin{equation}\label{eq:legc}
(t)_+=\sum_{\ell\ge0}c_\ell P_\ell(t),\qquad
c_\ell=\frac{2\ell+1}{2}\int_0^1tP_\ell(t)\dd t ,
\end{equation}
则 $c_0=\frac14$、$c_1=\frac12$、$c_2=\frac5{16}$、$c_3=0$、$c_4=-\frac3{32}$、$c_5=0$、
$c_6=\frac{13}{256}$、$c_8=-\frac{17}{512}$、$c_{10}=\frac{49}{2048}$、$c_{12}=-\frac{75}{4096}$，
且 \textbf{$c_\ell=0$ 对一切奇数 $\ell\ge3$}（因 $(t)_+=\frac t2+\frac{|t|}2$，而 $|t|$ 只有偶次项）。
令 $A_\ell:=\bigl(\sqrt{g_ig_j}\,P_\ell(\langle u_i,u_j\rangle)\bigr)_{ij}$，
则每个 $A_\ell\succeq0$（$P_\ell(\langle u,v\rangle)$ 是 $S^2$ 上的正定核，加法公式），
$P_+=\sum_\ell c_\ell A_\ell$，并且
\begin{equation}\label{eq:A01}
\lambda_{\max}(A_0)=\sum_ig_i=3,\qquad A_1=Q,\quad \lambda_{\max}(A_1)=1 .
\end{equation}
于是逐项放缩给出
\[
\rho\le\sum_{\ell:\,c_\ell>0}c_\ell\lambda_{\max}(A_\ell),
\]
而右端在冠军几何上截断到 $\ell=24$ 时等于 $1.585642$，比 $\rho_0$ 高 $0.304074$。
故这一族（含第~\ref{sec:nogo}~节的优占矩阵方案）不可能达到 $\rho_0$。
\end{proposition}

\begin{proof}
$c_\ell$ 由 \eqref{eq:legc}~直接积分（本轮用符号计算逐个求值，全部是有理数）。
奇数 $\ell\ge3$ 时 $tP_\ell(t)$ 是偶函数且 $\int_{-1}^1tP_\ell(t)\dd t=\int_0^1tP_\ell+\int_{-1}^0tP_\ell=0$
（用 $P_\ell$ 奇、$tP_\ell$ 偶），而 $\int_{-1}^1tP_\ell=0$ 对 $\ell\ne1$ 由正交性给出，
故 $\int_0^1tP_\ell=0$。$A_\ell\succeq0$ 与 $P_+=\sum_\ell c_\ell A_\ell$ 是逐条目的恒等式。
$A_0=\sqrt g\sqrt g{}^\top$ 给 $\lambda_{\max}=3$；$A_1=(\sqrt{g_ig_j}\langle u_i,u_j\rangle)=VV^\top=Q$。
\end{proof}

\begin{remark}[这条否证解释了一个数字]\label{rmk:harm54}
\eqref{eq:A01}~里 $\lambda_{\max}(A_0)=3$ 与 $\lambda_{\max}(A_1)=1$ 是\textbf{精确}的，
于是「只保留 $\ell\le1$」的逐项天花板是 $\frac14\cdot3+\frac12\cdot1=\frac54$，
恰好等于已证的 $\kappa_3(4)=\frac54$。这不是巧合式的装饰：
它说明 $\frac54$ 这一层信息完全来自核 $(t)_+$ 的\textbf{常数项与线性项}，
也就是「各向同性 $+$ 投影本身」这两件事；
$\rho_0-\frac54=0.031568$ 必须由 $\ell\ge2$ 的高次调和项提供。
但在冠军处这些高阶项的实际贡献是 $+0.051727$（其中 $\ell=6$ 一项就占 $+0.061846$，
$\ell=4,8,12$ 三项共 $-0.021199$），而且级数收敛很慢
（截断到 $\ell=24$ 时重构值 $1.279057$，误差 $-2.5\times10^{-3}$）。
\textbf{这就是逐项放缩失败的机制}：正负项交错、单项幅度大于总和，
任何「丢掉负项」的动作都会立刻越过 $\rho_0$。
\end{remark}

\subsection{Bloch 截面分解：主线的新形状}\label{sec:r19:sect}

到目前为止，主线一直是「在一个维数随 $N$ 增长的帧流形上求上确界」。
定理~\ref{thm:lamparam}~引入了 $5$ 个新参数（$\Lam\succeq0$、$\tr\Lam=1$，即一个三维 Bloch 体）。
本小节指出：因为帧集合在 $O(3)$ 下不变，\textbf{可以先对 $\Lam$ 的谱分类，再对帧取上确界}，
于是主线变成一个\textbf{二维单纯形}上的族。

\begin{theorem}[截面分解]\label{thm:section}
对 $\mu=(\mu_1,\mu_2,\mu_3)\in\Delta_2:=\{\mu_i\ge0,\ \sum_i\mu_i=1\}$（降序排列）定义
\begin{equation}\label{eq:section}
\kappa_3(\mu):=\sup_{N\ge3}\ \sup_{(g,U)\ \text{满足}\ \eqref{eq:frame18}}\ \rho_{g,U}(\operatorname{diag}\mu) .
\end{equation}
则
\begin{equation}\label{eq:secteq}
\kappa_3=\sup_{\mu\in\Delta_2}\kappa_3(\mu),
\end{equation}
且 $\kappa_3(R\Lam R^\top)=\kappa_3(\Lam)$ 对一切 $R\in O(3)$，
故 $\kappa_3(\mu)$ 只依赖 $\mu$ 作为\textbf{无序}三元组，即 $\kappa_3$ 是 $\Delta_2$ 上 $S_3$-不变的函数。
此外冠军满足
\begin{equation}\label{eq:mustar}
\mu(\tilde\Lam)=\Bigl(h_{\rm hub}^2,\ \tfrac{1-h_{\rm hub}^2}2,\ \tfrac{1-h_{\rm hub}^2}2\Bigr)
=(0.594608575866,\ 0.202695712067,\ 0.202695712067),
\end{equation}
是\textbf{单轴}的（后两个分量相等），且 $\kappa_3(\mu(\tilde\Lam))\ge\rho_0$。
\end{theorem}

\begin{proof}
乘积集上的上确界可换序：
$\sup_{(g,U)}\max_{\Lam}\rho_{g,U}(\Lam)=\sup_{(g,U),\Lam}\rho_{g,U}(\Lam)=\sup_{\Lam}\sup_{(g,U)}\rho_{g,U}(\Lam)$。
帧条件 \eqref{eq:frame18}~在 $u_i\mapsto Ru_i$ 下不变，且
$h_i(g,RU,\Lam)=\sqrt{(Ru_i)^\top\Lam(Ru_i)}=h_i(g,U,R^\top\Lam R)$、
$\langle Ru_i,Ru_j\rangle=\langle u_i,u_j\rangle$，故
$\rho_{g,RU}(\Lam)=\rho_{g,U}(R^\top\Lam R)$。
把 $R$ 换为 $R^\top$ 并对 $U$ 取上确界（$U\mapsto R^\top U$ 是帧集到自身的双射）即得不变性。
于是 $\kappa_3(\Lam)$ 是 $O(3)$-不变量，只依赖 $\operatorname{spec}(\Lam)$；
取 $\Lam=\operatorname{diag}\mu$ 得 \eqref{eq:secteq}。
\eqref{eq:mustar}~是引理~\ref{lem:lampos}；$\rho_{g^*,U^*}(\tilde\Lam)=\rho_0$
（定理~\ref{thm:lamparam}(ii)），故冠军是 $\mu(\tilde\Lam)$ 截面上的可行点。
\end{proof}

\begin{corollary}[两个特殊截面的显式形式]\label{cor:sectspecial}
\begin{enumerate}\itemsep3pt
\item[(i)] \textbf{中心截面} $\mu=(\frac13,\frac13,\frac13)$：$h_i\equiv\frac1{\sqrt3}$，故
\begin{equation}\label{eq:center}
\kappa_3\Bigl(\tfrac13,\tfrac13,\tfrac13\Bigr)=\frac13\sup_{(g,U)}\sum_{i,j}g_ig_j\bigl(\langle u_i,u_j\rangle\bigr)_+ ,
\end{equation}
即「各向同性测度 $\mu=\sum_ig_i\delta_{u_i}$ 的正部二次能量」除以 $3$。
此时 $\Lam$ 的自由度\textbf{完全消失}，问题变成一个纯测度论的极值问题。
\item[(ii)] \textbf{秩一截面} $\mu=(1,0,0)$，$\Lam=ww^\top$：$h_i=|\langle u_i,w\rangle|$，
令 $a_i:=\sqrt{g_i}\langle u_i,w\rangle$，则 $\sum_ia_i^2=\|w\|^2=1$ 且 $x_i(\Lam)=|a_i|$，故
$\rho_{g,U}(\Lam)=\sum_{ij}(Q_+)_{ij}|a_i||a_j|$。
\end{enumerate}
\end{corollary}

\begin{proposition}[截面地图，数值]\label{prop:sectmap}
对 $8$ 个截面 $\mu$ 与 $N\in\{6,8,10\}$，用 \eqref{eq:gradpsi}~之外的直接梯度上升
（对 $V\in\mathrm{St}(N,3)$，解析梯度）求 $\kappa_3(\mu)$ 的数值下界，得
\begin{center}\small
\begin{tabular}{lrrr}
\toprule
截面 $\mu$ & $N=6$ & $N=8$ & $N=10$ \\
\midrule
中心 $(\frac13,\frac13,\frac13)$ & $1.233669$ & $1.256733$ & $1.248877$ \\
单轴 $a=0.50$ & $1.249734$ & $1.249025$ & $1.245552$ \\
单轴 $a=0.5946=h_{\rm hub}^2$ & $1.250166$ & $1.270937$ & $1.273641$ \\
单轴 $a=0.70$ & $1.247226$ & $1.240466$ & $1.278518$ \\
单轴 $a=0.85$ & $1.209759$ & $1.222580$ & $1.261757$ \\
秩二 $(\frac12,\frac12,0)$ & $1.228572$ & $1.181850$ & $1.182550$ \\
秩一 $(1,0,0)$ & $1.150694$ & $1.175367$ & $1.166373$ \\
非对称 $(0.60,0.30,0.10)$ & $1.231907$ & $1.261213$ & $1.260940$ \\
\bottomrule
\end{tabular}
\end{center}
其中「单轴 $a$」指 $\mu=(a,\frac{1-a}2,\frac{1-a}2)$。
把冠军本身作为 $a=h_{\rm hub}^2$ 截面的可行点加入后，该截面的值为 $1.281568196872167=\rho_0$。
\textbf{观察}：（1）没有任何截面的数值上确界超过 $\rho_0$；
（2）最大值出现在单轴族 $a\approx0.6$--$0.7$，冠军恰在其中；
（3）秩一与秩二截面明显更低（最大 $1.228572$，余量 $0.052996$）；
（4）非对称截面 $(0.60,0.30,0.10)$ 也不超过 $1.2613$。
这些都是\textbf{数值}结论，不作为成果。
\end{proposition}

\begin{remark}[这条分解买到了什么]\label{rmk:sectread}
\textbf{买到三件事}。
其一，$\Lam$ 的 $8$ 个自由度被压缩成 $\Delta_2$ 的 $2$ 个：
旋转自由度是\textbf{冗余}的（定理~\ref{thm:section}~的不变性），
这是一个此前没有明确使用过的事实。
其二，$\Delta_2$ 的\textbf{边界}（秩 $\le2$ 的截面）与\textbf{中心}截面各自成为独立的、
形式上更简单的子问题；边界截面数值上余量很大（$\ge0.053$），
中心截面按 \eqref{eq:center}~根本不含 $\Lam$，因此最有希望先被解析解决。
其三，冠军落在\textbf{单轴}子族上，这不是偶然：
引理~\ref{lem:lampos}~说明单轴性是 $D_5$ 对称性的后果，
而冠军的极值性本身就来自 $D_5$ 对称。
\textbf{代价}是：$\kappa_3(\cdot)$ 是「一族凹函数的上确界」，本身既不凹也不凸，
所以不能用凹性把它压到某个特殊截面上；
排除非对称极大元（$S_3$-轨道大小为 $6$ 的那些）需要新的论证。
\end{remark}

\begin{corollary}[归约的精确形式]\label{cor:sectreduce}
主猜想 $\kappa_3\le\rho_0$ 等价于：对每个 $\mu\in\Delta_2$，$\kappa_3(\mu)\le\rho_0$。
由 $S_3$-对称性，只需对下列三类 $\mu$ 证明：
\textbf{(A)} 中心 $\mu=(\frac13,\frac13,\frac13)$；
\textbf{(B)} 单轴 $\mu=(a,\frac{1-a}2,\frac{1-a}2)$，$a\in[\frac13,1]$；
\textbf{(C)} 真正的非对称极大元（若存在）。
若能证明 $\kappa_3(\cdot)$ 在 $\Delta_2$ 上的极大元必落在 (A)$\cup$(B) 中，
则主猜想化归为\textbf{一个单参数族}。
\end{corollary}

\subsection{本轮状态、诚实标签与方法优先队列（第九版）}\label{sec:r19:status}

\textbf{已证（解析，含完整证明）}：
引理~\ref{lem:geo}（几何平均的仿射下确界表示）；
定理~\ref{thm:geo}（$\rho(\Lam)$ 的凹性与 $\inf_T\langle\Lam,S(T)\rangle$ 表示）；
推论~\ref{cor:gradpsi}（闭式梯度、恒半正定、Loewner 单调、Euler 恒等式）；
定理~\ref{thm:exactdual}（$S(T^*)=\rho I_3$ 对每个几何成立；对偶精确紧）；
推论~\ref{cor:argmax}（切向界 \eqref{eq:tangent}、$\tilde\Lam$ 是全局极大点、
极大点面 \eqref{eq:face}~及其秩判据；定理~\ref{thm:lamparam}(iii)~的「$\ge$」一侧在 $x>0$ 时无条件）；
引理~\ref{lem:lampos}（冠军 $\tilde\Lam$ 单轴且正定，$0<h_{\rm hub}^2<1$ 的精确证明）；
命题~\ref{prop:P1dead}（各向异性对偶证书不存在）；
命题~\ref{prop:ceilweyl}（Weyl$+$Frobenius 自举的天花板 $\varphi$ / $\frac32$）；
命题~\ref{prop:ceilharm}（Legendre 系数精确值、$A_\ell\succeq0$、$\lambda_{\max}(A_0)=3$、$\lambda_{\max}(A_1)=1$、
逐项天花板 $1.585642$）；
定理~\ref{thm:section}（截面分解）；推论~\ref{cor:sectspecial}（中心与秩一截面的显式形式）；
推论~\ref{cor:sectreduce}（三类归约）。

\textbf{条件性}：定理~\ref{thm:exactdual}~与推论~\ref{cor:argmax}~假设 $x>0$
（即 $Q_+$ 不可约 / 锐图连通）；一般情形的化归（按分块取最大值）在本报告中按
\textbf{条件性}处理，与第~\ref{sec:r18}~节 $g_i=0$、重根的退化情形是同一类技术债务。
推论~\ref{cor:sectreduce}~的 (C) 类（排除非对称极大元）尚未解决。

\textbf{数值（不作成果）}：
命题~\ref{prop:sectmap}（$8$ 个截面 $\times3$ 个阶数的地图）；
定理~\ref{thm:exactdual}~的 $13$ 个几何复核（$|S(T^*)-\rho I_3|_{\max}\le9.1\times10^{-15}$，
弱对偶最大违背 $0$）；推论~\ref{cor:argmax}~的 L-BFGS 复核（超出 $\le1.3\times10^{-15}$）；
命题~\ref{prop:ceilweyl}~的 $9$ 个阶数复核；命题~\ref{prop:ceilharm}~的 $\ell\le24$ 截断复核。
证书脚本：\texttt{证书\_凹性对偶与截面分解.py}（$23$ 项全绿）。

\textbf{否证}：命题~\ref{prop:P1dead}（各向异性对偶证书；可行集是单点 $\{\rho_0I_3\}$，
随机搜索最好者仍违反 $-4.2\times10^{-1}$）；
命题~\ref{prop:ceilweyl}（天花板 $\varphi=1.618034$，比 $\rho_0$ 高 $0.336466$）；
命题~\ref{prop:ceilharm}（天花板 $1.585642$，比 $\rho_0$ 高 $0.304074$）。
这三条与第~\ref{sec:r18}~节的命题~\ref{prop:refA}\ref{prop:refB}~一起，
把「凸化 / 对偶 / 逐项放缩」这一整片方法\textbf{全部}排除在主线路径之外。

\textbf{方法优先队列（第九版）}
\begin{enumerate}\itemsep3pt
\item \textbf{P1（队首，全局）——单轴截面归约}。目标：证明推论~\ref{cor:sectreduce}~的 (C) 类不存在，
即 $\kappa_3(\cdot)$ 在 $\Delta_2$ 上的极大元必为中心或单轴点。
入口（按优先级）：
(a) \textbf{对称临界点论证}：$\kappa_3$ 是 $S_3$-不变的；若它在 $\mu$ 处可微，
则 $\nabla\kappa_3(\mu)$ 必与稳定子群的不动点空间正交，
故稳定子群非平凡的点（中心、单轴）自动是临界点；需要的是\emph{排除}稳定子群平凡的极大元。
(b) \textbf{凹性的替代物}：$\kappa_3$ 是凹函数的上确界，本身不凹；
但对\textbf{固定}帧 $\rho_{g,U}(\cdot)$ 凹，故
$\kappa_3(\mu)\ge\int\rho_{g,U}(R\operatorname{diag}(\mu)R^\top)\dd R$ 对每个帧成立，
右端是 $\mu$ 的\emph{对称化凹}函数。若能证明对称化不改变上确界的位置，则 (C) 排除。
这一步目前\textbf{没有证明}，是第一障碍。
(c) 退一步：不求排除 (C)，而是直接对\textbf{单轴族}做解析（见 P2），
把 (C) 的处理留到最后。
\item \textbf{P2（全局）——中心截面的精确值}。由 \eqref{eq:center}，这是
$\frac13\sup\sum_{ij}g_ig_j(\langle u_i,u_j\rangle)_+$，
即「$S^2$ 上各向同性测度的正部能量」。它\textbf{不含} $\Lam$，是本轮之后形式最干净的子问题。
数值上 $\kappa_3(\frac13,\frac13,\frac13)\approx1.256733$（$N=8$），
小于 $\rho_0$ 达 $0.024836$。已知工具只能给到 $\Phi=1.309017$
（用 $(\langle u,v\rangle)_+\le\frac{|\langle u,v\rangle|+\langle u,v\rangle}2$ 与 $\lambda_{\max}(|Q|)\le\varphi$）。
需要的是一条\textbf{正部特有的}不等式：正部只在锐角对上计数，
而各向同性条件 $\sum_ig_iu_iu_i^\top=I_3$ 限制了锐角对的总权重。
\item \textbf{P3（全局）——秩一/秩二截面}。数值余量最大（$\ge0.053$），
推论~\ref{cor:sectspecial}(ii)~已把它写成 $x_i=|a_i|$、$\sum a_i^2=1$ 的形式，
应当是最容易先关掉的一类。
\item \textbf{P4（全局）——$\mathcal E_-$ 的下界}（第八版的 P3，未动）：
用引理~\ref{lem:Em2}~的显式 $\zeta_{ij}$ 在 \eqref{eq:star}~中给 $\frac12\mathcal E_-\ge\Phi-\tau_0$。
\item \textbf{P5（局部）——第~\ref{sec:r1517}~节的分支 I / 分支 II}（第八版的 P4）。
定理~\ref{thm:transitive}~给出分支 I 的新入口。
\item \textbf{P6（$N=6$ 收尾）}：第~\ref{sec:n6global}~节余下的 $20$ 个符号类。
\item \textbf{P7（副方向，下界）}：$\Lam_3(G)\ge\rho_0n-3$ 的可达到性，只差定理~\ref{thm:Fn}~的无穷族渐近。
\end{enumerate}

\textbf{已弹出}：第八版的 P1（构造对偶证书 $S$）——命题~\ref{prop:P1dead}~证明其可行集是单点；
第八版的 P2（解析证明凹性）——定理~\ref{thm:geo}~已完成。

\textbf{本轮把主线的形状改成了}：
\[
\kappa_3=\sup_{\mu\in\Delta_2}\kappa_3(\mu),\qquad
\kappa_3(\mu)=\sup_{N}\ \sup_{(g,U)}\ \rho_{g,U}(\operatorname{diag}\mu),
\]
其中 $\rho_{g,U}$ 对每个帧都是\textbf{凹}的、有\textbf{闭式梯度}的、
其内层对偶\textbf{精确紧}的泛函。$\Lam$ 的 $8$ 个自由度里 $6$ 个是冗余的旋转，
只剩 $\Delta_2$ 的 $2$ 个。冠军坐在单轴子族上，
中心截面则退化成一个不含 $\Lam$ 的经典测度极值问题 \eqref{eq:center}。
凸对偶与逐项放缩这两片方法已经被完整排除（命题~\ref{prop:P1dead}%
\ref{prop:ceilweyl}\ref{prop:ceilharm}），
下一步必须用\textbf{各向同性条件与正部之间的直接相互作用}，而不是再找一层松弛。


\section{通道型主函数方案的精确天花板}\label{sec:nogo}

\subsection{动机}

第~\ref{sec:stwanat} 节的注~\ref{rmk:loss} 指出逐项放缩注定失败。一个自然的反应是：
那就不要逐项放缩，而是把 $\Lam_3$ 直接写成 $\sum_{ij}c_ic_j(\langle u_i,u_j\rangle)_+$，
对\emph{正部} $(t)_+$ 本身找一个主函数，并把 $\|Q\|_1$ 与 $\one^\top Q\one$ 的信息\emph{联合}使用。
这一节证明：这条路也是死的，而且死得比想象的彻底。

\subsection{方案族}

沿用 \STW\ Theorem 4.1 的坐标：$Q=VV^\top\in P_3(n)$，行 $x_i=c_iu_i$，$c_i\ge0$，$u_i\in S^2$，
各向同性条件 $\sum_ic_i^2u_iu_i^\top=I_3$ 成立，故 $\sum_ic_i^2=3$，$q_{ij}=c_ic_j\langle u_i,u_j\rangle$。
由 $\Lam_3(G)=\sum_{i\ne j}(q_{ij})_+=\sum_{ij}(q_{ij})_+-3$ 与 $(q_{ij})_+=c_ic_j(\langle u_i,u_j\rangle)_+$，目标是给
\[
\Sigma(Q):=\sum_{i,j}c_ic_j\bigl(\langle u_i,u_j\rangle\bigr)_+
\]
一个 $\le Mn$ 型的界。

我们称下列六个量为\emph{通道}，并给出各自的最优线性界（每个界都对一切满足各向同性的配置成立）：
\begin{center}
{\small
\begin{tabular}{@{}lll@{}}
\toprule
通道核 $K(u,v)$ & $\sum_{ij}c_ic_jK(u_i,u_j)$ 的值 & 上界 \\
\midrule
$1$ & $C^2$，$C=\sum_ic_i$ & $\le 3n$ \\
$f_2(t)=t^2-\frac13$ & $X^2=\bigl\|\sum_ic_i\phi(u_i)\bigr\|_F^2$ & 与 $C^2$ 联合：$C^2+3X^2\le 3n$ \\
$t=\langle u,v\rangle$ & $\|y\|^2$，$y=\sum_ic_iu_i=\one^\top Q$ & $\le n$ \\
$u_1v_1$ & $y_1^2\le\|y\|^2$ & $\le n$ \\
$t\,u_1v_1$ & $\|z\|^2$，$z=\sum_ic_i(u_i)_1u_i$ & $\le C^2\le 3n$ \\
$P_\ell(t)$，$\ell\ge3$ & 由 Schoenberg 定理非负 & $\ge0$（故只有负系数可用） \\
\bottomrule
\end{tabular}}
\end{center}

第一行的 $C^2\le n\sum c_i^2=3n$ 是 Cauchy--Schwarz；第二行是 \STW\ Lemma 4.2（$r=3$）；
第三行是 $Q\preceq I$；第四行显然；第五行由 $\|z\|\le\sum_ic_i|(u_i)_1|\le C$。

\begin{definition}[通道型方案]\label{def:scheme}
称六元组 $(a,\beta,b,\gamma,\delta,\{\rho_\ell\}_{\ell\ge3})$（全部非负，只有有限个 $\rho_\ell$ 非零）
为一个\emph{可行方案}，若对一切 $u,v\in S^2$ 有逐点不等式
\begin{equation}\label{eq:scheme}
\bigl(\langle u,v\rangle\bigr)_+\ \le\
a+\beta\langle u,v\rangle+b\Bigl(\langle u,v\rangle^2-\tfrac13\Bigr)
+\gamma\,u_1v_1+\delta\,\langle u,v\rangle u_1v_1-\sum_{\ell\ge3}\rho_\ell P_\ell(\langle u,v\rangle).
\end{equation}
其\emph{代价}定义为
\[
\mathcal M(a,\beta,b,\gamma,\delta,\rho)=\max(3a,b)+(\beta+\gamma)+3\delta .
\]
\end{definition}

\begin{proposition}[已证]\label{prop:schemevalid}
若 $(a,\beta,b,\gamma,\delta,\rho)$ 可行，则对一切 $n$ 与一切 $n$ 阶图 $G$，
$\Lam_3(G)\le\mathcal M\cdot n-3$。
\end{proposition}

\begin{proof}
把 \eqref{eq:scheme} 代入 $u=u_i$、$v=u_j$，乘以 $c_ic_j\ge0$，对 $i,j$ 求和。逐通道使用上表：
\[
\Sigma(Q)\ \le\ aC^2+bX^2+\beta\|y\|^2+\gamma y_1^2+\delta\|z\|^2-\sum_{\ell\ge3}\rho_\ell(\ge0)
\ \le\ \max(3a,b)\,n+(\beta+\gamma)n+3\delta n .
\]
其中 $aC^2+bX^2\le\max(3a,b)\cdot n$ 来自在 $\{(u,v):u+3v\le3n,\ u,v\ge0\}$ 上最大化 $au+bv$；
$\beta\|y\|^2+\gamma y_1^2\le(\beta+\gamma)\|y\|^2\le(\beta+\gamma)n$。最后 $\Lam_3=\Sigma(Q)-3$。
\end{proof}

\begin{remark}
定义~\ref{def:scheme} 确实包含了 \STW\ 的证法。取
$a=\frac{a_3}2=\frac{\varphi}6$，$\beta=\frac12$，$b=\frac{b_3}2$，$\rho_4=\frac{\rho_3}2$，其余为零；
由 $(t)_+=\frac12(|t|+t)$ 与 \eqref{eq:maj}、\eqref{eq:fac} 知 \eqref{eq:scheme} 成立，
且 $\mathcal M=\max(3\cdot\frac\varphi6,\frac{b_3}2)+\frac12=\frac\varphi2+\frac12=\frac{1+\varphi}2=\Phi$，
精确复现 Proposition 6.1。
\end{remark}

\subsection{见证测度}

定义 $[-1,1]$ 上的一个三点非负测度
\begin{equation}\label{eq:lam}
\lambda=\frac{5-\sqrt5}{4}\,\delta_{-1/\sqrt5}\ +\ \frac{5+\sqrt5}{4}\,\delta_{1/\sqrt5}\ +\ \frac12\,\delta_{1}.
\end{equation}
三个权重都是正数（$\frac{5-\sqrt5}{4}=0.690983$，$\frac{5+\sqrt5}{4}=1.809017$，$\frac12$）。

\begin{lemma}[见证测度的四个矩，已证且精确]\label{lem:moments}
\[
\int\dd\lambda=3,\qquad \int t\,\dd\lambda(t)=1,\qquad \int t^2\,\dd\lambda(t)=1,\qquad
\int (t)_+\,\dd\lambda(t)=\frac{3+\sqrt5}{4}=\Phi .
\]
此外，对一切 $\ell\ge3$ 有 $\int P_\ell(t)\dd\lambda(t)\ge0$，且 $\ell\in\{3,4,8,14\}$ 时恰为 $0$。
\end{lemma}

\begin{proof}
记 $x_0=1/\sqrt5$。前三条是直接计算：
\[
\frac{5-\sqrt5}{4}+\frac{5+\sqrt5}{4}+\frac12=\frac52+\frac12=3,
\]
\[
\frac{5+\sqrt5}{4}\cdot\frac1{\sqrt5}-\frac{5-\sqrt5}{4}\cdot\frac1{\sqrt5}+\frac12
=\frac{2\sqrt5}{4\sqrt5}+\frac12=\frac12+\frac12=1,
\]
\[
\Bigl(\frac{5-\sqrt5}{4}+\frac{5+\sqrt5}{4}\Bigr)\cdot\frac15+\frac12=\frac{5/2}{5}+\frac12=1 .
\]
第四条：
\[
\int(t)_+\dd\lambda=\frac{5+\sqrt5}{4}\cdot\frac1{\sqrt5}+\frac12
=\frac{5+\sqrt5}{4\sqrt5}+\frac12=\frac{\sqrt5+1}{4}+\frac12=\frac{3+\sqrt5}{4}=\Phi .
\]

关于 Legendre 项。用 $P_\ell(-x)=(-1)^\ell P_\ell(x)$ 分奇偶：
\[
\int P_\ell\dd\lambda=\begin{cases}
\dfrac52P_\ell(x_0)+\dfrac12, & \ell\ \text{偶},\\[6pt]
\dfrac{\sqrt5}{2}P_\ell(x_0)+\dfrac12, & \ell\ \text{奇}.
\end{cases}
\]
\emph{奇数 $\ell\ge3$}：需要 $P_\ell(x_0)\ge-\frac{2}{\sqrt5}=-0.894427$。
由经典的 Laplace 型估计 $|P_\ell(x)|\le\sqrt{\frac{2}{\pi\ell}}\,(1-x^2)^{-1/4}$（$\ell\ge1$，$|x|<1$），
在 $x_0=1/\sqrt5$ 处 $(1-x_0^2)^{-1/4}=(4/5)^{-1/4}=1.057371$，故
$\ell\ge3$ 时 $|P_\ell(x_0)|\le\sqrt{2/(3\pi)}\times1.057371=0.487145<0.894427$，成立。

\emph{偶数 $\ell\ge4$}：需要 $P_\ell(x_0)\ge-\frac15$。
$\ell\ge18$ 时同一估计给 $|P_\ell(x_0)|\le\sqrt{2/(18\pi)}\times1.057371=0.198853<0.2$，成立。
剩下 $\ell\in\{4,6,8,10,12,14,16\}$ 用精确有理算术逐个核验：
\[
\begin{array}{c|ccccccc}
\ell & 4 & 6 & 8 & 10 & 12 & 14 & 16\\
\hline
P_\ell(x_0) & -\frac15 & \frac{41}{125} & -\frac15 & -\frac{131}{3125} & \frac{3301}{15625} & -\frac15 & \frac{3223}{78125}\\[3pt]
\int P_\ell\dd\lambda & 0 & \frac{33}{25} & 0 & \frac{247}{625} & \frac{3213}{3125} & 0 & \frac{9424}{15625}
\end{array}
\]
全部非负。奇数情形的前几项同样精确为零或正：
$\int P_3\dd\lambda=0$，$\int P_5\dd\lambda=\frac{18}{25}$，$\int P_7\dd\lambda=\frac{77}{125}$，
$\int P_9\dd\lambda=\frac{26}{125}$。
（以上全部由 \texttt{sympy} 以 $\Q(\sqrt5)$ 上的精确有理运算复核，见随附脚本
\texttt{证书\_zonal天花板定理.py}；另用浮点对 $\ell=3,\dots,4000$ 做了独立扫描，最小值为 $0$。）
\end{proof}

\subsection{主定理}

\begin{theorem}[通道型方案的天花板恰为 $\Phi$，已证]\label{thm:zonal}
设 $(a,\beta,b,\gamma,\delta,\rho)$ 是定义~\ref{def:scheme} 意义下的可行方案。则
\[
3a+\beta\ \ge\ \Phi=\frac{3+\sqrt5}{4}=1.309016994374947,
\]
从而 $\mathcal M(a,\beta,b,\gamma,\delta,\rho)\ge\Phi$。
因此，由命题~\ref{prop:schemevalid} 得到的任何形如 $\Lam_3(G)\le Mn-3$ 的界必有 $M\ge\Phi$，
\textbf{不可能改进 \STW\ 的 Proposition 6.1}。特别地，$M$ 不可能等于 $\rho_0=1.281568196872168$，
即通道型方案\textbf{不可能证明猜想~\ref{conj:A1}}。
\end{theorem}

\begin{proof}
在 \eqref{eq:scheme} 中依次取下列三对单位向量（都在 $e_2e_3$ 平面内，故第一个坐标为零）：
\[
(u,v)=\Bigl(e_3,\ -\tfrac1{\sqrt5}e_3+\tfrac2{\sqrt5}e_2\Bigr),\quad
\Bigl(e_3,\ \tfrac1{\sqrt5}e_3+\tfrac2{\sqrt5}e_2\Bigr),\quad (e_3,e_3),
\]
对应的 $\langle u,v\rangle$ 分别是 $-1/\sqrt5$、$1/\sqrt5$、$1$，即测度 $\lambda$ 的三个支撑点；
且因为 $u_1=v_1=0$，第四、第五通道在这三对上的值都恒为 $0$。
把三处不等式分别乘以 $\lambda$ 的对应权重并相加，得
\[
\int(t)_+\dd\lambda\ \le\ a\!\int\!\dd\lambda+\beta\!\int\! t\,\dd\lambda
+b\!\int\!\Bigl(t^2-\tfrac13\Bigr)\dd\lambda+0+0-\sum_{\ell\ge3}\rho_\ell\!\int\! P_\ell\dd\lambda .
\]
由引理~\ref{lem:moments}，$\int\dd\lambda=3$，$\int t\dd\lambda=1$，
$\int(t^2-\frac13)\dd\lambda=1-\frac13\cdot3=0$，且每个 $\int P_\ell\dd\lambda\ge0$ 而 $\rho_\ell\ge0$。故
\[
\Phi\ \le\ 3a+\beta+0-0\ =\ 3a+\beta .
\]
最后 $\mathcal M=\max(3a,b)+(\beta+\gamma)+3\delta\ge 3a+\beta\ge\Phi$，因为 $b,\gamma,\delta\ge0$。
\end{proof}

\begin{remark}[这条定理说了什么，没说什么]\label{rmk:scope}
\emph{说了什么}：只要证明的形状是「先对 $(\langle u,v\rangle)_+$（乃至对含 $u_1v_1$ 耦合的更一般核）
找一个逐点上界，再对若干全局通道各自使用一个 $\le\kappa n$ 型的线性界」，
那么无论把 Gegenbauer 展开做到多高阶、无论是否引入非球对称的耦合通道，
得到的常数都不可能低于 $\Phi$。这条结论的强度远超「$\gamma(3)=\varphi$ 所以逐项放缩饱和」：
它把\emph{联合}计费也一并封死了，而且证明只有三行、见证是显式的。

\emph{没说}什么：定理~\ref{thm:zonal} 完全没有用到自洽约束 $A=1\{\Pi_3(M)\ge0\}$
（定理~\ref{thm:selfcons}），也没有用到通道之间的\emph{联合}可行域
（它把 $(C^2,X^2,\|y\|^2,y_1^2,\|z\|^2)$ 当作各自独立地取到自己的上界）。
因此，逃逸定理~\ref{thm:zonal} 只有两种可能：
\textbf{(甲)} 使用自洽性／掩码组合结构；
\textbf{(乙)} 建立通道值与内积分布之间的非平凡\emph{联合}约束。
第~\ref{sec:methods} 节按这两条组织候选方法。
\end{remark}

\begin{remark}[见证的几何意义]\label{rmk:geo}
$\lambda$ 的支撑 $\{-1/\sqrt5,\,1/\sqrt5,\,1\}$ 恰是\textbf{正二十面体六条等角线系}的内收集
（公角 $\arccos\frac1{\sqrt5}=63.43^\circ$）。也就是说，
\emph{把 \STW\ 钉死在 $\Phi$ 的障碍是正二十面体，而不是我们的极值构造 $K_1\vee C_5$}。
冠军核心在通道型框架里是「不可见」的：它的 $(s,t)=(1.5865,0.9766)$ 两项都不接触各自的上确界。
这解释了为什么第~\ref{sec:nogo} 节的方案族无论怎么精化都够不到 $\rho_0$：
它在为另一个几何对象（正二十面体）优化。

同一现象在笔记的早期工作中已经以另一种形式出现过：
只用 $(\tr|G|,\tr|G|^2)=(3,3)$ 两个不变量的 Cauchy--Schwarz 包络，
也恰好被正二十面体 ETF 精确取到，从而任何只用这两个不变量的论证不可能低于 $\Phi$。
两条完全不同的通道死于同一个几何对象，是很强的结构性证据。
\end{remark}

\subsection{独立的数值确认}

在得到定理~\ref{thm:zonal} 之前，我们先把定义~\ref{def:scheme}（只含前三个通道与 $\ell\le8$ 的修正项）
离散化为一个线性规划：在 $[-1,1]$ 上取 $12001$ 个格点强制 \eqref{eq:scheme}，
目标是最小化 $\max(3a,b)+\beta$。求解结果如下。

\begin{center}
{\small
\begin{tabular}{@{}lrrrrr@{}}
\toprule
修正项集合 & 最优值 & 与 $\Phi$ 之差 & $a$ & $\beta$ & 接触点集合 \\
\midrule
$\{\ell=3,4\}$ & $1.309016840$ & $-1.5\times10^{-7}$ & $0.2696723$ & $0.5000000$ & $\{-1,\pm\frac1{\sqrt5},1\}$ \\
$\{\ell=3,\dots,6\}$ & $1.309016911$ & $-8.4\times10^{-8}$ & $0.2696723$ & $0.5000000$ & 同上 \\
$\{\ell=3,\dots,8\}$ & $1.309016911$ & $-8.4\times10^{-8}$ & $0.2696723$ & $0.5000000$ & 同上 \\
$\{\ell=3,\dots,12\}$ & $1.309016911$ & $-8.4\times10^{-8}$ & $0.2696723$ & $0.5000000$ & 同上 \\
\bottomrule
\end{tabular}}
\end{center}

三点说明。第一，负偏差 $10^{-7}$ 量级完全是格点离散化造成的（格点约束是可行域的松弛），
真值在 $[\Phi-1.5\times10^{-7},\ \Phi]$ 内，而定理~\ref{thm:zonal} 把它精确钉在 $\Phi$。
第二，线性规划\emph{独立地重新发现了 \STW\ 的解}：最优处 $a=\varphi/6=a_3/2=0.269672331$，$\beta=1/2$，
$b$ 与 $\rho_4$ 取值与 \STW\ 不同但代价相同，说明最优解构成一整面。
第三，接触点集合稳定地是 $\{-1,-1/\sqrt5,1/\sqrt5,1\}$，
与引理~\ref{lem:moments} 的见证支撑一致——这正是我们随后能猜出闭式 $\lambda$ 的原因。

\section{已经被排除的路线}\label{sec:dead}

这一节列出所有已经被\emph{定量}否证（而不是「试了没成功」）的路线。
列出来是必要的：一方面避免重复劳动，另一方面这些否定性结果本身刻画了问题的难度结构。
「天花板」一栏指该路线所能达到的最好常数；只要它 $\ge\Phi$，该路线就无法改进已发表纪录。

{\small
\begin{longtable}{@{}p{0.5cm}p{5.2cm}p{5.6cm}p{2.6cm}@{}}
\toprule
\# & 路线 & 定量结论 & 状态 \\
\midrule
\endfirsthead
\toprule
\# & 路线 & 定量结论 & 状态 \\
\midrule
\endhead
1 & 对称化 / 逐点多主函数 & 饱和于 $1.309017$，零改进 & 已否证 \\
2 & DL23 松弛量转移 & 链式界 $1.2904545$，但可用松弛只占所需的 $43.6\%$（差 $2.3$ 倍） & 已否证 \\
3 & 绝对值路线 & $\lambda_{\max}(|R|)=1.602534>2\rho_0-1=1.563136$（差 $+0.0234$） & 已否证 \\
4 & 纯投影常数 & $\gamma(3)=\varphi$ 精确（定理~\ref{thm:gam3}），天花板恰为 $\Phi$ & 已否证 \\
5 & trace / Frobenius 型 & 饱和于 $\varphi$；需 $\tr M^2\le2.233$ 而冠军为 $2.639$（差 $2.19$ 倍） & 已否证 \\
6 & 支撑--$6$ 归约 & 为假：$K_1\vee C_6$ 给 $\Phi_7-\max_{|S|=6}\Phi_6=+0.014281$ & 已否证 \\
7 & 双线性不等式 (P1) & 为假，含两个显式反例 & 已否证 \\
8 & 逐对放缩系列 (P2) & 天花板 $1.3452$（合法版）与 $\frac{1+\sqrt3}2=1.3660$（纯二阶矩版），均严格劣于 $\Phi$ & 已否证 \\
9 & 卦限覆盖 + 线性核删除 & 极限就是 $N\le24$（球面等径不等式迫使至少 $8$ 个卦限，$8\cdot\frac\pi2=4\pi$ 恰铺满） & 已否证 \\
10 & 球面覆盖 + $p$-variation & 支撑界形如 $N\le7m$、$m\ge4$，不可能优于 $N\le28$ & 已否证 \\
11 & 带隙条件单独使用 & 满足带隙的球冠图有无穷多，冠军只是其中之一 & 已否证 \\
12 & 配对松弛（形式 8） & 天花板 $\frac12(1+\sqrt3)=1.3660254$ \emph{恰为 Mohar 的 $k=3$ 界}；
实际最大值 $=4/3$（$K=3$，精确）与 $3/\sqrt5=1.3416408$（$K\ge4$），均 $>\Phi$ & 已否证 \\
13 & \textbf{通道型主函数方案（本报告定理~\ref{thm:zonal}）} &
\textbf{天花板恰为 $\Phi$}，见证为显式三点测度 \eqref{eq:lam} &
\textbf{本轮已否证} \\
\bottomrule
\end{longtable}}

第 12 条需要单独说明，因为它是本轮之前的主攻方向，投入最大。
考虑双线性松弛 $B(\nu_1,\nu_2)=\iint(\langle x,y\rangle)_+\dd\nu_1(x)\dd\nu_2(y)$，
其中 $\nu_i\in C_3=\{\nu:\int uu^\top\dd\nu=I_3\}$。由 $C_3$ 弱$*$紧凸、Dubins 极点定理（六个矩条件
$\Rightarrow$ 极点支撑 $\le7$）与两次线性极值原理，可得夹逼
\[
\rho_0=J(\nu^*)\ \le\ c_3^{\Sigma}\ \le\ \max\{B(e_1,e_2):e_1,e_2\in\operatorname{ext}C_3\}
\ \le\ \max\{B(\nu_1,\nu_2):\nu_1,\nu_2\in C_3\}.
\]
若能证明最右端 $=\rho_0$，则 A1$'$ 成立。但：

\begin{theorem}[配对松弛的天花板，已证]\label{thm:43A}
对一切 $\nu_1,\nu_2\in C_3$（不限原子数），
$B(\nu_1,\nu_2)=\frac12\iint|\langle u,w\rangle|+\frac12\langle\!\int u\dd\nu_1,\int w\dd\nu_2\rangle
\le\frac12(\sqrt3+1)=1.366025403784439$。
\end{theorem}

\begin{proof}
$\iint\langle u,w\rangle^2\dd\nu_1\dd\nu_2=\tr(M_1M_2)=\tr I_3=3$，Cauchy--Schwarz 给
$\iint|\langle u,w\rangle|\le\sqrt3$；对单位向量 $v$，$\langle\int u\dd\nu_1,v\rangle\le\sqrt{v^\top M_1v}=1$，
故 $\|\int u\dd\nu_1\|^2\le1$，从而第二个内积 $\le1$。
\end{proof}

这个天花板\emph{恰好}是 Mohar 的 $k=3$ 界，比 $\Phi$ 还高 $0.057008$。
原因是配对松弛只保留了 $(L,\alpha_1,\alpha_2)$ 三个二阶矩量，
而二阶矩方法的最优结果就是 Mohar 界。更进一步，上端确实被超过：

\begin{proposition}[显式反例，已证且精确]\label{prop:43B}
令 $R=\frac23J_3-I_3\in SO(3)$（$R^2=I$，$\det R=1$），
$e_1=\frac13\sum_{i=1}^3\delta_{\sqrt3e_i}$，$e_2=\frac13\sum_{i=1}^3\delta_{\sqrt3Re_i}$。
则 $e_1,e_2\in\operatorname{ext}C_3$，$\langle u_i,w_j\rangle\in\{-1,+2\}$（对角为 $-1$，非对角为 $+2$），
$\alpha_1=\alpha_2=1$，$L=\frac53$，且 $B(e_1,e_2)=\frac12(\frac53+1)=\frac43$。
\end{proposition}

\begin{theorem}[$K_1=K_2=3$ 时的精确最大值，已证]\label{thm:43C}
$\max\{B(\nu_1,\nu_2):\nu_1,\nu_2\in C_3,\ \text{各}\ 3\ \text{原子}\}=\frac43$。
归约为 $\max_{O\in O(3)}\sigma_1((O)_+)$；再用
$\max_{\|O\|\le1}\sigma_1((O)_+)=\max_{P\in\{0,1\}^{3\times3}}\max_{c,d\ge0}\|D_cPD_d\|_*$
把 $512$ 个符号模式归为 $36$ 个置换类做有限核验，全局最大 $=1.333333333333334$（与 $\frac43$ 差 $8.9\times10^{-16}$），
唯一取等模式是 $A(C_3)=J_3-I_3$，$c=d=\frac1{\sqrt3}(1,1,1)$，第二名 $\frac97$，间隔 $0.048$。
\end{theorem}

\begin{conjecture}[数值，非证明]
$K\ge4$ 时 $\max B=3/\sqrt5=1.3416407864998738$。
五组独立运行 $(K_1,K_2)=(4,4),(5,5),(6,6),(6,7),(7,7)$ 一致到 $8\times10^{-15}$，
$K=4$ 的见证为 $q=(4,3,4,4)/15$、$\alpha_2=121/125$、$\langle Eu,Ew\rangle=11/(5\sqrt5)$、$L=19/(5\sqrt5)$。
\end{conjecture}

于是完整的常数阶梯是
\[
\rho_0=1.2815682\ <\ \Phi=1.3090170\ <\ \tfrac43=1.3333333\ <\ \tfrac3{\sqrt5}=1.3416408\ <\ \tfrac{1+\sqrt3}2=1.3660254 .
\]

\begin{remark}[一个必须记录的数值教训]
早期一轮工作曾报告配对松弛的最大值为 $1.366758746$，看似「几乎取到天花板」。
该值是\emph{伪影}：它对应 $\alpha_2=1.0107>1$，违反定理~\ref{thm:43A} 中已证的 $\alpha\le1$，
且超出天花板 $7.3\times10^{-4}$。根因是白化步骤 $M^{-1/2}$ 的条件数爆炸。
修正办法是在搜索代码中加入可行性拒绝与审计门（每次输出 $\alpha_1,\alpha_2$ 与各向同性残差），
所有超出 $1$ 的解一律丢弃。这条经验值得写进论文的复现说明：
\textbf{任何违反已证不等式的数值「改进」都是伪影，不是发现。}
\end{remark}

\section{当前的障碍}\label{sec:obstacles}

\subsection{主要障碍}

\paragraph{障碍 0$''$（第 12 轮重写，现居首位）：$3.2128\times10^{-3}$ 的秩三缺口}
第~\ref{sec:rowvar} 节把主线的形状再次改变。原来的首位障碍（下面这条「障碍 0」）已经\textbf{被否证}：
命题~\ref{prop:chamberdead} 说明 $K_{2,2,2}$ 是腔掩码，数值扫描说明 $19/19$ 类都是腔掩码，
所以「证不存在腔掩码」这条几何路线能关闭的类数是 $0$。
现在的首位障碍是一条\emph{定量的}缺口：在 $Z_5$-对称、五角星图案的类内，
REL 松弛（四条 Cauchy--Schwarz $+$ 行能量恒等式）的天花板是 $1.284780995477$，
而定理~\ref{thm:conefam} 证明真正的秩三最大值是 $\rho_0=1.281568196872$，差 $3.2128\times10^{-3}$。
这个差的来源已经被精确定位：REL 只保留了\emph{逐行二次}信息 $a_i+r_i=g_i(1-g_i)$，
丢掉了引理~\ref{lem:rowid}(iii) 的三角形恒等式 \eqref{eq:tri}、四点 Gram 行列式 $\det G_4=0$
（$\R^3$ 中任意四点线性相关）以及等角线 $\le6$ 的约束。
它\emph{有利}的地方是：这些都是\textbf{等式}而非不等式，而且第~\ref{sec:conefamily} 小节已经证明
它们在等变层里退化为恒等式，所以必须\emph{先降对称性}它们才会咬人——
这给出一个明确的、可执行的第一步（第~\ref{sec:pq7} 小节的 P1(c)）。
它\emph{不利}的地方是：四点 Gram 行列式是 $q_{ij}$ 的\emph{四次}多项式，
而 REL 的变量是 $(g,\pi,r,w)$ 这些\emph{二次}量，把前者翻译成后者的语言没有现成办法。

\paragraph{障碍 0（第六版，\textbf{已被否证}，保留以记录演化）：$N=6$ 的 $19$ 个未闭合类，即一条几何不存在性陈述}
上一版把平面框架不等式 \eqref{eq:frame5} 列为首位障碍。定理~\ref{thm:bipglobal} 之后这条\textbf{不再是障碍}：
它被绕开而不是被证明，$\kappa_3(5)\le\rho_0$ 已经无条件（推论~\ref{cor:n5uncond}），
\eqref{eq:frame5} 与旧的 $(\ast_5)$ 一并降级为独立兴趣（注~\ref{rmk:bipvsframe}）。
现在的首位障碍是命题~\ref{prop:tab19} 列出的 $19$ 个 $6$ 阶掩码类。
它的\emph{有利}之处，也是本轮最重要的结构性收获：由注~\ref{rmk:tab19read}(iv)，
\textbf{不需要}对这 $19$ 类逐个求 $F$ 的上界，只需证明
\begin{equation*}
\Sigma\ \ne\ 1\{Q\ge0\}\qquad\text{对一切 }Q\in\mathcal P_3(6)
\end{equation*}
（即 $\Sigma\notin\mathcal S_6^{\mathrm{int}}$，不存在正权自洽见证），$\kappa_3(6)\le\rho_0$ 就随之成立。
其中 $17$ 类没有支配点，引理~\ref{lem:bdryrig} 已使这一步\emph{等价}；
引理~\ref{lem:chamber} 又把它翻译成「$\Sigma$ 不是任何 Parseval 框架的非锐角图」。
它\emph{不利}的地方有三条：(a) $\mathrm{Gr}(3,6)$ 是 $9$ 维紧流形，$17$ 个「开集为空」的断言要逐个给出理由；
(b) 符号模式 $1\{Q\ge0\}$ 在 $Q$ 的条目过零时\emph{不连续}，
所以「不存在」的证明必须处理退化面（这是上一轮翻转链失败的同一个原因，注~\ref{rmk:flipdead}）；
(c) 剩下 $2$ 类有支配点，引理~\ref{lem:bdryrig} 只能排除\emph{边界}见证，内部见证需要别的论证
（$K_{1,2,3}$ 可用引理~\ref{lem:mpsec} 的秩 $k$ 长期方程）。

\paragraph{障碍 0$'$（新，第二优先）：$k\ge3$ 时的尾和记账}
定理~\ref{thm:bipglobal} 的四步论证在完全多部、$k\ge3$ 时\textbf{断在第二步}，断点非常明确（注~\ref{rmk:mpopen}）：
$\operatorname{neg}(R)=k-1$（引理~\ref{lem:mpsec}，已证），故 Weyl 的指标平移从 $1$ 变成 $k-1$，
尾和里出现 $k-1$ 个负根而不是一个，「一个标量 $s$」的记账失效。
而且这条障碍\emph{不是}技术性的：注~\ref{rmk:mpopen} 的数值显示 $\min_p\sum_{i\ge4}\lambda_i=-0.2524$
（$K_{1,2,2}$），比 $-\frac14$ 还低 $2.4\times10^{-3}$，即 $k\ge3$ 的真值\textbf{严格高于} $\frac54$，
任何「$\frac54$ 型」的界都不可能推广过去。要走通，必须用秩 $k$ 长期方程 \eqref{eq:mpsec2}
直接研究根之和（Viète 型），而不是逐特征值放缩。

\begin{enumerate}[leftmargin=1.8em,itemsep=4pt]
\item \textbf{A1$'$ 未证，且没有产生新的无条件纪录。} 现役上界仍是 $\Phi n-3$。
缺口 $\Phi-\rho_0=0.0274488$ 一寸未动。
\emph{（2026-10-28 补注）}第~\ref{sec:domdual} 节没有改变这一条：冠军腔与四类近腔的全局上界
虽然已经\emph{证明}，但它们都是「条件性主定理~\ref{thm:assemble} 的前提」层面的进展，
不构成对任意图的无条件上界。缺口的\emph{性质}变了——
从「一批有限维非凸极值问题未解」变成「一条凸不等式（猜想~\ref{conj:A2}）在自洽掩码类上未证」。
\emph{（2026-10-28 深夜再补）}而这条凸不等式\textbf{已被否证}（命题~\ref{prop:A2no}）：
$H=\overline{\widehat{C_5}}$ 的 $7$ 原子自洽核心给出 $\tr((-M)_+)=0.288182>\tau$。
所以缺口的第二次表述也失效了，现在只剩第三次表述 (W1$''$)：$\kappa_3\le\rho_0$。
这实际上是一个\emph{好消息}：它说明「凸化」不是出路，
秩三结构 $\Pi_3$ 必须\textbf{原样保留}，而 $\kappa_3=\sup_V\rho((VV^\top)_+)$ 正是这样定义的。
\emph{（2026-10-28 深夜补注）}第~\ref{sec:kappa} 节同样没有改变这一条，但它把缺口的\emph{表述}又换了一次：
现在是「$\kappa_3\le\rho_0$」——一条关于秩三正交投影正部谱半径的纯线性代数不等式，
不再提及图、掩码或阶数。三次表述（非凸极值 $\to$ 凸不等式 A2 $\to$ 单常数 $\kappa_3$）
在逻辑上等价，但\emph{可用工具}完全不同：第三种表述允许使用
Stiefel 流形的一阶条件 \eqref{eq:stiefelkkt}、Perron 刚性恒等式组（命题~\ref{prop:perronrig}）
与 $N=6$ 的 $156$ 类符号穷举。
\item \textbf{定理~\ref{thm:zonal} 关闭了我们原先排在首位的技术路线。}
原计划的「路线 B：$O(2)$-不变正定核主函数」正是定义~\ref{def:scheme} 的一个实例
（它比 \eqref{eq:scheme} 更一般，但仍属通道型），故同样被封顶在 $\Phi$。
这是本轮最有价值的负面结果：它在我们投入数月之前排除了一个方向。
\item \textbf{自洽性还没有被真正用起来。} 定理~\ref{thm:selfcons} 告诉我们，
极值对象必须满足 $A=1\{\Pi_3(M)\ge0\}$，这是一条组合约束。
我们目前所有关于它的结论都是\emph{结构描述性}的（负图 $\Gamma^-=C_5$、无三角形、等权刚性），
还没有一条把它转化为\emph{定量}不等式的通道。这是全部难度的所在。
\emph{（2026-10-28 补注）}命题~\ref{prop:octnogo} 把这句话从定性升级为\textbf{定量}：
八面体掩码是 $\R^3$ 标架的符号掩码、满足定理~\ref{thm:P3} 的归一化、在 \eqref{eq:P3} 中目标值只有 $0.5$，
但它的 $\Omega$ 精确等于 $4/3>\rho_0$。也就是说，「$\R^3$ 标架可实现」这一整层几何信息
\emph{不足以}把优占对偶压到 $\rho_0$；必须用自洽性本身。反过来，
命题~\ref{prop:negmass}(i) 给出了自洽性的第一个\emph{定量}用法：
它把 $\Lam_3(M)$ 精确写成 $\frac12(\|Y\|_1+\one^\top Y\one)$，其中 $Y=D_{\sqrt p}\Pi_3(M)D_{\sqrt p}$——
这正是 \STW\ 的 $\eta_{\mathrm{mix}}$ 表达式，只是多了权 $p$ 与自洽性两道约束（注~\ref{rmk:etamix}）。
\item \textbf{支撑界缺口 (2a)：「极值测度的原子数 $\le6$」未证。}
无条件结果是 $N\le24$（Fenchel--Bunt 精化到 $N\le20$），
而定理~\ref{thm:selfcons} 要压到有限张表需要 $N\le8$（第 27 节定理 27.1：$SB(8)\Leftrightarrow c_3^{\Sigma}=\rho_0$）。
目前只有局部版本的定理 4.3--1。
\emph{（2026-10-28 深夜补注）}定理~\ref{thm:ladder} 与注~\ref{rmk:ladderinfo} 给了 (2a) 的第一条正面证据：
$\kappa_3(5)=1.2531973<\rho_0=\kappa_3(6)$，说明\textbf{六个原子恰好是达到最优常数的最小规模}；
而 $N=7,\dots,12$ 的最大值点全部退化（冠军加上一批权趋于零的冗余原子，
前六坐标与冠军一致到 $10^{-8}$），说明 $\kappa_3$ 的上确界\emph{不是}靠更多原子在 $N\to\infty$ 处逼近的。
这把 (2a) 从「一个孤立猜想」变成优先队列里的 P2：
证「$N\ge7$ 的任意投影可化约到某个 $\le6$ 个有效原子的构型而不增大 $\rho(Q_+)$」。
必须强调：$20$（或 $24$）\emph{不等于} $8$，把「支撑 $\le20$」写成「已证 A1$'$」是严格禁止的。
\item \textbf{第 6 条已否证路线（支撑--$6$ 归约为假）与 (2a) 之间的张力。}
$K_1\vee C_6$ 的反例说明「任意支撑 $7$ 的核心可以压到支撑 $6$」是假的；
(2a) 是更弱的陈述（只针对\emph{极值}测度），但两者的区别需要仔细论证，目前没有。
\item \textbf{（本轮新增）Perron 刚性系统缺一「第四条不等式」。}
命题~\ref{prop:perronrig} 把冠军处的全部一阶信息写成了八条\emph{精确}恒等式，
但注~\ref{rmk:rigiditylimit} 证明：这八条合起来给出的唯一新约束只有 $0\le\alpha\le1$，
再对 $\beta$ 单独用一次 Cauchy--Schwarz 也只能得到 $\rho\le\alpha+\sqrt{\alpha(1-\alpha)}\le\frac32$，
\textbf{远不足以}推出 $\rho\le\rho_0=1.2815682$。
缺口被精确定量了：冠军处 Cauchy--Schwarz 的余量是 $\alpha(1-\alpha)=0.022856959$，
而命题~\ref{prop:decouple} 说明 \STW\ 的总损失 $2(\Phi-\rho_0)=0.054897595$
恰好被劈成 $(1-\alpha)=0.023404741$ 与 $(\phi-\ell)=0.031492854$ 两半，
所以\emph{任何只吃掉其中一半的论证都必然失败}；需要的是一条\textbf{联合}不等式
$\alpha+\ell\le2\rho_0$（优先队列 P4）。这是本轮结束时最清晰、也最硬的一个障碍。
\item \textbf{（本轮新增）冠军在 $Q$ 层面的 $C_5$ 对称性未被证实。}
数值优化给出的 $N=6$ 最大值点对 $C_5$ 作用的偏差约 $2\times10^{-8}$，
与优化器本身的精度同阶；直接检验「rim 块的特征值结构」也不通过。
因此\textbf{不得}声称 $\kappa_3(6)$ 的最大值点具有 $C_5$ 闭式，
只能用「与冠军 $I+A(K_1\vee C_5)$ 的加权形式数值一致」这种表述。
\end{enumerate}

\subsection{数值证据的边界（必须诚实交代）}

本轮新增的数值条目及其强度，按可信度从高到低：
(i) \eqref{eq:frame5} 在 $937614$ 帧（Haar / 双轴含非等分 / 近退化 / 近等角四族）上 $0$ 违例，
$\sup=1.249993448627$；双轴族另有 $3\,600\,000$ 点的独立扫描，无一点超过等分点的值 $\frac54$。
(ii) 双轴 $120$ 个等分构型的 $\frac54I-B\succeq0$ 与「恰 $3$ 个取等」是\textbf{精确有理}结论，不属数值证据。
(iii) $N=6$ 的 $21$ 个例外：极大化点在边界（$\min_ip_i\le10^{-14}$）、逐类 $\max F=1.253197265$、
$19$ 类无支配点——均为数值。
(iv) $L(Z)\le1$ 的违例 $1583/937614$ 与 $\sup L=1.117916903$ 是数值观测，
但「$L\le1\Rightarrow\lambda_{\max}\le\frac54$」这一步是\textbf{定理}（命题~\ref{prop:schurcond}），
且逻辑一致性检查（$L\le1$ 却 $\lambda_{\max}>\frac54$ 的帧数 $=0$）通过。
(v) 双轴族上 $\varphi(t)$ 的\textbf{凹性}被否证（$360000$ 对中 $57149$ 对违例，$15.9\%$），
这是我本轮提出又自己推翻的一个中间论断，记录在案以免重复犯错。

我们把形式 9（见下）中的
$\mathrm{SR}_K=\max\{\lambda_{\max}((A^\top A)_+):A\in\R^{3\times K},\ AA^\top=I_3\}$
用随机起点加交替迭代加 Powell 精修做了数值封顶，结果为
$\mathrm{SR}_6=\rho_0+2.2\times10^{-15}$，
$\mathrm{SR}_7=\rho_0+2.9\times10^{-7}$，$\mathrm{SR}_8=\rho_0+2.5\times10^{-6}$，
$\mathrm{SR}_9=\rho_0+2.9\times10^{-6}$，
$\mathrm{SR}_{10}\dots\mathrm{SR}_{14}=\rho_0+(1.8\sim6.2)\times10^{-8}$。
超出量与「附加原子总权 $\sum_{k>6}c_k^2$」逐位同阶（$K\le9$ 为 $2.2\times10^{-7}$ 至 $2.3\times10^{-6}$，
$K=10\dots14$ 为 $1.4\times10^{-8}$ 至 $4.8\times10^{-8}$），
前六个坐标与冠军一致到 $10^{-6}$，故判定为「权趋于零的冗余原子」造成的病态摄动，真值应为 $\rho_0$。

\textbf{但这一整段只是数值证据，不是证明，而且它的价值有限。}理由有三：
第一，$\mathrm{SR}_K$ 是丢掉自洽约束后的松弛，即使对一切 $K\le20$ 证出 $\mathrm{SR}_K=\rho_0$，
也只在一个\emph{更强}的命题上得到结论，与 A1$'$ 之间还隔着自洽性；
第二，$K=10,\dots,14$ 的搜索强度（$142$ 个起点、交替迭代 $250$ 步、每档时限 $900$ 秒，
实际 $6\sim34$ 秒即在「冠军加零权列」处收敛）明显弱于 $K=6,\dots,9$
（$260$ 起点、\texttt{Powell} 精修至 \texttt{ftol}$=10^{-14}$、每档 $28\sim426$ 秒），
故 $K\ge10$ 各行只能算较弱的阴性证据；第三，$K=15,\dots,20$ 尚未计算，也没有任何区间算术证书。

这条线的正确定位是\emph{否证工具}：它的作用是排除错误方向（第 12 条路线正是被这类计算加精确反例一起否证的），
而不是推进上界。把它当作主线是方法论上的错误。

\section{候选方法与其评估}\label{sec:methods}

\subsection{从定理~\ref{thm:zonal} 反推出的必要条件}

注~\ref{rmk:scope} 给出了逃逸通道型框架的\emph{仅有}两种可能。把它写成一个判据：

\begin{proposition}[必要条件，已证]\label{prop:nec}
任何证明 $c_3^{\Sigma}\le\rho_0$ 的论证，必须至少使用下列信息之一：
\begin{enumerate}[label=(\Alph*),leftmargin=2em]
\item \textbf{自洽性}：$Q$ 不只是秩 $3$ 投影，而且是掩码矩阵 $D_w(A_H+I)D_w$
（其中 $A_H=1\{Q>0\}-I$）的前三特征投影；
\item \textbf{通道—分布的联合约束}：内积的经验分布 $\frac1n\#\{(i,j):\langle u_i,u_j\rangle\in\cdot\}$
与通道值 $(C^2,X^2,\|y\|^2,\dots)$ 之间存在非平凡的\emph{联合}限制，
而不能像通道型方案那样让它们各自独立地达到上界。
\end{enumerate}
\end{proposition}

\begin{proof}
若论证只用各向同性 \eqref{eq:iso} 与通道的独立线性界，则它必然可以改写成定义~\ref{def:scheme}
意义下的一个可行方案（主函数取为论证中实际使用的那一个），从而由定理~\ref{thm:zonal} 只能给出 $\ge\Phi$。
\end{proof}

这条判据的价值在于它是一个\emph{筛选器}：以后每提出一个新想法，先问它用的是 (A) 还是 (B)。
如果两者都不用，那它一定失败，不必尝试。第 12 条（配对松弛）与第 13 条（通道型方案）
都是没有通过这个筛选器的例子——配对松弛只用了二阶矩量（属于通道型），
所以它的天花板必然是 Mohar 界；这事后看来是显然的，但事前花了很大代价才发现。

\subsection{方向甲：自洽性 + 逐核精确 KKT}

这是唯一一条在同类问题上\emph{已经被验证能达到锐利}的架构：$k=2$ 时的
Kumar--Liu--Monterde--\allowbreak Pragada--Tait 定理 $\lambda_1+\lambda_2\le\frac87n$ 正是用图极限加凸分析做的，
而 \STW\ 的一般方法在 $k=2$ 处只能给 $\frac76n$。

我们已有的工具如下。

\begin{theorem}[形式 9：标架--Perron 形式，已证]\label{thm:form9}
\[
c_3^{\Sigma}=\max\Bigl\{\lambda_{\max}\bigl((A^\top A)_+\bigr):\
A\in\R^{3\times K},\ AA^\top=I_3,\ 3\le K\le20\Bigr\},
\]
其中 $(\cdot)_+$ 是逐元素正部，$A^\top A$ 是秩 $3$ 正交投影。
权是\emph{自由}的：双射 $a_k=\sqrt{p_k}\,u_k$、$c_k=\sqrt{p_k}$，且最优权自动等于 $(A^\top A)_+$ 的 Perron 向量。
参数个数是 $3K-6=54$（$K=20$）。
\end{theorem}

\begin{lemma}[Ky Fan 对偶，已证]\label{lem:kyfan}
$\max_{AA^\top=I_3}\tr(AMA^\top)=\mathrm{KyFan}_3(D_cMD_c)$ 型对偶成立，最优 $A$ 是前三特征向量。
取 $M=P=I+A_H$ 时，\textbf{逐核封顶（方向甲的主体）恰是形式 9 的精确对偶}，不是启发式近似。
由此得到一个「模式 $\leftrightarrow$ 权」的交替算法，其目标值单调不减、收敛到自洽不动点。
\end{lemma}

\begin{proposition}[冠军的 Perron 特征方程，已证且精确]\label{prop:perron}
冠军满足 $(G^*)_+\sqrt{p^*}=\rho_0\sqrt{p^*}$，其中 $G^*=(A^*)^\top A^*$。
数值上 $\|(G^*)_+\sqrt{p^*}-\rho_0\sqrt{p^*}\|_\infty=1.1\times10^{-16}$，$\langle\sqrt{p^*},\text{Perron}\rangle=1.0$，
且
\begin{multline*}
\mathrm{eig}\bigl(D_{\sqrt{p^*}}(I+A_{K_1\vee C_5})D_{\sqrt{p^*}}\bigr)\\
=\bigl[0.76203144,\ 0.25976838,\ 0.25976838,\\
-0.08312281,\ -0.09922269,\ -0.09922269\bigr],
\end{multline*}
前三大之和的 $\mathrm{KyFan}_3$ 与 $\rho_0$ 的 \texttt{sympy} 差为 $0$。
\end{proposition}

\textbf{方向甲的具体内容}：对每张支撑 $\le6$ 的无标号核心 $H$（共 $156$ 张，已由 hub 拆分命题从 $1044$ 张压缩而来），
在 $\|x\|=1$ 上求精确极大值 $\max_x\Lam_3(D_x(I+A_H)D_x)$，
办法是「对称对角化 + 一阶条件解析求解 + 区间算术认证符号模式」。
可行性支撑有三条：$k=2$ 的同一套流水线已端到端跑通并复现已发表定理；
符号模式的稳健性 $\min|\Pi_3|=0.1900$ 足够大，区间算术不会因舍入而翻转；
命题~\ref{prop:perron} 说明可以\emph{先解 Perron 方程再验自洽}，把搜索维度从「权 + 标架」降到「只有标架」，且每步有闭式。

\textbf{它的唯一实质缺口是 (2a)：「极值测度的原子数 $\le6$」}。没有它，$K\le20$ 无法压缩成有限的核心表。
这是一个明确的、有希望的研究问题：已有件包括负图 $\Gamma^-$ 的结构定理与等权刚性、
$\Gamma^-$ 无三角形、以及「第七原子最优权恰为零」的原子阶梯命题，但把它们拼成 (2a) 的论证目前没有。

\subsection{方向乙：通道—分布的联合约束（本轮新提出）}

定理~\ref{thm:zonal} 的见证 $\lambda$ 之所以能钉住 $\Phi$，是因为它让三个通道\emph{同时}达到各自的上界：
$\int\dd\lambda=3$ 对应 $C^2=3n$，$\int t\dd\lambda=1$ 对应 $\|y\|^2=n$，$\int t^2\dd\lambda=1$ 对应 $X^2$ 的余量。
但 $C^2=3n$ 当且仅当所有 $c_i$ 相等。于是自然的想法是\textbf{按 $C^2/n$ 分情形}。

\emph{等权子问题。}当所有 $c_i$ 相等时，$\sum_ic_i^2u_iu_i^\top=I_3$ 变成
$\sum_iu_iu_i^\top=\frac n3I_3$，即 $\{u_i\}$ 是 $\R^3$ 中的\emph{单位范数紧框架}（UNTF）；
同时 $\Sigma(Q)/n=\frac3{n^2}\sum_{ij}(\langle u_i,u_j\rangle)_+$。记
\[
m_{\rm eq}=\sup_{n\ge3}\ \max_{\{u_i\}\ \mathrm{UNTF}}\ \frac3{n^2}\sum_{i,j}\bigl(\langle u_i,u_j\rangle\bigr)_+ .
\]
这是 $(P_3)$ 在等权限制下的版本，故 $m_{\rm eq}\le c_3^{\Sigma}\le\Phi$。
两条平凡的观察：
(i) \textbf{复制单调性}——把 UNTF 复制一份得到 $2n$ 个向量的 UNTF，目标值不变，故 $m_{\rm eq}(2n)\ge m_{\rm eq}(n)$；
(ii) 正二十面体（$n=6$）给 $m_{\rm eq}\ge\frac3{36}(6+15/\sqrt5)=1.0590$，
而我们的一次初步搜索在 $n=6$ 处找到 $1.245304$，故 $m_{\rm eq}\ge1.245304$。

\emph{定量收益估计。}冠军核心在团吹起下的 $c_i$ 只取两个值
$\sqrt{q_{\rm hub}}=\sqrt{a^*V(\text{hub},\text{hub})}=0.575973$ 与
$\sqrt{q_{\rm cyc}}=\sqrt{b^*V(\text{cyc},\text{cyc})}=0.730511$，于是尺度无关的比值
\[
\frac{C^2}{3n}=\frac{(\mathbb E c)^2}{\mathbb E[c^2]}
=\frac{\bigl(a^*\sqrt{q_{\rm hub}}+5b^*\sqrt{q_{\rm cyc}}\bigr)^2}{(a^*)^2V(\text{hub},\text{hub})+5(b^*)^2V(\text{cyc},\text{cyc})}
=0.992343569348608
\]
（\texttt{sympy} 精确计算），即 $3-C^2/n=0.0229692919542$。
若能在 $C^2/n\le3-\epsilon$ 的子类上把天花板按 $\Phi(3-\epsilon)/3$ 的比例压缩，
则在 $\epsilon=0.02297$ 处得到 $1.29899459654$，比 \STW\ 改进 $0.7657\%$，
但仍比 $\rho_0$ 高 $0.01743$。

\textbf{评估}：方向乙\emph{有可能}产出一个新的无条件纪录（约 $1.299$），但\emph{不可能}证到 $\rho_0$；
而且「按 $C^2/n$ 分情形」要成为证明，需要对\emph{两个}子情形都有界，
其中 $C^2/n$ 接近 $3$ 的那一支恰恰是最难的（它逼近等权子问题 $m_{\rm eq}$）。
所以方向乙的正确定位是：\textbf{一个可以独立发表的部分结果（改进 STW 约 $0.77\%$），
以及一个把主问题分解为「等权」与「非等权」两支的手段}，而不是主线的替代。

\emph{一条必须记录的负面数值经验。}我们对 $m_{\rm eq}$ 做的第一轮梯度搜索（$n=6,8,9,12,16,18,24$，
各 $160$ 个随机起点）给出 $1.2453,0.9073,0.9197,0.8846,0.8344,0.8857,0.8119$。
\textbf{这组数不可信}：$n\ge8$ 的结果与复制单调性 (i) 直接矛盾（$m_{\rm eq}(12)$ 不可能小于 $m_{\rm eq}(6)$），
根因是搜索陷入了「近均匀框架」的局部极小（$n\to\infty$ 时目标趋于
$3\iint(\langle u,v\rangle)_+\dd\sigma\dd\sigma=3/4=0.75$，与 $0.81\sim0.89$ 一致），
而 $(t)_+$ 在大圆 $\langle u,v\rangle=0$ 上不可微，朴素梯度在此失效。
本报告不采信任何 $m_{\rm eq}$ 的数值上界；要重做，必须用光滑化（如 $(t)_+\to\sqrt{t^2+\eta^2}$ 再令 $\eta\to0$）
或以复制构造为起点的分层搜索。

\subsection{方向丙：三角 SOS / Fejér--Riesz}

定理~\ref{thm:champ}(d) 的 $O(2)$ 归约把问题化为：固定标量轮廓 $a=(a_I)$ 后，
角度 $\theta\in S^1$ 是唯一的连续自由度，而冠军是 $5$ 重对称（$\theta$ 间距 $2\pi/5$）。
于是待证不等式可以写成 $5$ 次单位根上的非负三角多项式，Fejér--Riesz 定理保证它有 SOS 表示，
并且由常数域 $\Q(\sqrt5,\sqrt7)$ 的分析，系数有望在该域上\emph{精确}求出并用 \texttt{sympy} 判等。

\textbf{必须注意它与定理~\ref{thm:zonal} 的关系。}如果 SOS 证书最终被组织成
「$(t)_+$ 的逐点上界 + 若干全局通道的线性界」，那它就落在通道型方案内，天花板仍是 $\Phi$，白做。
只有当证书是\emph{直接}针对 $\sum p_ip_j(\langle z_i,z_j\rangle)_+\le\rho_0$
在约束 $\sum p_iz_iz_i^\top=I_3$ 下的一个约束多项式非负性证书
（即形如「目标 $-\rho_0+$ 约束的 SOS 组合 $\ge0$」的 Putinar/Schmüdgen 型证书）时，
它才使用了 (B) 类信息而不受定理~\ref{thm:zonal} 限制。这个区别必须在动手前想清楚。

一条硬约束（已实测）：把 $a_Ia_J$ 用均值 $1/n$ 代替的「常数移位代理」在 $\|Q\|_1$ 上有 $-3.97\%$
的相对误差，\emph{比整个 $2.0969\%$ 的缺口还大}。所以标量轮廓 $a$ 必须精确保留，
「先做球面平均化再处理角度」这条最自然的简化路线被直接否掉。

\subsection{方向丁：图极限—凸分析（KLMPT 架构的移植）}

$k=2$ 的锐利结果 $\lambda_1+\lambda_2\le\frac87n$ 是 KLMPT 用图极限与凸分析证的，
而它的极值核心 $P_3=K_1\vee2K_1$ 只有三个原子、负图 $K_2$ 是二部的。
移植到 $k=3$ 的障碍已在第~3.3 节说明：$\Gamma^-=C_5$ 是奇圈，切换结构失效，目标函数不再是完全平方。
可行的切入点是绕过切换、直接对 $\Lam_3$ 的 Ky Fan 表示做凸性分析：
$\Lam_3(M)$ 关于 $M$ 是凸的，而 $M=D_w(A_H+I)D_w$ 关于 $w$ 是\emph{二次}的，
故 $\psi_3(H)=\max_{w\in\Delta}\Lam_3(M(H,w))$ 是一个凸函数在单纯形上的最大化，极值在边界或临界点。
定理~\ref{thm:champ}(b) 的一阶条件 $a^2+24ab-31b^2=0$ 正是这条路线的产物。
尚缺的是把「每张核心的临界点都已算出」升级为「全局上确界在有限张核心上取到」——又回到 (2a)。

\subsection{其他考虑过但不推荐的想法}

\begin{itemize}[leftmargin=1.4em,itemsep=2pt]
\item \textbf{矩阵值（而非标量值）主函数}：把 \eqref{eq:maj} 的右端换成一个矩阵值核。
这属于通道型方案的推广，若其计费仍是「逐通道 $\le\kappa n$」，则定理~\ref{thm:zonal} 的见证
（取 $u_1=v_1=0$ 使所有含 $u_1,v_1$ 的通道归零）仍然适用，故无收益。
\item \textbf{二阶 lift / 矩量矩阵 SOS}：把 $(c_i,u_i)$ 的经验测度提升到二阶矩量矩阵再做 SDP。
理论上这属于 (B) 类，但维度随原子数平方增长，在 $K\le20$ 时不可行；
且若不引入自洽性，SDP 的极限仍会收敛到通道型松弛的值。
\item \textbf{稀疏图侧}：\STW\ 在 \S6.4 中承认其方法只对稠密图有效，稀疏侧用迹恒等式另做。
但 $c_3^{\Sigma}$ 是稠密渐近常数，稀疏侧与本课题无关。
\item \textbf{把 $k=3$ 化成 $k=2$ 的已有结果}：$\Lam_3=\Lam_2+\lambda_3$，
但 $\lambda_3$ 的最好界（Tang 的 $\lambda_3\le n/3-1$，已被 \STW\ 合并取代）太弱，相加后远劣于 $\Phi$。
\end{itemize}

\subsection{优先级建议}

\begin{enumerate}[leftmargin=1.8em,itemsep=2pt]
\item \textbf{(2a) 支撑 $\le6$}（方向甲的唯一缺口）。这是把无限问题变有限的唯一开关，
一旦证出，方向甲就是一套机械的、可用区间算术完整认证的流程，主定理随之得证。
\item \textbf{方向乙的等权子问题 $m_{\rm eq}$}。即使主线受阻，它也能产出一个可发表的独立结果
（改进 \STW\ 约 $0.77\%$ 的无条件上界），适合作为毕业论文的保底章节。需要先修正数值方法。
\item \textbf{方向丙的约束 SOS 证书}。产物形态最适合答辩（一个可现场复核的代数证书），
但必须先确认它不落入定理~\ref{thm:zonal}。
\item \textbf{下界一侧的整理}。定理~\ref{thm:lower} 与~\ref{thm:champ} 已是完整的、精确的、不依赖 A1$'$ 的成果，
可以直接作为毕业论文的一个独立章节交付，不需要等待上界。
\end{enumerate}

\subsection{本轮新增的候选方向（庚、戊、己）}

在方向甲至丁之外，本轮又考虑了三条新方向。它们与甲—丁的区别在于：
不再试图构造更好的主函数（定理~\ref{thm:zonal} 已经封死这条路），
而是直接处理定理~\ref{thm:zonal} 指出的那两类缺失信息。

\textbf{方向庚：攻掩码，而不是攻支撑（建议取代 (2a) 成为主线缺口）。}
方向甲的缺口 (2a) 是「极值测度的原子数 $\le6$」。但重新审视定理~\ref{thm:selfcons} 会发现，
真正决定问题的是\emph{掩码} $A_H=1\{Q>0\}-I$ 的组合类型，而不是原子个数：
冠军的负图 $\Gamma^-=C_5$，而我们已经知道 $\Gamma^-$ 无三角形、且满足等权刚性。
于是提出一个比 (2a) 更弱、因而更有希望的新猜想：
\emph{任何达到值 $\ge\rho_0$ 的自洽构型，其负图必为 $C_5$ 的吹起（允许孤立顶点与零权原子）。}
若此猜想成立，掩码就被完全钉死为 $K_1\vee C_5$ 的符号模式，
问题立即退化为「固定符号模式下的有限维极值」，而后者已由定理~\ref{thm:champ} 解出。
攻它的可行路径是按 $\Gamma^-$ 的极小禁子结构分情形：每一情形是一个低维多项式极值问题，
可用 SOS 证书加区间算术逐块认证上界 $<\rho_0$。
与 (2a) 相比，这条路的分情形对象是\emph{图}而不是\emph{测度}，
可用的结构信息多一个量级（三角形自由、奇圈长度、度序列都能用），这是它更有希望的原因。

\textbf{方向戊：$\varepsilon$-稳定性定理 + 计算机辅助有限覆盖。}
不直接证 $c_3^{\Sigma}\le\rho_0$，改证稳定性：存在显式 $\varepsilon_0,C>0$，
使任何值 $\ge\rho_0-\varepsilon$（$\varepsilon<\varepsilon_0$）的自洽构型距离冠军吹起族不超过 $C\sqrt\varepsilon$。
证明分两段。「近段」是冠军处的二阶展开：常数全部在 $\Q(\sqrt5,\sqrt7)$ 中，
Hessian 在对称轨道（$D_5$ 与 $O(2)$ 的商）横向上的负定性可用区间算术认证；
符号模式稳健性 $\min|\Pi_3|=0.1900$ 保证展开邻域内掩码不变。
「远段」需要把「距冠军 $\ge\eta$」的构型空间分成有限块并逐块给出显式上界 $<\rho_0$，
这是真正的难点：块数可能很大，但每块是一次凸可行性检查，可并行、可认证。
这种「解析主部 + 计算机辅助有限核验」的证明形态在本领域有被接受的先例：
第五大投影常数的确定 [DFFL22] 与八维、二十四维球堆积 [Via17, CKM$^+$17] 都是这一结构。
风险在于远段块数爆炸，以及 Hessian 商掉对称后是否仍有零方向（目前未验证）。

\textbf{方向己：交替动力系统的全局吸引性。}
引理~\ref{lem:kyfan} 的「模式 $\leftrightarrow$ 权」交替映射使目标单调不减。
若能证明 (i) 任何极限点都是自洽不动点；(ii) 不动点集模对称后有限；
(iii) 除冠军外每个不动点的值都 $<\rho_0$ 且局部为鞍或排斥，
则从任意初值出发的迭代都收敛到冠军，冠军从而是全局最大。
(ii) 可由不动点方程组的多项式性质经半代数几何得到有限性；
(iii) 的缺口是不动点个数未知。即使不能完成，(i)+(ii) 本身是一个干净的结构定理，
可以单独作为论文的一节。

\textbf{优先级（更新）。} 庚的掩码刚性猜想取代 (2a) 成为主线第一缺口；
戊作为主线证明形态的蓝图；等权子问题 $m_{\rm eq}$ 的正确重算（第~7.3 节，
产出改进 \STW\ 约 $0.77\%$ 的无条件上界）保持为保底发表件；己作为结构副产品。
四者共享同一套基础设施：定理~\ref{thm:form9} 的形式 9、区间算术认证管线、
以及命题~\ref{prop:nec} 的 (A)/(B) 筛选器。

\section{诚实声明与结论状态}\label{sec:honest}

本节把全文的结论按证明强度重新排一遍，使读者不必回到正文逐条核对，
就能准确知道课题现在站在什么位置、哪些东西可以引用、哪些不可以。

\subsection{可以引用的结果}

第一类是下界一侧的完整结果，也是目前成熟度最高的部分。
定理~\ref{thm:lower} 给出 $\liminf_{n}M_3(n)/n\ge\rho_0$，
其中 $\rho_0=(9+\sqrt5+2\sqrt{35})/18$，小数展开为 $1.281568196872$，
见证是 $K_1\vee C_5$ 的一个六原子团吹起。
定理~\ref{thm:champ} 进一步给出该极值核心的完整闭式刻画：
最优权为 $s=1.586541134689$、$t=0.976595259055$，
由一阶条件 $a^2+24ab-31b^2=0$ 与归一化条件解出；
对应的系数是 $a^*=(91-25\sqrt7)/126$ 与 $b^*=(7+5\sqrt7)/126$，
其比值 $a^*/b^*=5\sqrt7-12$，五个核函数值全部落在 $\Q(\sqrt5,\sqrt7)$ 中。
这两条不依赖任何未证猜想，其中的数值等式已由 \texttt{sympy} 在相应数域上判定为恒等式。

第二类是归约链与结构定理。定理~\ref{thm:reduce} 把「对一切 $n$、一切图」的陈述
压缩成一个有限维极值问题；定理~\ref{thm:gam3} 给出 $\gamma(3)=\varphi$ 的精确值；
定理~\ref{thm:selfcons} 给出极值对象必须满足的自洽条件 $A=1\{\Pi_3(M)\ge0\}$；
定理~\ref{thm:caseab} 给出 Case A / Case B 二分，并把全部难度定位到 $t\in(0.9451,0.995]$ 这条窄带上。
链条本身是严格的，不严格的只是它末端的猜想~\ref{conj:A1}。

第七类（第六版新增，也是本轮最强的一条）是\textbf{完全二部掩码的全局闭合}。
定理~\ref{thm:bipglobal} 说：对任意 $a,b\ge1$、任意阶数 $N=a+b\ge4$、任意权 $p\in\Delta_{N-1}$，
$\Lam_3(D_{\sqrt p}(I+A(K_{a,b}))D_{\sqrt p})\le\frac54$，
取等当且仅当 $a,b\ge2$、$p$ 恰有四个非零分量、每部各两个、都等于 $\frac14$；
故 $F(\Sigma_{K_{a,b}})=\frac54$（$a,b\ge2$）是\textbf{精确值}。
这一条\emph{不含任何数值成分、不含任何未证前提}，而且是对\emph{一切}阶数同时成立的全局陈述，
可以直接作为毕业论文的一个独立定理章节。它的证明链条只有四步，全部初等：
引理~\ref{lem:biprank}（秩二分解 \eqref{eq:rank2} $+$ Sylvester 惯性律 $\Rightarrow\operatorname{neg}(M)\le1$）、
引理~\ref{lem:bipweyl}（$R\succeq-\sqrt{\alpha\beta}e_-e_-^\top$ $+$ Weyl $\Rightarrow\lambda_i\ge q^\downarrow_{i+1}$）、
引理~\ref{lem:bipsec}（秩二长期方程 \eqref{eq:secular} $+$ 单调性 $\Rightarrow$ \eqref{eq:mono}）、
引理~\ref{lem:core4}（核心引理 $s\le\frac14+\sigma_4$，用 \eqref{eq:Delta}、\eqref{eq:philin}、Titu、AM--GM）。
其中引理~\ref{lem:core4} 的等号条件是逐步反推出来的（四个非零权全为 $\frac14$、$r=2$），
这一点值得强调：$\frac54$ 不是猜出来再验证的，而是被取等条件\emph{唯一确定}的。
配套的三条已证结果是引理~\ref{lem:mpsec}（完全多部的秩 $k$ 长期方程 \eqref{eq:mpsec} 与
$\operatorname{inertia}(R)=(1,k-1)$，含 $k=2$ 的特例）、推论~\ref{cor:n5uncond}
（$\kappa_3(5)\le\rho_0$ \textbf{无条件}，$F(\Sigma_{K_{2,3}})=\frac54$ 精确）、
以及命题~\ref{prop:tab19}（$N=6$ 剩下 $19$ 类的完整表，含 graph6 码，可复算）。
注~\ref{rmk:bipvsframe} 给出了这条路线与 Grassmannian 路线的难度对比，
注~\ref{rmk:mpopen}、注~\ref{rmk:k1b}、注~\ref{rmk:tab19read} 分别精确指出了三处\textbf{失效边界}
（$k\ge3$ 的尾和退化、$\min(a,b)=1$ 族真值 $\frac65$、$19$ 类的结构分布）——
这三条是「已证结果的适用范围」的一部分，引用时不应省略。

第三类是本轮的主要新贡献，即定理~\ref{thm:zonal}：
定义~\ref{def:scheme} 意义下的整个通道型主函数方案族，其可达到的最好常数恰好是 $\Phi$。
这条定理的证明是初等的，见证测度 \eqref{eq:lam} 是显式的三点测度，
因此它的可靠性不依赖任何数值计算；随附脚本的作用只是对引理~\ref{lem:moments} 中
「对一切 $\ell\ge3$ 有 $\int P_\ell\dd\lambda\ge0$」这一唯一需要分档讨论的部分做机器复核。
命题~\ref{prop:nec} 把它转化为一条可操作的筛选判据。

第四类是对旧路线（配对松弛）的完整否证：定理~\ref{thm:43A} 的天花板 $\frac12(1+\sqrt3)$、
命题~\ref{prop:43B} 的显式反例、定理~\ref{thm:43C} 的三原子精确最大值 $\frac43$，
共同说明该松弛的值严格大于 $\Phi$，从而这条投入最大的路线不可能证明猜想~\ref{conj:A1}。

第五类是本轮补上的近段与腔面结果。定理~\ref{thm:w3} 把定理~\ref{thm:locrig} 的局部刚性升级为
带显式三次余项常数的完整近段稳定性不等式（$C_3=2.521337009637$、适用半径 $R=0.024942655$），
其三阶导张量由 $50$ 位截断幂级数精确求出、经五点中心差分与 $210$ 方向线性方程组残差双重背书；
第~\ref{sec:wallcert} 节把三类单条目腔面的 argmax 认证为严格局部最大（Weyl 加 $50$ 位）；
引理~\ref{lem:cover} 与定理~\ref{thm:assemble} 把整条归约链装配成一个条件性主定理，
明确指出主线仅剩 W1-全局与 W2 两个具名缺口。
这些结果的\emph{局部}与\emph{条件}性质必须连同结论一起引用。

第六类是本轮的主要正面贡献，也是成熟度仅次于第一类的部分：优占对偶。
定理~\ref{thm:domdual} 建立恒等式 $F(\Sigma)\le G(\Sigma)=\Omega(\Sigma)$，
其中 $\Omega$ 是半定规划、$G$ 是核范数形式，等号由强对偶给出；证明是初等的（Ky Fan 单调性 $+$ Lagrange 对偶），
不依赖任何数值计算。定理~\ref{thm:champglobal} 在冠军掩码上把这个松弛\textbf{精确}求出：
$F(\Sigma_0)=G(\Sigma_0)=\Omega(\Sigma_0)=\rho_0$，证书 $T$ 只有三个根式常数，
八个恒等式与两个正性条件全部由 \texttt{sympy} 在 $\Q(\sqrt5,\sqrt7)$ 上判定为 $0$ 或正；
配合定理~\ref{thm:lower} 的 $\ge\rho_0$，这是一个\textbf{双向夹紧}的定理，不含数值成分。
推论~\ref{cor:chamberclosed} 由此闭合了 W1-全局的冠军腔部分。
定理~\ref{thm:wallsexact} 给四类近腔掩码配上精确有理数证书（Sylvester 判据，$63$ 个主子式逐个在 $\Q$ 上判非负），
其中 class2、class3 的腔面全局最大值由「数值」升级为「已证且精确」。
命题~\ref{prop:octnogo} 与命题~\ref{prop:negmass} 是配套的：
前者给出一条精确的否定性结论（$\Omega(\Sigma_{\mathrm{oct}})=4/3$），
后者曾给出 A1$'$ 的一个凸的、尖锐的充分形式（猜想~\ref{conj:A2}）——
但该猜想\textbf{已在本轮被否证}，见下（第七类）。
引用第六类结果时请注意其\textbf{适用范围}：它覆盖的是掩码，不是任意图；
「$\Omega(A)\le\rho_0$ 对一切自洽 $A$」仍未证。

第七类是本轮的全局重构，其中\textbf{已证}与\textbf{数值}两部分必须分开引用。
已证的部分有四条：
(i) 定理~\ref{thm:kappa} 的全局重构与 $\Lam_3(G)\le\kappa_3\,n-3$（证明是自足的，不依赖定理~\ref{thm:reduce} 的规范化段落）；
(ii) 推论~\ref{cor:kappamain}：主定理 $\iff\kappa_3\le\rho_0$；
(iii) 命题~\ref{prop:perronrig} 的 (a)--(h) 全部恒等式，以及命题~\ref{prop:decouple} 的解耦恒等式
（两者都是纯代数推导，证书 A/B 段以 $10^{-9}\sim10^{-8}$ 逐条核验，$64$ 项全绿；
残差的唯一来源是 Perron 方程的数值残差 $1.8\times10^{-11}$，而 $Q$ 本身满足 $\|Q^2-Q\|_{\max}=4.4\times10^{-16}$）；
(iv) 定理~\ref{thm:ladder} 中 $\kappa_3(3)=1$、$\kappa_3(4)=\frac54$、$\kappa_3(5)=\frac{9+4\sqrt6}{15}$
三个\emph{精确}值（\texttt{sympy} 判定为恒等式；$\kappa_3(4)$ 含三类符号分布的完整不等式证明）。
\textbf{数值}的部分是：$N=6,\dots,14$ 处 $\kappa_3(N)=\rho_0$
（$N=6,\dots,12$ 由 SCF 与 Stiefel 上升两套独立求解器同时给出，KKT 残差 $\le3\times10^{-9}$、
与闭式差 $\le1.6\times10^{-15}$，证书 F 段独立复跑 $N=3,\dots,8$；$N=13,14$ \emph{只有} SCF 一套，
证据强度低一档），以及第~\ref{sec:exh7} 小节重跑后的 $F$ 表
（$1252$ 类全部跑完，$\max F=\rho_0$、$\#\{F>\rho_0+10^{-9}\}=0$、自洽 $115$ 个）。
前者是本课题最强的数值证据——它在\emph{全局}（而非单个腔面）上做极大化，
且用两条独立算法（自洽场 SCF 与流形上升）得到同一常数——但它仍然\textbf{不是证明}，
因为没有区间算术证书、也没有对全部 $156$ 个符号类穷举临界点。
此外，注~\ref{rmk:normcaution} 是一条\emph{自我否证}：定理~\ref{thm:reduce} 原证明的规范化步骤有误，
虽然结论已被 (i) 独立重证，但原证明在修复前不应被单独引用。

第八类是本轮的\textbf{精确否证}结果，成熟度与第一类同级（因为它是一条\emph{定理}，不是猜想）。
命题~\ref{prop:A2no}：猜想~\ref{conj:A2}（负质量界）为假。
显式反例是 $A=I_7+A_H$、$H=\overline{\widehat{C_5}}$
（$\widehat{C_5}=C_5$ 的两个相邻顶点各挂一片叶子；$H$ 有 $14$ 边、$\omega=4$、$\alpha=2$、$|\mathrm{Aut}|=2$），
$p$ 为 $F(A)$ 的极大化点，$\tr((-M)_+)=0.288182336567654\ldots>\tau=0.281568196872168\ldots$。
证明\textbf{完全无浮点}：$M$ 相似于有理矩阵 $A\,\mathrm{diag}(\hat p)$（$\hat p$ 是 $p$ 的精确有理化，
$|\hat p-p|_\infty\le2.8\times10^{-17}$），故特征多项式在 $\Q$ 上精确可算；
对它自实现 Sturm 序列做二分隔离与精化（宽度 $<10^{-32}$），得 $\tr((-M)_+)$ 的严格有理区间；
$\tau$ 的严格有理区间由整数开方 \texttt{isqrt} 给出（宽 $1.7\times10^{-46}$，
包含性由 \texttt{sympy} 在 $\Q(\sqrt5,\sqrt{35})$ 上精确判定）；两区间不相交。
自洽性由 $80$ 位 \texttt{mpmath} 认证（谱投影三条残差 $\le5.6\times10^{-82}$，符号余量 $4.9\times10^{-2}$），
另有 $2000$ 次 $10^{-9}$ 扰动的鲁棒性体检（自洽性破坏 $0$ 次）。
配套的两条已证结果是注~\ref{rmk:A2why}（失效机制：A2 $=$ 删去 Ky Fan 对偶的迹约束 $\tr Z=N-3$，
等号恰在 $\mathrm{npos}=3$ 时成立）与命题~\ref{prop:A2sharp}（修复形式 A2$^\sharp$ 与 A1$'$ 等价）。
注~\ref{rmk:A2strata} 的分层表（$115$ 个自洽掩码按 $\mathrm{npos}$ 分为 $3/103/9$）属于\textbf{数值}，
但它对每个掩码都是可复算的确定量，不含任何猜测成分。
\textbf{必须强调}：这条否证\emph{不}影响 A1$'$——反例的 $\Lam_3=1.257711<\rho_0$。

\subsection{本报告没有主张的内容}

把边界划清楚与说明成果同样重要。

我们\textbf{没有}改进 $\Lam_3(G)\le\Phi n-3$ 这条无条件纪录，缺口 $\Phi-\rho_0=0.0274488$ 一寸未动；
现役的最好上界仍然属于 \STW。猜想~\ref{conj:A1}（即 $c_3^{\Sigma}=\rho_0$）\textbf{未证}。
具体地，第~\ref{sec:assemble} 节条件性主定理~\ref{thm:assemble} 的两个前提——
W1-全局与 W2——\textbf{均未完全 discharged}。W1-全局的当前精确状态是：
冠军腔\emph{已证}（定理~\ref{thm:champglobal}，$F(\Sigma_0)=\rho_0$，闭式证书，无数值成分）；
class1、class2、class3 三类单条目腔面与 class12 一类双条目腔面\emph{已证} $<\rho_0$
（定理~\ref{thm:wallsexact}，精确有理数证书，余量 $7.5\times10^{-3}\sim5.1\times10^{-2}$）；
$N\le7$ 的 $1252$ 个同构类中 $832$ 个已证（第~\ref{sec:exh7} 小节）。
\textbf{未证}的是：$\Omega>\rho_0$ 的那批掩码（$N\le7$ 中 $420$ 个，最大 $\Omega=1.3733$；
其中包括双条目近腔 class13，$\Omega=1.30667$）、$N\ge8$ 的全部情形、
以及最关键的——「自洽 $\Rightarrow\Omega\le\rho_0$」这条蕴含关系（猜想~\ref{conj:A2} 是它的一个更强的凸形式）。
W2（远腔度序列型上界）仍只有否证性数值扫描，没有引理。
定理~\ref{thm:w3} 的近段稳定性只在半径 $R=0.024942655$ 的冠军腔内邻域成立，\emph{不是}全局上界；
冠军腔的全局上界是另一条定理（推论~\ref{cor:chamberclosed}），两者不可混用。
第~\ref{sec:exh7} 小节中对 $420$ 个掩码计算的 $F(A)$ 值全部是\textbf{数值}，不构成证明；
只有由 $\Omega(A)\le\rho_0$ 排除的那 $832$ 个是已证的。
支撑界 (2a)（极值测度的原子数 $\le6$）\textbf{未证}，目前无条件的只有 $N\le20$
（由 $N\le24$ 经 Fenchel--Bunt 精化），而定理~\ref{thm:selfcons} 要压成有限的核心表需要 $N\le8$；
必须强调 $20\ne8$，把「支撑 $\le20$」表述为「已证 A1$'$」是严格禁止的。
关于等权子问题 $m_{\rm eq}$ 的\textbf{一切数值上界都不被采信}，理由见第~7.3 节：
$n\ge8$ 的搜索结果与已证的复制单调性直接矛盾，属于优化伪影；
目前可用的只有下界 $m_{\rm eq}\ge1.0590$（正二十面体，精确）与初步的 $1.245304$（数值，未认证）。
第~6.2 节的 $\mathrm{SR}_K$ 计算只是\textbf{阴性证据}，不构成任何命题的证明，不应被引用为成果，
其中 $K=15,\dots,20$ 各档尚未计算。
第~5 节表格中标注「数值」或出现在猜想环境中的条目（例如 $K\ge4$ 时 $\max B=3/\sqrt5$）一律不是定理。
关于第~\ref{sec:kappa} 节，我们\textbf{没有}主张：$\kappa_3\le\rho_0$（这正是待证的 (W1$''$)）；
$\kappa_3(N)=\rho_0$ 对任何 $N\ge6$ 的\emph{证明}（只有流形上升的数值认证）；
$N=6$ 最大值点的 $C_5$ 闭式（对称性偏差 $2\times10^{-8}$，与优化器精度同阶，见障碍第 6 条）；
以及「Perron 刚性恒等式组足以推出 $\rho\le\rho_0$」——注~\ref{rmk:rigiditylimit} 明确否证了这一点，
它最多给到 $\frac32$。我们也\textbf{没有}主张 $\frac1N\|Q_+\|_1$、行和型或优占型上界能达到 $\rho_0$：
冠军处 $\frac1N\|Q_+\|_1=1.279704891<\rho_0$，这一族被\emph{定量}排除。
最后，定理~\ref{thm:reduce} 的规范化步骤在修复前\textbf{不得}单独引用（注~\ref{rmk:normcaution}）。
关于 A2 的否证，我们\textbf{没有}主张：A1$'$ 为假（反例的 $\Lam_3$ 比 $\rho_0$ 低 $2.4\times10^{-2}$，
A1$'$ 完好）；$\mathrm{npos}\le3$ 档上 A2$'$ 的\emph{证明}（只有 $N\le7$ 全部 $106$ 个自洽掩码的穷举核查，
$\max\tr((-M)_+)$ 恰为 $\tau$、由冠军饱和，这是\textbf{数值}）；
以及 A2$^\sharp$ 比 A1$'$ 好证——命题~\ref{prop:A2sharp} 说明二者\emph{等价}，
A2$^\sharp$ 的唯一价值是指认了正确的对象（尾谱和而非负质量）。

关于第~\ref{sec:bipglobal} 节（第六版新增），我们\textbf{没有}主张以下几条，请引用时注意边界：
(i) \eqref{eq:frame5} \textbf{没有被证明}。定理~\ref{thm:bipglobal} 使它变得\emph{不必要}，
而不是使它成立；$B(Z)$ 与 $M_p$ 是两个不同族的矩阵，前者的 $\frac54$ 由命题~\ref{prop:biaxsharp} 知其 sharp，
但它的\emph{证明}仍然缺失。
(ii) $F(\Sigma_{K_{1,b}})=\frac65$ 只有\textbf{下界}是精确的（$K_{1,3}$ 的有理见证），
上界未证；定理~\ref{thm:bipglobal} 给的是 $\frac54$，松 $0.05$（注~\ref{rmk:k1b} 精确指出了失效点）。
(iii) $k\ge3$ 的完全多部掩码（$K_{2,2,2}$、$K_{1,2,3}$、$K_{1,2,2}$、$K_{2,2,3}$ 等）
的 $F$ 值全部是\textbf{数值}：$\sup\Lam_3\in[1.2500767598,\,1.2530756675]$，
虽严格低于 $\rho_0$（余量 $\ge2.85\times10^{-2}$），但没有任何证明；
引理~\ref{lem:mpsec} 只提供了长期方程与惯性两条\emph{工具}，不提供界。
(iv) 「$\kappa_3(5)$ 的取等类是 $K_{1,2,2}$」是\textbf{数值}推测（$\sup\Lam_3=1.2530756675$
与 $\frac{9+4\sqrt6}{15}=1.2531972647$ 差 $1.2\times10^{-4}$，尚未收敛到闭式），不得当作结论。
(v) 命题~\ref{prop:tab19} 表中的 $\max F$ 列与注~\ref{rmk:tab19read}(ii) 的两个未识别闭式
（$1.251351067$、$1.216091742$）都是\textbf{数值}；$\Omega$ 列是已证的 SDP 上界。
(vi) \eqref{eq:atom6}（$N=6$ 的 $19$ 类极大化点在边界）\textbf{未证}，
而且它的自然推广（对一切 $6$ 阶掩码）是\emph{假}的——冠军类就是反例。
(vii) \eqref{eq:s6int} 只是单向蕴含（$\Leftarrow$），本报告\textbf{没有}主张 $\kappa_3(6)\le\rho_0$。

\subsection{数值证据的使用纪律}

全文使用四个标签：「已证」「精确」「数值」「否证」。它们不是修辞，而有明确的操作含义——
「已证」要求本报告或所引文献中含有完整论证；
「精确」要求相关数值等式已由符号计算判定为恒等（差化简后恰为 $0$）；
「数值」表示只有浮点证据；「否证」表示已有反例。
所有标注「数值」的条目都不得作为成果引用。

这条纪律来自一次具体的教训。早期一轮工作曾报告配对松弛的最大值为 $1.366758746$，
几乎取到定理~\ref{thm:43A} 的天花板 $\frac12(1+\sqrt3)=1.366025404$，
当时被当作「主线接近完成」的证据。复核后发现它是伪影：
对应的解满足 $\alpha_2=1.0107>1$，而定理~\ref{thm:43A} 的证明中已经严格建立了 $\alpha\le1$；
根因是白化步骤 $M^{-1/2}$ 的条件数爆炸。
修正办法是在所有搜索代码中加入可行性拒绝门与审计输出（每次打印 $\alpha_1,\alpha_2$ 与各向同性残差），
超出 $1$ 的解一律丢弃。由此得到本报告反复使用的一条原则：
\textbf{任何违反已证不等式的数值「改进」都是伪影，不是发现。}

\subsection{第 12 轮新增结论的证明状态标注}\label{sec:honest12}

按第~\ref{sec:honest} 节的纪律，第~\ref{sec:rowvar} 节的结论逐条标注如下。

\textbf{已证（含完整论证，且不依赖任何未证猜想）}：
命题~\ref{prop:chamberdead}（$\Q$ 上精确，无浮点）；
引理~\ref{lem:rowid} 的三条恒等式（含三角形恒等式 \eqref{eq:tri}）；
命题~\ref{prop:fourcs} 的四条 Cauchy--Schwarz 不等式 $[f],[g],[h_1],[h_2]$；
定义~\ref{def:rel} 的松弛 REL 与 $\kappa_3(N)\le\mathrm{REL}(N)$；
定理~\ref{thm:LN}（$L(N)$ 闭式，含 AM--GM 取等点的可行性验证）；
命题~\ref{prop:fracdeg}（$n^*$ 的整性刻画与 Pell 例外集 $\mathcal E$ 的显式两支递推）；
定理~\ref{thm:unattain}（例外阶数上的强正则图障碍；(b) 支为符号谱计算，(a) 支为\emph{纯整数}判别式检验，
覆盖 $N\le7.1\times10^8$ 的 $13$ 个例外阶数）；命题~\ref{prop:ladder6}（$N=6$ 整数化阶梯的六个闭式）；
命题~\ref{prop:twovalue}（两值秩三投影的谱分类，全部由符号特征多项式给出）；
引理~\ref{lem:z5class}（$Z_5$-等变秩三投影的完全分类，特征多项式恒等式 $L^3(L-1)^3$）；
定理~\ref{thm:conefam}（$\rho(s)$ 闭式、唯一临界点、$\rho(s^*)=\rho_0$、精确 Perron 权重、ETF 识别）；
定理~\ref{thm:convex}（$\Phi$ 的凸性与 $\operatorname{conv}(\mathcal P_3(N))=\mathcal S_N$）；
引理~\ref{lem:perp2}（垂直二标架分解）；命题~\ref{prop:corank1} 与推论~\ref{cor:k23}（$\kappa_2(3)=\frac87$）；
引理~\ref{lem:slope1}（$\rho=\frac{\alpha+\ell}2$）。

\textbf{精确（数值等式已由符号计算判定为恒等，差化简后恰为 $0$）}：
冠军处四条 CS 同时取等（残差 $0$，非 $10^{-16}$ 量级的近似）；
$\|Q_-\|_F^2=\frac{16+5\sqrt7}{81}=\rho_0\tau$；
$\lambda^*(\text{冠军数据})=\rho_0$；$L(4)=\frac54$ 且被 $Q=I_4-vv^\top$ 取到；
$\rho(\frac12)=\frac{5+\sqrt5+\sqrt{30}}{10}$（ETF）；$\pi_h=\frac{91-25\sqrt7}{126}$、$\pi_c=\frac{7+5\sqrt7}{126}$；
$\mathrm{spec}(Q_+)$ 与 $\mathrm{spec}(Q_-)$ 的四个根式；
$\kappa_2(3)$ 见证的特征多项式 $\frac1{343}(7L-8)(7L-5)(7L-1)$；
$6$ 阶会议矩阵的 $S^2=5I_6$。

\textbf{数值（只有浮点证据，不得作为成果引用）}：
$\mathrm{REL}(N)=L(N)$ 对 $N=5,6,7,8,10,12$（差分进化，吻合到 $\le4\times10^{-6}$；\emph{猜想}，非定理）；
$Z_5$-对称层八个图案类的天花板（其中 $6/8$ 类被 REL 单独排除，冠军所在类为 $1.284780995477$）；
$19/19$ 未闭合类都是腔掩码（符号违背量 $\le5\times10^{-19}$）；
$\alpha=1$ 层的 $\max\beta=0.258695$；通道线性规划的 $8701$ 个构型与 HiGHS 解；
$\Phi$ 的凸性在 $N=5,6,7$ 上共 $24000$ 对随机点无 Jensen 违背（\emph{注}：定理~\ref{thm:convex}(a) 本身已证，
这条数值只是独立交叉核验）；
$E_\perp=\frac12x^\top|B|x=0.519536756454917$（未找到闭式）。

\textbf{否证}：命题~\ref{prop:chamberdead} 否证了 \eqref{eq:chamber6} 与第六版队首 P1；
注~\ref{rmk:tab19read}(iv) 的「等价」降级为单向蕴含；
不等式 $[g]$ 的取等\emph{不是}冠军的刻画（$N=4$ 上允许 $\rho=1.4186$，见第~\ref{sec:miscid} 小节 (6)）。

\subsection{一句话总结}

下界一侧已经完整：$\rho_0$ 是精确的，极值核心已闭式刻画，这部分可以立即写成毕业论文的独立章节。
上界一侧，本轮把「极值核心的全局最优性」从猜想变成了定理：
$F(\Sigma_0)=\Omega(\Sigma_0)=\rho_0$ 已证且精确，证书是三个根式常数与八个符号恒等式（定理~\ref{thm:champglobal}），
四类近腔掩码另有精确有理数证书（定理~\ref{thm:wallsexact}）。
方法论上的收获是把非凸问题\textbf{提升}到一对原对偶半定变量 $(T,Y)$ 上（定理~\ref{thm:domdual}），
这与此前被否证的通道型主函数路线（定理~\ref{thm:zonal}）、被证明零间隙的球面松弛（注~\ref{rmk:tight}）
构成一组完整的「哪条路走得通、哪条走不通」的图谱。
剩下的一步，本报告先后给了三种表述，前两种\textbf{都已被否证}：
「逐点主函数 $+$ 逐通道计费」被定理~\ref{thm:zonal} 封顶在 $\Phi$；
「负质量界 $\tr((-M)_+)\le\tau$」（猜想~\ref{conj:A2}）被命题~\ref{prop:A2no} 用一个 $7$ 原子自洽核心
精确推翻（超出 $\tau$ 达 $6.6\times10^{-3}$）。两次失败的原因\emph{相同}：
它们都删去了 $\Pi_3$ 的秩三结构——前者删在通道分解上，后者删在 Ky Fan 对偶的迹约束 $\tr Z=N-3$ 上
（注~\ref{rmk:A2why}）。命题~\ref{prop:octnogo} 则从另一侧保证了这一步绕不开自洽性。
于是本轮把主线换成第三种表述，也是唯一保留秩三结构的那种：
主定理现在\emph{等价}于单一常数不等式 $\kappa_3=\sup_N\sup_{V\in\mathrm{St}(N,3)}\rho((VV^\top)_+)\le\rho_0$
（推论~\ref{cor:kappamain}），它不含图、不含掩码、不含阶数，一阶条件有闭式 \eqref{eq:stiefelkkt}，
且 $N=3,4,5$ 三档已被\emph{精确}求出（定理~\ref{thm:ladder}）、
$N=6,\dots,14$ 已数值认证为 $\rho_0$（$6\le N\le12$ 两套独立求解器，$13,14$ 一套）。
同时我们也精确地知道\emph{为什么还不能收工}：Perron 刚性系统本身只能给到 $\frac32$
（注~\ref{rmk:rigiditylimit}），\STW\ 的损失被劈成两半而必须\emph{联合}处理（命题~\ref{prop:decouple}），
缺的是一条形如 $\alpha+\ell\le2\rho_0$ 的第四不等式。
至于无条件纪录，方向乙的等权子问题仍可望把常数压到 $1.299$ 左右（相对 \STW\ 改进约 $0.77\%$），
但它到不了 $\rho_0$；要证到 $\rho_0$，现在唯一在列的路径是 (W1$''$)（$\kappa_3\le\rho_0$）
配合支撑界 (2a) 或 $N\ge7$ 的有效原子化约（优先队列 P2）。

\medskip
\noindent\textbf{（第六版增补）一句话总结的更新。}
上面这段是第五版写的，其中「$N=3,4,5$ 三档已被精确求出」与「唯一在列的路径是 (W1$''$) 配合支撑界 (2a)」两处需要修正。
现在的状态是：$\kappa_3(3)=1$、$\kappa_3(4)=\frac54$、$\kappa_3(5)=\frac{9+4\sqrt6}{15}$ 三档\textbf{精确}（定理~\ref{thm:ladder}），
并且 $\kappa_3(N)\le\rho_0$ 对 $N=3,4,5$ 已\textbf{无条件}成立（推论~\ref{cor:k34uncond}、\ref{cor:n5uncond}）——
第六版之前 $N=5$ 那一档是条件性的，条件 \eqref{eq:frame5} 现已被定理~\ref{thm:bipglobal} 绕开。
同时我们拿到了第一条\emph{跨阶数}的定理：一切完全二部掩码、任意阶数，$\Lam_3\le\frac54$（定理~\ref{thm:bipglobal}），
取等构型完全刻画。这改变了主线的形状——剩下的路不再是「证一条越来越难的不等式」，
而是「把一条纯几何的不存在性陈述 \eqref{eq:chamber6} 对 $17$ 个符号模式逐个建立」（优先队列 P1），
外加 $2$ 个有支配点的类（P2）与 $k\ge3$ 的尾和改进（P3）。
至于毕业论文的写作，现在有三个可以\textbf{独立成章}的完整结果：
下界一侧的 $\rho_0$ 与极值核心闭式刻画（定理~\ref{thm:lower}、\ref{thm:champ}）、
冠军腔的全局最优性（定理~\ref{thm:champglobal}，原对偶半定证书）、
以及完全二部掩码的全局闭合（定理~\ref{thm:bipglobal}，初等证明）；
三者都不依赖任何未证猜想，且方法论互不重叠（团吹起 $+$ 根式恒等式 / 半定对偶 / 秩二长期方程 $+$ 初等不等式）。

\medskip
\noindent\textbf{（第七版增补）一句话总结的更新。}
第六版把主线定为「对 $17$ 个符号模式逐个建立 \eqref{eq:chamber6}」。第 12 轮把这条路线\textbf{否证}了：
$19/19$ 类都是腔掩码（命题~\ref{prop:chamberdead}，$K_{2,2,2}$ 有精确有理证书），
所以「腔掩码不存在性」这条路能关闭的类数是 $0$，第六版增补中「唯一在列的路径」的说法作废。
现在的状态是：\textbf{主线从「逐类排除」换成「一条不等式链的逐级收紧」}，
链条是 $1.366025$（Mohar）$>1.309017$（$\Phi$，等于 REL 的自由上确界）$>1.300000$（$\frac{13}{10}$，加度的整性）
$>1.284781$（加 $Z_5$ 对称）$>1.281568$（$\rho_0$，加秩三）。
前两级已经\textbf{定理化}（定理~\ref{thm:LN} 与命题~\ref{prop:ladder6}），
第三级是数值天花板，第四级由定理~\ref{thm:conefam} 精确闭合。
于是「还差什么」第一次有了一个\emph{定量且定位清楚}的答案：差 REL 与秩三之间的 $3.2\times10^{-3}$，
缺的原料是三角形恒等式 \eqref{eq:tri} 与四点 Gram 行列式（第~\ref{sec:diag} 小节的 (A)(B) 两项）。
另外，$\rho_0$ 现在不再只是一个见证：它是 $Z_5$-等变秩三投影弧的唯一极大值（定理~\ref{thm:conefam}），
$\sqrt{35}$ 来自判别式 $140$，$\sqrt7$ 来自 Perron 权重的二次方程——两个根式都有了出处。
至于毕业论文的章节结构，第七版之后可以独立成章的结果增加到\textbf{五}个：
上面三个，加上「行方差机制与天花板 $L(N)$」（定理~\ref{thm:LN}，含 Mohar 与 \STW\ 两个常数的统一重导），
以及「$Z_5$-等变弧的完全分类与 $\rho_0$ 的闭式导出」（定理~\ref{thm:conefam}，含 ETF 的意外识别）。
这两章都是纯数学的、自包含的、不需要任何数值验证即可宣读。

\appendix
\section{复现说明与产物清单}\label{app:repro}

\subsection{运行环境}

全部计算在 Windows 上完成，解释器为 \texttt{C:\textbackslash Python314\textbackslash python.exe}，
依赖 \texttt{numpy}、\texttt{scipy}、\texttt{sympy} 三个库。
统一的调用方式是把 \texttt{-X utf8} 与 \texttt{-u} 一起打开：前者避免 Windows 控制台默认 GBK
编码在打印含根号或希腊字母的中间结果时报错，后者保证长任务的输出实时可见。

\begin{center}
\fbox{\texttt{C:\textbackslash Python314\textbackslash python.exe -X utf8 -u <脚本名>}}
\end{center}

本报告由 \LaTeX\ 源码编译，编译器为 MiKTeX 的 \texttt{xelatex}，文档类为 \texttt{ctexart}。
因为含目录与交叉引用，需要连续运行两遍；若日志中仍出现 \texttt{undefined references}，运行第三遍即可。

\subsection{随本报告分发的文件}

{\small
\begin{longtable}{@{}p{6.9cm}p{6.4cm}p{2.0cm}@{}}
\toprule
文件 & 内容 & 状态 \\
\midrule
\endfirsthead
\toprule
文件 & 内容 & 状态 \\
\midrule
\endhead
\texttt{进展报告\_前三大特征值之和的精确上界.tex} & 本报告源码（保留的 \LaTeX\ 版本） & --- \\
\texttt{进展报告\_前三大特征值之和的精确上界.pdf} & 本报告编译产物 & --- \\
\texttt{Sivashankar-Tang-Wakhare\_}\newline
\texttt{2608.02429\_GraphEigenvalues}\newline
\texttt{AndProjectionConstants.pdf} &
\STW\ 原文，arXiv:2608.02429v1，2026 年 8 月 3 日，CC BY 4.0；第~\ref{sec:stwanat} 节的解剖对象 & 他人成果 \\
\texttt{证书\_zonal天花板定理.py} & 定理~\ref{thm:zonal} 见证测度的四个矩与 Legendre 各档非负性的 \texttt{sympy} 复核 & 本报告 \\
\texttt{证书\_负行挤压恒等式.py} & 定理~\ref{thm:squeeze} 四条恒等式与冠军处精确饱和表的复核 & 本报告 \\
\texttt{证书\_冠军局部刚性.py} & 定理~\ref{thm:locrig} 的 $50$ 位 Hessian、$\kappa$、谱隙与腔半径认证 & 本报告 \\
\texttt{证书\_三阶稳定性.py} & 定理~\ref{thm:w3} 的三阶导张量、$C_3$、近段半径 $R$ 与有限差复核 & 本报告 \\
\texttt{证书\_腔面局部认证.py} & 第~\ref{sec:wallcert} 节三类单条目 argmax 的 Weyl 局部认证 & 本报告 \\
\texttt{证书\_双条目模式扫描.py} & 第~\ref{sec:wallcert} 节 P1--P6 双条目模式上界（数值） & 本报告 \\
\texttt{证书\_腔面全局归约.py} & 定理~\ref{thm:wallreduce} 球面归约、KKT、紧致松弛否定与否证扫描的复核 & 本报告 \\
\texttt{证书\_优占对偶全局上界.py} & 第~\ref{sec:domdual} 节全部结论：定理~\ref{thm:domdual} 的三列互核、定理~\ref{thm:champglobal} 的八个符号恒等式、定理~\ref{thm:wallsexact} 的 $63$ 个有理主子式、命题~\ref{prop:octnogo} 的 $4/3$、命题~\ref{prop:negmass} 的负质量饱和 & 本报告 \\
\texttt{证书\_极值阶梯与Perron刚性.py} &
第~\ref{sec:kappa} 节全部可证结论，分六段 A--F 共 $64$ 项检查：A/B Perron 刚性与恒等式组 (a)--(h)、
C $\kappa_3(3)=1$ 与 $\kappa_3(4)=\frac54$（含三类符号分布的 \texttt{sympy} 不等式证明）、
D $\kappa_3(5)=\frac{9+4\sqrt6}{15}$（$C_{4v}$ Ansatz 单参数闭式 $+$ 临界点 $u^*=\frac{9-\sqrt6}{30}$）、
E 阶梯 $F(I+A(H))$ 的数值复核（$H=K_3,C_4,C_5,K_4,W_5,W_6,W_7$）、
F \textbf{独立}复核：Stiefel 流形直接上升求 $\max\rho((VV^\top)_+)$（$N=3,\dots,8$）与 Stiefel-KKT \eqref{eq:stiefelkkt} 残差 &
本报告 \\
\texttt{证书\_A2反例与修复.py} &
第~\ref{sec:A2no} 小节全部结论，分五段 A--E 共 $34$ 项检查：A 反例图 $H=\overline{\widehat{C_5}}$ 的识别
（边集、度序列、$\omega,\alpha$、$|\mathrm{Aut}|=2$ 及对合）、B $400$ 起点重跑 $F(A)$ 与全部余量、
C \textbf{精确否证}（有理化 $\hat p$ $\to$ 有理 charpoly $\to$ 自实现 Sturm 序列隔离精化 $\to$
\texttt{isqrt} 给 $\tau$ 的有理区间 $\to$ 两区间不相交，全程无浮点）、
D $80$ 位 \texttt{mpmath} 认证 $\mathrm{sign}(\Pi_3)=2A-J$、E 按 $\mathrm{npos}$ 分档核查修复形式 A2$'$ &
本报告 \\
\texttt{证书\_N5特征多项式闭合.py} &
第~\ref{sec:assemble6} 节之前 $N=5$ 缺口的\textbf{无条件}闭合（特征多项式因式分解路线），分六段 A--F 共 $45$ 项检查：
$\chi$-因式分解 $\det(\lambda I-M)=\prod_i(\lambda-p_i)+N_AN_B$ 的符号核验、
根分布（唯一负根 $\lambda_5=-s$、$\lambda_4\ge p_{\min}$）、核心引理 $s\le\frac14+p_{\min}$
（Jensen $+$ 显式非负分解）、$F(\Sigma_{K_{2,3}})=\frac54$ 精确、$\kappa_3(5)\le\rho_0$ 无条件、
取等点刻画（$p=(\frac14,\frac14,\frac14,\frac14,0)$，零点在 $B$ 部，共 $3$ 个） &
本报告（第五版新增，运行约 $96$ s，输出 \texttt{\_n5chi.json}） \\
\texttt{证书\_完全二部全局闭合.py} &
定理~\ref{thm:bipglobal}、引理~\ref{lem:biprank}--\ref{lem:mpsec}、推论~\ref{cor:n5uncond}、
注~\ref{rmk:mpopen}、注~\ref{rmk:k1b} 的全部内容，分八段 A--H 共 $52$ 项检查：
A 秩二分解与惯性（符号恒等式 $+$ $48000$ 组随机权）、B 长期方程与根分布（$40000$ 组）、
C Weyl 下界与尾和界（$60000$ 组）、D 核心引理（\textbf{全符号}：$\Delta$ 恒等式、$k=3,4$ 链式、
五个 $r$ 分支的 Titu 缺口分子、$2-\Xi(m)$ 的因式分解 $2m(1+14m)$、AM--GM；另 $150000$ 组数值双保险）、
E 装配与取等点穷举（$9$ 个构型）、F 与旧结果的关系、
G 完全多部秩 $k$ 长期方程（$13$ 个构型符号核验）与 $\operatorname{inertia}(R)=(1,k-1)$、
H $\min(a,b)=1$ 族的 $\frac65$ 精确见证 &
本报告（第六版新增，运行约 $137$ s，输出 \texttt{\_kmnbi.json}、\texttt{\_kmn7.json}） \\
\texttt{\_kmn5.json}\newline\texttt{\_kmn6.json}\newline\texttt{\_kmn7.json}\newline\texttt{\_kmnbi.json}\newline\texttt{\_n5chi.json} &
上述两份证书的原始输出。其中 \texttt{\_kmn7.json} 逐类记录 $30$ 个完全（多）部掩码的
$\sup\Lam_3$、$\sup\tr((-M)_+)$、$\min\sum_{i\ge4}\lambda_i$、$\max\operatorname{neg}(M)$ 与极大化点，
是注~\ref{rmk:mpopen} 全部数值的来源 &
本报告（数据） \\
\texttt{\_danger23.json}\newline\texttt{\_exc23\_struct.json}\newline\texttt{\_exc19\_open.json}\newline\texttt{\_tab19.tex} &
例外类的结构数据。\texttt{\_danger23.json} 是 $23$ 个 $\Omega>\rho_0$ 的类（$1$ 个 $N=5$、$22$ 个 $N=6$）的原始记录；
\texttt{\_exc23\_struct.json} 由 graph6 码解码后补算「是否完全多部（及各部大小）」「是否含诱导 $C_5$」；
\texttt{\_exc19\_open.json} 是扣除冠军、$K_{2,4}$、$K_{3,3}$ 之后剩下的 $19$ 类；
\texttt{\_tab19.tex} 是命题~\ref{prop:tab19} 表格的 \LaTeX\ 行 &
本报告（数据） \\
\texttt{\_core.py} &
共用核心：$F(A)=\max_{p\in\Delta}\Lam_3(D_{\sqrt p}AD_{\sqrt p})$ 的稳健优化器（多起点 L-BFGS-B
$+$ softmax 解析梯度 $+$ KKT 诊断）。上列三个证书脚本都 \texttt{import \_core}，故\textbf{必须}随本报告一并分发 &
本报告 \\
\texttt{\_kd.py} &
Stiefel 流形直接上升：在 $\{V:V^\top V=I_3\}$ 上用解析梯度 \eqref{eq:gradfix} 极大化 $\rho((VV^\top)_+)$，
QR 回缩 $+$ 步长回溯。是定理~\ref{thm:ladder} 中 $N\ge6$ 各档\textbf{数值}结论的第二套独立求解器 &
本报告 \\
\texttt{\_kappa\_profile.json}\newline\texttt{\_kappa\_direct.json} &
定理~\ref{thm:ladder} 的原始输出：$\kappa_3(N)$、$\alpha$、$\beta$、$\mathrm{diag}(Q)$、Collatz 行和与 KKT 残差。
前者为自洽场 SCF（$N=3,\dots,14$），后者为 Stiefel 上升（$N=3,\dots,12$）——
正文中「两套独立求解器」与「$N=13,14$ 仅一套」的说法即以此二文件为准 &
本报告（数据） \\
\texttt{\_scan7F.json} &
第~\ref{sec:exh7} 小节 $1252$ 个同构类的完整原始记录（每类：$F(A)$、$\Omega$、$N$、$\mathrm{npos}$、
$\tr((-M)_+)$、自洽标志、$\alpha,\beta,\rho(Q_+)$、三种残差、graph6 码）。
注~\ref{rmk:A2strata} 的分层表、命题~\ref{prop:A2no} 反例的定位、
以及「$F=\rho_0$ 的 $17$ 个掩码全部含诱导 $C_5$」这条结构观察都直接由它算出 &
本报告（数据） \\
\bottomrule
\end{longtable}}

\subsection{证书脚本核对了什么}

\texttt{证书\_优占对偶全局上界.py} 分五块，全部可独立复跑（约 $3$ 分钟）。
(A) 对六个掩码同时用三条独立算法算 $F$（解析梯度 L-BFGS-B 多起点）、$G$（Nelder-Mead/Powell 最大化 $\tr(\cdot)_+$）、
$\Omega$（\texttt{cvxpy}+SCS，$\varepsilon=10^{-12}$），核对 $F\le G=\Omega$；
$G$ 与 $\Omega$ 的差在 $\pm10^{-13}$ 内，是对定理~\ref{thm:domdual}(ii) 强对偶的独立数值背书。
(B) 用 \texttt{sympy} 在 $\Q(\sqrt5,\sqrt7)$ 上把定理~\ref{thm:champglobal} 的八个恒等式化简为 $0$、
两个正性条件判为正，并用 \texttt{numpy.linalg.eigh} 交叉核对 $T$ 与 $T-\Sigma_0$ 的最小特征值（$\sim10^{-15}$）与谱。
(C) 对四类近腔掩码，把 SDP 解舍入到分母 $2^{40}$ 的二进有理数、对角加 $2^{-20}$ 裕度，
再用 \texttt{fractions.Fraction} 精确算出 $T$ 与 $T-\Sigma$ 各 $63$ 个主子式并逐个判非负（Sylvester）；
实际最小主子式量级 $10^{-25}>0$，说明裕度足以吸收舍入。
(D) 算八面体掩码：$\Omega=1.333333333333341$（与 $4/3$ 差 $7\times10^{-15}$），
并核对下界 $1+\tr((-\Sigma)_+)/6=4/3$ 与之相符。
(E) 冠军处：$\Lam_3(M^*)-\rho_0=0$、$\tr((-M^*)_+)-\tau=5.6\times10^{-17}$、
$\frac12(\|Y^*\|_1+\one^\top Y^*\one)-\rho_0=0$、互补松弛残差 $5.6\times10^{-17}$ 与 $3.3\times10^{-16}$、
$\mathrm{sgn}(Y^*)=2\Sigma_0-J$ 逐条目相符。
\textbf{注意}：脚本中 $F(A)$ 的数值一律只作交叉核对，报告里凡标注「已证」的结论都由 (B)(C) 的符号/有理算术给出。

\texttt{证书\_zonal天花板定理.py} 做四件事，全部在 $\Q(\sqrt5)$ 上用精确有理运算完成。
第一步输出见证测度 $\lambda$ 的三个质量 $\frac{5-\sqrt5}4$、$\frac{5+\sqrt5}4$、$\frac12$ 并验证其非负。
第二步验证三条矩条件 $\sum_i\lambda_i=3$、$\sum_i\lambda_it_i=1$、$\sum_i\lambda_it_i^2=1$；
后两条正是把方案 \eqref{eq:scheme} 中的 $\beta$ 通道与 $b$ 通道「钉死」所必需的，
其中 $\int(t^2-\frac13)\dd\lambda=1-\frac13\cdot3=0$ 使 $b$ 的取值完全不影响下界。
第三步验证目标值 $\sum_i\lambda_i(t_i)_+=\frac{\sqrt5}4+\frac34=\Phi$，
判据是 \texttt{simplify(差)\,==\,0} 而非浮点比较。
第四步对 $\ell=3,\dots,300$ 逐档验证 $\sum_i\lambda_iP_\ell(t_i)\ge0$，
把结果用 \texttt{nsimplify} 化到 $\Q(\sqrt5)$ 后判定符号；
其中 $\ell\in\{3,4,8,14\}$ 四档恰为 $0$，这与引理~\ref{lem:moments} 的解析论证一致。

需要说明分工：第四步在逻辑上\emph{不是}定理的证明，定理的证明由引理~\ref{lem:moments} 完成——
奇 $\ell\ge3$ 用 Laplace 型估计 $|P_\ell(x)|\le\sqrt{2/(\pi\ell)}\,(1-x^2)^{-1/4}$
（在 $x_0=1/\sqrt5$ 处给出 $0.487145<0.894427$），偶 $\ell\ge18$ 用同一估计（$0.198853<0.2$），
中间有限档 $\ell\in\{4,6,\dots,16\}$ 用精确有理算术逐个核验。
脚本的作用是把解析估计与手工核验之间的所有档位机器扫完，
并另以浮点对 $\ell=3,\dots,4000$ 做了独立扫描（最小值为 $0$）。
该脚本运行数分钟，主要开销在 \texttt{nsimplify}。

\texttt{证书\_冠军局部刚性.py} 在 $50$ 位下用三阶截断幂级数与基无关的二阶摄动公式重算冠军处 $9\times9$ Hessian，
舍入到 $20$ 位后以 Weyl 单调性给出 $\lambda_{\max}(\mathrm{Hess})\le-\kappa$（$\kappa=0.020962946762952704$），
并打印谱隙 $0.588930287016$ 与腔半径 $0.189960218562$，与独立高精度有限差在 $8$ 位吻合。
\texttt{证书\_三阶稳定性.py} 用同套级数机器（外加加边正则化 $L+v_0v_0^\top$，见注~\ref{rmk:bordered}）
求出全部 $165$ 个三阶导 $T_\tau$，输出 $\max|T|$、$\|T\|_F$、算子范数、$C_3$ 与近段半径 $R=3\kappa/C_3$，
并以五点中心差分（吻合到 $2.3\times10^{-10}$）与 $210$ 方向线性方程组残差（$1.6\times10^{-48}$）双重复核。
\texttt{证书\_腔面局部认证.py} 对三类单条目腔面的数值 argmax 重算梯度与 Hessian 并 Weyl 认证负定
（$\lambda_{\max}\le-0.0015\sim-0.0164$）。\texttt{证书\_负行挤压恒等式.py} 复核定理~\ref{thm:squeeze} 的四条恒等式与冠军处饱和表。
\texttt{证书\_双条目模式扫描.py} 为 P1--P6 双条目模式上界的数值来源（\textbf{数值}，非证明）。

\texttt{证书\_极值阶梯与Perron刚性.py} 分六段，输出末行给出「总计 $64$ 项检查，全部通过: True」。
(A) 在冠军 $W_6=K_1\vee C_5$ 上核验 Perron 与自洽临界点方程：$\Lam_3=\rho_0$（差 $<10^{-13}$）、
$\mathrm{sign}_+(Q)=\Sigma$ 逐条目相符、$Q_+x=\Lam_3x$ 残差 $1.8\times10^{-11}$、
$\mathrm{diag}(\Pi_3M\Pi_3)=\Lam_3p$ 残差 $7.8\times10^{-12}$。
(B) 逐条核验命题~\ref{prop:perronrig} 的 (a)--(h)：$\Lam_3=\alpha+\beta$、$\alpha=\|Qx\|^2$、
$\delta=\beta^2+\alpha(1-\alpha)$、$U^\top d=(\Lam_3-1)U^\top x$、
$QE(I-Q)=(Qx)((I-Q)x)^\top$ 且秩为 $1$（次奇异值 $<10^{-12}$）、$\tr((-M)_+)=\tau$、
$\ell=\alpha+2\beta=\Lam_3+\beta$、$Qx\perp(I-Q)x$、$\langle d,(I-Q)x\rangle=\rho(1-\alpha)$、
$Qd=(\Lam_3-1)Qx$、Collatz 包络 $[1.2587229,\,1.3846149]\ni\Lam_3$，
以及命题~\ref{prop:decouple} 的解耦恒等式与「\STW\ 损失精确劈成两半」
（$(1-\alpha)+(\phi-\ell)=0.054897595=2(\Phi-\rho_0)$，逐位相符）。
(C)(D) 用 \texttt{sympy} 把 $\kappa_3(3),\kappa_3(4),\kappa_3(5)$ 的闭式与全部辅助不等式
（如 $|P|=|M|=2$ 时 $4A(1-A)\le A(1-A)+\frac34\iff-\frac34(A-\frac12)^2\le0$；
$|P|=1$ 时 $32A^2-24A+7\ge0$，判别式 $-320<0$；$D(u^*)=\frac32$）判定为恒等式。
(E) 对七个图 $H$ 重算 $F(I+A(H))$ 并与闭式比较（差 $\le10^{-14}$）。
(F) 这是\textbf{独立于 (A)--(E) 的第二次复核}：不读任何缓存、从随机 $V\in\mathrm{St}(N,3)$ 出发，
用投影梯度上升直接极大化 $\rho((VV^\top)_+)$，$N=3,\dots,8$ 各 $60$ 个起点，
输出与闭式的差 $\le1.6\times10^{-15}$ 及 Stiefel-KKT \eqref{eq:stiefelkkt} 残差 $\le3\times10^{-9}$；
末行另做一次否证性体检（未发现 $\rho(Q_+)>\rho_0$）。
\textbf{注意}本段对 $N\ge6$ 的等式判定属于\textbf{数值}，只说明流形上升找不到更大的值，不构成证明。

\texttt{证书\_A2反例与修复.py} 是本报告第一个\textbf{纯精确算术}的否证证书，运行约 $14$ 秒，
末行输出「总计 $34$ 项检查，全部通过: True」。它的 C 段值得单独说明，因为这是全报告中
唯一一条\emph{完全不含浮点}的不等式证明：先用 $2^{-60}$ 网格把最优权有理化为 $\hat p$
（$|\hat p-p|_\infty\le2.8\times10^{-17}$，归一化后 $\sum\hat p_i=1$ 精确成立），
再利用相似性 $D_{\sqrt{\hat p}}^{-1}\hat M D_{\sqrt{\hat p}}=A\,\mathrm{diag}(\hat p)$
把含 $\sqrt{\hat p_i}$ 的对称矩阵换成\textbf{有理}矩阵，
于是特征多项式可在 $\Q$ 上精确算出（清分母后整数系数，首项 $99$ 位十进制）；
接着用\emph{自实现}的 Sturm 序列（纯 \texttt{Fraction}，不依赖 \texttt{sympy} 的实根隔离内部 API）
在 Cauchy 界 $R=2$ 内计数得 $7$ 个实根、二分隔离出 $7$ 个互不相交的有理区间、精化到宽度 $<10^{-32}$，
从中读出 $2$ 负 $1$ 零 $4$ 正与 $\tr((-M)_+)$ 的严格有理区间；
$\tau=\frac{-9+\sqrt5+2\sqrt{35}}{18}$ 的严格有理区间由整数开方 \texttt{isqrt} 到 $10^{-45}$ 得到
（宽 $1.7\times10^{-46}$），包含性用 \texttt{sympy} 在 $\Q(\sqrt5,\sqrt{35})$ 上\emph{精确}判定。
两个有理区间不相交（间距 $6.6\times10^{-3}$），故结论是定理。
交叉核对：Sturm 区间中点与 \texttt{numpy} 谱逐位相符（最大差 $3.3\times10^{-16}$）。

\texttt{证书\_N5特征多项式闭合.py}（第五版）与 \texttt{证书\_完全二部全局闭合.py}（第六版）是本轮的两份新证书，
末行分别输出「$45$ 项检查，全部通过: True」（$96.3$ s）与「$52$ 项检查，全部通过: True」（$137.1$ s），
\texttt{fail} 列表为空。两者的分工是：前者把 $N=5$ 这一个缺口按\emph{特征多项式因式分解}关掉，
后者把同一个论证\emph{推广到一切完全二部掩码与一切阶数}；后者的 $N=5$ 特例
（$\sigma_4(p)=p_{\min}$，核心引理退化为 $s\le\frac14+p_{\min}$）与前者的 C 段完全一致，
但证法短得多（Titu $+$ AM--GM 两行，而不是按 $p_{\min}$ 落在 $A/B$ 部分情形做齐次三次 SOS 分解），
因此前者保留下来作为独立复核。

\texttt{证书\_完全二部全局闭合.py} 的 D 段（核心引理）值得单独说明，因为它是\textbf{全符号}的、
不含任何数值搜索：$\phi''(z)=-2\sigma/(z+\sigma)^3$ 与 $\phi(z)-z/\sigma=-z^2/(\sigma(z+\sigma))$
都被 \texttt{sympy} 化简为恒等式；二元 Jensen 缺陷 $\Delta(u,v)$ 的分子被展开为 $\sigma(u-v)^2$、
分母系数被逐个判正；$k=4$ 的链式分解 $4\phi(\bar z)-\sum\phi(z_i)=\Delta(z_1,z_2)+\Delta(z_3,z_4)+2\Delta(\frac{z_1+z_2}2,\frac{z_3+z_4}2)$
与 $k=3$ 的补均值原子版本都被验证残差\emph{恰为} $0$；$g_c(z)=c-c^2\sigma/(z+c\sigma)$ 对 $c=1,2,3,4$ 逐个核对；
Titu 的五个分支 $r=0,\dots,4$ 的缺口分子用一个自实现的递归非负判据
（常数非负 / 偶次幂 / 系数全非负 / 乘积的每个因子非负）判定，输出即 \eqref{eq:titugap} 的五式；
$2-\Xi(m)$ 被 \texttt{factor} 成 $2m(1+14m)/[(2+3m)(1+4m)]$，分子分母分别判非负与判正。
D 段末尾另有 $150000$ 组随机权（$16$ 个 $(a,b)$、$N$ 最大 $10$）的数值双保险：
$\sup\bigl[h(-(\frac14+\sigma_4))-1\bigr]=-9.9\times10^{-5}\le0$、$\sup(s-\sigma_4)=0.2499105517\le\frac14$。
E 段穷举 $9$ 个「每部各两个 $\frac14$」的取等构型（$(a,b)$ 取 $(2,2)$ 到 $(5,5)$），
逐个核对 $\Lam_3=\frac54$、$s=\frac14$、$\sigma_4=0$、尾和 $=-\frac14$；
并单独认证 $K_{2,4}$ 与 $K_{3,3}$ 两类（$N=6$ 例外）的 $F=\frac54$ 精确。
G 段对 $13$ 个完全多部构型（$k=2,3,4$、$N=4,\dots,8$）符号核验清分母形式的行列式公式与
长期方程 $\sum_t\frac1{1-A_t}=k-1$，残差恒为 $0$，并对 $k=2,\dots,5$ 核验
$B=D_c(J_k-I_k)D_c$、二次型改写与 $\operatorname{spec}(J_k-I_k)=\{-1,k-1\}$；
H 段给出 $K_{1,3}$ 配 $p=(\frac25,\frac15,\frac15,\frac15)$ 的\textbf{精确}谱
$\{-\frac15,\frac15,\frac15,\frac45\}$（\texttt{sympy} 在 $\Q$ 上算特征多项式并求根），
即注~\ref{rmk:k1b} 的 $\frac65$ 下界。

\begin{sloppypar}
\texttt{\_exc23\_struct.json}、\texttt{\_exc19\_open.json} 与 \texttt{\_tab19.tex}
由 \texttt{证书\_例外类结构标注.py} 生成：它从 \texttt{\_danger23.json}
的 graph6 码解码出邻接矩阵（\texttt{ord(ch)-63} 标准解码，不依赖 \texttt{networkx} 的读入），
再判定「补图的每个连通分支是否都是团」（$\iff$ 原图是完全多部，并输出各部大小）与
「是否含诱导 $C_5$」（$5$ 元子集上恰 $5$ 条边且 $2$-正则），
最后扣除已闭合的三类（冠军、$K_{2,4}$、$K_{3,3}$）得到 $19$ 行表格。
该脚本运行不到 $1$ 秒，是命题~\ref{prop:tab19} 表格与注~\ref{rmk:tab19read}(i)(ii) 结构计数的唯一来源。
\end{sloppypar}

\medskip
\noindent\textbf{两条运行注意。}
其一，\texttt{证书\_极值阶梯与Perron刚性.py} 与 \texttt{证书\_优占对偶全局上界.py} 都通过
\texttt{sys.path.insert(0, <上级目录>)} 导入根目录下的 \texttt{\_core.py}，
因此请\emph{在课题根目录}或以脚本自身路径调用，勿把证书脚本单独拷出。
其二，\texttt{\_core.py} 中的 softmax 梯度曾在 2026-10-28 被修正：
早期写成 $g-\langle p,g\rangle\one$（仅在 $p$ 均匀时正确），正确式为 $g-(\sum g)\,p=g-\Lam_3p$，
其中 $g=\mathrm{diag}(\Pi_3M\Pi_3)$；有限差分核验到 $4\times10^{-10}$。
所有依赖 $F(A)$ 的数值（第~\ref{sec:exh7} 小节）都必须在修正后重跑，否则会得到系统性偏低的值。

\subsection{未随本报告分发的中间产物}

课题目录下另有大量以 \texttt{\_} 开头的临时脚本与文本转储
（文献抓取、图极限数值搜索、SDP/LP 离散化、区间算术试验等）。
它们不属于本报告的证明链条，且多数已被后续版本取代，故未纳入版本库的本次提交。
本轮有两个\textbf{例外}：\texttt{\_core.py}（已被两个证书脚本导入，见上表）
与 \texttt{\_kd.py}（Stiefel 流形直接上升的实现，证书 F 段的同源算法，
其梯度公式 $\nabla_V=2(I-VV^\top)GV$、$G=\mathbf 1\{Q\ge0\}\odot xx^\top$
即报告中的 \eqref{eq:gradfix}），二者随本次提交一并入库以便复跑；
此外，本轮把三份\textbf{数据文件}也一并入库，因为正文直接引用它们作为数值结论的出处：
\texttt{\_kappa\_profile.json}（SCF，$N=3,\dots,14$）、\texttt{\_kappa\_direct.json}（Stiefel 上升，$N=3,\dots,12$）、
\texttt{\_scan7F.json}（$1252$ 类的完整穷举记录）。
有了这三份文件，注~\ref{rmk:A2strata} 的分层表、「$F=\rho_0$ 的 $17$ 个掩码全部含诱导 $C_5$、
全部 $\mathrm{npos}=3$、负质量全部恰为 $\tau$」这条结构观察、
以及定理~\ref{thm:ladder} 中 $N\ge6$ 各档的数值，都可以用几行 \texttt{numpy} 直接复核，\emph{不必}重跑数小时的优化。
其余 \texttt{\_sd*.py}、\texttt{\_chk*.py}、\texttt{\_fix*.py}、\texttt{\_pp*.py}、\texttt{\_*.npy}
仍属临时产物，不入库。
若需要复核某一具体数值（例如第~7.3 节的 $C^2/(3n)=0.992343569348608$），
可按本附录「运行环境」一节的命令重算：该值由冠军权 $a^*,b^*$ 与 $V(\text{hub},\text{hub})$、
$V(\text{cyc},\text{cyc})$ 的闭式直接代入得到，\texttt{sympy} 化简后即为此数，不依赖任何优化过程。

\noindent{\small\textit{说明：正文以方括号缩写键指代下列文献，例如 \STW\ 对应 [STW26]，「DL23」对应 [DL23]，「KLMPT」对应 [KLM$^+$26]。第~\ref{sec:stwanat} 节解剖 \STW\ 的证明时用到投影常数一脉的文献（[Gr\"u60]、[KTJ94]、[CL09]、[DFFL22]、[DL23]），一并列出以便查阅。}}
\vspace{4pt}

\begin{thebibliography}{99}
\small

\bibitem[CKM$^+$17]{ckm17}
H.~Cohn, A.~Kumar, S.~Miller, D.~Radchenko, and M.~Viazovska,
\newblock The sphere packing problem in dimension 24,
\newblock \textit{Annals of Mathematics} \textbf{185} (2017), no.~3, 1017--1033.

\bibitem[CL09]{cl09}
B.~L.~Chalmers and G.~Lewicki,
\newblock Three-dimensional subspace of $\ell_\infty^5$ with maximal projection constant,
\newblock \textit{Journal of Functional Analysis} \textbf{257} (2009), no.~2, 553--592.

\bibitem[DFFL22]{dffl22}
B.~Der\k{e}gowska, M.~Fickus, S.~Foucart, and B.~Lewandowska,
\newblock On the value of the fifth maximal projection constant,
\newblock \textit{Journal of Functional Analysis} \textbf{283} (2022), no.~10, 109634.

\bibitem[DL23]{dl23}
B.~Der\k{e}gowska and B.~Lewandowska,
\newblock A simple proof of the Gr\"unbaum conjecture,
\newblock \textit{Journal of Functional Analysis} \textbf{285} (2023), no.~2, 109950.

\bibitem[DMS19]{dms19}
K.~C.~Das, S.~A.~Mojallal, and S.~Sun,
\newblock On the sum of the $k$ largest eigenvalues of graphs and symmetric matrices,
\newblock \textit{Linear Algebra and its Applications} \textbf{569} (2019), 175--194.

\bibitem[EMNA08]{emna08}
J.~B.~Ebrahimi, B.~Mohar, V.~Nikiforov, and A.~S.~Ahmady,
\newblock On the sum of two largest eigenvalues of a symmetric matrix,
\newblock \textit{Linear Algebra and its Applications} \textbf{429} (2008), no.~11--12, 2781--2787.

\bibitem[Gr\"u60]{gru60}
B.~Gr\"unbaum,
\newblock Projection constants,
\newblock \textit{Transactions of the American Mathematical Society} \textbf{95} (1960), no.~3, 451--465.

\bibitem[HW26]{hw26}
J.~Huang and W.~Wei,
\newblock The exact maximum of the spectral sum of graphs,
\newblock arXiv preprint arXiv:2607.23081 (2026).

\bibitem[KLM$^+$26]{klmpt26}
H.~Kumar, L.~Liu, H.~Monterde, S.~Pragada, and M.~Tait,
\newblock Maximum spectral sum of graphs,
\newblock arXiv preprint arXiv:2604.00512 (2026).

\bibitem[KTJ94]{ktj94}
H.~K\"onig and N.~Tomczak-Jaegermann,
\newblock Norms of minimal projections,
\newblock \textit{Journal of Functional Analysis} \textbf{119} (1994), no.~2, 253--280.

\bibitem[Moh09]{moh09}
B.~Mohar,
\newblock On the sum of $k$ largest eigenvalues of graphs and symmetric matrices,
\newblock \textit{Journal of Combinatorial Theory, Series B} \textbf{99} (2009), no.~2, 306--313.

\bibitem[Nik15]{nik15}
V.~Nikiforov,
\newblock Extrema of graph eigenvalues,
\newblock \textit{Linear Algebra and its Applications} \textbf{482} (2015), 158--190.

\bibitem[Sch42]{sch42}
I.~J.~Schoenberg,
\newblock Positive definite functions on spheres,
\newblock \textit{Duke Mathematical Journal} \textbf{9} (1942), no.~1, 96--108.

\bibitem[STW26]{stw26}
V.~Sivashankar, Q.~Tang, and T.~Wakhare,
\newblock Graph eigenvalues and projection constants,
\newblock arXiv preprint arXiv:2608.02429 (2026).

\bibitem[Tan26]{tan26}
Q.~Tang,
\newblock A sharp upper bound on the third adjacency eigenvalue of a graph,
\newblock arXiv preprint arXiv:2603.21181 (2026).

\bibitem[Via17]{via17}
M.~Viazovska,
\newblock The sphere packing problem in dimension 8,
\newblock \textit{Annals of Mathematics} \textbf{185} (2017), 991--1015.

\end{thebibliography}

\end{document}
